% Options for packages loaded elsewhere
\PassOptionsToPackage{unicode}{hyperref}
\PassOptionsToPackage{hyphens}{url}
\PassOptionsToPackage{dvipsnames,svgnames,x11names}{xcolor}
\documentclass[
  magyar,
  11pt]{article}
\usepackage{xcolor}
\usepackage[margin=2.5cm]{geometry}
\usepackage{amsmath,amssymb}
\setcounter{secnumdepth}{-\maxdimen} % remove section numbering
\usepackage{iftex}
\ifPDFTeX
  \usepackage[T1]{fontenc}
  \usepackage[utf8]{inputenc}
  \usepackage{textcomp} % provide euro and other symbols
\else % if luatex or xetex
  \usepackage{unicode-math} % this also loads fontspec
  \defaultfontfeatures{Scale=MatchLowercase}
  \defaultfontfeatures[\rmfamily]{Ligatures=TeX,Scale=1}
\fi
\usepackage{lmodern}
\ifPDFTeX\else
  % xetex/luatex font selection
  \setmainfont[]{DejaVu Serif}
  \setsansfont[]{DejaVu Sans}
  \setmonofont[]{DejaVu Sans Mono}
\fi
% Use upquote if available, for straight quotes in verbatim environments
\IfFileExists{upquote.sty}{\usepackage{upquote}}{}
\IfFileExists{microtype.sty}{% use microtype if available
  \usepackage[]{microtype}
  \UseMicrotypeSet[protrusion]{basicmath} % disable protrusion for tt fonts
}{}
\makeatletter
\@ifundefined{KOMAClassName}{% if non-KOMA class
  \IfFileExists{parskip.sty}{%
    \usepackage{parskip}
  }{% else
    \setlength{\parindent}{0pt}
    \setlength{\parskip}{6pt plus 2pt minus 1pt}}
}{% if KOMA class
  \KOMAoptions{parskip=half}}
\makeatother
\usepackage{longtable,booktabs,array}
\usepackage{calc} % for calculating minipage widths
% Correct order of tables after \paragraph or \subparagraph
\usepackage{etoolbox}
\makeatletter
\patchcmd\longtable{\par}{\if@noskipsec\mbox{}\fi\par}{}{}
\makeatother
% Allow footnotes in longtable head/foot
\IfFileExists{footnotehyper.sty}{\usepackage{footnotehyper}}{\usepackage{footnote}}
\makesavenoteenv{longtable}
\usepackage{graphicx}
\makeatletter
\newsavebox\pandoc@box
\newcommand*\pandocbounded[1]{% scales image to fit in text height/width
  \sbox\pandoc@box{#1}%
  \Gscale@div\@tempa{\textheight}{\dimexpr\ht\pandoc@box+\dp\pandoc@box\relax}%
  \Gscale@div\@tempb{\linewidth}{\wd\pandoc@box}%
  \ifdim\@tempb\p@<\@tempa\p@\let\@tempa\@tempb\fi% select the smaller of both
  \ifdim\@tempa\p@<\p@\scalebox{\@tempa}{\usebox\pandoc@box}%
  \else\usebox{\pandoc@box}%
  \fi%
}
% Set default figure placement to htbp
\def\fps@figure{htbp}
\makeatother
\ifLuaTeX
\usepackage[bidi=basic,shorthands=off,]{babel}
\else
\usepackage[bidi=default,shorthands=off,]{babel}
\fi
\ifPDFTeX
\else
\babelfont{rm}[]{DejaVu Serif}
\fi
\ifLuaTeX
  \usepackage{selnolig} % disable illegal ligatures
\fi
\setlength{\emergencystretch}{3em} % prevent overfull lines
\providecommand{\tightlist}{%
  \setlength{\itemsep}{0pt}\setlength{\parskip}{0pt}}
\usepackage{amsmath}
\usepackage{amssymb}
\usepackage{mathtools}
\providecommand{\boldsymbol}[1]{\mathbf{#1}}
\setcounter{secnumdepth}{0}
\usepackage{bookmark}
\IfFileExists{xurl.sty}{\usepackage{xurl}}{} % add URL line breaks if available
\urlstyle{same}
\hypersetup{
  pdftitle={Tudományos számítások},
  pdfauthor={Michael T. Heath},
  pdflang={hu-HU},
  colorlinks=true,
  linkcolor={teal},
  filecolor={Maroon},
  citecolor={Blue},
  urlcolor={Blue},
  pdfcreator={LaTeX via pandoc}}

\title{Tudományos számítások}
\usepackage{etoolbox}
\makeatletter
\providecommand{\subtitle}[1]{% add subtitle to \maketitle
  \apptocmd{\@title}{\par {\large #1 \par}}{}{}
}
\makeatother
\subtitle{Egy bevezető áttekintés --- Átdolgozott második kiadás (magyar
fordítás, folyamatban)}
\author{Michael T. Heath}
\date{}

\begin{document}
\maketitle

{
\setcounter{tocdepth}{3}
\tableofcontents
}
\section{Előszó a Classics
kiadáshoz}\label{elux151szuxf3-a-classics-kiaduxe1shoz}

Megtiszteltetés számomra, hogy a SIAM a Classics sorozatában jelenteti
meg a könyvemet. Ennek az új kiadásnak az előkészítése során jelentősen
átdolgoztam a könyv leginkább frissítésre szoruló részeit, elsősorban a
rendelkezésre álló szoftverekre mutató hivatkozásokat és az
irodalomjegyzéket. A könyv tartalmának nagyobbik része azonban
időtállónak bizonyult, és nagyrészt érintetlen maradt. Szeretném azt
hinni, hogy ez a könyv korszerű szemléletének köszönhető, de
valószínűbb, hogy inkább az alapozó jellegéből adódik: a numerikus
analízis és a tudományos számítások legalapvetőbb fogalmait és
módszereit tárgyalja, amelyek viszonylag lassan változnak (Newton, Euler
és Gauss nem publikált sokat az utóbbi időben!). A legjelentősebb
technikai változás talán az, hogy a Lagrange-interpoláció baricentrikus
megfogalmazását beépítettem a főszövegbe --- korábban ez csak egy
feladatban szerepelt.

A könyv eredeti megjelenése óta számos új irányzat bontakozott ki,
amelyek közül néhányra ez az új kiadás utal is, míg mások túlmutatnak a
könyv keretein. Pedagógiai szempontból talán a legfontosabb trend a
Python mint programozási nyelv és számítási környezet térnyerése, amely
az oktatásban és a szakmai gyakorlatban egyaránt jelentős hatást
gyakorolt. Ennek megfelelően a könyv szoftverhivatkozásait
kiegészítettem a megfelelő Python-szoftverekre mutató hivatkozásokkal. A
szoftverfronton másik figyelemre méltó új fejlemény a Chebfun, amely
több hagyományos megközelítést ötvöz innovatív módon, és ezáltal új,
mélyreható számítási lehetőségeket kínál. A nagyléptékű tudományos
számításokban a párhuzamosság egyre erőteljesebb kihasználásának
tendenciája töretlenül folytatódott, ám a párhuzamos algoritmusok komoly
tárgyalása már egy egész újabb könyvet igényelne.

Egy másik trend a sztochasztikus módszerek növekvő alkalmazása ---
például a lineáris algebrában és az optimalizálásban ---, amelyet
gyakran a gyorsan fejlődő területek, mint a gépi tanulás és az
adattudomány alkalmazásai motiválnak. A könyv már eredetileg is
tartalmazott egy rövid fejezetet a sztochasztikus módszerekről, amely
felkészít a más forrásokból való további tanulmányozásra. A
bizonytalanság kvantifikációja (UQ) iránti érdeklődés is ugrásszerűen
megnőtt az utóbbi években. A létezés, egyértelműség, kondicionáltság és
stabilitás következetes hangsúlyozásával a könyv szilárd fogalmi alapot
nyújt a számítási eredmények megbízhatóságának megítéléséhez, de a
legújabb UQ-technikák részletes tárgyalása ismét csak meghaladja a
kereteinket.

A korábbi kiadás megjelenése óta eltelt időben egy újdonság, hogy
elérhetővé vált mintegy 100 oktatási modul, amelyeket több diákkal
együttműködve készítettem. Ezek a Java-appletek dinamikusan szemléltetik
a tudományos számítások számos fogalmát és módszerét, és a könyvet
kísérő interaktív példákként szolgálnak. Alkalmasak oktatói tantermi
bemutatókra, illetve a diákok önálló tanulásához is. A modulok elérhetők
webböngészőn keresztül, illetve letölthetők egy önálló változatban,
amely a felhasználó saját számítógépén futtatható. A modulok szabadon
beszerezhetők a www.siam.org/books/cl80 címen.

Az eredeti előszóban megköszönt két mentorom (Alston Householder és Gene
Golub) mellett szeretném ezt az alkalmat arra is használni, hogy
elismerésemet fejezzem ki több szerzőnek, akiknek gondos és vonzó írói
stílusát következetesen inspirálónak és követendőnek találtam: Nick
Higham, Pete Stewart, Gil Strang, Nick Trefethen és Charlie Van Loan.
Végül ismét köszönetet szeretnék mondani feleségem, Mona türelméért és
megértéséért, aki mellett „nyugdíjazásom'' ellenére továbbra is dolgozom
--- illő, hogy ezzel az új kiadással végre ő is saját dedikációs oldalt
kapjon!

Michael T. Heath

\section{Előszó}\label{elux151szuxf3}

Ez a könyv átfogó áttekintést nyújt a numerikus módszerekről olyan
hallgatók és a gyakorlatban dolgozó szakemberek számára, akiknek
matematikai feladatokat kell megoldaniuk a számítási irányultságú
tudományterületeken. A hagyományos numerikus analízis könyvektől abban
különbözik, hogy a bemutatott algoritmusok mögött rejlő motivációra és
gondolatokra helyezi a hangsúlyt, nem pedig azok részletes elemzésére.
Igyekszem általános képet adni azokról a technikákról, amelyek az egyes
nagyobb problémakategóriákban rendelkezésre állnak, ideértve a helyes
feladatmegfogalmazást és az eredmények értelmezését, ugyanakkor ahol
csak lehetséges, a szakmailag megírt matematikai szoftverek használatát
szorgalmazom a megoldás előállítására. A könyv inkább a matematikai
szoftverek potenciális felhasználóinak szól, mint az ilyen szoftverek
lehetséges készítőinek. Remélem, hogy felhívhatom az olvasó figyelmét
arra, hogy milyen kérdések lényegesek a megfelelő módszerek és
szoftverek kiválasztásánál és körültekintő alkalmazásánál.

A University of Illinois-on ez a könyv egy átfogó, egy féléves numerikus
módszerek kurzus tankönyveként szolgál, és több célt is betölt:

\begin{itemize}
\tightlist
\item
  Befejező kurzusként végzős alapképzéses hallgatóknak, főként
  informatika, matematika és mérnöki szakokon
\item
  Alapozó kurzusként olyan mesterképzéses hallgatóknak informatika és
  matematika szakon, akik tudományos számításokra kívánnak szakosodni
\item
  Átfogó kurzusként olyan informatika mesterképzéses hallgatóknak, akik
  nem kívánnak tudományos számításokra szakosodni
\item
  Gyakorlati képzésként olyan természettudományos és mérnöki
  mesterképzéses hallgatóknak, akiknek numerikus módszereket és
  szoftvereket kell használniuk a kutatásukban. Ez egy kötelező
  alapkurzus a Mérnöki Kar által támogatott interdiszciplináris
  Computational Science and Engineering mesterképzési programban.
\end{itemize}

Hogy ezt a sokszínű hallgatói kört kiszolgáljam, a kurzus és a könyv
előkövetelményeit a lehető legalacsonyabb szintre szorítottuk: alapvető
jártasság a lineáris algebrában, a többváltozós kalkulusban és némi
ismeret a differenciálegyenletekről. Numerikus módszerekben való
előzetes jártasságot nem feltételezünk. A könyv ugyanakkor meglehetősen
kifinomult szemléletet alkalmaz, ezért a hallgatónál (vagy olvasónál)
bizonyos fokú érettség ajánlott. Az akadémiai környezeten túl remélem,
hogy a könyv referenciaként is hasznos lesz a gyakorló mérnökök és
tudósok számára, akiknek gyors áttekintésre lehet szükségük egy adott
számítási feladatról, valamint a megoldásához rendelkezésre álló
módszerekről és szoftverekről.

Noha a könyv hangsúlyozza a matematikai szoftverek használatát, néhány
más szoftverközpontú munkával ellentétben maga nem biztosít szoftvert,
és nem összpontosít konkrét szoftvercsomagokra, könyvtárakra vagy
környezetekre. Ehelyett minden problémakategóriához hivatkozásokat ad
olyan konkrét rutinokra, amelyek nyilvánosan elérhető tárolókból,
valamint a jelentősebb kereskedelmi könyvtárakból és csomagokból
szerezhetők be. Sok akadémiai és ipari számítási környezetben az ilyen
szoftverek már telepítve vannak, mindenesetre hivatkozások szerepelnek a
nyilvánosan hozzáférhető, az interneten keresztül szabadon elérhető
szoftverekre is. A könyvben szereplő számítógépes feladatok nem függenek
egy konkrét szoftver- vagy programnyelv-választástól.

A könyv felépítésének főbb elemei a következők:

Fejezetek: A könyv minden fejezete egy-egy jelentős számítási
problémakört tárgyal. A könyv első fele főként algebrai feladatokkal,
míg a második fele a deriváltakat és integrálokat érintő analitikus
feladatokkal foglalkozik. Az első két fejezet alapvető a könyv többi
részéhez, a további fejezetek azonban az oktató preferenciája szerint
különféle sorrendben vehetők. Pontosabban, a fejezetek közötti főbb
függőségek nagyjából a következők:

\begin{longtable}[]{@{}
  >{\raggedright\arraybackslash}p{(\linewidth - 10\tabcolsep) * \real{0.1273}}
  >{\raggedright\arraybackslash}p{(\linewidth - 10\tabcolsep) * \real{0.1818}}
  >{\raggedright\arraybackslash}p{(\linewidth - 10\tabcolsep) * \real{0.1273}}
  >{\raggedright\arraybackslash}p{(\linewidth - 10\tabcolsep) * \real{0.1818}}
  >{\raggedright\arraybackslash}p{(\linewidth - 10\tabcolsep) * \real{0.1273}}
  >{\raggedright\arraybackslash}p{(\linewidth - 10\tabcolsep) * \real{0.2545}}@{}}
\toprule\noalign{}
\begin{minipage}[b]{\linewidth}\raggedright
Fejezet
\end{minipage} & \begin{minipage}[b]{\linewidth}\raggedright
Függ ettől
\end{minipage} & \begin{minipage}[b]{\linewidth}\raggedright
Fejezet
\end{minipage} & \begin{minipage}[b]{\linewidth}\raggedright
Függ ettől
\end{minipage} & \begin{minipage}[b]{\linewidth}\raggedright
Fejezet
\end{minipage} & \begin{minipage}[b]{\linewidth}\raggedright
Függ ettől
\end{minipage} \\
\midrule\noalign{}
\endhead
\bottomrule\noalign{}
\endlastfoot
2 & 1 & 6 & 1, 2, 3, 5 & 10 & 1, 2, 5, 7, 9 \\
3 & 1, 2 & 7 & 1, 2 & 11 & 1, 2, 7, 9, 10 \\
4 & 1, 2, 3 & 8 & 1, 2, 7 & 12 & 1, 2, 7 \\
5 & 1, 2 & 9 & 1, 2, 5, 7 & 13 & 1 \\
\end{longtable}

Így a 7., 8., 12. és 13. fejezet sokkal korábban, míg a 3., 4. és 6.
fejezet sokkal később is tárgyalható, mint ahogy a könyvben szerepelnek.
Például a 3., 7. és 12. fejezet mindegyike valamilyen adatillesztést
tartalmaz, így kívánatos lehet egy egységként tárgyalni őket. Egy másik
példa: a lineáris egyenletrendszerek iterációs módszerei a parciális
differenciálegyenletekről szóló 11. fejezetben találhatók, mert innen
származnak a legfontosabb motivációs példák, de ennek az anyagnak nagy
része közvetlenül a 2. fejezet lineáris egyenletrendszerekre vonatkozó
direkt módszerei után is tárgyalható. Megjegyezzük, hogy a
sajátértékeket a könyv többi részében szabadon használjuk, így célszerű
a 4. fejezetet viszonylag korán feldolgozni, hacsak a hallgatók már nem
ismerik e téma alapjait.

A könyvben több mint elegendő anyag van egy teljes féléves kurzushoz,
így egy féléves kurzusban körültekintő válogatásra lehet szükség.
Például a véletlen számokról és a sztochasztikus szimulációról szóló 13.
fejezet csak érintőlegesen kapcsolódik a könyv többi részéhez, és
nyilvánvaló jelölt az elhagyásra (véletlenszám-generátorokat azonban a
könyv számos feladatában használunk). A teljes könyv két negyedéves
(quarter) vagy két féléves kurzus alatt dolgozható fel.

Példák: Szinte minden bevezetett fogalmat és módszert egy vagy több
példa illusztrál. Ezek a példák a viszonylag tömör általános tárgyalás
kiegészítését szolgálják, és a szöveg lényeges részeként olvasandók. A
példákat a lehető legegyszerűbben tartottam (olykor a túlzott
leegyszerűsítés kockázatával), hogy az olvasó könnyen követhesse őket.
Tapasztalatom szerint egy alaposan megértett egyszerű példa rendszerint
többet segít, mint egy realisztikusabb, de nehezebben követhető.

Szoftverek: Az egyes problémakategóriákhoz tartozó rendelkezésre álló
szoftverek listája ésszerűen átfogóra törekszik. Nem kíséreltem meg
kiemelni a „legjobb'' szoftvert egy adott feladathoz, részben azért,
mert rendszerint egyetlen csomag sem felsőbbrendű minden tekintetben,
részben pedig azért, hogy helyet adjak az eltérő szoftverelérhetőségnek
és programozási nyelv-választásnak, amely a különböző olvasóknál
érvényes lehet. Minden idézett szoftver legalább szakszerűen megírt, és
némelyik egészen kiváló.

Feladatok: A könyv sok feladatot tartalmaz, amelyek három kategóriába
sorolhatók:

\begin{itemize}
\tightlist
\item
  Áttekintő kérdések, amelyek rövid választ igénylő kérdések, és az
  alapvető fogalmi megértés ellenőrzésére szolgálnak
\item
  Feladatok, amelyek valamivel több gondolkodást, hosszabb válaszokat és
  esetleg kézi számításokat igényelnek
\item
  Számítógépes feladatok, amelyek némi programozást igényelnek, és
  gyakran meglévő szoftverek használatát is magukban foglalják.
\end{itemize}

Az áttekintő kérdések az olvasó öntesztelését szolgálják. Szándékosan
tartalmaznak némi ismétlést, hogy rögzítsék a kulcsfontosságú
gondolatokat, és megerősítsék az anyag elsajátításába vetett bizalmat. A
hosszabb feladatok írásbeli házi feladatok céljára alkalmasak.
Némelyikük egyszerű példákkal való kézi számítást igényel, míg másokat
úgy terveztünk, hogy kiegészítsék a főszövegből kihagyott levezetések és
bizonyítások részleteit. Ez utóbbiak különösen hasznosak lehetnek, ha a
könyvet elméletibb jellegű kurzushoz használjuk. A számítógépes
feladatok lehetőséget biztosítanak a gyakorlati tapasztalatszerzésre az
ajánlott szoftverek használatában, az egyes kategóriákhoz tartozó
tipikus feladatok megoldására. Egyes feladatok általánosak, mások
viszont közvetlenül kapcsolódnak a különböző tudományos és mérnöki
diszciplínák konkrét alkalmazásaihoz.

Változások a második kiadásban. Minden fejezet most motivációs
tárgyalással és egy vagy több szemléltető példával kezdődik, amelyet az
adott típusú feladat megoldásainak létezéséről, egyértelműségéről és
kondicionáltságáról szóló tárgyalás követ. A cél az, hogy a hallgató
jobban megértse, miért fontos az adott feladat, és hogy hogyan ismerheti
fel a feladat „jó'' vagy „rossz'' megfogalmazását, mielőtt a megoldó
algoritmusokat tárgyalnánk. A fő algoritmusokat most formálisan
megfogalmaztuk és a könnyebb hivatkozhatóság érdekében megszámoztuk. Az
irodalomjegyzéket naprakésszé tettük, a történeti jegyzeteket pedig
némileg kibővítettük. Az 1. fejezet előreható és hátraható hibáról,
valamint azok kapcsolatáról szóló tárgyalását kibővítettük és
tisztáztuk. A szingulárisérték-felbontásról szóló anyag nagy részét a 4.
fejezetből a 3. fejezetbe helyeztük át, ahol az alkalmazásai
kényelmesebben illeszkednek. A 4. fejezet sajátértékalgoritmusokról
szóló tárgyalását kibővítettük, több motivációval és részlettel ---
különösen a QR-iterációról ---, valamint néhány további módszerrel. A 6.
fejezet feltételes optimalizálásról szóló tárgyalását lényegesen
kibővítettük. A differenciálegyenletekről szóló fejezeteket némileg
átszerveztük, és a spektrális módszerek tárgyalását kibővítettük. A
gyors Fourier-transzformációról szóló 12. fejezetet átszerveztük és
letisztultabbá tettük néhány felesleges anyag törlésével.

Szeretném megköszönni mentoraim hatását, akik először bevezettek a
numerikus számítás váratlan bájaiba: Alston Householdernek és Gene
Golubnak. Hálás vagyok azért a bőséges visszajelzésért, amelyet a könyv
első kiadását használó hallgatóktól és oktatóktól kaptam. Az első kiadás
előzetes lektorai a következők voltak: Alan George, University of
Waterloo; Dianne O\textquotesingle Leary, University of Maryland; James
Ortega, University of Virginia; John Strikwerda, University of
Wisconsin; és Lloyd N. Trefethen, Oxford University. Az első kiadás
második kiadásra való előkészítésének lektorai a következők voltak:
Thomas Coleman, Cornell University; Robert Funderlic, North Carolina
State University; Thomas Hagstrom, University of New Mexico; Ramon
Moore, Ohio State University; Mark Pernarowski, Montana State
University; Linda Petzold, University of California at Santa Barbara; és
Brian Suchomel, University of Minnesota. Köszönöm valamennyi lektor
felbecsülhetetlen javaslatait. Ezen felül külön köszönetet szeretnék
mondani kollégáimnak --- Joerg Liesennek, Paul Saylornak és Eric de
Sturlernek, mindhármuknak a University of Illinois-ról ---, akik
mindegyike elolvasta az átdolgozott kézirat egy részét vagy egészét, és
felbecsülhetetlen visszajelzést adott. Köszönöm Melchior Franznak és
Justin Winklernek a tördeléssel kapcsolatos hasznos tanácsait, amelyek
elengedhetetlenek voltak a nyomdakész példány előkészítéséhez. Végül
mélyen nagyra értékelem feleségem, Mona türelmét és megértését e könyv
megírásának és átdolgozásának számtalan órája alatt. Nagy örömmel és
hálával ajánlom neki ezt a könyvet.

Michael T. Heath

A könyv honlapja: www.siam.org/books/cl80 A szerző honlapja:
heath.cs.illinois.edu

\section{Jelölések}\label{jeluxf6luxe9sek}

A könyvben használt jelölésrendszer meglehetősen szabványos, és kevés
magyarázatra szorul. Szabadon használunk vektor- és mátrixjelölést:
mátrixokat általában nagybetűs félkövér szedéssel, vektorokat kisbetűs
félkövér szedéssel, skalárokat pedig normál szedéssel jelölünk. Az
iterációs és komponensindexeket alsó indexszel jelöljük, általában
\(i\)-től \(n\)-ig. Például egy \(\boldsymbol{x}\) vektornak és egy
\(\boldsymbol{A}\) mátrixnak rendre \(x_i\) és \(a_{ij}\) elemei vannak.
Azokban a ritkább esetekben, amikor iterációs index és komponensindex
egyaránt szükséges, az iterációt zárójeles felső index jelzi, például
\(x_i^{(k)}\) egy sorozat \(k\)-adik vektorának \(i\)-edik komponensét
jelenti. Egyébként \(x_i\) az \(\boldsymbol{x}\) vektor \(i\)-edik
komponensét jelöli, míg \(\boldsymbol{x}_k\) egy sorozat \(k\)-adik
vektorát.

Az egyszerűség kedvéért elsősorban valós vektorokkal és mátrixokkal
foglalkozunk, noha az általunk tárgyalt elmélet és algoritmusok nagy
része kevés változtatással (vagy akár változtatás nélkül) átvihető a
komplex esetre. A valós számok halmazát \(\mathbb{R}\) jelöli, az
\(n\)-dimenziós valós euklideszi teret \(\mathbb{R}^n\), a valós
\(m \times n\) mátrixok halmazát pedig \(\mathbb{R}^{m \times n}\). A
megfelelő komplex megfelelők jelölése rendre \(\mathbb{C}\),
\(\mathbb{C}^n\) és \(\mathbb{C}^{m \times n}\).

Egy vektor vagy mátrix transzponáltját felső \(T\) index jelzi,
konjugált transzponáltját pedig felső \(H\) index (a Hermitian transpose
rövidítése). Hacsak másként nem jelezzük, minden vektort oszlopvektornak
tekintünk; egy sorvektort egy oszlopvektor explicit transzponálásával
adunk meg. A tördelési kényelem kedvéért egy oszlopvektor komponenseit
néha a megfelelő sorvektor transzponáltjaként jelöljük, például
\(\mathbf{x} = \begin{bmatrix} x_1 & x_2 \end{bmatrix}^T\). Két
\(n\)-dimenziós \(\mathbf{x}\) és \(\mathbf{y}\) vektor belső szorzata
(más néven skaláris szorzat) a mátrixszorzás speciális eseteként áll
elő, ezért \(\mathbf{x}^T\mathbf{y}\)-nal jelöljük (vagy
\(\mathbf{x}^H\mathbf{y}\)-nal a komplex esetben). Hasonlóképpen két
\(n\)-dimenziós \(\mathbf{x}\) és \(\mathbf{y}\) vektor külső szorzata,
amely egy \(n \times n\)-es mátrix, \(\mathbf{x}\mathbf{y}^T\) alakban
írható. Az \(n\)-edrendű egységmátrixot \(\mathbf{I}_n\) jelöli (vagy
csak \(\mathbf{I}\), ha az \(n\) dimenzió a szövegkörnyezetből
egyértelmű), \(i\)-edik oszlopát pedig \(\mathbf{e}_i\). A zérusmátrixot
\(\mathbf{O}\), a zérusvektort \(\mathbf{0}\), a zérusskalárt pedig
\(0\) jelöli. A \(d_1, \ldots, d_n\) átlóbeli elemekkel rendelkező
diagonális mátrixot \(\operatorname{diag}(d_1, \ldots, d_n)\) jelöli.
Vektorok vagy mátrixok közötti egyenlőtlenségek elemenként
értelmezendők. Az \(m \times n\)-es \(\mathbf{A}\) mátrix oszlopai által
kifeszített alteret \(\mathbb{R}^m\)-ben, azaz
\(\{\mathbf{A}\mathbf{x}: \mathbf{x} \in \mathbb{R}^n\}\)-t,
\(\operatorname{span}(\mathbf{A})\) jelöli.

Egyváltozós \(f(t)\) függvény közönséges deriváltját \(df/dt\) vagy
\(f'(t)\) jelöli. Több változóban értelmezett függvény --- például
\(u(x,y)\) --- parciális deriváltjait \(\partial u/\partial x\) alakban,
vagy egyes szövegkörnyezetekben alsó indexszel, például \(u_x\) alakban
jelöljük. A gradiensvektorok, a Jacobi-mátrixok és a Hesse-mátrixok
jelölését szükség szerint vezetjük be. Minden logaritmus természetes
alapú (alap \(e \approx 2.718\)), hacsak más alapot kifejezetten nem
jelölünk. A \(\approx\) szimbólumot hagyományos értelemben vett közelítő
egyenlőségre, a \(\cong\) szimbólumot pedig kifejezetten legkisebb
négyzetes közelítésekre tartjuk fenn.

A numerikus algoritmusok számítási költségét, avagy számítási
komplexitását általában a szükséges aritmetikai műveletek számával
mérjük. Hagyományosan a numerikus analízisben csak a szorzásokat
számolták (esetleg az osztásokat és négyzetgyökvonásokat is), mert a
szorzás általában lényegesen drágább volt az összeadásnál vagy
kivonásnál, és mert a legtöbb algoritmusban a szorzások jellemzően
hasonló számú összeadással párosulnak (például két vektor belső
szorzatának számításánál). Újabban az összeadás és a szorzás közötti
költségkülönbség nagyrészt eltűnt (sőt, sok modern mikroprocesszor
egyetlen szorzás-összeadás utasítással el tudja végezni a kapcsolt
szorzást és összeadást). A számítógépgyártók és -felhasználók szeretik a
legmagasabb lehetséges teljesítményt hirdetni, ezért egyre gyakoribb,
hogy minden egyes aritmetikai műveletet számolnak. Mivel azonban egyes
műveletszámok nagyon jól ismertek a hagyományos gyakorlat szerint, ebben
a könyvben rendszerint csak a szorzásokat számoljuk. Az értelem
egyértelművé tétele végett hozzátesszük a „és hasonló számú összeadás''
fordulatot, vagy külön jelezzük, ha mindkét műveletet számoljuk.

A műveletszám és a közelítések pontosságának számszerűsítéséhez gyakran
a „nagy O'' jelölést használjuk egy függvény nagyságrendjének, illetve
domináns tagjának megadására. Műveletszám esetén arra vagyunk
kíváncsiak, hogyan viselkedik a függvény, ha a feladat mérete ---
mondjuk \(n\) --- nagy lesz. Azt mondjuk, hogy

\[f(n) = \mathcal{O}(g(n))\]

(„\(f\) nagy O-ja \(g\)-nek'', vagy „\(f\) \(g\) nagyságrendű''), ha
létezik olyan \(C\) pozitív állandó, amelyre

\[|f(n)| \le C|g(n)|\]

minden elég nagy \(n\)-re teljesül. Például

\[2n^3 + 3n^2 + n = \mathcal{O}(n^3),\]

mert ahogy \(n\) nagy lesz, az \(n^3\)-nál kisebb rendű tagok relatíve
jelentéktelenné válnak. Pontossági becslésnél egy olyan \(h\) mennyiség
(például egy lépésköz vagy rácstávolság) kicsivé válását vizsgáljuk. Azt
mondjuk, hogy

\[f(h) = \mathcal{O}(g(h)),\]

ha létezik olyan \(C\) pozitív állandó, amelyre

\[|f(h)| \leq C|g(h)|\]

minden elég kicsi \(h\)-ra teljesül. Például

\[\frac{1}{1-h} = 1 + h + h^2 + h^3 + \dots = 1 + h + \mathcal{O}(h^2),\]

mert ahogy \(h\) kicsi lesz, a \(h^2\)-en túli elhagyott tagok relatíve
jelentéktelenné válnak. Vegyük észre, hogy a két definíció ekvivalens,
ha \(h = 1/n\).

\section{Tudományos számítások}\label{tudomuxe1nyos-szuxe1muxedtuxe1sok}

\subsection{1.1 Bevezetés}\label{bevezetuxe9s}

Ennek a könyvnek a tárgyát hagyományosan numerikus analízisnek nevezik.
A numerikus analízis olyan algoritmusok tervezésével és elemzésével
foglalkozik, amelyek számos területen --- különösen a
természettudományokban és a mérnöki tudományokban --- felmerülő
matematikai feladatokat oldanak meg. Ezért a numerikus analízist az
utóbbi időben tudományos számításoknak is nevezik. A tudományos
számításokat a számítástudomány legtöbb más ágától az különbözteti meg,
hogy folytonos --- nem pedig diszkrét --- mennyiségekkel foglalkozik.
Olyan függvényekkel és egyenletekkel dolgozik, amelyek alapul szolgáló
változói --- idő, távolság, sebesség, hőmérséklet, sűrűség, nyomás,
feszültség és hasonlók --- természetüknél fogva folytonosak.

A folytonos matematika legtöbb feladata (például szinte bármely feladat,
amely deriváltakat, integrálokat vagy nemlinearitást tartalmaz) még
elvben sem oldható meg véges számú lépésben egzaktul, ezért ezeket egy
(elméletileg végtelen) iteratív eljárással kell megoldani, amely végül
egy megoldáshoz konvergál. A gyakorlatban természetesen nem iterálunk
örökké, hanem csak addig, amíg a válasz megközelítőleg helyes, azaz
gyakorlati célokra „elég közel'' van a kívánt eredményhez. Így a
tudományos számítások egyik legfontosabb aspektusa a gyorsan konvergáló
iteratív algoritmusok megtalálása és a kapott közelítés pontosságának
felmérése. Ha a konvergencia kellően gyors, akkor még néhány olyan
feladat is --- amely véges algoritmussal megoldható, mint például a
lineáris algebrai egyenletrendszerek --- bizonyos esetekben előnyösebben
oldható meg iterációs módszerekkel, amint azt látni fogjuk.

Másodsorban a tudományos számításokat az jellemzi, hogy foglalkoznia
kell a közelítések hatásaival. Sok megoldási technika különféle típusú
közelítések egész sorozatát foglalja magában. Még a használt aritmetika
is csak közelítő, hiszen a digitális számítógépek nem tudnak minden
valós számot egzaktul ábrázolni. A jó algoritmusok szokásos
tulajdonságain --- például a hatékonyságon --- túl a numerikus
algoritmusoknak a lehető legmegbízhatóbbnak és legpontosabbnak kell
lenniük az útközben alkalmazott különféle közelítések ellenére.

\subsubsection{1.1.1 Számítási
feladatok}\label{szuxe1muxedtuxe1si-feladatok}

Ahogy a név is sugallja, a tudományos számításokban felmerülő sok
feladat a természettudományokból és a mérnöki tudományokból származik,
ahol a végső cél valamely természeti jelenség megértése vagy egy eszköz
tervezése. A számítógépes szimuláció egy fizikai rendszer vagy folyamat
számítógépes ábrázolása és utánzása. A számítógépes szimuláció
nagymértékben elmélyítheti a tudományos megértést, mert lehetővé teszi
olyan helyzetek vizsgálatát, amelyeket kizárólag elméleti, megfigyelési
vagy kísérleti úton nehéz vagy lehetetlen volna vizsgálni. Az
asztrofizikában például két ütköző fekete lyuk részletes viselkedése túl
bonyolult ahhoz, hogy elméletileg meghatározzuk, és lehetetlen
közvetlenül megfigyelni vagy laboratóriumban reprodukálni. Számítógépes
szimulációjához mindössze egy megfelelő matematikai leírásra (ebben az
esetben Einstein általános relativitáselméleti egyenleteire), az
egyenletek numerikus megoldására szolgáló algoritmusra és egy elegendően
nagy számítógépre van szükség, amelyen az algoritmus implementálható.

A számítógépes szimuláció nemcsak egzotikus vagy egyébként
hozzáférhetetlen helyzetek feltárására hasznos, hanem a „normális''
forgatókönyvek szélesebb körének vizsgálatát is lehetővé teszi, mint
amennyit ésszerű költséggel és idővel egyébként vizsgálni lehetne. A
mérnöki tervezésben a számítógépes szimuláció lehetővé teszi, hogy sok
tervezési alternatívát sokkal gyorsabban, olcsóbban és biztonságosabban
próbáljunk ki, mint a fizikai prototípusokat használó hagyományos
„építsd meg és teszteld'' módszerekkel. Ebben az értelemben a
számítógépes szimuláció virtuális prototípus-készítés néven is ismertté
vált. A gépjárműbiztonság javításánál például a törésteszt számítógépen
sokkal olcsóbb és veszélytelenebb, mint a valóságban, így a lehetséges
tervezési paraméterek tere sokkal alaposabban feltárható egy optimális
terv kidolgozásához.

A számítógépes szimuláció átfogó problémamegoldási folyamata általában a
következő lépéseket tartalmazza:

\begin{itemize}
\tightlist
\item
  \begin{enumerate}
  \def\labelenumi{\arabic{enumi}.}
  \tightlist
  \item
    Az érdeklődési körbe tartozó fizikai jelenség vagy rendszer
    matematikai modelljének kidolgozása --- általában valamilyen típusú
    egyenletekkel megfogalmazva.
  \end{enumerate}
\item
  \begin{enumerate}
  \def\labelenumi{\arabic{enumi}.}
  \setcounter{enumi}{1}
  \tightlist
  \item
    Algoritmusok kidolgozása az egyenletek numerikus megoldására.
  \end{enumerate}
\item
  \begin{enumerate}
  \def\labelenumi{\arabic{enumi}.}
  \setcounter{enumi}{2}
  \tightlist
  \item
    Az algoritmusok implementálása számítógépes szoftverben.
  \end{enumerate}
\item
  \begin{enumerate}
  \def\labelenumi{\arabic{enumi}.}
  \setcounter{enumi}{3}
  \tightlist
  \item
    A szoftver futtatása számítógépen, a fizikai folyamat numerikus
    szimulációjára.
  \end{enumerate}
\item
  \begin{enumerate}
  \def\labelenumi{\arabic{enumi}.}
  \setcounter{enumi}{4}
  \tightlist
  \item
    A számított eredmények valamilyen érthető formában --- például
    grafikus megjelenítéssel --- való bemutatása.
  \end{enumerate}
\item
  \begin{enumerate}
  \def\labelenumi{\arabic{enumi}.}
  \setcounter{enumi}{5}
  \tightlist
  \item
    A számított eredmények értelmezése és validálása, szükség esetén az
    előző lépések bármelyikének megismétlésével.
  \end{enumerate}
\end{itemize}

Az 1. lépést gyakran matematikai modellezésnek nevezik. Ehhez az
érintett tudományos vagy mérnöki terület speciális ismerete, valamint az
alkalmazott matematikában való jártasság szükséges. A 2. és 3. lépés ---
numerikus algoritmusok és szoftverek tervezése, elemzése, implementálása
és használata --- a tudományos számítások fő témája, így e könyvé is.
Noha a 2. és 3. lépésre összpontosítunk, lényeges, hogy mindezek a
lépések --- a probléma megfogalmazásától egészen az eredmények
értelmezéséig és validálásáig --- megfelelően történjenek, hogy az
eredmények értelmesek és hasznosak legyenek. A tudományos számítások
elveit és módszereit meglehetősen általános szinten lehet tanulmányozni,
amint azt látni fogjuk, de a konkrét feladat konkrét forrását és az
eredmények felhasználási módját mindig szem előtt kell tartani, mivel
mindegyik szempont hat a többire és fordítva. Például az eredeti
feladatmegfogalmazás erősen befolyásolhatja a numerikus eredmények
pontosságát, ami viszont kihat az eredmények értelmezésére és
validálására.

Egy matematikai feladatot jól kitűzöttnek mondunk, ha létezik megoldása,
az egyértelmű, és folytonosan függ a feladat adataitól. Az utóbbi
feltétel azt jelenti, hogy a feladat adatainak kis megváltozása nem okoz
hirtelen, aránytalan változást a megoldásban; ez a tulajdonság különösen
fontos a numerikus számításokban, ahol --- amint hamarosan látni fogjuk
--- az ilyen perturbációk általában elkerülhetetlenek. A jól kitűzöttség
erősen kívánatos a fizikai rendszerek matematikai modelljeiben, de ez
nem mindig érhető el. Például egy fizikai rendszer belső szerkezetének
csak külső megfigyelésekből való következtetése --- például
tomográfiában vagy szeizmológiában --- gyakran eredendően rosszul
kitűzött matematikai feladatokra vezet, amelyekben egymástól jelentősen
eltérő belső konfigurációk megkülönböztethetetlen külső megjelenéssel
rendelkeznek.

Még ha egy feladat jól kitűzött is, a megoldás akkor is nagyon
érzékenyen (bár folytonosan) reagálhat a feladat adataiban bekövetkező
perturbációkra. Ahhoz, hogy felmérjük az ilyen perturbációk hatásait, a
folytonosság kvalitatív fogalmán túl kell lépnünk, és meg kell
határoznunk a feladat érzékenységének mennyiségi mértékét. Ezenfelül
ügyelnünk kell arra is, hogy az adott feladat numerikus megoldására
használt algoritmus ne tegye az eredményeket érzékenyebbé annál, mint
ami az alapul szolgáló feladatban eleve benne rejlik (a „ne árts!''
hippokratészi eskü a numerikus matematikusokra éppúgy vonatkozik, mint
az orvosokra). Ez a követelmény vezet el a stabil algoritmus fogalmához.
Ezeket az általános fogalmakat és kérdéseket ebben a fejezetben vezetjük
be, majd a későbbi fejezetekben részletesen tárgyaljuk a konkrét típusú
számítási feladatokhoz.

\subsubsection{1.1.2 Általános
stratégia}\label{uxe1ltaluxe1nos-stratuxe9gia}

Egy adott számítási feladat megoldásának keresésében --- amely alapvető
általános stratégiaként végigvonul a könyvön --- az a cél, hogy egy
nehéz feladatot egy könnyebbre cseréljünk, amelynek ugyanaz a megoldása,
vagy legalábbis egy közeli rokon megoldása. Ennek a megközelítésnek a
példái:

\begin{itemize}
\tightlist
\item
  Végtelen dimenziós terek helyettesítése véges dimenziós terekkel
\item
  Végtelen folyamatok helyettesítése véges folyamatokkal, például
  integrálok vagy végtelen sorok helyettesítése véges összegekkel, vagy
  deriváltak véges differenciákkal
\item
  Differenciálegyenletek helyettesítése algebrai egyenletekkel
\item
  Nemlineáris feladatok helyettesítése lineáris feladatokkal
\item
  Magasabb rendű rendszerek helyettesítése alacsonyabb rendű
  rendszerekkel
\item
  Bonyolult függvények helyettesítése egyszerű függvényekkel, például
  polinomokkal
\item
  Általános mátrixok helyettesítése egyszerűbb alakú mátrixokkal
\end{itemize}

Például egy nemlineáris differenciálegyenlet-rendszer megoldásához
először egy nemlineáris algebrai egyenletrendszerrel helyettesíthetjük,
majd a nemlineáris algebrai rendszert egy lineáris algebrai rendszerrel,
majd a lineáris rendszer mátrixát olyan speciális alakú mátrixra
cserélhetjük, amellyel a megoldás könnyen számítható. E folyamat minden
lépésében meg kell győződnünk arról, hogy a megoldás nem változott meg,
vagy legalábbis a valódi megoldástól valamilyen előírt tűréshatáron
belül marad.

Ahhoz, hogy ez az általános stratégia egy adott feladatra működjön, a
következőkre van szükség:

\begin{itemize}
\tightlist
\item
  Egy olyan alternatív feladat vagy feladatosztály, amely könnyebben
  megoldható
\item
  Az adott feladatnak olyan transzformációja egy ilyen alternatív típusú
  feladattá, amely valamilyen értelemben megőrzi a megoldást
\end{itemize}

Így erőfeszítéseink nagy része olyan megfelelő feladatosztályok
azonosítására irányul, amelyeknek egyszerű megoldásuk van, valamint az
ilyen osztályokba vezető, megoldástartó transzformációk megtalálására.

Ideális esetben a transzformált feladat megoldása azonos az eredeti
feladat megoldásával, de ez nem mindig lehetséges. Az utóbbi esetben a
megoldás csak közelíti az eredeti feladatét, de a pontosság általában
tetszőlegesen jóvá tehető további munka és tárhely árán. Ezért
elsődleges feladatunk az ilyen közelítő megoldás pontosságának becslése
és a valódi megoldáshoz való konvergencia igazolása a határértékben.

\subsection{1.2 Közelítések a tudományos
számításokban}\label{kuxf6zeluxedtuxe9sek-a-tudomuxe1nyos-szuxe1muxedtuxe1sokban}

\subsubsection{1.2.1 A közelítés
forrásai}\label{a-kuxf6zeluxedtuxe9s-forruxe1sai}

A számítási tudományban a közelítésnek, illetve a pontatlanságnak számos
forrása van. Néhány közelítés már a számítás megkezdése előtt felléphet:

\begin{itemize}
\tightlist
\item
  Modellezés: A vizsgált feladat vagy rendszer bizonyos fizikai vonásait
  leegyszerűsíthetjük vagy elhagyhatjuk (pl. súrlódás, viszkozitás,
  légellenállás).
\item
  Empirikus mérések: A laboratóriumi műszereknek véges a pontossága.
  Pontosságukat tovább korlátozhatja a kis mintaelemszám, vagy a kapott
  leolvasásokat véletlen zaj, illetve szisztematikus torzítás is
  terhelheti. Például még a fontos fizikai állandók --- mint Newton
  gravitációs állandója vagy Planck állandója --- legalaposabb
  méréseiből is jellemzően legfeljebb nyolc vagy kilenc értékes jegyet
  kapunk, a legtöbb laboratóriumi mérés pedig ennél jóval pontatlanabb.
\item
  Korábbi számítások: A bemenő adatok egy korábbi számítási lépésből
  származhatnak, amelynek eredményei csak közelítők voltak.
\end{itemize}

Az imént felsorolt közelítések általában kívül esnek a befolyásunkon, de
fontos szerepet játszanak annak meghatározásában, hogy egy számítástól
milyen pontosság várható. Figyelmünket elsősorban azokra a közelítésekre
fordítjuk, amelyekre van valamilyen befolyásunk. Ezek a számítás közben
fellépő szisztematikus közelítések a következők:

• Csonkítás vagy diszkretizáció: egy matematikai modell bizonyos
jellemzőit elhagyhatjuk vagy leegyszerűsíthetjük (pl. deriváltakat véges
differenciákkal helyettesítünk, vagy egy végtelen sornak csak véges sok
tagját használjuk).

• Kerekítés: legyen szó kézi számításról, zsebszámológépről vagy
digitális számítógépről, a valós számok és a rajtuk végzett aritmetikai
műveletek ábrázolása végső soron véges mennyiségű precízióra
korlátozódik, így általában nem egzakt.

Egy számítás végső eredményeinek pontossága tükrözheti ezen közelítések
bármelyikének vagy akár mindegyikének kombinációját, a fellépő
perturbációkat pedig a megoldandó feladat természete vagy a használt
algoritmus --- netán mindkettő --- fel is erősítheti. Az ilyen
közelítéseknek a numerikus algoritmusok pontosságára és stabilitására
gyakorolt hatásának vizsgálatát hagyományosan hibaelemzésnek nevezik.

\textbf{1.1. Példa. Közelítések.} A Föld felszínét ki lehetne számítani
az

\[A = 4\pi r^2\]

képlet alapján egy \(r\) sugarú gömb felszínére. Ennek a képletnek a
használata a számításban számos közelítést foglal magában:

\begin{itemize}
\tightlist
\item
  A Földet gömbként modellezzük, ami a tényleges alakjának idealizálása.
\item
  Az \(r \approx 6370\) km sugárérték empirikus mérések és korábbi
  számítások kombinációján alapul.
\item
  A \(\pi\) értékét egy végtelen határértékfolyamat adja meg, amelyet
  valahol csonkolni kell.
\item
  A bemenő adatok numerikus értékei, valamint a rajtuk végzett
  aritmetikai műveletek eredményei egy számítógépen vagy számológépen
  kerekítve vannak.
\end{itemize}

A számított eredmény pontossága mindezektől a közelítésektől függ.

\subsubsection{1.2.2 Abszolút hiba és relatív
hiba}\label{abszoluxfat-hiba-uxe9s-relatuxedv-hiba}

Egy hiba jelentősége nyilvánvalóan összefügg a mért vagy számított
mennyiség nagyságával. Például 1-es hiba sokkal kisebb jelentőségű a
Föld népességének számolásánál, mint egy telefonfülke lakóinak
számolásánál. Ez motiválja az \emph{abszolút hiba} és a \emph{relatív
hiba} fogalmát, amelyeket a következőképpen definiálunk:

\[\text{abszolút hiba} = \text{közelítő érték} - \text{valódi érték},\]

\[\text{relatív hiba} = \frac{\text{abszolút hiba}}{\text{valódi érték}}.\]

Egyes szerzők az abszolút hibát a fenti különbség abszolút értékeként
definiálják, de mi az abszolút értéket (vagy vektorok és mátrixok esetén
a normát) csak akkor vesszük kifejezetten, amikor kizárólag a hiba
nagyságára van szükségünk. Megjegyezzük, hogy a relatív hiba nincs
definiálva, ha a valódi érték nulla.

A relatív hiba százalékban is kifejezhető, amely egyszerűen a relatív
hiba 100-zal szorozva. Így például egy 0,1 abszolút hiba a 10 valódi
értékhez viszonyítva 0,01, azaz 1 százalékos relatív hibának felel meg.
Egy teljesen hibás közelítés legalább 1, azaz legalább 100 százalékos
relatív hibának felelne meg, ami azt jelenti, hogy az abszolút hiba
akkora, mint maga a valódi érték.

A relatív hiba egy másik értelmezése: ha egy közelítő érték relatív
hibája körülbelül \(10^{-p}\), akkor tizedestört alakjában körülbelül
\(p\) helyes értékes jegye van (az értékes jegyek az első nem nulla
számjegy és az azt követő összes jegy). Ezzel kapcsolatban érdemes
különbséget tenni a \emph{precízió} és a \emph{pontosság} között: a
precízió azon jegyek számára utal, amelyekkel a számot kifejezzük, a
pontosság pedig a kívánt mennyiség közelítésében helyes értékes jegyek
számára (azaz a relatív hibára). Például a 3,252603764690804 egy nagyon
precíz szám, de a \(\pi\) közelítéseként nem túl pontos. Amint hamarosan
látni fogjuk, az, hogy egy mennyiséget adott precízióval számítunk,
önmagában nem jelenti, hogy az eredmény erre a precízióra pontos is
lesz.

Az abszolút és relatív hiba közötti összefüggés egy hasznos kifejezése a
következő:

\[\text{közelítő érték} = (\text{valódi érték}) \times (1 + \text{relatív hiba}).\]

Természetesen általában nem ismerjük a valódi értéket; ha ismernénk, nem
kellene bajlódnunk a közelítésével. Ezért a hibát rendszerint csak
becsüljük vagy megkorlátozzuk, nem pedig egzaktul kiszámítjuk, hiszen a
valódi érték ismeretlen. Ugyanezen okból a relatív hibát gyakran a
közelítő értékhez, nem pedig a valódi értékhez viszonyítva veszik,
eltérően a fenti definíciótól.

\subsubsection{1.2.3 Adathiba és számítási
hiba}\label{adathiba-uxe9s-szuxe1muxedtuxe1si-hiba}

Amint láttuk, egyes hibák a bemenő adatoknak tulajdoníthatók, míg mások
a későbbi számítási folyamatok eredményei. Bár ez a megkülönböztetés nem
mindig egyértelmű (a kerekítés például érintheti a bemenő adatokat és a
későbbi számítási eredményeket is), mégis segít megérteni a
közelítéseknek a numerikus számításokra gyakorolt összhatását.

A legtöbb realisztikus feladat többdimenziós, de az egyszerűség kedvéért
ebben a fejezetben csak egydimenziós feladatokat tekintünk; a definíciók
és eredmények magasabb dimenziókra való kiterjesztése egyértelmű,
általában csak az abszolút értékek megfelelő normákra cserélését igényli
(lásd a 2.3.1. szakaszt). Egy tipikus egydimenziós feladat felfogható
egy függvény --- mondjuk \(f: \mathbb{R} \to \mathbb{R}\) --- értékének
számításaként, amely egy adott bemenő értéket egy kimenő eredményre
képez. Jelölje a bemenet valódi értékét \(x\), így a kívánt valódi
eredmény \(f(x)\). Tegyük fel, hogy pontatlan bemenettel, mondjuk
\(\hat{x}\)-szel kell dolgoznunk, és csak egy közelítést --- mondjuk
\(\hat{f}\)-et --- tudunk számítani a függvényre. Akkor azzal a jól
ismert matematikai trükkel, hogy ugyanazt a mennyiséget hozzáadjuk és
kivonjuk, úgy, hogy az összeg változatlan marad, kapjuk:

\[\begin{aligned}
\text{teljes hiba} &= \hat{f}(\hat{x}) - f(x) \\
&= \bigl(\hat{f}(\hat{x}) - f(\hat{x})\bigr) + \bigl(f(\hat{x}) - f(x)\bigr) \\
&= \text{számítási hiba} + \text{propagált adathiba}.
\end{aligned}\]

Ebben az összegben az első tag a pontos és a közelítő függvény
különbsége ugyanarra a bemenetre, tehát tiszta \emph{számítási hibának}
tekinthető. A második tag a pontos függvényértékek különbsége a bemenet
hibája miatt, tehát tiszta \emph{propagált adathibának} tekinthető.
Megjegyezzük, hogy az algoritmus választásának nincs hatása a propagált
adathibára.

\textbf{1.2. Példa. Adathiba és számítási hiba.} Tegyük fel, hogy nincs
számítógépünk vagy számológépünk, és „gyors és piszkos'' közelítésre van
szükségünk a \(\sin(\pi/8)\) értékére. Először a \(\pi\) valamilyen
értékére van szükségünk a bemenethez. Röviden fontolóra vesszük a
\(\pi \approx 22/7\) klasszikus, iskoláskorunkból megmaradt közelítést,
de úgy döntünk, túl sok munka volna átváltani a kívánt tizedestört
formátumra, és inkább a \(\pi \approx 3\) egyszerű „bibliai''
közelítésnél maradunk, így a tényleges bemenetünk \(3/8\). A
függvényérték kiszámításához felidézzük a kalkulusból, hogy kis
argumentumokra jó közelítést ad a Taylor-sor első tagja, ami \(\sin(x)\)
esetén egyszerűen \(x\). Végeredményünk tehát

\[\sin(\pi/8) \approx \sin(3/8) \approx 3/8 = 0{,}3750.\]

A fenti közelítések közül az első --- egy perturbált \(\hat{x} = 3/8\)
bemenet használata a valódi \(x = \pi/8\) helyett --- propagált
adathibát okoz: még ha egzaktul értékelnénk is ki a szinuszfüggvényt,
hibás eredményt kapnánk, mert hibás bemenetet használtunk. A második
közelítés egy számítási hibát jelent, noha a „számításunk'' ebben az
esetben csupán a bemenet lemásolása volt! (A számítási hibák gyakran
ilyen „mulasztásos hibák'', bár általában nem ennyire szélsőségesek.) A
korábban bevezetett jelölésben az \(\hat{f}(x) = x\) csonkolt
matematikai kifejezést használtuk a valódi \(f(x) = \sin(x)\) függvény
helyett, ami azt jelenti, hogy helyes bemenettel is helytelen eredményt
kaptunk volna. A teljes pontosságunkat e két közelítés kombinációja
határozza meg.

Később --- miután hozzáférünk egy zsebszámológéphez --- megállapítjuk,
hogy négy tizedesjegyre a helyes válasz

\[\sin(\pi/8) \approx 0{,}3827,\]

így a teljes hiba

\[\hat{f}(\hat{x}) - f(x) \approx 0{,}3750 - 0{,}3827 = -0{,}0077.\]

Észrevéve, hogy a perturbált bemenetre a helyes válasz

\[f(\hat{x}) = \sin(3/8) \approx 0{,}3663,\]

látjuk, hogy a pontatlan bemenet használata által okozott propagált
adathiba

\[f(\hat{x}) - f(x) = \sin(3/8) - \sin(\pi/8) \approx 0{,}3663 - 0{,}3827 = -0{,}0164.\]

A végtelen sor csonkolásából adódó számítási hiba pedig

\[\hat{f}(\hat{x}) - f(\hat{x}) = 3/8 - \sin(3/8) \approx 0{,}3750 - 0{,}3663 = 0{,}0087.\]

Ez a két hiba összege adja a megfigyelt teljes hibát. Ebben a konkrét
példában a két hibának ellentétes az előjele, tehát részben kioltják
egymást; más körülmények között azonos előjelűek is lehetnek, és ekkor
felerősítik egymást. Ennél a bemenetnél a propagált adathiba és a
számítási hiba nagyságrendileg hasonló --- körülbelül kétszeres
szorzóval tér el ---, de bármelyik hibaforrás dominánssá válhat más
bemenő értékek esetén. Ha a \(\pi\)-re és a \(\sin(x)\)-re ugyanezeket a
közelítéseket használjuk, a propagált adathiba sokkal kisebb
bemeneteknél lenne domináns, míg a számítási hiba sokkal nagyobb
bemeneteknél (Miért?). A teljes hiba csökkentésére használhatnánk
pontosabb értéket \(\pi\)-re (ami csökkentené a propagált adathibát),
vagy pontosabb matematikai leírást \(\sin(x)\)-re (pl. a végtelen sor
több tagját), ami csökkentené a számítási hibát.

\subsubsection{1.2.4 Csonkítási hiba és kerekítési
hiba}\label{csonkuxedtuxe1si-hiba-uxe9s-kerekuxedtuxe9si-hiba}

A számítási hiba (azaz a számítás során elkövetett hiba) tovább bontható
\emph{csonkítási} (vagy \emph{diszkretizációs}) \emph{hibára} és
\emph{kerekítési hibára}:

\begin{itemize}
\tightlist
\item
  A csonkítási hiba a valódi eredmény (a tényleges bemenetre) és az
  adott algoritmus egzakt aritmetikával adott eredménye közötti
  különbség. Olyan közelítésekből ered, mint egy végtelen sor
  csonkolása, deriváltak véges differenciákkal való helyettesítése, vagy
  egy iterációs sorozat konvergencia előtti megszakítása.
\item
  A kerekítési hiba az adott algoritmus egzakt aritmetikával adott
  eredménye és ugyanezen algoritmus véges precíziójú, kerekített
  aritmetikával adott eredménye közötti különbség. A valós számok és a
  rajtuk végzett aritmetikai műveletek ábrázolásának pontatlanságából
  ered, amelyet az 1.3. szakaszban részletesen tárgyalunk.
\end{itemize}

A definíció szerint tehát a számítási hiba egyszerűen a csonkítási hiba
és a kerekítési hiba összege. Az 1.2. példában a bemenetet kerekítettük,
de a számítás során nem volt kerekítési hiba, így a számítási hiba
kizárólag a végtelen sor egyetlen tagjának használatából eredő
csonkítási hibából állt. Több tag használata a sorban csökkentette volna
a csonkítási hibát, de valószínűleg némi kerekítési hibát vezetett volna
be a sor kiértékeléséhez szükséges aritmetikába. A csonkítási hiba és a
kerekítési hiba közötti ilyen kompromisszumok nem szokatlanok.

\textbf{1.3. Példa. Véges differencia-közelítés.} Egy differenciálható
\(f: \mathbb{R} \to \mathbb{R}\) függvényre tekintsük az első derivált
véges differencia-közelítését:

\[f'(x) \approx \frac{f(x+h) - f(x)}{h}.\]

Taylor tétele szerint

\[f(x+h) = f(x) + f'(x)h + f''(\theta)h^{2}/2\]

valamilyen \(\theta \in [x, x+h]\)-ra, így a véges differencia-közelítés
csonkítási hibája felülről korlátozható \(Mh/2\)-vel, ahol \(M\) egy
\(|f''(t)|\)-re adott korlát \(x\) közelében. Feltéve, hogy a
függvényértékek hibája \(\epsilon\)-nal korlátozott, a véges
differencia-képlet kiértékelésének kerekítési hibája \(2\epsilon/h\)-val
korlátozható. A teljes számítási hiba így felülről korlátozható két
függvény összegével:

\[\frac{Mh}{2} + \frac{2\epsilon}{h},\]

ahol az első tag \(h\) csökkenésével csökken, a második nő. Ezért \(h\)
lépésköz megválasztásánál kompromisszum áll fenn a csonkítási hiba és a
kerekítési hiba között. Ezt a függvényt \(h\) szerint deriválva és a
deriváltat nullával egyenlővé téve látható, hogy a teljes számítási
hibára adott korlát minimumát

\[h = 2\sqrt{\epsilon/M}\]

esetén veszi fel. Egy tipikus példa az 1.1. ábrán látható, ahol a véges
differencia-közelítés teljes számítási hibáját --- a csonkítási és
kerekítési hibára adott külön-külön korlátokkal együtt --- az
\(f(x) = \sin(x)\) függvényre \(x = 1\)-ben ábrázoljuk \(h\) lépésköz
függvényében, \(M = 1\) mellett és \(\epsilon \approx 10^{-16}\) értékű
számítógépen. Látjuk, hogy a teljes hiba valóban minimumot ér el
\(h \approx 10^{-8} \approx \sqrt{\epsilon}\) környékén. A teljes hiba
nagyobb \(h\) értékekre a növekvő csonkítási hiba miatt, kisebb \(h\)
értékekre pedig a növekvő kerekítési hiba miatt nő.

\pandocbounded{\includegraphics[keepaspectratio]{_page_8_Figure_7.jpeg}}

1.1. ábra: A véges differencia-közelítés számítási hibája adott
lépésközre.

A csonkítási hiba csökkenthető lenne pontosabb véges differencia-formula
használatával, például a középpontos differencia-közelítéssel (lásd a
8.6.1. szakaszt):

\[f'(x) \approx \frac{f(x+h) - f(x-h)}{2h}.\]

A kerekítési hiba magasabb precíziójú aritmetikával csökkenthető, ha
rendelkezésre áll.

Bár egy adott számításban a csonkítási hiba és a kerekítési hiba
egyaránt fontos szerepet játszhat, általában egyik vagy másik válik a
teljes számítási hiba domináns tényezőjévé. Nagyjából elmondható, hogy a
kerekítési hiba rendszerint a tisztán algebrai, véges megoldási
algoritmusú feladatokban dominál, míg a csonkítási hiba az integrálokat,
deriváltakat vagy nemlinearitást tartalmazó feladatokban, amelyek
gyakran elméletileg végtelen megoldási folyamatot igényelnek.

A különböző hibatípusok között tett megkülönböztetéseink fontosak a
numerikus algoritmusok viselkedésének és a pontosságukat befolyásoló
tényezők megértéséhez, de az egyes hibatípusokat általában nem szükséges
--- és nem is mindig lehetséges --- pontosan kvantifikálni. Valójában,
amint hamarosan látni fogjuk, gyakran előnyös, ha az összes hibát
összevonjuk és a bemenő adatok hibájára visszavezetjük.

\subsubsection{1.2.5 Előreható hiba és hátraható
hiba}\label{elux151rehatuxf3-hiba-uxe9s-huxe1trahatuxf3-hiba}

Egy számított eredmény minősége egyéb tényezők mellett a bemenő adatok
minőségétől is függ. Ha például a bemenő adatok csak --- mondjuk ---
négy értékes jegyre pontosak, akkor legfeljebb négy értékes jegyet
várhatunk a számított eredményben is, bárhogy is végezzük a számítást. A
híres számítástechnikai szentencia, a „szemét megy be, szemét jön ki''
(garbage in, garbage out) ezt a megfigyelést viszi el logikus végletéig.
Ezért a számított eredmények minőségének értékelésekor nem szabad
figyelmen kívül hagynunk a bemenő adatoknak azok bizonytalansági
szintjén belüli perturbációi lehetséges hatásait.

\textbf{1.4. Példa. Az adathiba hatásai.} Tegyük fel, hogy egy ország
népességét szeretnénk előrejelezni tíz év múlva. Először egy matematikai
modellre van szükségünk, amely leírja a népesség időbeli változásait.
Feltéve, hogy a születések és a halálozások egyaránt arányosak a
mindenkori népességgel, a

\[P(t + \Delta t) = P(t) + (B - D)P(t)\Delta t\]

egyszerű modellhez jutunk, ahol \(P(t)\) a népességet jelöli \(t\)
időpontban, \(\Delta t\) egy adott időintervallum (mondjuk egy év),
\(B\) és \(D\) pedig a születési és halálozási arány (például a nettó
növekedési ráta \(B - D = 0{,}04\), azaz évi 4 százalék lehet).
Használhatjuk ezt a diszkrét modellt egy rögzített \(\Delta t\)
időintervallummal, vagy vehetjük a \(\Delta t \to 0\) határátmenetet,
hogy a

\[dP(t)/dt = (B - D)P(t)\]

differenciálegyenlethez jussunk, amelynek megoldása a jól ismert
exponenciális növekedési törvény:

\[P(t) = P(0) \exp\bigl((B - D)t\bigr).\]

Mindkét esetben a modell nyilvánvalóan csak közelítése a valóságnak;
például elhanyagoltuk a bevándorlást és a korlátozott kapacitás bármely
hatását. Bár e szemléltetésben az ilyen modellezési hibákat figyelmen
kívül hagyjuk, szinte minden valós tudományos feladatban való
jelenlétüknek mérsékelnie kell a további számításokban elérhető
pontossággal szembeni elvárásainkat. Akár a diszkrét, akár a folytonos
modellt használjuk népesség-előrejelzéshez, ismernünk kell a jelenlegi
népességet, valamint a születési és halálozási arányokat. Közismert,
hogy egy népszámlálásnál nehéz mindenkit megszámolni anélkül, hogy
valakit kihagynánk, így a kiindulási népesség csak korlátozott
pontossággal ismert. Ennek megfelelően a kiindulási népességet
rendszerint „kerek számokkal'' (azaz csak néhány értékes jeggyel) adják
meg, utalva az érték bizonytalansági szintjére. Ezt nem szabad
kerekítési hibának tekinteni, mert a kapott érték éppúgy lehet helyes,
mint a precízebb érték, hiszen az elhagyott jegyek kétségesek vagy
értelmetlenek voltak. Hasonlóképpen a születési és halálozási arányok is
--- amelyek sok diszkrét eseményre vett átlagok --- csak korlátozott
pontossággal ismertek.

A modellünk bemenetében levő bizonytalanságok szükségszerűen bizonyos
bizonytalanságot okoznak a kapott népesség-előrejelzésben is.
Következésképpen úgy tekinthetjük a modellünket, hogy a bemenő tér egy
„elmosódott'' tartományát a kimenő tér egy „elmosódott'' tartományára
képezi le (ez utóbbi az összes lehetséges bemenet-kombináció
bizonytalansági tartományán belül adódó összes lehetséges eredmény
halmaza). A modell implementálása során némi számítási hiba (csonkítási
vagy kerekítési hiba) léphet fel, de amíg az előállított eredmény a
bemenő adatok bizonytalanságának megfelelő kimenő tér elmosódott
tartományában marad, addig az eredmény aligha kifogásolható. Másként
fogalmazva: bármely eredmény --- bárhogy is kaptuk --- amely egzakt
eredménye egy ugyanolyan valószínűséggel helyes bemenetnek, mint amilyet
ténylegesen használtunk, annyira jó válasz, amennyire egyáltalán jogunk
van várni.

Most formalizáljuk ezeket a fogalmakat, ismét az egyszerűség kedvéért
egydimenziós feladatokra összpontosítva. Tegyük fel, hogy egy
\(y = f(x)\) függvényérték kiszámítása a cél, ahol
\(f: \mathbb{R} \to \mathbb{R}\), de helyette egy közelítő \(\hat{y}\)
értéket kapunk. A számított és valódi értékek közötti eltérést,
\(\Delta y = \hat{y} - y\), \emph{előreható hibának} nevezzük. A
számított eredmény minőségének értékelésének egyik módja, hogy
megpróbáljuk becsülni ennek az előreható hibának a relatív nagyságát,
ami a konkrét körülményektől függően lehet könnyű vagy nehéz. Általában
azonban egy számításban a hibák előrefelé való propagálásának elemzése
--- a későbbiekben látható okokból --- gyakran nehéz. Sőt, a minden
szakaszban feltételezett legrosszabb esetek gyakran nagyon pesszimista
korláthoz vezetnek a teljes hibára nézve.

Egy alternatív megközelítés, ha a kapott közelítő megoldást úgy
tekintjük, mint egy módosított feladat egzakt megoldását, és
megkérdezzük: mekkora módosítás szükséges az eredeti feladathoz ahhoz,
hogy éppen a ténylegesen kapott eredményt adja? Más szóval: mekkora
adathibára volna szükség a kiindulási bemenetben, hogy megmagyarázza a
végső számított eredmény összes hibáját? Formálisabban a
\(\Delta x = \hat{x} - x\) mennyiséget --- ahol \(f(\hat{x}) = \hat{y}\)
--- \emph{hátraható hibának} nevezzük, amelynek relatív nagyságát a
hátraható hibaelemzésben próbáljuk megbecsülni. Ebből a nézőpontból egy
adott feladatra adott közelítő megoldás akkor jó, ha egy „közeli''
feladat egzakt megoldása (azaz a relatív hátraható hiba kicsi). Ha pedig
a közeli feladat a bemenő adatok bizonytalanságán belül van, akkor a
\(\hat{y}\) számított megoldás valójában lehet „a'' valódi megoldás,
amennyire tudhatjuk (vagy tudhatnánk, tekintve a bemenet
bizonytalanságát), tehát aligha kifogásolható.

Ezeket az összefüggéseket sematikusan (nem méretarányosan) az 1.2. ábra
szemlélteti, ahol \(x\) és \(f\) rendre az egzakt bemenet és az egzakt
függvény, \(\hat{f}\) a ténylegesen számított közelítő függvény,
\(\hat{x}\) pedig egy olyan bemeneti érték, amelyre az egzakt függvény
éppen ezt a számított eredményt adná. Megjegyezzük, hogy a
\(\hat{f}(x) = f(\hat{x})\) egyenlőség éppen \(\hat{x}\) választásából
adódik; valójában ez a követelmény definiálja \(\hat{x}\)-et. Az 1.2.6.
szakaszban számszerűsítjük az előreható és a hátraható hiba közötti
kapcsolatot.

\pandocbounded{\includegraphics[keepaspectratio]{_page_11_Figure_3.jpeg}}

1.2. ábra: Sematikus ábra az előreható és a hátraható hiba
szemléltetésére.

\textbf{1.5. Példa. Előreható és hátraható hiba.} A \(y = \sqrt{2}\)
közelítéseként a \(\hat{y} = 1{,}4\) érték abszolút előreható hibája

\[|\Delta y| = |\hat{y} - y| = |1{,}4 - 1{,}41421\ldots| \approx 0{,}0142,\]

azaz körülbelül 1 százalékos relatív előreható hiba. A hátraható hiba
meghatározásához megfigyeljük, hogy \(\sqrt{1{,}96} = 1{,}4\), tehát az
abszolút hátraható hiba

\[|\Delta x| = |\hat{x} - x| = |1{,}96 - 2| = 0{,}04,\]

azaz 2 százalékos relatív hátraható hiba.

\textbf{1.6. Példa. Hátraható hibaelemzés.} Tegyük fel, hogy egyszerű
közelítést keresünk a \(y = f(x) = \cos(x)\) koszinuszfüggvényre az
\(x = 1\) helyen. A koszinuszfüggvényt a

\[\cos(x) = 1 - \frac{x^2}{2!} + \frac{x^4}{4!} - \frac{x^6}{6!} + \cdots\]

végtelen sor adja meg, így fontolóra vehetjük a sor csonkolását, mondjuk
két tag után, hogy a

\[\hat{y} = \hat{f}(x) = 1 - x^2/2\]

közelítést kapjuk. Ekkor a közelítés előreható hibája

\[\Delta y = \hat{y} - y = \hat{f}(x) - f(x) = 1 - x^2/2 - \cos(x).\]

A hátraható hiba meghatározásához meg kell találnunk azt a \(\hat{x}\)
bemenő értéket \(f\)-re, amely a ténylegesen kapott \(\hat{y}\) kimenő
értéket adja, azaz amelyre \(\hat{f}(x) = f(\hat{x})\). A
koszinuszfüggvény esetén ezt az értéket

\[\hat{x} = \arccos\bigl(\hat{f}(x)\bigr) = \arccos(\hat{y})\]

adja meg. Így \(x = 1\) esetén:

\[y = f(1) = \cos(1) \approx 0{,}5403,\]

\[\hat{y} = \hat{f}(1) = 1 - 1^2/2 = 0{,}5,\]

\[\hat{x} = \arccos(\hat{y}) = \arccos(0{,}5) \approx 1{,}0472,\]

\[\text{előreható hiba} = \Delta y = \hat{y} - y \approx 0{,}5 - 0{,}5403 = -0{,}0403,\]

\[\text{hátraható hiba} = \Delta x = \hat{x} - x \approx 1{,}0472 - 1 = 0{,}0472.\]

Az előreható hiba arra utal, hogy a pontosság meglehetősen jó, mert a
kimenet közel van ahhoz, amit számítani szerettünk volna; a hátraható
hiba pedig arra utal, hogy a pontosság jó, mert a kapott kimenet helyes
egy csak kissé perturbált bemenetre. A következőkben látni fogjuk,
hogyan kapcsolódik az előreható és a hátraható hiba egymáshoz
számszerűen.

\subsubsection{1.2.6 Érzékenység és
kondicionáltság}\label{uxe9rzuxe9kenysuxe9g-uxe9s-kondicionuxe1ltsuxe1g}

Egy pontatlan megoldás nem feltétlenül egy rosszul átgondolt algoritmus
miatt pontatlan, hanem lehet eredendő tulajdonsága a megoldandó
feladatnak is. Még egzakt számítással is előfordulhat, hogy a feladat
megoldása nagyon érzékeny a bemenő adatok perturbációira. Az érzékenység
kvalitatív fogalma és annak kvantitatív mértéke --- a kondicionáltság
--- a propagált adathibával foglalkozik, azaz a bemenő adatok
perturbációinak a megoldásra gyakorolt hatásaival.

Egy feladatot érzéketlennek, vagy \emph{jól kondicionáltnak} mondunk, ha
a bemenő adatokban bekövetkező adott relatív változás a megoldásban
ésszerűen arányos relatív változást okoz. Egy feladatot érzékenynek,
vagy \emph{rosszul kondicionáltnak} mondunk, ha a megoldásban
bekövetkező relatív változás sokkal nagyobb lehet, mint a bemenő
adatokban bekövetkező változás. Aki valaha érezte, hogy a zuhany a
hőmérséklet-szabályzó legkisebb elállítására fagyosból forróvá vagy
fordítva vált, az első kézből tapasztalta az érzékeny rendszert: a hatás
aránytalan az okkal.

Pontosabban egy feladat \emph{kondíciószámát} a megoldásban bekövetkező
relatív változás és a bemenetben bekövetkező relatív változás
hányadosaként definiáljuk. Egy feladat rosszul kondicionált, vagyis
érzékeny, ha kondíciószáma sokkal nagyobb 1-nél. Egy függvény
kiértékelésének korábbi példájában használt jelöléssel:

\[\text{kondíciószám} = \frac{\bigl|(f(\hat{x}) - f(x))/f(x)\bigr|}{\bigl|(\hat{x} - x)/x\bigr|} = \frac{|(\hat{y} - y)/y|}{|(\hat{x} - x)/x|} = \frac{|\Delta y/y|}{|\Delta x/x|}.\]

Az arány számlálójában és nevezőjében rendre a relatív előreható és a
relatív hátraható hibát ismerhetjük fel, így ez az összefüggés így is
átfogalmazható:

\[|\text{relatív előreható hiba}| = \text{kondíciószám} \times |\text{relatív hátraható hiba}|.\]

Így a kondíciószám úgy értelmezhető, mint egy „erősítési tényező'',
amely az előreható és a hátraható hibát kapcsolja össze. Ha egy feladat
rosszul kondicionált (azaz kondíciószáma nagy), akkor a relatív
előreható hiba (a megoldás relatív perturbációja) akkor is nagy lehet,
ha a relatív hátraható hiba (a bemenet relatív perturbációja) kicsi.

Általában a kondíciószám a bemenettől függően változik, és a
gyakorlatban pontosan általában úgysem ismerjük. Ezért gyakran meg kell
elégednünk valamilyen bemeneti tartomány feletti maximális kondíciószám
durva becslésével vagy felső korlátjával, így a hátraható és az
előreható hiba közötti összefüggés közelítő egyenlőtlenséggé válik:

\[|\text{relatív előreható hiba}| \lesssim \text{kondíciószám} \times |\text{relatív hátraható hiba}|,\]

amely felülről korlátozza a legrosszabb esetben fellépő előreható hibát,
de nem feltétlenül valósul meg minden bemenet esetén. Ezen összefüggés
alapján a kondíciószám lehetővé teszi az --- általában leginkább érdekes
--- előreható hiba korlátozását a --- gyakran könnyebben becsülhető ---
hátraható hiba függvényében.

Analízissel közelíthetjük egy differenciálható
\(f: \mathbb{R} \to \mathbb{R}\) függvény kiértékelésének
kondíciószámát:

\[\text{abszolút előreható hiba} = f(x + \Delta x) - f(x) \approx f'(x)\Delta x,\]

így

\[\text{relatív előreható hiba} = \frac{f(x + \Delta x) - f(x)}{f(x)} \approx \frac{f'(x)\Delta x}{f(x)},\]

és ezért

\[\text{kondíciószám} \approx \left| \frac{f'(x)\Delta x/f(x)}{\Delta x/x} \right| = \left| \frac{x\,f'(x)}{f(x)} \right|.\]

Így a kimenő függvényérték relatív hibája lehet sokkal nagyobb vagy
sokkal kisebb is a bemenet relatív hibájánál, a szóban forgó függvény
tulajdonságaitól és a bemenet konkrét értékétől függően.

Egy adott feladat \emph{inverz feladata} annak meghatározása, hogy
melyik bemenet adna adott kimenetet. Egy \(y = f(x)\) függvény
kiértékelésének inverz feladata --- jelöléssel \(x = f^{-1}(y)\) ---
adott \(y\) értékre olyan \(x\) érték meghatározása, amelyre
\(f(x) = y\). A definícióból látható, hogy az inverz feladat
kondíciószáma az eredeti feladaté reciproka. Következésképpen, ha a
kondíciószám közel van 1-hez, akkor mind a feladat, mind az inverze jól
kondicionált. Ha azonban a kondíciószám 1-nél sokkal nagyobb vagy sokkal
kisebb, akkor rendre a feladat vagy az inverze rosszul kondicionált.

Felidézzük a kalkulusból, hogy ha \(g\) az inverz függvény
\(g(y) = f^{-1}(y)\), és \(x\) és \(y\) olyan értékek, amelyekre
\(y = f(x)\), akkor \(g'(y) = 1/f'(x)\), feltéve hogy \(f'(x) \neq 0\).
Így a \(g\) inverz függvény kondíciószáma

\[\text{kondíciószám} \approx \left| \frac{y\,g'(y)}{g(y)} \right| = \left| \frac{f(x)\,(1/f'(x))}{x} \right| = \left| \frac{f(x)}{x\,f'(x)} \right|,\]

ami az eredeti \(f\) függvény kondíciószámának reciproka.

\textbf{1.7. Példa. Kondíciószám.} Tekintsük az \(f(x) = \sqrt{x}\)
függvényt. Mivel \(f'(x) = 1/(2\sqrt{x})\),

\[\text{kondíciószám} \approx \left| \frac{xf'(x)}{f(x)} \right| = \left| \frac{x/(2\sqrt{x})}{\sqrt{x}} \right| = \frac{1}{2}.\]

Ez azt jelenti, hogy a bemenet adott relatív változása a kimenetben
körülbelül fele akkora relatív változást okoz. Vagyis a relatív
előreható hiba nagyságrendileg fele a relatív hátraható hibának, ahogy
ugyanerre a feladatra az 1.5. példában is láttuk. Így a négyzetgyökvonás
feladata igen jól kondicionált. Megjegyezzük, hogy az inverz feladat,
\(g(y) = y^2\), kondíciószáma \(|y\,g'(y)/g(y)| = |y(2y)/y^2| = 2\), ami
a négyzetgyökvonás kondíciószámának reciproka, ahogy vártuk.

\textbf{1.8. Példa. Érzékenység.} Tekintsük a tangensfüggvényt,
\(f(x) = \tan(x)\). Mivel \(f'(x) = \sec^2(x) = 1 + \tan^2(x)\),

\[\text{kondíciószám} \approx \left| \frac{xf'(x)}{f(x)} \right| = \left| \frac{x(1 + \tan^2(x))}{\tan(x)} \right| = \left| x \left( \frac{1}{\tan(x)} + \tan(x) \right) \right|.\]

Így \(\tan(x)\) nagyon érzékeny \(\pi/2\) bármely egész többszörösének
közelében, ahol értéke végtelenné válik. Például \(x = 1{,}57079\)-re a
kondíciószám körülbelül \(2{,}48275 \times 10^5\). Ennek hatásának
szemléltetésére kiértékeljük a függvényt két közeli pontban:

\[\tan(1{,}57079) \approx 1{,}58058 \times 10^5, \qquad \tan(1{,}57078) \approx 6{,}12490 \times 10^4,\]

és valóban látjuk, hogy a kimenet relatív változása --- körülbelül
\(1{,}58\) --- körülbelül negyedmilliószor nagyobb, mint a bemenet
relatív változása, amely körülbelül \(6{,}37 \times 10^{-6}\). Másfelől
a \(g(y) = \arctan(y)\) inverz függvény kondíciószáma
\(|y\,g'(y)/g(y)| = |y\,(1/(1+y^2))/\arctan(y)|\). A
\(y = 1{,}58058 \times 10^5\)-re ez a kondíciószám körülbelül
\(4{,}0278 \times 10^{-6}\), ami a tangens függvény megfelelő \(x\)-re
vett kondíciószámának reciproka. Így \(\arctan(y)\) rendkívül érzéketlen
ebben a pontban.

Az így definiált kondíciószámot néha \emph{relatív} kondíciószámnak
nevezik, mivel a bemenet és a kimenet relatív változásaira, illetve
relatív hibáira nézve van definiálva. Ez rendszerint a legmegfelelőbb,
de nincs értelmezve, ha a bemeneti \(x\) vagy a kimeneti \(y\) érték
nulla. Ilyen esetekben az \emph{abszolút} kondíciószám --- a megoldás
abszolút változásának és a bemenet abszolút változásának hányadosa,
\(|\Delta y|/|\Delta x|\) --- az érzékenység megfelelő mértéke. Ilyen
helyzet merül fel például a \emph{gyökkeresésben}: egy
\(f: \mathbb{R} \to \mathbb{R}\) függvényhez olyan \(x^*\) értéket
keresünk, amelyre \(f(x^*) = 0\) (lásd az 5. fejezetet). Az \(f(x)\)
függvény kiértékelése egy ilyen \(x^*\) gyök közelében abszolút
kondíciószáma

\[\frac{|\Delta y|}{|\Delta x|} \approx |f'(x^*)|,\]

a gyök meghatározásának inverz feladata --- azaz olyan \(x\) bemenő
érték keresése, amely \(y = f(x) = 0\)-t ad --- abszolút kondíciószáma
pedig \(1/|f'(x^*)|\), feltéve hogy \(f'(x^*) \neq 0\).

\subsubsection{1.2.7 Stabilitás és
pontosság}\label{stabilituxe1s-uxe9s-pontossuxe1g}

Egy számítási algoritmus \emph{stabilitásának} fogalma analóg egy
matematikai feladat kondicionáltságával abban, hogy mindkettő a
perturbációk hatásairól szól. A különbség az, hogy a stabilitás a
számítási hibának egy algoritmus által számított eredményre gyakorolt
hatására vonatkozik, míg a kondicionáltság az adathibának a feladat
megoldására gyakorolt hatására. Egy algoritmus \emph{stabil}, ha a
kiszámított eredmény viszonylag érzéketlen a számítás során alkalmazott
közelítések miatti perturbációkra. A hátraható hibaelemzés szemszögéből
egy algoritmus akkor stabil, ha az általa adott eredmény egy közeli
feladat egzakt megoldása, azaz a számítás során fellépő perturbációk
hatása nem rosszabb, mint egy kis mennyiségű adathiba hatása lenne az
adott feladat bemenetében. Ezen definíció szerint egy stabil algoritmus
pontosan a helyes eredményt adja egy közel helyes feladatra. Sok --- de
nem minden --- hasznos algoritmus stabil ebben az erős értelemben. Egyes
szövegkörnyezetekben hasznos egy gyengébb stabilitási fogalom: az, hogy
az algoritmus közel helyes eredményt ad egy közel helyes feladatra.

A \emph{pontosság} a számított megoldás és a vizsgált feladat valódi
megoldása közötti közelséget jelenti. Egy algoritmus stabilitása
önmagában nem garantálja, hogy a számított eredmény pontos: a pontosság
egyszerre függ a feladat kondicionáltságától és az algoritmus
stabilitásától. A stabilitás azt mondja meg, hogy a kapott megoldás
egzakt egy közeli feladatra, de ennek a közeli feladatnak a megoldása
nem feltétlenül van közel az eredeti feladat megoldásához, hacsak a
feladat jól kondicionált nem. Így pontatlanság eredhet akár egy stabil
algoritmusnak egy rosszul kondicionált feladatra való alkalmazásából,
akár egy instabil algoritmusnak egy jól kondicionált feladatra való
alkalmazásából. Ezzel szemben ha egy stabil algoritmust egy jól
kondicionált feladatra alkalmazunk, akkor pontos megoldáshoz jutunk.

\subsection{1.3 Számítógépes
aritmetika}\label{szuxe1muxedtuxf3guxe9pes-aritmetika}

Ahogy már említettük, a tudományos számításokban elkerülhetetlenül
jelentkező egyik közelítés a valós számok számítógépes ábrázolásában
rejlik. Ebben a szakaszban részletesebben megvizsgáljuk azokat a véges
precíziójú aritmetikai rendszereket, amelyeket a digitális
számítógépeken végzett tudományos számítások túlnyomó többségében
használunk.

\subsubsection{1.3.1 Lebegőpontos
számok}\label{lebegux151pontos-szuxe1mok}

Egy digitális számítógépen a matematika \(\mathbb{R}\) valós
számrendszerét egy \emph{lebegőpontos számrendszerrel} ábrázoljuk. Az
alapötlet a tudományos jelölésrendszerre hasonlít, amelyben egy nagyon
nagy vagy nagyon kicsi abszolút értékű számot egy mérsékelt méretű szám
és a tíz megfelelő hatványának szorzataként fejezünk ki. Például a 2347
és a 0,0007396 rendre \(2{,}347 \times 10^3\) és
\(7{,}396 \times 10^{-4}\) alakban írható. Ebben a formátumban a
tizedespont mozog, vagyis „lebeg", ahogy a \(10\) hatványa változik.
Formálisan egy \(F\) lebegőpontos számrendszert négy egész szám
jellemez:

\[\begin{aligned}
\beta &\quad \text{a számrendszer alapja (radix)}\\
p &\quad \text{pontosság}\\
[L, U] &\quad \text{exponens-tartomány}
\end{aligned}\]

Bármely \(x \in F\) lebegőpontos szám alakja

\[x = \pm \left(d_0 + \frac{d_1}{\beta} + \frac{d_2}{\beta^2} + \dots + \frac{d_{p-1}}{\beta^{p-1}}\right) \beta^E,\]

ahol \(d_i\) olyan egész szám, amelyre

\[0 \le d_i \le \beta - 1, \quad i = 0, \dots, p - 1,\]

\(E\) pedig olyan egész szám, amelyre

\[L \leq E \leq U.\]

A zárójelben lévő részt, amelyet a \(d_0 d_1 \cdots d_{p-1}\) \(p\)
darab \(\beta\)-alapú számjegyből álló sorozat ábrázol,
\emph{mantisszának} vagy \emph{szignifikándnak} nevezzük, \(E\)-t pedig
az \(x\) lebegőpontos szám \emph{exponensének} vagy
\emph{karakterisztikájának}. A mantissza \(d_1 d_2 \cdots d_{p-1}\)
része a \emph{törtrész}. Egy számítógépen az előjelet, az exponenst és a
mantisszát egy adott lebegőpontos szó külön mezőiben tárolják, amelyek
mindegyike rögzített szélességű. A nulla egyedileg ábrázolható úgy, hogy
mind a mantisszája, mind az exponense nulla, vagy úgy, hogy a törtrésze
nulla és az exponens egy speciális értéket vesz fel.

Manapság a legtöbb számítógép bináris (\(\beta = 2\)) aritmetikát
használ, de más alapok is előfordultak a múltban: például hexadecimális
(\(\beta = 16\)) az IBM mainframe-eken, illetve \(\beta = 3\) egy
balsorsú orosz számítógépen. Az oktális (\(\beta = 8\)) és a
hexadecimális jelölést gyakran használjuk kényelmes rövidítésként a
bináris számok hármas, illetve négyes bitcsoportokba rendezett
leírására. Nyilvánvaló okokból a decimális (\(\beta = 10\)) aritmetika
népszerű a zsebszámológépekben. Az emberi interakció megkönnyítése
érdekében egy számítógép általában a bevitt számértékeket decimális
jelölésről, a kimeneten pedig decimális jelölésre alakítja át,
függetlenül attól, milyen alapot használ belsőleg. Néhány tipikus
lebegőpontos rendszer paramétereit az 1.1. táblázat tartalmazza, amely
jól illusztrálja a mezőszélességekből adódó kompromisszumot a pontosság
és az exponens-tartomány között. Például ugyanolyan 64 bites
szóhosszúság mellett a Cray rendszer szélesebb exponens-tartománnyal
rendelkezik, mint az IEEE dupla pontosság, de ezért kisebb pontossággal
fizet.

Az IEEE szabvány szerinti egyszeres pontosságú (SP) és dupla pontosságú
(DP) bináris lebegőpontos rendszerek ma messze a legfontosabbak. Szinte
egyetemesen elterjedtek a személyi számítógépekben és munkaállomásokon,
valamint számos mainframe- és szuperszámítógépen is. Az IEEE szabványt
gondosan úgy tervezték, hogy kiküszöbölje a korábbi, gyártó-specifikus
lebegőpontos megvalósítások számos anomáliáját és kétértelműségét, és
ezzel nagyban elősegítette a hordozható és megbízható numerikus
szoftverek fejlesztését. Emellett lehetővé teszi a kivételes helyzetek
--- például a nullával való osztás --- értelmes és következetes
kezelését.

\begin{longtable}[]{@{}lllll@{}}
\toprule\noalign{}
Rendszer & \(\beta\) & \(p\) & \(L\) & \(U\) \\
\midrule\noalign{}
\endhead
\bottomrule\noalign{}
\endlastfoot
IEEE SP & 2 & 24 & \(-126\) & \(127\) \\
IEEE DP & 2 & 53 & \(-1022\) & \(1023\) \\
Cray & 2 & 48 & \(-16\,383\) & \(16\,384\) \\
HP számológép & 10 & 12 & \(-499\) & \(499\) \\
IBM mainframe & 16 & 6 & \(-64\) & \(63\) \\
\end{longtable}

1.1. táblázat: Tipikus lebegőpontos rendszerek paraméterei.

\subsubsection{1.3.2 Normalizáltság}\label{normalizuxe1ltsuxe1g}

Egy lebegőpontos rendszert \emph{normalizáltnak} mondunk, ha a vezető
\(d_0\) számjegy mindig nullától különböző, kivéve ha az ábrázolt szám
nulla. Így egy normalizált lebegőpontos rendszerben egy adott nem nulla
lebegőpontos szám \(m\) mantisszájára mindig fennáll, hogy

\[1 \leq m < \beta.\]

(Egy másik konvenció szerint \(d_0\) mindig nulla, ebben az esetben egy
lebegőpontos szám akkor normalizált, ha \(d_1 \neq 0\), és ekkor
\(\beta^{-1} \leq m < 1\).) A lebegőpontos rendszereket általában
normalizálják, mert

\begin{itemize}
\tightlist
\item
  Így minden szám ábrázolása egyértelmű.
\item
  Nem pazarolunk el számjegyeket vezető nullákra, és ezzel
  maximalizáljuk a pontosságot.
\item
  Egy bináris (\(\beta = 2\)) rendszerben a vezető bit mindig \(1\), így
  nem kell tárolni, és ezáltal egy adott mezőszélességen belül egy
  bittel több pontosságot nyerünk.
\end{itemize}

\subsubsection{1.3.3 A lebegőpontos rendszerek
tulajdonságai}\label{a-lebegux151pontos-rendszerek-tulajdonsuxe1gai}

Egy lebegőpontos számrendszer véges és diszkrét. Egy adott lebegőpontos
rendszerben a normalizált lebegőpontos számok száma

\[2(\beta - 1)\,\beta^{p-1}\,(U - L + 1) + 1,\]

mivel az előjelre két választás van, a mantissza vezető számjegyére
\(\beta - 1\), a mantissza további \(p - 1\) számjegyének mindegyikére
\(\beta\), az exponensre pedig \(U - L + 1\) lehetséges érték. A
\(+1\)-et azért adjuk hozzá, mert a szám lehet nulla is.

Létezik egy legkisebb pozitív normalizált lebegőpontos szám,

\[\text{UFL (alulcsordulási szint)} = \beta^L,\]

amelyben a vezető számjegy \(1\), a mantissza többi számjegye \(0\), az
exponens pedig a legkisebb lehetséges érték. Létezik egy legnagyobb
lebegőpontos szám,

\[\text{OFL (túlcsordulási szint)} = \beta^{U+1}\bigl(1 - \beta^{-p}\bigr),\]

amelyben a mantissza minden számjegye \(\beta - 1\), az exponens pedig a
legnagyobb lehetséges érték. Az adott lebegőpontos rendszerben nem
ábrázolható OFL-nél nagyobb szám, sem UFL-nél kisebb pozitív szám. A
lebegőpontos számok nem egyenletesen oszlanak el az értékkészletükben,
hanem csak egymást követő \(\beta\)-hatványok között egyenletes a közük.

\textbf{1.9. Példa. Lebegőpontos rendszer.} Egy példa lebegőpontos
rendszert szemléltet az 1.3. ábra, ahol a skálajelek a \(\beta = 2\),
\(p = 3\), \(L = -1\), \(U = 1\) paraméterű rendszer mind a 25
lebegőpontos számát jelzik. Ebben a rendszerben a legnagyobb szám
\(\text{OFL} = (1{,}11)_2 \times 2^1 = (3{,}5)_{10}\), a legkisebb
pozitív normalizált szám pedig
\(\text{UFL} = (1{,}00)_2 \times 2^{-1} = (0{,}5)_{10}\). Ez egy nagyon
apró, szemléltető célú „játék" rendszer, mégis jellemző a lebegőpontos
rendszerekre általában: kellően nagy nagyítás mellett minden normalizált
lebegőpontos rendszer lényegében így néz ki --- szemcsés és egyenetlen
közű.

\pandocbounded{\includegraphics[keepaspectratio]{_page_18_Figure_4.jpeg}}

1.3. ábra: Példa lebegőpontos számrendszer.

\subsubsection{1.3.4 Kerekítés}\label{kerekuxedtuxe9s}

Azokat a valós számokat, amelyek egy adott lebegőpontos rendszerben
pontosan ábrázolhatók, \emph{gépi számoknak} nevezzük. Ha egy adott
\(x\) valós szám nem ábrázolható pontosan lebegőpontos számként, akkor
valamilyen „közeli" lebegőpontos számmal kell közelítenünk. Egy \(x\)
valós szám lebegőpontos közelítését \(\mathrm{fl}(x)\)-szel jelöljük.
Annak a folyamatát, amely során egy adott \(x\) valós számhoz egy közeli
\(\mathrm{fl}(x)\) lebegőpontos számot választunk, \emph{kerekítésnek}
nevezzük, az ilyen közelítés által bevezetett hibát pedig
\emph{kerekítési hibának}. Két gyakran használt kerekítési szabály:

\begin{itemize}
\tightlist
\item
  Csonkítás: az \(x\) szám \(\beta\)-alapú kifejtését a \((p-1)\)-edik
  számjegy után elvágjuk. Mivel \(\mathrm{fl}(x)\) a következő
  lebegőpontos szám \(x\)-től a nulla felé haladva, ezt a szabályt néha
  \emph{nulla felé kerekítésnek} is nevezik.
\item
  Kerekítés a legközelebbihez: \(\mathrm{fl}(x)\) az \(x\)-hez
  legközelebbi lebegőpontos szám; döntetlen esetén azt a lebegőpontos
  számot választjuk, amelynek az utolsó tárolt számjegye páros. Ez
  utóbbi tulajdonság miatt ezt a szabályt \emph{páros felé kerekítésnek}
  is nevezik.
\end{itemize}

A legközelebbihez kerekítés a legpontosabb és torzítatlan, de valamivel
költségesebb helyesen megvalósítani. Egyes régebbi rendszerek olcsóbban
megvalósítható szabályokat használtak, például csonkítást, de az IEEE
szabványos rendszerekben a legközelebbihez kerekítés az alapértelmezett
szabály.

\textbf{1.10. Példa. Kerekítési szabályok.} Az alábbi decimális számok
két számjegyre kerekítése mindkét szabály szerint a következő eredményt
adja:

\begin{longtable}[]{@{}llllll@{}}
\toprule\noalign{}
Szám & Csonkítás & Legközelebbi & Szám & Csonkítás & Legközelebbi \\
\midrule\noalign{}
\endhead
\bottomrule\noalign{}
\endlastfoot
1,649 & 1,6 & 1,6 & 1,749 & 1,7 & 1,7 \\
1,650 & 1,6 & 1,6 & 1,750 & 1,7 & 1,8 \\
1,651 & 1,6 & 1,7 & 1,751 & 1,7 & 1,8 \\
1,699 & 1,6 & 1,7 & 1,799 & 1,7 & 1,8 \\
\end{longtable}

Egy további, gyakran figyelmen kívül hagyott hibaforrás a lebegőpontos
számok bevitelekor és kiírásakor szokásosan elvégzett decimális-bináris
és bináris-decimális átváltásokban rejlik. Ezekre az átváltásokra az
IEEE szabvány nem terjed ki; a szabvány csak a belső aritmetikai
műveleteket rögzíti. Helyesen kerekített bevitel és kiírás elérhető
ésszerű költséggel, de nem minden számítógépes rendszerben valósítják
meg. A helyesen kerekített bináris-decimális és decimális-bináris
átváltásokra hatékony, hordozható rutinok --- a \texttt{dtoa} és a
\texttt{strtod} --- érhetők el a Netlibből (lásd az 1.4.1. szakaszt).

\subsubsection{1.3.5 Gépi pontosság}\label{guxe9pi-pontossuxe1g}

Egy lebegőpontos rendszer pontosságát olyan mennyiség jellemzi, amelyet
\emph{egységnyi kerekítési hibának} (unit roundoff), \emph{gépi
pontosságnak} vagy \emph{gépi epszilonnak} is neveznek. Ennek értéke ---
amelyet \(\epsilon_{\text{mach}}\)-kal jelölünk --- az alkalmazott
kerekítési szabálytól függ. Csonkítással végzett kerekítés esetén

\[\epsilon_{\text{mach}} = \beta^{1-p},\]

a legközelebbihez való kerekítés esetén pedig

\[\epsilon_{\text{mach}} = \tfrac{1}{2}\,\beta^{1-p}.\]

Az egységnyi kerekítési hiba azért fontos, mert korlátozza azt a
\emph{relatív hibát}, amellyel egy lebegőpontos rendszer normalizált
tartományában bármely nem nulla \(x\) valós szám ábrázolható:

\[\left| \frac{\mathrm{fl}(x) - x}{x} \right| \le \epsilon_{\text{mach}}.\]

Az egységnyi kerekítési hiba egy másik jellemzése, amellyel néha
találkozhatunk, hogy \(\epsilon_{\text{mach}}\) a legkisebb olyan
\(\epsilon\) szám, amelyre

\[\mathrm{fl}(1 + \epsilon) > 1,\]

ám ez nem teljesen ekvivalens az előbbi definícióval, ha a páros felé
kerekítés szabályát alkalmazzuk. Egy további, néha használt definíció
szerint \(\epsilon_{\text{mach}}\) az \(1\) és az utána következő
lebegőpontos szám közötti távolság, ám ez eltérhet az előző két
definíciótól. Részleteikben különbözhetnek, de
\(\epsilon_{\text{mach}}\) mindhárom definíciójának ugyanaz az alapvető
szándéka: a lebegőpontos rendszer relatív szemcsézettségének mérése.

Az 1.9. példa szemléltető játékrendszerében csonkítás esetén
\(\epsilon_{\text{mach}} = 0{,}25\), legközelebbihez kerekítés esetén
pedig \(\epsilon_{\text{mach}} = 0{,}125\). Az IEEE bináris lebegőpontos
rendszerekre \(\epsilon_{\text{mach}} = 2^{-24} \approx 10^{-7}\)
egyszeres pontosság mellett, és
\(\epsilon_{\text{mach}} = 2^{-53} \approx 10^{-16}\) dupla pontosság
mellett. Ezért azt mondjuk, hogy az IEEE egyszeres, illetve dupla
pontosságú lebegőpontos rendszerek körülbelül 7, illetve 16 decimális
értékes jegy pontosságúak.

Bár mindkettő „kicsi", az egységnyi kerekítési hibát nem szabad
összetéveszteni az alulcsordulási szinttel. Az
\(\epsilon_{\text{mach}}\) egységnyi kerekítési hibát a lebegőpontos
rendszer mantisszamezejének számjegyszáma határozza meg, míg az UFL
alulcsordulási szintet az exponensmező számjegyszáma. Minden gyakorlati
lebegőpontos rendszerben fennáll, hogy

\[0 < \text{UFL} < \epsilon_{\text{mach}} < \text{OFL}.\]

\subsubsection{1.3.6 Szubnormális számok és fokozatos
alulcsordulás}\label{szubnormuxe1lis-szuxe1mok-uxe9s-fokozatos-alulcsorduluxe1s}

Az 1.3. ábrán bemutatott játék lebegőpontos rendszerben érezhető rés van
a nulla körül. Ez a rés --- amely valamilyen mértékben minden
lebegőpontos rendszerben jelen van --- a normalizáltság következménye: a
legkisebb lehetséges mantissza \(1{,}00\dots\), a legkisebb lehetséges
exponens pedig \(L\), így nincs lebegőpontos szám a nulla és \(\beta^L\)
között. Ha feladjuk a normalizáltsági követelményt, és megengedjük, hogy
a vezető számjegy nulla legyen (de csak amikor az exponens a minimumon
áll), akkor a nulla körüli rést további lebegőpontos számokkal
„tölthetjük ki". A szemléltető játékrendszerünkben ez az enyhítés hat
további lebegőpontos számot eredményez, amelyek közül a legkisebb
pozitív \((0{,}01)_2 \times 2^{-1} = (0{,}125)_{10}\), ahogy azt az 1.4.
ábra mutatja.

\pandocbounded{\includegraphics[keepaspectratio]{_page_20_Figure_6.jpeg}}

1.4. ábra: Példa lebegőpontos rendszer szubnormális számokkal.

Az így a rendszerhez hozzáadott további számokat \emph{szubnormális}
vagy \emph{denormalizált} lebegőpontos számoknak nevezzük. Noha
hasznosan kiterjesztik az ábrázolható abszolút értékek tartományát, a
szubnormális számoknak eredendően kisebb a pontosságuk, mint a
normalizáltaknak, mert a törtrészükben kevesebb értékes számjegy
szerepel. Különösen fontos, hogy a tartomány ilyen módon történő
kiterjesztése nem csökkenti az \(\epsilon_{\text{mach}}\) egységnyi
kerekítési hibát.

Az ilyen módon kiterjesztett lebegőpontos rendszerről azt mondjuk, hogy
\emph{fokozatos alulcsordulást} valósít meg, mivel az ábrázolható
abszolút értékek alsó tartományát kiterjeszti, ahelyett hogy a minimális
exponens érték alullépésekor a nullára csordulna. Az IEEE szabvány
tartalmazza a szubnormális számokat és a fokozatos alulcsordulást. A
fokozatos alulcsordulást az exponensmező egy speciális, fenntartott
értékén keresztül valósítják meg, mivel a vezető bináris számjegyet nem
tárolják, így azt nem lehet a denormalizáltság jelzésére használni.

\subsubsection{1.3.7 Kivételes
értékek}\label{kivuxe9teles-uxe9rtuxe9kek}

Az IEEE lebegőpontos szabvány két további speciális értéket biztosít a
kivételes helyzetek jelzésére:

\begin{itemize}
\tightlist
\item
  \texttt{Inf}, amely a végtelenséget jelöli, egy véges szám nullával
  való osztásakor áll elő, például \(1/0\).
\item
  \texttt{NaN}, amely a \emph{nem szám} (not a number) rövidítése,
  definiálatlan vagy határozatlan művelet eredménye, például \(0/0\),
  \(0 \cdot \text{Inf}\) vagy \(\text{Inf}/\text{Inf}\).
\end{itemize}

Az \texttt{Inf} és a \texttt{NaN} az IEEE aritmetikában az exponensmező
speciális, fenntartott értékein keresztül van megvalósítva.

Az, hogy az \texttt{Inf} és a \texttt{NaN} támogatott-e a felhasználói
szinten egy adott számítási környezetben, a nyelvtől, a fordítótól és a
futásidejű rendszertől függ. Ha elérhetők, akkor ezek a mennyiségek
segíthetnek olyan szoftvert tervezni, amely a kivételes helyzeteket
kecsesen kezeli, és nem szakítja meg hirtelen a programot. A MATLAB-ban
(lásd az 1.4.2. szakaszt) például, ha \texttt{Inf} vagy \texttt{NaN} lép
fel, azokat értelmesen továbbterjeszti a számítás során (pl.
\(1 + \text{Inf} = \text{Inf}\)). Ennek ellenére érdemes, ha lehetséges,
teljesen elkerülni az ilyen kivételes helyzeteket. Amellett, hogy
felhívják a felhasználó figyelmét az aritmetikai kivételekre, ezek a
speciális értékek jelzőként is hasznosak lehetnek, mivel egyetlen
törvényes numerikus értékkel sem téveszthetők össze. Például a
\texttt{NaN} használható egy tömb még definiálatlan részének
megjelölésére.

\subsubsection{1.3.8 Lebegőpontos
aritmetika}\label{lebegux151pontos-aritmetika}

Két lebegőpontos szám összeadásánál vagy kivonásánál a mantisszák
összeadása vagy kivonása előtt az exponenseknek meg kell egyezniük. Ha
kezdetben nem egyeznek meg, akkor az egyik szám mantisszáját addig kell
eltolni, amíg az exponensek megegyeznek. Egy ilyen eltolás során a
(abszolút értékben) kisebb szám néhány záró számjegye a mantisszamező
rögzített szélességén túlra kerül, így az aritmetikai művelet helyes
eredménye nem ábrázolható pontosan a lebegőpontos rendszerben. Sőt, ha
az abszolút értékek közötti különbség túl nagy, akkor a kisebb szám
teljes mantisszája a mezőszélességen túlra kerülhet, és így az eredmény
egyszerűen a két operandus közül a nagyobbik lesz. Másképpen fogalmazva:
ha két \(p\) számjegyű szám valódi összege \(p\)-nél több számjegyet
tartalmaz, akkor a \(p\) számjegyre kerekítéskor a felesleges számjegyek
elvesznek, és a legrosszabb esetben a kisebb abszolút értékű operandus
teljesen eltűnhet.

Két lebegőpontos szám szorzásához nem szükséges, hogy exponenseik
megegyezzenek --- az exponenseket egyszerűen összeadjuk, a mantisszákat
pedig összeszorozzuk. Két \(p\) számjegyű mantissza szorzata azonban
általában akár \(2p\) számjegyet is tartalmazhat, így a helyes eredmény
ismét csak nem ábrázolható pontosan a lebegőpontos rendszerben, és
kerekíteni kell.

\textbf{1.11. Példa. Lebegőpontos aritmetika.} Tekintsünk egy
\(\beta = 10\) és \(p = 6\) paraméterű lebegőpontos rendszert. Ha
\(x = 1{,}92403 \times 10^2\) és \(y = 6{,}35782 \times 10^{-1}\), akkor
a lebegőpontos összeadás --- legközelebbihez kerekítés mellett --- az
\(x + y = 1{,}93039 \times 10^2\) eredményt adja. Vegyük észre, hogy
\(y\) utolsó két számjegye nem befolyásolja az eredményt. Még kisebb
exponens esetén \(y\) egyáltalán nem hatott volna az eredményre.
Hasonlóképpen, a lebegőpontos szorzás az
\(x \cdot y = 1{,}22326 \times 10^2\) eredményt adja, ami a valódi
szorzat számjegyeinek felét eldobja.

Két lebegőpontos szám osztása is adhat olyan eredményt, amely nem
ábrázolható pontosan. Például az \(1\) és a \(10\) pontosan ábrázolható
bináris lebegőpontos számként, de hányadosuk, az \(1/10\) nem véges
bináris kifejtésű, így nem bináris lebegőpontos szám.

Minden fent említett esetben egy lebegőpontos aritmetikai művelet
eredménye eltérhet attól, amit a megfelelő valós aritmetikai művelet
adna ugyanazokra az operandusokra, mert a helyes valós eredmény
ábrázolásához nincs elegendő pontosság. A valós eredmény azért is lehet
nem ábrázolható, mert az exponense a lebegőpontos rendszer rendelkezésre
álló tartományán kívül esik (túlcsordulás vagy alulcsordulás). A
túlcsordulás általában súlyosabb probléma, mint az alulcsordulás, abban
az értelemben, hogy egy lebegőpontos rendszerben nincs jó közelítés a
tetszőlegesen nagy számokra, míg a nulla gyakran ésszerű közelítés a
tetszőlegesen kis számokra. Ezért sok számítógépes rendszerben a
túlcsordulás bekövetkezése végzetes hibával megszakítja a programot, az
alulcsordulást viszont csendben nullára állíthatja a rendszer a
végrehajtás megzavarása nélkül.

\textbf{1.12. Példa. Sor összegzése.} E kérdések szemléltetéseként a

\[\sum_{n=1}^{\infty} \frac{1}{n}\]

végtelen sornak véges az összege lebegőpontos aritmetikában, noha a
valós sor divergens. Első ránézésre azt gondolhatnánk, hogy ez azért
van, mert \(1/n\) előbb-utóbb alulcsordul, vagy a részletösszeg
előbb-utóbb túlcsordul --- ami valóban meg is történik. De mielőtt ezek
bármelyike bekövetkezne, a részletösszeg megszűnik változni, amint
\(1/n\) a részletösszeghez képest elhanyagolhatóvá válik, azaz amikor
\(1/n < \epsilon_{\text{mach}} \sum_{k=1}^{n-1} (1/k)\), így az összeg
véges (lásd az 1.8. számítógépes feladatot).

Amint megjegyeztük, két lebegőpontos számon elvégzett valós aritmetikai
művelet eredménye nem feltétlenül lebegőpontos szám. Ha a számítógépbe
olyan szám kerül, amely lebegőpontos számként nem ábrázolható pontosan,
vagy egy későbbi aritmetikai művelet ilyet állít elő, akkor azt a
korábban megadott kerekítési szabályok valamelyikével kerekíteni kell,
hogy lebegőpontos számot kapjunk. Mivel a lebegőpontos számok nem
egyenletes közűek, az ilyen közelítésben fellépő abszolút hiba nem
egyenletes, a relatív hibát azonban az \(\epsilon_{\text{mach}}\)
egységnyi kerekítési hiba korlátozza.

Ideális esetben
\(x \mathbin{\mathrm{flop}} y = \mathrm{fl}(x \mathbin{\mathrm{op}} y)\)
(azaz a lebegőpontos aritmetikai műveletek helyesen kerekített eredményt
adnak), és sok számítógép --- például az IEEE lebegőpontos szabványnak
megfelelőek --- el is éri ezt az ideált, amíg
\(x \mathbin{\mathrm{op}} y\) a lebegőpontos rendszer tartományán belül
van. Ennek ellenére a valós aritmetika néhány ismert törvénye nem
feltétlenül érvényes a lebegőpontos rendszerben. Különösen, a
lebegőpontos összeadás és szorzás kommutatív, de nem asszociatív.
Például ha \(\epsilon\) az \(\epsilon_{\text{mach}}\) egységnyi
kerekítési hibánál egy kicsit kisebb pozitív lebegőpontos szám, akkor
\((1 + \epsilon) + \epsilon = 1\), de \(1 + (\epsilon + \epsilon) > 1\).
Az, hogy a lebegőpontos aritmetika nem teljesíti a valós aritmetika
szokásos törvényeit, egyik oka annak, hogy az előreható hibaelemzés
nehéz lehet. A hátraható hibaelemzés egyik előnye, hogy lehetővé teszi a
valós aritmetika használatát az elemzésben.

A kerekítési hibaelemzés rendszerint a következő szabványos modellre
épül a lebegőpontos aritmetikára vonatkozóan:

\[\mathrm{fl}(x \mathbin{\mathrm{op}} y) = (x \mathbin{\mathrm{op}} y)(1 + \delta),\]

ahol \(|\delta| \le \epsilon_{\text{mach}}\), és \(\mathrm{op}\) a
\(+\), \(-\), \(\times\), \(/\) szabványos aritmetikai műveletek
valamelyike. Átrendezve látjuk, hogy ez a modell értelmezhető a relatív
előreható hibára vonatkozó állításként:

\[\frac{|\mathrm{fl}(x \mathbin{\mathrm{op}} y) - (x \mathbin{\mathrm{op}} y)|}{|(x \mathbin{\mathrm{op}} y)|} = |\delta| \le \epsilon_{\text{mach}}.\]

Értelmezhető a hátraható hibára vonatkozó állításként is. Például
\(\mathrm{op} = +\) esetén

\[\mathrm{fl}(x + y) = (x + y)(1 + \delta) = x(1 + \delta) + y(1 + \delta),\]

ami azt jelenti, hogy a lebegőpontos összeadás eredménye egzakt olyan
\(x\) és \(y\) operandusokra, amelyek mindegyike \(|\delta|\) relatív
mértékben perturbált.

\textbf{1.13. Példa. Kerekítési hibaelemzés.} Tekintsük az \(x(y + z)\)
egyszerű számítást. Lebegőpontos aritmetikában

\[\mathrm{fl}(y + z) = (y + z)(1 + \delta_1), \quad |\delta_1| \le \epsilon_{\text{mach}},\]

és így

\[\begin{aligned}
\mathrm{fl}\bigl(x(y + z)\bigr) &= \bigl(x\bigl((y + z)(1 + \delta_1)\bigr)\bigr)(1 + \delta_2), \quad |\delta_2| \le \epsilon_{\text{mach}} \\
&= x(y + z)(1 + \delta_1 + \delta_2 + \delta_1\delta_2) \\
&\approx x(y + z)(1 + \delta_1 + \delta_2) \\
&= x(y + z)(1 + \delta), \quad |\delta| = |\delta_1 + \delta_2| \le 2\epsilon_{\text{mach}}.
\end{aligned}\]

Amint korábban, ez a korlát értelmezhető az előreható hiba (a számított
és a kívánt eredmény közötti eltérés) vagy a hátraható hiba
(operandusperturbációk) szempontjából is, és meglehetősen pesszimista
lehet. Például \(\delta_1\) és \(\delta_2\) ellentétes előjelű lehet, és
így kiolthatják egymást, aminek eredményeként a teljes hiba sokkal
kisebb is lehet, mint amit a legrosszabb esetű elemzés sugall.
Bonyolultabb számítások hasonló elemzése általában
\(\epsilon_{\text{mach}}\) ennek megfelelően nagyobb többszöröseit
tartalmazó hibakorlátokhoz vezet. Szerencsére \(\epsilon_{\text{mach}}\)
olyan kicsi, hogy a gyakorlatban rendkívül ritkán romlik az eredmény
pusztán a számítás mennyisége miatt súlyosan.

\subsubsection{1.3.9 Számjegykioltás}\label{szuxe1mjegykioltuxe1s}

A kerekítés nem az egyetlen szükséges rossz a véges precíziójú
aritmetikában. Két azonos előjelű és hasonló abszolút értékű \(p\)
számjegyű szám kivonása \(p\)-nél kevesebb értékes számjegyű eredményt
ad, ezért mindig pontosan ábrázolható (feltéve, hogy a két szám abszolút
értéke nem több mint kétszeresen tér el egymástól). Ennek oka, hogy a
két szám vezető számjegyei kioltják egymást (azaz különbségük nulla).
Például ismét \(\beta = 10\)-t és \(p = 6\)-ot véve, ha
\(x = 1{,}92403 \times 10^2\) és \(z = 1{,}92275 \times 10^2\), akkor
\(x - z = 1{,}28000 \times 10^{-1}\) adódik, amely mindössze három
értékes jegyével pontosan ábrázolható.

Bár az eredmény pontos, az ilyen kioltás gyakran mégis jelentős
információvesztést jelenthet. A probléma az, hogy az operandusok
kerekítés vagy más korábbi hibák miatt gyakran bizonytalanok, ezért a
különbségben megjelenő relatív bizonytalanság nagy lehet. Más szóval: ha
két közel azonos szám csak a kerekítési hiba pontosságáig helyes, akkor
a különbségük már csak a kerekítési hibát tartalmazza.

Egyszerű példaként: ha \(\epsilon\) az \(\epsilon_{\text{mach}}\)
egységnyi kerekítési hibánál kicsit kisebb pozitív szám, akkor
\((1 + \epsilon) - (1 - \epsilon) = 1 - 1 = 0\) lebegőpontos
aritmetikában, ami az utolsó kivonás tényleges operandusaira nézve
helyes, de a teljes számítás valódi eredménye, \(2\epsilon\), teljesen
elveszett. Maga a kivonás nem hibás: csupán jelzi a már korábban
bekövetkezett információvesztést.

Természetesen az információvesztés nem mindig teljes, de tény marad,
hogy a kioltás során elveszett számjegyek a legfontosabb, vezető
számjegyek, míg a kerekítés során elveszett számjegyek a legkevésbé
fontos, záró számjegyek. Emiatt a hatás miatt egy kis mennyiség nagyok
különbségeként történő számítása általában rossz ötlet, mert a
kerekítési hiba valószínűleg dominálni fogja az eredményt. Például egy
váltakozó előjelű sor összegzése, mint

\[e^x = 1 + x + \frac{x^2}{2!} + \frac{x^3}{3!} + \cdots\]

\(x < 0\) esetén, katasztrofális kioltás miatt végzetesen rossz
eredményt adhat (lásd az 1.9. számítógépes feladatot).

\textbf{1.14. Példa. Kioltás.} A kioltás nemcsak a számítógépes
aritmetikában okoz problémát, hanem bármely olyan helyzetben, ahol csak
korlátozott pontosság érhető el --- például empirikus mérések vagy
laboratóriumi kísérletek során. Például Manhattan és Staten Island
közötti távolság meghatározása a két hely Los Angelestől mért
távolságainak segítségével nagyon rossz eredményt adna, hacsak ez utóbbi
távolságokat nem ismerjük rendkívül nagy pontossággal.

Másik példaként a fizikusok évek óta próbálják a hélium atom teljes
energiáját első elvekből kiszámítani Monte-Carlo-módszerekkel. E
számítások pontosságát nagyrészt a felhasznált véletlen próbák száma
határozza meg. Ahogy gyorsabb számítógépek válnak elérhetővé és a
számítási technikák finomodnak, az elérhető pontosság is javul. A teljes
energia a mozgási és a potenciális energia összege, amelyeket
külön-külön számítanak ki, és ellentétes előjelűek. Így a teljes energia
különbségként áll elő, és kioltás miatt szenved. Az 1.2. táblázat évek
során kapott értékek sorozatát mutatja be (ezeket az adatokat Dr. Robert
Panoff bocsátotta rendelkezésre). Ebben az időszakban a mozgási és a
potenciális energiára számított értékek csak legfeljebb 6 százalékkal
változtak, a teljes energiára adott becslés azonban 144 százalékkal. A
korábbi számítások egy-két értékes jegye a későbbi kivonásban teljesen
elveszett.

\begin{longtable}[]{@{}llll@{}}
\toprule\noalign{}
Év & Mozgási & Potenciális & Teljes \\
\midrule\noalign{}
\endhead
\bottomrule\noalign{}
\endlastfoot
1971 & 13,0 & \(-14{,}0\) & \(-1{,}0\) \\
1977 & 12,76 & \(-14{,}02\) & \(-1{,}26\) \\
1980 & 12,22 & \(-14{,}35\) & \(-2{,}13\) \\
1985 & 12,28 & \(-14{,}65\) & \(-2{,}37\) \\
1988 & 12,40 & \(-14{,}84\) & \(-2{,}44\) \\
\end{longtable}

1.2. táblázat: A hélium atom teljes energiájára számított értékek.

\textbf{1.15. Példa. Másodfokú megoldóképlet.} A kioltás és egyéb
numerikus buktatók nem feltétlenül igényelnek hosszú számítási
sorozatot. Például a másodfokú egyenlet gyökeire szolgáló szabványos
képlet használata tele van numerikus csapdákkal. Ahogy minden iskolás
gyermek megtanulja, az

\[ax^2 + bx + c = 0\]

másodfokú egyenlet két megoldását a megoldóképlet adja:

\[x = \frac{-b \pm \sqrt{b^2 - 4ac}}{2a}.\]

Az együtthatók bizonyos értékei esetén e képlet naiv alkalmazása
lebegőpontos aritmetikában túlcsordulást, alulcsordulást vagy
katasztrofális kioltást okozhat.

Ha például az együtthatók nagyon nagyok vagy nagyon kicsik, akkor
\(b^2\) vagy \(4ac\) túlcsordulhat vagy alulcsordulhat. A túlcsordulás
lehetősége az együtthatók átskálázásával elkerülhető, például úgy, hogy
mindhárom együtthatót elosztjuk a legnagyobb abszolút értékű
együtthatóval. Az ilyen átskálázás nem változtatja meg a másodfokú
egyenlet gyökeit, ám most a legnagyobb együttható már \(1\), és \(b^2\)
vagy \(4ac\) számításakor nem léphet fel túlcsordulás. Az ilyen
átskálázás nem szünteti meg az alulcsordulás lehetőségét, de megelőzi a
felesleges alulcsordulást, amely egyébként bekövetkezhetne, ha mindhárom
együttható nagyon kicsi.

A \(-b\) és a négyzetgyök közötti kioltás elkerülhető az egyik gyök
kiszámításával az

\[x = \frac{2c}{-b \mp \sqrt{b^2 - 4ac}}\]

alternatív képlet szerint, amelyben az előjelek a szabványos képletével
ellentétesek. A négyzetgyökön belüli kioltás azonban nehezen kerülhető
el nagyobb pontosság használata nélkül (ha a diszkrimináns az
együtthatókhoz képest kicsi, akkor a két gyök közel van egymáshoz, és a
feladat eredendően rosszul kondicionált).

Szemléltetésképpen négy számjegyű decimális aritmetikát ---
legközelebbihez kerekítéssel --- használunk az \(a = 0{,}05010\),
\(b = -98{,}78\), \(c = 5{,}015\) együtthatójú másodfokú egyenlet
gyökeinek számítására. Összehasonlításul a helyes gyökök, tíz értékes
jegyre kerekítve:

\[1971{,}605916 \quad \text{és} \quad 0{,}05077069387.\]

A diszkrimináns négy számjegyű aritmetikában

\[b^2 - 4ac = 9757 - 1{,}005 = 9756,\]

így

\[\sqrt{b^2 - 4ac} = 98{,}77.\]

A szabványos megoldóképlet ekkor a

\[\frac{98{,}78 \pm 98{,}77}{0{,}1002} = 1972 \quad \text{és} \quad 0{,}09980\]

gyököket adja. Az első gyök a helyesen kerekített négy számjegyű
eredmény, a másik viszont teljesen hibás, körülbelül 100 százalékos
hibával. A ludas a kioltás --- nem abban az értelemben, hogy a végső
kivonás helytelen volna (sőt, az egzaktul pontos), hanem abban, hogy a
vezető számjegyek kioltása csak a korábbi kerekítési hibákat hagyta meg.
Az alternatív megoldóképlet a

\[\frac{10{,}03}{98{,}78 \mp 98{,}77} = 1003 \quad \text{és} \quad 0{,}05077\]

gyököket adja. Ismét egy teljesen pontos és egy teljesen hibás gyököt
kaptunk, de minden esetben az előzővel ellentétes gyök romlott el. A
magyarázat ismét a kioltás, de az eltérő előjelminta miatt az ellenkező
gyök szennyeződik. Általában az egyes gyökök számításához azt a képletet
kell választanunk, amelyik elkerüli ezt a kioltást --- \(b\) előjelétől
függően.

\textbf{1.16. Példa. Szórás.} Egy véges \(x_i\), \(i = 1, \dots, n\)
valós sorozat átlagának definíciója

\[\bar{x} = \frac{1}{n} \sum_{i=1}^{n} x_i,\]

a szórásé pedig

\[\sigma = \left[ \frac{1}{n-1} \sum_{i=1}^{n} (x_i - \bar{x})^2 \right]^{1/2}.\]

Ezen képletek használata két menetet igényel az adatokon: egyet az
átlag, egy másikat a szórás számítására. A hatékonyság érdekében csábító
lehet a matematikailag ekvivalens

\[\sigma = \left[ \frac{1}{n-1} \left( \sum_{i=1}^{n} x_i^2 - n\bar{x}^2 \right) \right]^{1/2}\]

képletet használni a szórás számítására, hiszen így az összeg és a
négyzetösszeg is egyetlen menetben számítható.

Sajnos az egymenetes képlet végén fellépő egyetlen kioltás numerikusan
gyakran sokkal rombolóbb, mint a kétmenetes képlet összes kioltása
együttvéve. A probléma az, hogy az egymenetes képletben kivonandó két
mennyiség általában viszonylag nagy és közel azonos, így a különbségben
fellépő relatív hiba nagy lehet (sőt, az eredmény akár negatív is lehet,
és ilyenkor a négyzetgyökvonás is meghiúsul).

\textbf{1.17. Példa. Maradékok számítása.} Tegyük fel, hogy az
\(ax = b\) skaláris lineáris egyenletet oldjuk meg az \(x\)
ismeretlenre, és kaptunk egy \(\hat{x}\) közelítő megoldást. A válaszunk
minőségének egyik mérőszámaként kiszámítjuk az \(r = b - a\hat{x}\)
maradékot. Lebegőpontos aritmetikában

\[\mathrm{fl}(a\hat{x}) = a\hat{x}(1 + \delta_1), \quad |\delta_1| \le \epsilon_{\text{mach}},\]

így

\[\begin{aligned}
\mathrm{fl}(b - a\hat{x}) &= \bigl(b - a\hat{x}(1 + \delta_1)\bigr)(1 + \delta_2), \quad |\delta_2| \le \epsilon_{\text{mach}} \\
&= (r - a\hat{x}\delta_1)(1 + \delta_2) \\
&= r(1 + \delta_2) - a\hat{x}\delta_1 - a\hat{x}\delta_1\delta_2 \\
&\approx r(1 + \delta_2) - \delta_1 b.
\end{aligned}\]

Ám \(\delta_1 b\) akár \(\epsilon_{\text{mach}} b\) nagyságrendű is
lehet, ami pedig \(r\) nagyságrendű lehet --- ez azt jelenti, hogy a
maradék kiszámításában fellépő hiba akár 100 százalékos vagy annál
nagyobb is lehet. Ezért az \(r\) maradék értelmes kiszámításához nagyobb
pontosság szükséges lehet.

Szemléltetésként tegyük fel, hogy három számjegyű decimális aritmetikát
használunk, és legyen \(a = 2{,}78\), \(b = 3{,}14\),
\(\hat{x} = 1{,}13\). Három számjegyű aritmetikában
\(2{,}78 \times 1{,}13 = 3{,}14\), így a három számjegyű maradék \(0\).
Ez az eredmény egyfelől megnyugtató, mert azt mutatja, hogy a
felhasznált precízió mellett a lehető legjobbat tettük, másfelől viszont
egyetlen értékes jegyet sem ad a valódi maradékból, amely
\(3{,}14 - 2{,}78 \times 1{,}13 = 3{,}14 - 3{,}1414 = -0{,}0014\). Ami
történt: \(b\) és \(a\hat{x}\) tervezés szerint csaknem egyenlők, így a
különbségük csak kerekítési hibát tartalmaz. A feladat korábbi elemzése
szerint

\[3{,}14 = \mathrm{fl}(a\hat{x}) = a\hat{x}(1 + \delta_1) = 3{,}1414(1 + \delta_1),\]

amelyből \(\delta_1 = -0{,}0014/3{,}1414 \approx -0{,}00044566\). A
későbbi kivonásban nincs kerekítési hiba (mint a kioltás esetén
általában), így \(\delta_2 = 0\). A kapott számított maradék

\[\mathrm{fl}(b - a\hat{x}) \approx r(1 + \delta_2) - \delta_1 b = r - \delta_1 b \approx -0{,}0014 - (-0{,}0014) = 0.\]

Általánosságban a maradék pontos kiszámításához akár a munkaprecíziónak
a kétszeresére is szükség lehet.

\subsubsection{1.3.10 Más aritmetikai
rendszerek}\label{muxe1s-aritmetikai-rendszerek}

Manapság a tudományos számítások túlnyomó többségében IEEE szabvány
szerinti lebegőpontos aritmetikát használunk. Mivel szinte minden modern
mikroprocesszor hardverébe be van építve, az IEEE lebegőpontos
aritmetika rendkívül gyors, az IEEE dupla pontosság pedig a legtöbb
gyakorlati célra több mint elegendő pontosságot tesz lehetővé. Időnként
előfordulhatnak azonban olyan körülmények, amelyekben nagyobb
pontosságra van szükség --- például egy elkerülhetetlenül nagyon
érzékeny feladat megoldásához, vagy esetleg a puszta számítási mennyiség
miatt (bár ehhez rendszerint valóban hatalmas számítási mennyiség
szükséges). Ilyen esetekben a rendelkezésre álló pontosság
kiterjesztésére több alternatíva áll rendelkezésre, ám siettünk
hozzátenni, hogy a kiterjesztett pontosság használata nem
helyettesítheti a gondos feladatmegfogalmazást és a megbízható megoldási
technikákat.

Néhány számítógépes rendszer négyszeres pontosságot biztosít
hardveresen, és ez valószínűleg egyre elterjedtebbé válik, ahogy a
processzorok egyre gyorsabbak, a memóriák pedig egyre nagyobbak lesznek.
Több szoftvercsomag is elérhető, amely többszörös pontosságú
lebegőpontos aritmetikát biztosít (lásd az 1.4.3. szakaszt). Egyes
szimbolikus számítási környezetek --- például a Maple és a Mathematica
(lásd az 1.4.2. szakaszt) --- tetszőleges pontosságú aritmetikát
biztosítanak abban az értelemben, hogy a pontosság szükség szerint nő,
hogy a számítás során végig egzakt eredményeket tartsunk fenn. Az ilyen
szoftveres megközelítéseknek azonban jelentős hátrányai vannak
sebességben és memóriában: a szoftveresen végzett aritmetika
nagyságrendekkel lassabb lehet, mint a megfelelő hardveres lebegőpontos
műveletek, és nagy mennyiségű memóriára lehet szükség a kapott hosszú
számjegysorozatok tárolásához. Még ha a pontosság abban az értelemben
tetszőleges is, hogy változhat, végső soron az adott számítógépen
rendelkezésre álló véges memória korlátozza.

Bármilyen pontosságot is használunk, a numerikus analízis egyik fő
kérdése, hogy egy adott számításban valójában milyen pontosságot érünk
el. Ahogy láttuk ebben a fejezetben, és látni fogjuk a könyv hátralévő
részében is, a hagyományos hibaelemzés a megoldandó feladat, a használt
algoritmus, valamint annak az aritmetikának a tulajdonságain alapuló
matematikai elemzésen nyugszik, amellyel az algoritmust megvalósítjuk.
Az \emph{intervallumanalízis} alternatív megközelítés, amely a
bizonytalanságok továbbterjedését magába az aritmetikai rendszerbe építi
be, és ezzel legalábbis elvben automatikusan meghatározza a pontosságot.
Az alapötlet az, hogy minden numerikus mennyiséget egy zárt
intervallumként --- mondjuk \([a, b]\) --- ábrázolunk, amelyet az adott
mennyiség összes lehetséges értékének halmazaként értelmezünk. Az
intervallum hossza ekkor a mennyiség bizonytalanságát, azaz lehetséges
hibáját képviseli (ha nincs bizonytalanság, akkor \(a = b\)).

Két intervallumon végzett aritmetikai művelet eredménye a megfelelő
valós aritmetikai művelet összes lehetséges eredményének halmaza, a két
intervallumból vett operandusokkal:

\[[a, b] \mathbin{\mathrm{op}} [c, d] = \{x \mathbin{\mathrm{op}} y : x \in [a, b] \text{ és } y \in [c, d]\},\]

ahol \(\mathrm{op}\) a \(+\), \(-\), \(\times\), \(/\) műveletek
valamelyike. Az eredményül kapott halmaz maga is intervallum, amelynek
végpontjai a következő szabályok szerint számíthatók:

\[\begin{aligned}
[a, b] + [c, d] &= [a + c, \; b + d], \\
[a, b] - [c, d] &= [a - d, \; b - c], \\
[a, b] \times [c, d] &= \bigl[\min(ac, ad, bc, bd), \; \max(ac, ad, bc, bd)\bigr], \\
[a, b] / [c, d] &= [a, b] \times [1/d, \; 1/c], \quad \text{feltéve, hogy } 0 \notin [c, d].
\end{aligned}\]

Lebegőpontos aritmetikában megvalósítva az intervallum végpontjainak
kiszámítását megfelelően kell kerekíteni annak érdekében, hogy a
számított intervallum tartalmazza a megfelelő egzakt intervallumot (a
bal végpontot \(-\infty\) felé, a jobb végpontot \(+\infty\) felé
kerekítjük; az ilyen \emph{irányított kerekítés} az IEEE aritmetikában
opcióként elérhető).

A bizonytalanságok (köztük a kezdeti bemenet bizonytalanságainak)
továbbterjesztése egy számításon keresztül könnyű útnak tűnhet az
automatikus hibaelemzéshez, a gyakorlatban azonban a kapott
intervallumok gyakran túl szélesek ahhoz, hogy hasznosak legyenek ---
nagyjából ugyanabból az okból, amiért az előreható hibaelemzés is
gyakran túlzottan pesszimista korlátokat ad. Ha a számítást rögzített
precízióval dolgozó, egymás után alkalmazott függvények kompozíciójának
tekintjük, akkor a végső intervallum szélessége az összes függvény
kondíciószámának szorzatával arányos. Ha ezek a függvények számosak vagy
rosszul kondicionáltak --- vagy mindkettő ---, akkor a végső intervallum
valószínűleg túl széles lesz ahhoz, hogy hasznos információt nyújtson.
Egy másik probléma, hogy az intervallumanalízis figyelmen kívül
hagyhatja a számítási változók közötti függéseket --- például ha ugyanaz
a változó többször is megjelenik egy számításban ---, és ezeket
különálló bizonytalansági forrásokként kezelheti, ami ismét feleslegesen
széles intervallumokhoz vezethet. Triviális példa:
\([a, b] - [a, b] = [a - b, \; b - a]\) ahelyett, hogy \([0, 0]\) volna.

Az intervallumaritmetika által kapott intervallumszélességek kordában
tarthatók, ha változó precízióval valósítjuk meg. Az úgynevezett
\emph{tartomány-aritmetikában} (range arithmetic) például a szokásos
lebegőpontos aritmetikát két módon bővítjük ki: a mantissza tetszőleges
hosszúságú lehet egy előre meghatározott precízióhatárig, és egy további
\(r\) \emph{tartományjegyet} tartunk fenn, amely a mantissza
bizonytalanságát (\(\pm r\)) jelzi. A tartományszámokon végzett
aritmetikai műveletek a korábban megadott szokásos
intervallumaritmetikai szabályokat követik. A számítások
előrehaladtával, ha egy adott érték bizonytalansága nő, akkor a
mantisszájának hosszát ennek megfelelően csökkentjük. Ha a végső
eredmény mantisszája a pontossági követelmény teljesítéséhez túl kevés
számjegyet tartalmaz, akkor a precízióhatárt növeljük, és az egész
számítást annyiszor ismételjük, ahányszor szükséges a pontossági
követelmény teljesüléséhez. Ily módon sok (de nem minden) számítási
feladat megoldható tetszőlegesen előírt pontosságra, ám gyakran
megjósolhatatlan (és potenciálisan nagy) futásidő- és memóriaigény árán.

Több szoftvercsomag is elérhető, amely intervallum- vagy
tartomány-aritmetikát valósít meg (lásd az 1.4.3. szakaszt). Az említett
hátrányok korlátozták az intervallum- és tartomány-aritmetika általános
tudományos számításokban való használatát, ám az intervallumos módszerek
ennek ellenére bizonyos típusú feladatokra --- például nemlineáris
egyenletrendszerek többdimenziós garantált megoldási módszereire és
globális optimalizálási feladatokra --- a kevés elérhető lehetőség
egyikét nyújtják. Bár az intervallumos módszerek részletes bemutatása
meghaladja e könyv kereteit, a megfelelő helyeken jelezni fogjuk az
intervallum-alapú szoftverek elérhetőségét.

\subsubsection{1.3.11 Komplex aritmetika}\label{komplex-aritmetika}

Ahogyan az \(m/n\) alakú törtek --- ahol \(m\) és \(n\) egész számok ---
az egész számokat kiegészítve a racionális számokat alkotják, az
algebrai számok (például \(\sqrt{2}\)) és a transzcendens számok
(például \(\pi\)) pedig a racionális számokat kiegészítve a valós
számokat, úgy a valós számok \(\sqrt{-1}\)-gyel kiegészítve a
\emph{komplex} számokat alkotják. Ebben a könyvben a komplex számokat
explicit módon elsősorban a sajátérték-feladatok (4. fejezet) és a
Fourier-transzformációk (12. fejezet) tárgyalásánál használjuk. Máshol
az egyszerűség kedvéért nagyrészt a valós esetre szorítkozunk, ám az
általunk tárgyalt algoritmusok többsége a komplex esetre is
alkalmazható.

Ahogy a matematikában szokásos, ebben a könyvben is a \(\sqrt{-1}\)
képzetes egységet \(i\)-vel jelöljük, noha egyes mérnöki tudományágak
ehelyett \(j\)-t használnak. Egy adott \(z\) komplex szám kifejezhető a

\[z = x + iy\]

alakban, ahol az \(x = \mathrm{Re}(z)\) és \(y = \mathrm{Im}(z)\) valós
számok rendre \(z\) valós és képzetes részei. Az összes komplex szám
halmazát \(\mathbb{C} = \{x + iy : x, y \in \mathbb{R}\}\) jelöli. Egy
\(z\) komplex szám \emph{valós}, ha \(\mathrm{Im}(z) = 0\), és
\emph{képzetes} (vagy tisztán képzetes), ha \(\mathrm{Re}(z) = 0\).

\(\mathbb{C}\)-t elképzelhetjük egy kétdimenziós Descartes-féle síkként
--- ez a \emph{komplex számsík} (vagy más néven \emph{Argand-diagram})
---, amelynek két valós koordinátáját egy adott komplex szám valós és
képzetes része adja meg. Konvenció szerint a valós rész a vízszintes, a
képzetes rész pedig a függőleges tengelyen szerepel a komplex számsíkban
(lásd az 1.5. ábrát). Az \(i = \sqrt{-1}\)-vel való szorzás a komplex
számsíkban \(\pi/2\)-vel való elforgatásnak felel meg. Hasonlóképpen az
\(i^2 = -1\)-gyel való szorzás \(\pi\)-vel való elforgatásnak, az
\(i^4 = 1\)-gyel való szorzás pedig \(2\pi\)-vel való elforgatásnak
(azaz az \(1\), ahogy várható, a multiplikatív egység).

\pandocbounded{\includegraphics[keepaspectratio]{_page_30_Figure_8.jpeg}}

1.5. ábra: Komplex számsík a \(z\) komplex szám és \(\bar{z}\) komplex
konjugáltjának Descartes- és polárkoordinátás ábrázolásával.

Két \(z_1 = x_1 + iy_1\) és \(z_2 = x_2 + iy_2\) komplex szám szabványos
aritmetikai műveletei a következőképpen definiáltak:

\[\begin{aligned}
z_1 + z_2 &= (x_1 + x_2) + i(y_1 + y_2), \\
z_1 - z_2 &= (x_1 - x_2) + i(y_1 - y_2), \\
z_1 \times z_2 &= (x_1 x_2 - y_1 y_2) + i(x_1 y_2 + x_2 y_1), \\
z_1 / z_2 &= \left(\frac{x_1 x_2 + y_1 y_2}{x_2^2 + y_2^2}\right) + i\left(\frac{x_2 y_1 - x_1 y_2}{x_2^2 + y_2^2}\right), \quad \text{feltéve, hogy } z_2 \neq 0.
\end{aligned}\]

Megjegyezzük, hogy két komplex számon végzett aritmetikai művelet
kettőtől tizenegy valós aritmetikai műveletet igényel --- a konkrét
komplex művelettől függően ---, így a komplex aritmetika lényegesen
költségesebb lehet, mint a valós aritmetika.

Egy \(z = x + iy\) komplex szám \emph{komplex konjugáltját} \(\bar{z}\)
jelöli, és

\[\bar{z} = x - iy.\]

Egy komplex szám és komplex konjugáltja tehát a valós tengelyre nézve
szimmetrikus (egymás tükörképei) (lásd az 1.5. ábrát). Bármely
\(z = x + iy\) komplex számra

\[z\bar{z} = (x + iy)(x - iy) = x^2 + y^2\]

mindig valós és nemnegatív, így egy \(z = x + iy\) komplex szám
\emph{modulusa} (más néven \emph{abszolút értéke}) definiálható:

\[|z| = \sqrt{z\bar{z}} = \sqrt{x^2 + y^2},\]

amely éppen az \((x, y)\) pont és az origó \((0, 0)\) közötti szokásos
euklideszi távolság a komplex számsíkban.

Egy \(z\) komplex szám egy másik hasznos ábrázolása a polárkoordinátás
alak:

\[z = x + iy = r(\cos\theta + i\sin\theta),\]

ahol \(r = |z|\) és \(\theta = \arctan(y/x)\) (lásd az 1.5. ábrát). Az
Euler-azonosságot használva,

\[e^{i\theta} = \cos\theta + i\sin\theta,\]

megkapjuk az exponenciális ábrázolást:

\[z = r e^{i\theta}.\]

Ez utóbbi különösen fontos lesz a Fourier-transzformációk tárgyalásánál
(lásd a 12. fejezetet). Vegyük észre, hogy az exponenciális ábrázolásban
a komplex szorzás és osztás egyszerűbb alakú:

\[\begin{aligned}
z_1 \times z_2 &= r_1 r_2\, e^{i(\theta_1 + \theta_2)}, \\
z_1 / z_2 &= (r_1 / r_2)\, e^{i(\theta_1 - \theta_2)}.
\end{aligned}\]

Egy számítógépen egy komplex számot jellemzően két szokásos lebegőpontos
értékkel ábrázolunk, amelyek a valós és a képzetes résznek felelnek meg.
A komplex aritmetika támogatása nyelvenként és rendszerenként változó; a
Fortran minden változata például kifejezett támogatást ad a komplex
adattípusokhoz és aritmetikához, míg a C és leszármazottai jellemzően
nem. A MATLAB (lásd az 1.4.2. szakaszt) zökkenőmentesen támogatja a
komplex aritmetikát; sőt a MATLAB elsődleges numerikus adattípusa
komplex. A komplex aritmetikát a Python is támogatja --- a képzetes
számok \texttt{j} utótagúak, az ezekhez tartozó függvények pedig a
\texttt{cmath} modulban találhatók.

\subsection{1.4 Matematikai szoftver}\label{matematikai-szoftver}

Ez a könyv a numerikus analízis és a tudományos számítások témáinak
széles körét tárgyalja. Minden téma alapvető aspektusairól szót ejtünk,
ám arra nincs módunk, hogy bármelyikkel részletesen foglalkozzunk.
Ahhoz, hogy érdekes számítási feladatokat tudjunk megoldani, gyakran
szakemberek által írt matematikai szoftverekre támaszkodunk. Ha az
algoritmikus részleteket az ilyen szoftverekre bízzuk, a figyelmünket a
helyes feladatmegfogalmazásra és az eredmények értelmezésére
összpontosíthatjuk. Minden feladattípusnál csak a legalapvetőbb
algoritmusokat tekintjük át, elsősorban az általuk nyújtott rálátás
kedvéért --- hogy megalapozottan tudjunk megfelelő módszert választani
és azt körültekintően alkalmazni. Elsődleges célunk, hogy a matematikai
szoftverek értő felhasználói legyünk, ne pedig alkotói.

Mielőtt néhány konkrét forrást megemlítenénk a jó matematikai
szoftverekhez, foglaljuk össze azokat a kívánatos tulajdonságokat,
amelyekkel az ilyen szoftvereknek rendelkezniük kell --- nem valamilyen
fontossági sorrendben:

\begin{itemize}
\tightlist
\item
  \textbf{Megbízhatóság:} könnyű feladatokra mindig helyesen működik.
\item
  \textbf{Robusztusság:} nehéz feladatokra rendszerint működik, ám
  amikor mégis csődöt mond, azt kecsesen és tájékoztató jelleggel teszi.
\item
  \textbf{Pontosság:} a feladat és a bemenő adatok által indokolt
  mértékben pontos eredményt ad, lehetőleg az elért pontosság
  becslésével együtt.
\item
  \textbf{Hatékonyság:} a végrehajtási idő és a tárhely közel áll a
  megoldandó feladat lehető legkisebb igényéhez.
\item
  \textbf{Hordozhatóság:} kis vagy semmilyen módosítással alkalmazkodik
  új számítási környezetekhez.
\item
  \textbf{Karbantarthatóság:} könnyen érthető és módosítható.
\item
  \textbf{Használhatóság:} kényelmes és jól dokumentált felhasználói
  felülettel rendelkezik.
\item
  \textbf{Alkalmazhatóság:} feladatok széles körét képes megoldani.
\end{itemize}

Nyilvánvaló, hogy ezek a tulajdonságok gyakran ütköznek egymással, és
ritka az olyan szoftver, amely mindegyiknek megfelel. Ez a lista azonban
segít a matematikai szoftverek felhasználóinak felismerni, milyen
tulajdonságokat érdemes keresni, a fejlesztőknek pedig ad néhány
követendő célt.

\paragraph{1.4.1 Matematikai
szoftverkönyvtárak}\label{matematikai-szoftverkuxf6nyvtuxe1rak}

Itt sorolunk fel néhány széles körben elérhető, általános célú
matematikai szoftverforrást. Ennek a könyvnek minden fejezete végén
konkrét rutinokat ismertetünk adott típusú feladatokhoz --- mind ezekből
az általános könyvtárakból, mind specializáltabb csomagokból.

\begin{itemize}
\tightlist
\item
  \textbf{FMM:} Forsythe, Malcolm és Moler \emph{Computer Methods for
  Mathematical Computations} című könyvéhez {[}152{]} tartozó szoftver.
  Elérhető a Netlibből.
\item
  \textbf{GAMS:} \emph{Guide to Available Mathematical Software},
  amelyet a National Institute of Standards and Technology (NIST) állít
  össze. \texttt{gams.nist.gov}
\item
  \textbf{GSL (GNU Scientific Library):} Átfogó matematikai
  szoftverkönyvtár C és C++ nyelvhez. \texttt{www.gnu.org/software/gsl/}
\item
  \textbf{HSL (Harwell Subroutine Library):} Az angliai Harwell
  Laboratóriumban kifejlesztett szoftver. \texttt{www.hsl.rl.ac.uk}
\item
  \textbf{IMSL (International Mathematical and Statistical Libraries):}
  Átfogó matematikai szoftverkönyvtár. A Rogue Wave Software, Inc.
  kereskedelmi terméke.
  \texttt{www.roguewave.com/products-services/imsl-numerical-libraries}
\item
  \textbf{KMN:} Kahaner, Moler és Nash \emph{Numerical Methods and
  Software} {[}262{]} könyvéhez tartozó szoftver.
\item
  \textbf{NAG (Numerical Algorithms Group):} Átfogó matematikai
  szoftverkönyvtár. A NAG Inc. kereskedelmi terméke.
  \texttt{www.nag.com}
\item
  \textbf{NAPACK:} Hager \emph{Applied Numerical Linear Algebra}
  {[}220{]} könyvéhez tartozó szoftver. A lineáris algebrán kívül
  rutinokat tartalmaz nemlineáris egyenletekhez, feltétel nélküli
  optimalizáláshoz és gyors Fourier-transzformációkhoz is. Elérhető a
  Netlibből.
\item
  \textbf{Netlib:} Sokféle forrásból származó ingyenes szoftver.
  \texttt{www.netlib.org}
\item
  \textbf{NR (Numerical Recipes):} Press, Teukolsky, Vetterling és
  Flannery \emph{Numerical Recipes} {[}377{]} könyvéhez tartozó
  szoftver. Elérhető Fortran-, C- és C++-változatban.
  \texttt{numerical.recipes}
\item
  \textbf{NUMAL:} Az amszterdami Mathematisch Centrumban kifejlesztett
  szoftver. Algol- és Fortran-változatban is létezik, de leginkább Lau
  \emph{A Numerical Library for Scientists and Engineers} {[}297, 298{]}
  könyveiből érhető el C és Java változatban.
\item
  \textbf{NUMERALGO:} A \emph{Numerical Algorithms} folyóiratban
  megjelent szoftverek. Elérhető a Netlibből.
\item
  \textbf{PORT:} A Bell Laboratories-ben kifejlesztett szoftver. Egyes
  rutinok szabadon elérhetők a Netlibből, más részek csak kereskedelmi
  licenc alatt.
\item
  \textbf{SciPy/NumPy:} Átfogó, nyílt forráskódú matematikai
  szoftverkönyvtár Pythonhoz. \texttt{www.scipy.org}
\item
  \textbf{SLATEC:} Egy USA-beli kormányzati laboratóriumokból álló
  konzorcium által összeállított szoftver. Elérhető a Netlibből.
\item
  \textbf{SOL:} Optimalizálási és kapcsolódó feladatokhoz készült
  szoftver a Stanford University Systems Optimization Laboratoryjából.
  \texttt{web.stanford.edu/\textasciitilde{}saunders/brochure/brochure.html}
\item
  \textbf{TOMS:} Az \emph{ACM Transactions on Mathematical Software}
  (korábban \emph{Collected Algorithms of the ACM}) folyóiratban
  megjelent szoftverek. Az algoritmusokat szám (megjelenési sorrendben)
  és név szerint egyaránt azonosítjuk. Elérhető a Netlibből.
\end{itemize}

\paragraph{1.4.2 Tudományos számítási
környezetek}\label{tudomuxe1nyos-szuxe1muxedtuxe1si-kuxf6rnyezetek}

Az imént felsorolt szoftverkönyvtárak olyan szubrutinokat tartalmaznak,
amelyeket a felhasználó által írt --- rendszerint egy hagyományos
programozási nyelvben, például Fortranban vagy C-ben készült ---
programoknak kell meghívniuk. A tudományos számítások egyre népszerűbb
alternatívája az olyan interaktív környezetek használata, amelyek
erőteljes, kényelmesen elérhető, beépített matematikai képességeket
kínálnak, gyakran kifinomult grafikával és egy nagyon magas szintű
programozási nyelvvel párosítva, amely új algoritmusok gyors
prototípus-készítésére alkalmas.

Az egyik leggyakrabban használt ilyen számítási környezet a MATLAB,
amely a The MathWorks, Inc. (\texttt{www.mathworks.com}) kereskedelmi
terméke. A MATLAB (MATrix LABoratory) egy interaktív rendszer, amely
kiterjedt matematikai képességeket --- különösen lineáris algebrában
---, nagy teljesítményű tudományos vizualizációt, magas szintű
programozási nyelvet és többféle választható „toolboxot" ötvöz, amelyek
speciális képességeket nyújtanak egy-egy alkalmazási területhez, például
jelfeldolgozáshoz, képfeldolgozáshoz, szabályozáshoz,
rendszeridentifikációhoz, optimalizáláshoz vagy statisztikához. A MATLAB
sokféle számítógépre elérhető, szakmai és hallgatói kiadásban egyaránt.
Hasonló képességeket nyújt az Octave, amely szabadon elérhető
(\texttt{www.gnu.org/software/octave}). További, egyre népszerűbb
alternatíva a Python, amely általános célú magas szintű programozási
nyelvet, interaktív végrehajtási környezetet, kiterjedt matematikai
szoftverkönyvtárakat (pl. SciPy) és átfogó ábrázolási képességeket
(Matplotlib) kínál.

Az interaktív számítási környezetek egy másik családja elsősorban
szimbolikus (nem pedig numerikus) számításon alapul, amit gyakran
\emph{számítógépes algebrának} neveznek. Ezek a csomagok --- az Axiom,
Derive, Magma, Maple, Mathematica, Maxima és Reduce --- sok ugyanolyan
matematikai és grafikus képességet nyújtanak, és ezenfelül szimbolikus
differenciálást, integrálást, egyenletmegoldást, polinommanipulációt és
hasonlókat kínálnak, valamint tetszőleges pontosságú aritmetikát is. A
MATLAB egy „Symbolic Math" toolbox révén szintén rendelkezik szimbolikus
számítási képességekkel, Pythonhoz pedig a SymPy könyvtár kínál hasonló
lehetőségeket.

Mivel a MATLAB és a Python ezek közül a környezetek közül a
leggyakrabban használtak a könyvben tárgyalt feladattípusokhoz, az egyes
feladatkategóriákhoz rendelkezésre álló szoftverek összegzéseiben a
hagyományos szoftverkönyvtárak mellett konkrét MATLAB- és
Python-függvényeket --- akár az alapkörnyezetből, akár a kiegészítő
toolboxokból és könyvtárakból --- is felsorolunk.

\paragraph{1.4.3 Kiterjesztett aritmetikai
csomagok}\label{kiterjesztett-aritmetikai-csomagok}

A többszörös pontosságú aritmetikát biztosító szoftvercsomagok közé
tartoznak a TOMS-ból származó MP(\#524), FM(\#693), MPFUN(\#719) és
ZM(\#786); a CLN (Class Library for Numbers); a GMP (GNU Multiple
Precision Arithmetic Library); a Python \texttt{mpmath} könyvtára;
valamint az NR {[}377{]} és {[}378{]} rutinjai. Intervallum- vagy
tartomány-aritmetikát megvalósító szoftvercsomagok és rendszerek közé
tartoznak a TOMS-ból származó INTLIB(\#737, \#763); a PROFIL/BIAS
{[}278{]}; az Arb (\texttt{arblib.org}); az INTLAB toolbox a MATLAB-hoz;
valamint az {[}1{]}-hez tartozó szoftver.

\paragraph{1.4.4 Gyakorlati tanácsok a
szoftverhasználathoz}\label{gyakorlati-tanuxe1csok-a-szoftverhasznuxe1lathoz}

Ez a szakasz gyakorlati tanácsokat tartalmaz a könyvben hivatkozott
szoftverek beszerzésére és használatára, különös tekintettel az egyes
fejezetek végén található számítógépes feladatok alapján készíthető
programozási házi feladatokra. A számítógépes feladatok egyik konkrét
szoftvertől vagy programozási nyelvtől sem függnek, így több lehetőség
is rendelkezésre áll. Egy adott esetben a legjobb választás a
felhasználó tapasztalatától, erőforrásaitól és céljaitól függ.

A hivatkozott szoftverek sokféle forrásból származnak: nagy kereskedelmi
könyvtárakból, például az IMSL-ből és a NAG-ból; nyilvános, szabadon
hozzáférhető tárolókból, például a Netlibből; valamint tudományos
számítási környezetekből, például a MATLAB-ból és a Pythonból. Sok
akadémiai és ipari számítási központban és laboratóriumban az ilyen
szoftverek reprezentatív mintája már telepítve és használatra kész
állapotban rendelkezésre áll. Mindenesetre e könyv összes számítógépes
feladatához bőséges mennyiségű szoftver szerezhető be ingyenesen az
interneten keresztül, vagy jelképes összegért egyéb forrásokból (például
a kísérő tankönyvekből). A megfelelő szoftver megtalálása, letöltése és
telepítése hasznos, valós tapasztalatot jelent, és az ennek során
megszerzett készségek fontos gyakorlati kiegészítői a könyvben
tanultaknak.

Talán a legfontosabb választás a programozási nyelvé. Ahogy az 1.4.2.
szakaszban már tárgyaltuk, a MATLAB (vagy az Octave) és a Python
környezetei nemcsak magas szintű programozási nyelvet nyújtanak, hanem
viszonylag zökkenőmentes hozzáférést adnak kiterjedt matematikai
szoftverkönyvtárakhoz is. Az ilyen környezetek felhasználói több előnyt
is élvezhetnek. A felhasználói programok általában sokkal rövidebbek
lesznek az explicit deklarációk, a tárhelykezelés és sok ciklus
kiiktatása miatt. További bónusz a beépített grafika, amely feleslegessé
teszi ennek külön utófeldolgozási fázisban történő elvégzését.

A hagyományos programozási nyelvek közül a Fortran volt a tudományos
számítások eredeti nyelve, és a régi szoftverkönyvtárak és alkalmazási
kódok nagy része Fortranban íródott. E könyv mellett a
Fortran-felhasználó előnyére válik, hogy a legszélesebb választék áll
rendelkezésre szoftverekben, és kompatibilis a létező alkalmazási kódok
túlnyomó többségével. Emellett, mivel a Fortran viszonylag szigorú
nyelv, és a hozzá tartozó fordítók sok év optimalizálási munkáján vannak
túl, a Fortran egyes számítógépes rendszereken gyakran hatékonyabb
futtatható kódot állít elő.

A C sokoldalúbb és kifejezőbb nyelv, mint a Fortran, és a C (valamint
változatai, például a C++ és a Java) gyakrabban szerepel kezdő
programozási tárgyak tananyagában. Ha ezzel a könyvvel C-t szeretnél
használni, bőséges mennyiségű szoftver áll rendelkezésedre. Ezenfelül
sok számítógépes rendszeren viszonylag egyszerű Fortran-rutinokat
C-programokból meghívni. A főbb eltérések, amelyekre figyelni kell: a
rutinok neve kissé módosulhat (gyakran a szokásos név elé vagy után
kerül egy aláhúzás); minden argumentumot Fortran szubrutinokba cím
szerint (azaz mutatóként C-ben) kell átadni; a C és a Fortran ellentétes
konvenciókat használnak a tömbök tárolására (a C mátrixokat sor-, a
Fortran mátrixokat oszlopfolytonosan tárolja). Végül a Fortran
forráskódja automatikusan át is alakítható C-re a Netlibből elérhető
\texttt{f2c} konverterrel, így az --- például az interneten ---
beszerzett Fortran-rutinok könnyen használhatók C-programokkal.

A könyv egyes számítógépes feladatai grafikus kimenetet igényelnek. A
szükséges ábrák előállítására több lehetőség is rendelkezésre áll, a
használt számítási környezettől függően. Az olyan interaktív
környezetekhez, mint a MATLAB vagy a Python, kiterjedt grafikai
képességek automatikusan elérhetők. Hagyományos programozási nyelvekkel
való használatra szabadon elérhető ábrázoló csomagok többek között a
\texttt{gnuplot} (\texttt{www.gnuplot.info}), a \texttt{pgplot}
(\texttt{www.astro.caltech.edu/\textasciitilde{}tjp/pgplot}) és a
\texttt{plplot} (\texttt{plplot.sourceforge.net}). Kifinomultabb és
erőteljesebb tudományos vizualizációs rendszerek is elérhetők, ám ezek
képességei messze meghaladják a könyv számítógépes feladataihoz
szükséges, viszonylag egyszerű ábrázolásokat.

Egy másik fontos programozási szempont a teljesítmény. A mai
mikroprocesszor-alapú számítógépes rendszerek teljesítménye gyakran
kritikusan függ a memóriahierarchia (regiszterek, gyorsítótár, RAM,
lemez stb.) megfontolt kihasználásától, mind a programozó, mind az
optimalizáló fordító részéről. Ezért nemcsak a megfelelő algoritmust
kell megválasztani, hanem gondosan meg is kell valósítani, hogy
maximalizáljuk az adatok újrafelhasználását, miközben azok a
memóriahierarchia gyorsabb elérési idejű részeiben vannak. A programozás
teljesítményszempontjairól lásd {[}121, 188, 201{]}. Szerencsére az
ilyen programozás részletei általában rejtve maradnak a felhasználó
előtt a könyvben ajánlott könyvtári rutinok belsejében. Ez csak egyike a
létező, szakemberek által írt tudományos számítási szoftverek
használatából adódó számos előnynek, ahol csak lehetséges.

Ha olyan tudományos számítási környezetet használsz, mint a MATLAB vagy
a Python, tudnod kell, hogy jelentős teljesítménykülönbségek lehetnek a
beépített műveletek --- amelyek általában nagyon gyorsak --- és az
általad explicit módon programozottak között, amelyek a tolmács
(interpreter) üzemmód és a memóriakezelés többletköltsége miatt
általában sokkal lassabbak. Ezért ilyen körülmények között óvatosnak
kell lenni a teljesítmény-összehasonlításokban. Például lehet, hogy
egyik algoritmus elvileg rosszabb a másiknál, mégis jobban teljesít,
mert hatékonyabban hasznosítja a gyors beépített műveleteket.

\subsection{1.5 Történeti jegyzetek és további
olvasnivaló}\label{tuxf6rtuxe9neti-jegyzetek-uxe9s-tovuxe1bbi-olvasnivaluxf3}

A ma numerikus analízisnek vagy tudományos számításoknak nevezett
tudományág jóval megelőzi a modern számítógépek megjelenését. A ma is
használatos fogalmak és módszerek közül sokat a huszadik század előtti
óriások --- Newton (1642--1727), Euler (1707--1783), Lagrange
(1736--1813), Laplace (1749--1827), Legendre (1752--1833), Gauss
(1777--1855), Cauchy (1789--1857), Jacobi (1804--1851), Adams
(1819--1892), Csebisev (1821--1894), Hermite (1822--1901), Laguerre
(1834--1886) és mások --- fogalmaztak meg először, és nevük végigvonul
ezen a könyvön. A fő feladat akkor, mint most is, a fizikában, a
csillagászatban, a geodéziában és más tudományterületeken felmerülő
matematikai feladatok közelítő megoldására szolgáló hatékony módszerek
megtalálása volt. Sőt a számítási erőforrások hatékony felhasználása még
kritikusabb, ha ceruzával, papírral és emberi agyerővel (esetleg
zsebszámológéppel) dolgozunk, mint egy modern, nagy sebességű
számítógéppel.

Nagyrészt a modern számítógépek egyszerűen megnövelték a megoldható
feladatok méretét. Emellett szükségessé tették a kerekítési hibák
gondosabb elemzését és kontrollját is, mert a számítást már nem ember
végzi, aki szükség szerint további pontossággal dolgozhat. Nem kétséges
azonban, hogy a digitális számítógépek kifejlődése adta a lökést a
numerikus analízis virágzó, termékeny és lendületesen fejlődő területté
válásához, amely lehetővé tette, hogy a számítás ma a modern
természettudomány és mérnöki gyakorlat minden területén mindenütt jelen
legyen. Sőt a számítás az elmélet és a kísérlet mellett a tudományos
ismeretek és a mérnöki gyakorlat előrehaladásának egyenrangú és
nélkülözhetetlen partnerévé vált {[}267{]}.

A numerikus analízis korai történetéről lásd {[}193{]}-at; a tudományos
számításoknak a számítógépes korszakban való újabb fejlődéséről lásd
{[}57, 63, 346{]}-ot. A numerikus analízis irodalma --- a tankönyvektől
a kutatási monográfiákon át a folyóiratokig --- jóval terjedelmesebb
annál, semhogy itt megfelelően tárgyalni lehetne. Ebben a szövegben
megpróbáljuk megfelelően megadni az elismerést a bemutatott fő
gondolatok eredetének (legalábbis azoknak, amelyek nem nyilvánvalók már
a névről), és (rendszerint másodlagos) forrásokat idézünk a további
olvasáshoz, ám ezek a hivatkozások és ajánlások semmiképpen sem
teljesek. Túl sok kiváló általános numerikus analízis tankönyv létezik
ahhoz, hogy mindegyiket megemlítsük, de közülük sok ma is érdemes
olvasmány (beleértve a régebbieket is, amelyek közül sokat olcsó
utánnyomásokban is újra kiadtak). Csak a tárgyalásunkhoz leginkább
közvetlenül kapcsolódókat idézzük.

A legtöbb numerikus analízis tankönyv általánosan tárgyalja a
hibaelemzést. A kerekítési hibák elemzésének alapművei közé tartozik az
{[}500{]} --- értékes meglátások kincsesháza. Szerzője, James H.
Wilkinson kulcsszerepet játszott a hátraható hibaelemzés fogalmának
kidolgozásában és elterjesztésében, és számos híres „számítási
ellenpéldáért" is felelős, amelyek váratlan feladatokban és
algoritmusokban fellépő numerikus instabilitásokat tárnak fel {[}502,
503{]}. Hasonló szellemben készült újabb munka a {[}237{]}. A
hibaelemzés automatizálására --- beleértve az intervallumaritmetikát ---
különféle megközelítésekről lásd {[}1, 6, 213, 254, 329, 337, 338, 348,
469{]}-et. Hibaelemzésre szolgáló MATLAB toolboxot tárgyal {[}69{]}. A
kondicionáltság általános fogalmának egy korai tárgyalása {[}385{]}-ben
található.

A számítógépes aritmetika általános tárgyalásai közé tartoznak {[}54,
282, 283, 287, 356, 363{]}. A lebegőpontos aritmetika részletes
tárgyalásához lásd {[}345, 423{]}-at; rövidebb áttekintéshez
{[}279{]}-et. A lebegőpontos aritmetika szabványosítására irányuló
erőfeszítést és a létrejött szabvány magas színvonalát nagyrészt William
Kahan inspirálta, aki számos ismert számítási ellenpéldáért is felelős.
Az IEEE lebegőpontos szabvány {[}247{]}-ben található. A lebegőpontos
aritmetikáról és az IEEE szabványról szóló hasznos bemutatók közé
tartozik {[}189, 361{]}.

A matematikai szoftver mint a numerikus analízis és az informatika
részterületének megjelenéséről lásd a {[}80{]} áttekintést és a {[}42,
88, 131, 230, 251, 386, 387{]} gyűjteményeket. A szakmai minőségű
szoftvereket (nem csak szemléltető kódrészleteket) tartalmazó talán
legkorábbi numerikus módszerek tankönyv a {[}406{]} volt, amely
hangvételében, stílusában és tartalmában hasonló Forsythe, Malcolm és
Moler nagy hatású könyvéhez {[}152{]}, amely ezt a megközelítést
népszerűsítette. További, bevezető szintű tárgyalást és példákat nyújtó
tankönyvek közé tartoznak {[}13, 18, 85, 123, 208, 214, 240, 262, 331,
353, 380, 407, 428, 472, 481{]}. Haladóbb, általános tankönyvek közé
tartoznak {[}17, 87, 92, 93, 111, 172, 196, 248, 275, 358, 371, 379,
402, 429, 434, 448, 459, 474{]}. A numerikus számítás buktatóinak
felismeréséhez és elkerüléséhez gyakorlati tanácsokat nyújt {[}3, 4,
127, 151, 422{]}.

Strang könyvei {[}435, 439, 440, 441{]} kiváló hátteret és rálátást
nyújtanak az alkalmazott matematika számos olyan területére, amelyek
általában is --- és különösen e könyv szinte minden fejezete
szempontjából --- fontosak a numerikus számításhoz. A tudományos
programozás elemi bevezetéseihez lásd {[}265, 480, 482, 517, 518{]}-at.
Numerikus szoftverek tervezéséhez, megvalósításához és teszteléséhez ---
nem pusztán a használatukhoz --- lásd {[}43, 80, 127, 328, 354{]}-et.
További számítógépes gyakorlatok és projektek találhatók {[}89, 94, 130,
155, 167, 170, 200, 203, 205, 294, 353, 418{]}-ban.

\subsubsection{Áttekintő
kérdések}\label{uxe1ttekintux151-kuxe9rduxe9sek}

\begin{itemize}
\tightlist
\item
  \textbf{1.1.} Igaz vagy hamis: Egy feladat rosszul kondicionált, ha a
  megoldása nagyon érzékeny a feladat adatainak kis változásaira.
\item
  \textbf{1.2.} Igaz vagy hamis: Nagyobb precíziójú aritmetika
  használata egy rosszul kondicionált feladatot jobban kondicionálttá
  tesz.
\item
  \textbf{1.3.} Igaz vagy hamis: Egy feladat kondicionáltsága függ a
  megoldására használt algoritmustól.
\item
  \textbf{1.4.} Igaz vagy hamis: Egy jó algoritmus a feladat
  kondíciójától függetlenül pontos megoldást ad.
\item
  \textbf{1.5.} Igaz vagy hamis: Egy feladat megoldására választott
  algoritmusnak nincs hatása a propagált adathibára.
\item
  \textbf{1.6.} Igaz vagy hamis: Egy stabil algoritmus jól kondicionált
  feladatra alkalmazva szükségképpen pontos megoldást ad.
\item
  \textbf{1.7.} Igaz vagy hamis: Ha két valós szám lebegőpontos számként
  pontosan ábrázolható, akkor a rajtuk végzett valós aritmetikai művelet
  eredménye is ábrázolható lebegőpontos számként.
\item
  \textbf{1.8.} Igaz vagy hamis: A lebegőpontos számok egyenletesen
  oszlanak el az értékkészletükön.
\item
  \textbf{1.9.} Igaz vagy hamis: A lebegőpontos összeadás asszociatív,
  de nem kommutatív.
\item
  \textbf{1.10.} Igaz vagy hamis: Egy lebegőpontos számrendszerben az
  alulcsordulási szint az a legkisebb pozitív szám, amely az 1-hez
  hozzáadva perturbálja az 1-et.
\item
  \textbf{1.11.} Igaz vagy hamis: Az IEEE dupla pontosságú lebegőpontos
  aritmetika mantisszája pontosan kétszer olyan hosszú, mint az IEEE
  egyszeres pontosságúé.
\item
  \textbf{1.12.} Mely három tulajdonság jellemzi a jól kitűzött
  feladatot?
\item
  \textbf{1.13.} Sorolj fel három hibaforrást a tudományos
  számításokban.
\item
  \textbf{1.14.} Fejtsd ki a különbséget a csonkítási (vagy
  diszkretizációs) hiba és a kerekítési hiba között.
\item
  \textbf{1.15.} Fejtsd ki a különbséget az abszolút és a relatív hiba
  között.
\item
  \textbf{1.16.} Fejtsd ki a különbséget a számítási hiba és a propagált
  adathiba között.
\item
  \textbf{1.17.} Fejtsd ki a különbséget a precízió és a pontosság
  között.
\item
  \textbf{1.18.} (a) Mit értünk egy feladat kondicionáltságán? (b)
  Függ-e ez a megoldásra használt algoritmustól? (c) Függ-e ez a
  megoldáshoz használt aritmetika precíziójától?
\item
  \textbf{1.19.} Ha egy számítási feladat kondíciószáma \(1\), ez jó
  vagy rossz? Miért?
\item
  \textbf{1.20.} Fejtsd ki a különbséget a relatív és az abszolút
  kondíciószám között.
\item
  \textbf{1.21.} Mi az inverz feladat? Hogyan viszonyul egymáshoz egy
  feladat és az inverz feladat kondicionáltsága?
\item
  \textbf{1.22.} (a) Mit értünk egy számított eredmény hátraható
  hibáján? (b) A hátraható hibaelemzés szerint mikor tekinthető jónak
  egy adott feladat közelítő megoldása?
\item
  \textbf{1.23.} Tegyük fel, hogy egy adott feladatot egy adott
  algoritmussal oldunk meg. Az alábbiak mindegyikéről mondd meg, hogy
  függ-e az algoritmus stabilitásától, és miért: (a) propagált adathiba;
  (b) a számított eredmény pontossága; (c) a feladat kondicionáltsága.
\item
  \textbf{1.24.} (a) Fejtsd ki a különbséget az előreható és a hátraható
  hiba között. (b) Hogyan kapcsolódik az előreható és a hátraható hiba
  mennyiségileg egymáshoz?
\item
  \textbf{1.25.} Adott lebegőpontos számrendszer esetén írd le szavakkal
  a gépi számok eloszlását a valós egyenes mentén.
\item
  \textbf{1.26.} Lebegőpontos aritmetikában általában melyik károsabb:
  az alulcsordulás vagy a túlcsordulás? Miért?
\item
  \textbf{1.27.} Lebegőpontos aritmetikában az alábbi műveletek közül,
  két pozitív lebegőpontos operandusra, melyek okozhatnak túlcsordulást?
  (a) összeadás; (b) kivonás; (c) szorzás; (d) osztás.
\item
  \textbf{1.28.} Lebegőpontos aritmetikában az alábbi műveletek közül,
  két pozitív lebegőpontos operandusra, melyek okozhatnak
  alulcsordulást? (a) összeadás; (b) kivonás; (c) szorzás; (d) osztás.
\item
  \textbf{1.29.} Sorolj fel két okot, amiért a lebegőpontos
  számrendszereket rendszerint normalizálják.
\item
  \textbf{1.30.} Egy lebegőpontos rendszerben mely mennyiség határozza
  meg egy adott valós szám gépi számmal való ábrázolásakor fellépő
  maximális relatív hibát?
\item
  \textbf{1.31.} (a) Fejtsd ki a különbséget a „nulla felé kerekítés" és
  a „legközelebbihez kerekítés" szabályai között egy lebegőpontos
  rendszerben. (b) E két szabály közül melyik pontosabb? (c) Mekkora
  számszerű különbséget okoz ez az \(\epsilon_{\text{mach}}\) egységnyi
  kerekítési hibában?
\item
  \textbf{1.32.} Egy \(p\) számjegyű bináris lebegőpontos rendszerben
  legközelebbihez kerekítés mellett mekkora az
  \(\epsilon_{\text{mach}}\) egységnyi kerekítési hiba értéke?
\item
  \textbf{1.33.} Egy fokozatos alulcsordulást (szubnormális számokat)
  engedélyező lebegőpontos rendszerben minden szám ábrázolása egyértelmű
  marad-e? Miért?
\item
  \textbf{1.34.} Egy lebegőpontos rendszerben két gépi szám szorzata
  rendszerint pontosan ábrázolható a lebegőpontos rendszerben? Miért?
\item
  \textbf{1.35.} Egy lebegőpontos rendszerben két nem nulla gépi szám
  hányadosa mindig pontosan ábrázolható a lebegőpontos rendszerben?
  Miért?
\item
  \textbf{1.36.} (a) Adj példát, amely megmutatja, hogy a lebegőpontos
  összeadás nem feltétlenül asszociatív. (b) Adj példát, amely
  megmutatja, hogy a lebegőpontos szorzás nem feltétlenül asszociatív.
\item
  \textbf{1.37.} (a) Milyen körülmények között lép fel számjegykioltás
  egy lebegőpontos rendszerben? (b) A kioltás fellépése azt jelenti-e,
  hogy az azt kiváltó konkrét művelet valódi eredménye nem ábrázolható
  pontosan a lebegőpontos rendszerben? Miért? (c) Miért általában rossz
  a kioltás?
\item
  \textbf{1.38.} Adj példát olyan számra, amelynek a tízes
  számrendszerbeli alakja véges (azaz véges számú nem nulla számjegyet
  tartalmaz), de a bináris alakja nem.
\item
  \textbf{1.39.} Adj példákat olyan lebegőpontos aritmetikai
  műveletekre, amelyek a \texttt{Inf} és a \texttt{NaN} kivételes
  értékek mindegyikét előállítják.
\item
  \textbf{1.40.} Egy \(\beta\) alapú, \(p\) precíziójú, legközelebbihez
  kerekítést használó lebegőpontos rendszerben mi a maximális relatív
  hiba bármely nem nulla valós szám ábrázolásában, a rendszer
  tartományán belül?
\item
  \textbf{1.41.} Fejtsd ki, miért gyakran rossz a két hasonló abszolút
  értékű szám kivonásakor fellépő kioltás, annak ellenére, hogy az
  eredmény az operandusokra nézve pontosan helyes lehet.
\item
  \textbf{1.42.} Tegyük fel, hogy egy tízes alapú lebegőpontos
  rendszerrel dolgozunk, amelynek gépi precíziója
  \(\epsilon_{\text{mach}} = 10^{-5}\), exponens-tartománya pedig
  \(\pm 20\). Mi az alábbi lebegőpontos aritmetikai műveletek eredménye?
  (a) \(1 + 10^{-7}\); (b) \(1 + 10^{3}\); (c) \(1 + 10^{7}\); (d)
  \(10^{10} + 10^{3}\); (e) \(10^{10} / 10^{-15}\); (f)
  \(10^{-10} \times 10^{-15}\).
\item
  \textbf{1.43.} Egy olyan lebegőpontos rendszerben, amelynek
  alulcsordulási szintje \(\text{UFL} = 10^{-38}\), az alábbi számítások
  közül melyekben lép fel alulcsordulás?

  \begin{enumerate}
  \def\labelenumi{(\alph{enumi})}
  \tightlist
  \item
    \(a = \sqrt{b^2 + c^2}\), ahol \(b = 1\), \(c = 10^{-25}\).
  \item
    \(a = \sqrt{b^2 + c^2}\), ahol \(b = c = 10^{-25}\).
  \item
    \(u = (v \times w)/(y \times z)\), ahol \(v = 10^{-15}\),
    \(w = 10^{-30}\), \(y = 10^{-20}\), \(z = 10^{-25}\). Azokban az
    esetekben, ahol alulcsordulás lép fel, ésszerű-e egyszerűen nullára
    állítani az alulcsorduló mennyiséget?
  \end{enumerate}
\item
  \textbf{1.44.} (a) Fejtsd ki szavakkal a különbséget egy lebegőpontos
  rendszer \(\epsilon_{\text{mach}}\) egységnyi kerekítési hibája és
  \(\text{UFL}\) alulcsordulási szintje között. E két mennyiség közül
  (b) melyik függ csak a mantissza mezőjében található számjegyek
  számától? (c) melyik függ csak az exponens mezőjében található
  számjegyek számától? (d) melyik nem függ a használt kerekítési
  szabálytól? (e) melyiket nem befolyásolja a szubnormális számok
  engedélyezése?
\item
  \textbf{1.45.} Legyen \(x_k\) pozitív számok szigorúan csökkenő, véges
  sorozata (azaz \(x_k > x_{k+1}\) minden \(k\)-ra). Ha a számokat
  tetszőleges általunk választott sorrendben vehetjük, milyen sorrendben
  kellene összegezni a sorozatot a kerekítési hiba minimalizálása
  érdekében?
\item
  \textbf{1.46.} Példa-e a kioltás a kerekítési hibára? Miért?
\item
  \textbf{1.47.} (a) Magyarázd meg, miért lehet véges az összege
  lebegőpontos aritmetikában egy olyan divergens végtelen sornak, mint
\end{itemize}

\[\sum_{n=1}^{\infty} \frac{1}{n}.\]

\begin{enumerate}
\def\labelenumi{(\alph{enumi})}
\setcounter{enumi}{1}
\tightlist
\item
  Melyik ponton fognak a részletösszegek megszűnni változni?
\end{enumerate}

\begin{itemize}
\tightlist
\item
  \textbf{1.48.} Lebegőpontos aritmetikában, ha egy pozitív tagú
  konvergens végtelen sor
\end{itemize}

\[S = \sum_{i=1}^{\infty} x_i\]

összegét számítod természetes sorrendben, milyen megállási kritériumot
alkalmaznál a lehető legnagyobb pontosság eléréséhez a legkevesebb
taggal? - \textbf{1.49.} Fejtsd ki, miért nehéz lebegőpontos
aritmetikában pontosan kiértékelni egy váltakozó előjelű végtelen sort,
például

\[e^x = 1 + x + \frac{x^2}{2!} + \frac{x^3}{3!} + \cdots\]

\(x < 0\) esetén. - \textbf{1.50.} Ha \(f\) egy valós változós valós
értékű függvény, akkor a deriváltjának

\[f'(x) \approx \frac{f(x+h) - f(x)}{h}\]

véges differencia közelítésében a csonkítási hiba \(h \to 0\) esetén
nullához tart. Ha lebegőpontos aritmetikát használunk, sorolj fel két
tényezőt, amely korlátozza, mennyire kicsi \(h\) értéket tudunk a
gyakorlatban használni. - \textbf{1.51.} Sorolj fel legalább két módot,
ahogy a

\[x = \frac{-b \pm \sqrt{b^2 - 4ac}}{2a}\]

másodfokú megoldóképlet kiértékelése lebegőpontos aritmetikában
numerikus nehézségekbe ütközhet.

\subsubsection{Feladatok}\label{feladatok}

\begin{itemize}
\tightlist
\item
  \textbf{1.1.} Az átlagos emberi testhőmérsékletet szokás \(98{,}6\)
  Fahrenheit-foknak megadni, amit vélhetnénk egy nagy népesség
  átlagaként, három értékes jegyre kerekítve. Valójában azonban a
  \(98{,}6\) egyszerűen a \(37\) Celsius-fok Fahrenheit-megfelelője,
  amely csak két értékes jegyre pontos.

  \begin{enumerate}
  \def\labelenumi{(\alph{enumi})}
  \tightlist
  \item
    Mi a maximális relatív hiba az elfogadott értékben, feltéve, hogy az
    \(\pm 0{,}05\,^\circ\text{F}\) pontosságú?
  \item
    Mi a maximális relatív hiba az elfogadott értékben, feltéve, hogy az
    \(\pm 0{,}5\,^\circ\text{C}\) pontosságú?
  \end{enumerate}
\item
  \textbf{1.2.} Mi a közelítő abszolút és relatív hiba \(\pi\)
  közelítésére a következő mennyiségek mindegyikével?

  \begin{enumerate}
  \def\labelenumi{(\alph{enumi})}
  \tightlist
  \item
    \(3\); (b) \(3{,}14\); (c) \(22/7\).
  \end{enumerate}
\item
  \textbf{1.3.} Ha \(a\) egy \(t\) valódi értékű mennyiség közelítő
  értéke, és \(a\) relatív hibája \(r\), akkor e fogalmak definícióiból
  bizonyítsuk be, hogy \(a = t(1 + r)\).
\item
  \textbf{1.4.} Tekintsük a \(\sin(x)\) függvény kiértékelésének
  feladatát --- különösen a propagált adathibát, azaz a függvényérték
  hibáját, amelyet az \(x\) argumentum \(h\) perturbációja okoz.

  \begin{enumerate}
  \def\labelenumi{(\alph{enumi})}
  \tightlist
  \item
    Becsüld meg a \(\sin(x)\) kiértékelésének abszolút hibáját.
  \item
    Becsüld meg a \(\sin(x)\) kiértékelésének relatív hibáját.
  \item
    Becsüld meg a feladat kondíciószámát.
  \item
    Az \(x\) argumentum mely értékeire nagyon érzékeny ez a feladat?
  \end{enumerate}
\item
  \textbf{1.5.} Tekintsük az \(f: \mathbb{R}^2 \to \mathbb{R}\)
  függvényt, amelyet \(f(x, y) = x - y\) definiál. A bemenet \((x, y)\)
  méretét \(|x| + |y|\)-nal mérve, és feltételezve, hogy
  \(|x| + |y| \approx 1\) és \(x - y \approx \epsilon\), mutasd meg,
  hogy \(\operatorname{cond}(f) \approx 1/\epsilon\). Mit mondhatsz a
  kivonás érzékenységéről?
\item
  \textbf{1.6.} A szinuszfüggvényt a
\end{itemize}

\[\sin(x) = x - \frac{x^3}{3!} + \frac{x^5}{5!} - \frac{x^7}{7!} + \cdots\]

végtelen sor adja meg. (a) Mi az előreható és a hátraható hiba, ha a
szinuszfüggvényt a sor első tagjával közelítjük --- azaz
\(\sin(x) \approx x\) --- \(x = 0{,}1\), \(0{,}5\) és \(1{,}0\) esetén?
(b) Mi az előreható és a hátraható hiba, ha a szinuszfüggvényt a sor
első két tagjával közelítjük --- azaz \(\sin(x) \approx x - x^3/6\) ---
\(x = 0{,}1\), \(0{,}5\) és \(1{,}0\) esetén? - \textbf{1.7.} Egy
lebegőpontos számrendszert négy egész szám jellemez: a \(\beta\) alap, a
\(p\) precízió, valamint az exponens-tartomány \(L\) és \(U\) alsó,
illetve felső határai. (a) Ha \(\beta = 10\), mik \(p\) és \(U\)
legkisebb, illetve \(L\) legnagyobb értékei, hogy a \(2365{,}27\) és a
\(0{,}0000512\) egyaránt pontosan ábrázolható legyen egy normalizált
lebegőpontos rendszerben? (b) Hogyan változna a válaszod, ha a rendszer
nem normalizált, azaz fokozatos alulcsordulást engedélyezünk? -
\textbf{1.8.} Egy \(p = 6\) decimális precíziójú lebegőpontos
rendszerben legyen \(x = 1{,}23456\) és \(y = 1{,}23579\). (a) Hány
értékes jegyet tartalmaz a \(y - x\) különbség? (b) Ha a lebegőpontos
rendszer normalizált, mi a minimális exponens-tartomány, amely esetén
\(x\), \(y\) és \(y - x\) mindegyike pontosan ábrázolható? (c) Az
exponens-tartománytól függetlenül pontosan ábrázolható-e a \(y - x\)
különbség, ha fokozatos alulcsordulás megengedett? Miért? -
\textbf{1.9.} (a) Négy számjegyű decimális aritmetikával és az 1.1.
példában adott képlettel számítsd ki a Föld felszínét \(r = 6370\) km
mellett. (b) Ugyanezzel a képlettel és precízióval számítsd ki a felszín
különbségét, ha a sugár értékét 1 km-rel növeljük. (c) Mivel
\(dA/dr = 8\pi r\), a felszín változása \(8\pi r h\)-val közelíthető,
ahol \(h\) a sugár változása. Ezzel a képlettel --- még mindig négy
számjegyű aritmetikával --- számítsd ki a felszín különbségét a sugár 1
km-es növekedése miatt. Hogyan viszonyul az e közelítő képlettel kapott
érték a (b)-ben az „egzakt" képlettel kapott értékhez? (d) Határozd meg,
melyik az előző két válasz közül a közelebb álló a helyeshez úgy, hogy
mindkét számítást magasabb precízióval --- mondjuk hat számjegyű
decimális aritmetikával --- megismétled. (e) Magyarázd meg az (a)--(d)
pontokban kapott eredményeket. (f) Próbáld ki ezt a feladatot
számítógépen. Milyen kicsi \(h\) sugárváltoztatás esetén lép fel ugyanez
a jelenség? Próbáld ki mind egyszeres, mind dupla pontossággal, ha
rendelkezésre áll. - \textbf{1.10.} Tekintsük az

\[\frac{1}{1-x} - \frac{1}{1+x}\]

kifejezést, feltételezve, hogy \(x \neq \pm 1\). (a) Az \(x\) értékek
mely tartományára nehéz e kifejezést pontosan számítani lebegőpontos
aritmetikában? (b) Add meg a tagok olyan átrendezését, amely az (a)-beli
\(x\)-tartományra lebegőpontos aritmetikában pontosabb. - \textbf{1.11.}
Ha \(x \approx y\), akkor azt várnánk, hogy a \(\log(x) - \log(y)\)
számítás kioltást szenved. Másfelől \(\log(x) - \log(y) = \log(x/y)\),
és az utóbbi nem tartalmaz kioltást. Jelenti-e ez azt, hogy
\(\log(x/y)\) számítása valószínűleg jobb eredményt ad?
(\emph{Útmutatás:} Milyen értékre érzékeny a logaritmusfüggvény?) -
\textbf{1.12.} (a) A matematikailag ekvivalens

\[x^2 - y^2 \quad \text{és} \quad (x - y)(x + y)\]

kifejezések közül melyik értékelhető ki pontosabban lebegőpontos
aritmetikában? Miért? (b) \(x\) és \(y\) mely --- egymáshoz viszonyított
--- értékeire van lényeges különbség a két kifejezés pontosságában? -
\textbf{1.13.} Egy \(n\)-dimenziós \(\boldsymbol{x}\) vektor euklideszi
normáját az

\[\|\boldsymbol{x}\|_2 = \left(\sum_{i=1}^n x_i^2\right)^{1/2}\]

képlet definiálja. Hogyan kerülheted el a túlcsordulást és a káros
alulcsordulást ebben a számításban? - \textbf{1.14.} Egy \([a, b]\)
intervallum \(m\) középpontjának lebegőpontos aritmetikában történő
számításához a következő két képlet közül melyik jobb, és miért? Mikor?
(\emph{Útmutatás:} Találj ki példákat, amelyekben a képlet által adott
„középpont" az \([a, b]\) intervallumon kívülre esik.) (a)
\(m = (a + b)/2{,}0\) (b) \(m = a + (b - a)/2{,}0\) - \textbf{1.15.} Adj
konkrét példákat, amelyek megmutatják, hogy a lebegőpontos összeadás nem
asszociatív, a következő lebegőpontos rendszerek mindegyikében: (a) az
1.9. példa játék lebegőpontos rendszere; (b) IEEE egyszeres pontosságú
lebegőpontos aritmetika. - \textbf{1.16.} Magyarázd meg, hogyan
különbözhetnek a gyakorlatban az \(\epsilon_{\text{mach}}\) egységnyi
kerekítési hiba 1.3.5. szakaszban adott különféle definíciói.
(\emph{Útmutatás:} Tekintsd az 1.9. példa játék lebegőpontos
rendszerét.) - \textbf{1.17.} Legyen \(x\) egy adott nem nulla
lebegőpontos szám egy normalizált rendszerben, és legyen \(y\) egy
szomszédos --- szintén nem nulla --- lebegőpontos szám. (a) Mi a
minimális lehetséges távolság \(x\) és \(y\) között? (b) Mi a maximális
lehetséges távolság \(x\) és \(y\) között? - \textbf{1.18.} Hány
normalizált gépi szám van egy IEEE egyszeres pontosságú lebegőpontos
rendszerben? Hány további gépi számot nyerünk, ha megengedjük a
szubnormálisakat? - \textbf{1.19.} Egy IEEE egyszeres pontosságú
lebegőpontos rendszerben mi a legnagyobb gépi szám \(\text{OFL}\), és a
legkisebb pozitív normalizált gépi szám \(\text{UFL}\) értéke? Hogyan
változnak a válaszaid, ha a szubnormálisak megengedettek? -
\textbf{1.20.} Mi a \(0{,}1\) decimális tört IEEE egyszeres pontosságú
bináris lebegőpontos ábrázolása (a) csonkítással? (b) legközelebbihez
kerekítéssel? - \textbf{1.21.} (a) Egy lebegőpontos rendszerben az
\(\epsilon_{\text{mach}}\) egységnyi kerekítési hiba szükségképpen gépi
szám-e? (b) Létezhet-e olyan lebegőpontos rendszer, amelyben
\(\epsilon_{\text{mach}} < \text{UFL}\)? Ha igen, adj példát. -
\textbf{1.22.} Tegyük fel, hogy az \(ax^2 + bx + c = 0\) másodfokú
egyenletet oldjuk meg \(a = 1{,}22\), \(b = 3{,}34\), \(c = 2{,}28\)
mellett, egy \(\beta = 10\), \(p = 3\) normalizált lebegőpontos
rendszerben. (a) Mi a diszkrimináns \(b^2 - 4ac\) számított értéke? (b)
Mi a diszkrimináns helyes értéke valós (egzakt) aritmetikában? (c) Mi a
diszkrimináns számított értékének relatív hibája? - \textbf{1.23.}
Tegyünk fel egy \(\beta = 10\), \(p = 3\), \(L = -98\) normalizált
lebegőpontos rendszert. (a) Mi az \(\text{UFL}\) alulcsordulási szint
értéke ebben a rendszerben? (b) Ha \(x = 6{,}87 \times 10^{-97}\) és
\(y = 6{,}81 \times 10^{-97}\), mi \(x - y\) eredménye? (c) Mi lenne
\(x - y\) eredménye, ha a rendszer megengedné a fokozatos
alulcsordulást? - \textbf{1.24.} Tekintsük a következő állítást: ha két
azonos előjelű lebegőpontos szám, \(x\) és \(y\) legfeljebb a \(\beta\)
alap szerinti tényezővel tér el (azaz \(1/\beta \le x/y \le \beta\)),
akkor a különbségük, \(x - y\) pontosan ábrázolható a lebegőpontos
rendszerben. Mutasd meg, hogy ez az állítás igaz \(\beta = 2\) esetén,
de adj ellenpéldát \(\beta > 2\) esetére. - \textbf{1.25.} Néhány
mikroprocesszor rendelkezik egy \texttt{mpyadd(a,\ b,\ c)} utasítással
(szorzás-összeadás), amely egyszeres hosszúságú bemeneteket vesz, és
hozzáadja \(c\)-t \(a\) és \(b\) dupla hosszúságú szorzatához, mielőtt
normalizálná és egyszeres hosszúságú eredményt adna vissza. Hogyan lehet
egy ilyen utasítást felhasználni dupla pontosságú szorzatok
kiszámítására úgy, hogy dupla hosszúságú változókat egyáltalán nem
használunk (azaz \(a\) és \(b\) dupla hosszúságú szorzata két egyszeres
hosszúságú változóban --- például \(s\)-ben és \(t\)-ben --- tárolódik)?
- \textbf{1.26.} Igazold, hogy az 1.15. példában adott alternatív
másodfokú megoldóképlet valóban helyes gyököket ad a másodfokú
egyenletre (egzakt aritmetikában). - \textbf{1.27.} Adj részletes
magyarázatot a szórás számításához használt egymenetes képlet numerikus
gyengeségére az 1.16. példában adott kétmenetes képlethez képest.

\subsubsection{Számítógépes
feladatok}\label{szuxe1muxedtuxf3guxe9pes-feladatok}

\begin{itemize}
\tightlist
\item
  \textbf{1.1.} Írj programot, amely kiszámítja Stirling közelítésének
\end{itemize}

\[n! \approx \sqrt{2\pi n}\,\left(n/e\right)^n\]

abszolút és relatív hibáját \(n = 1, \dots, 10\) esetén. Az abszolút
hiba nő vagy csökken \(n\) növekedtével? És a relatív hiba? -
\textbf{1.2.} Írj programot, amely közelítő értékeket határoz meg az
\(\epsilon_{\text{mach}}\) egységnyi kerekítési hibára és az
\(\text{UFL}\) alulcsordulási szintre, és teszteld egy valós
számítógépen. (\emph{Opcionális:} Meg tudod határozni a géped
\(\text{OFL}\) túlcsordulási szintjét is? Ez trükkösebb, mert egy
tényleges túlcsordulás végzetes lehet.) Írasd ki a kapott értékeket
decimálisan, és próbáld megállapítani a használt lebegőpontos rendszer
mantissza- és exponensmezőjének bitszámát. - \textbf{1.3.} A legtöbb
lebegőpontos rendszerben az egységnyi kerekítési hibára gyors közelítés
kapható az

\[\epsilon_{\text{mach}} \approx |3 \cdot (4/3 - 1) - 1|\]

kifejezés kiértékelésével. (a) Magyarázd meg, miért működik ez a trükk.
(b) Próbáld ki különféle számítógépeken (egyszeres és dupla pontossággal
egyaránt) és számológépeken, hogy megerősítsd a működését. (c) Működne-e
ez a trükk egy \(\beta = 3\) alapú lebegőpontos rendszerben? -
\textbf{1.4.} Írj programot, amely a természetes logaritmus \(e\)
matematikai alapját a

\[e = \lim_{n \to \infty} (1 + 1/n)^n\]

definícióból számítja ki. Konkrétan számítsd ki \((1 + 1/n)^n\)-t
\(n = 10^k\)-ra, \(k = 1, 2, \dots, 20\). Ha a használt programnyelvben
nincs hatványozási operátor, használhatod az ekvivalens

\[(1 + 1/n)^n = \exp\bigl(n\,\log(1 + 1/n)\bigr)\]

képletet, ahol az \texttt{exp} és a \texttt{log} beépített függvények.
Határozd meg a közelítéseid hibáját az \texttt{exp(1)} értékkel való
összehasonlítással. A hiba mindig csökken \(n\) növekedtével? Magyarázd
meg az eredményeidet. - \textbf{1.5.} (a) Tekintsük az

\[f(x) = (e^x - 1)/x\]

függvényt. L\textquotesingle Hôpital szabályával mutasd meg, hogy
\(\lim_{x \to 0} f(x) = 1\). (b) Ellenőrizd ezt az eredményt
kísérletileg úgy, hogy írsz egy programot, amely \(f(x)\)-et számítja
\(x = 10^{-k}\)-ra, \(k = 1, \dots, 15\). Az eredményeid egyeznek az
elméleti várakozásokkal? Magyarázd meg, miért. (c) Végezd el a (b)-beli
kísérletet ismét, ezúttal a matematikailag ekvivalens

\[f(x) = (e^x - 1)/\log(e^x)\]

megfogalmazást használva, a jelzett módon, minden egyszerűsítés nélkül.
Ha ez jobban működik, meg tudod magyarázni, miért? - \textbf{1.6.}
Tegyük fel, hogy \(n + 1\) egyenletesen elhelyezett pontot kell
generálnunk az \([a, b]\) intervallumon, \(h = (b - a)/n\) lépésközzel.
(a) Lebegőpontos aritmetikában a következő két módszer közül melyik
jobb, és miért?

\[x_0 = a, \quad x_k = x_{k-1} + h, \quad k = 1, \dots, n\]

vagy

\[x_k = a + kh, \quad k = 0, \dots, n.\]

\begin{enumerate}
\def\labelenumi{(\alph{enumi})}
\setcounter{enumi}{1}
\tightlist
\item
  Írj programot, amely mindkét módszert implementálja, és találj egy
  példát --- mondjuk \(a = 0\) és \(b = 1\) mellett ---, amely
  szemlélteti a különbséget közöttük.
\end{enumerate}

\begin{itemize}
\tightlist
\item
  \textbf{1.7.} (a) Írj programot, amely egy függvény deriváltjának
  közelítő értékét a
\end{itemize}

\[f'(x) \approx \frac{f(x+h) - f(x)}{h}\]

véges differencia képlettel számítja. Teszteld a programot a \(\tan(x)\)
függvényen \(x = 1\)-nél. Határozd meg a hibát a beépített \(\sec(x)\)
függvény négyzetével való összehasonlítással. Ábrázold a hiba nagyságát
\(h\) függvényében \(h = 10^{-k}\)-ra, \(k = 0, \dots, 16\). Mind
\(h\)-t, mind a hiba nagyságát logaritmikus skálán érdemes használni.
Van a hiba nagyságának minimuma? Hogyan viszonyul a hozzá tartozó \(h\)
érték az 1.3. példában levezetett
\(h \approx \sqrt{\epsilon_{\text{mach}}}\) hüvelykujj-szabályhoz? (b)
Ismételd meg a feladatot a

\[f'(x) \approx \frac{f(x+h) - f(x-h)}{2h}\]

középpontos differencia közelítéssel. - \textbf{1.8.} Tekintsük a

\[\sum_{n=1}^{\infty} \frac{1}{n}\]

végtelen sort. (a) Bizonyítsd be, hogy a sor divergens.
(\emph{Útmutatás:} Csoportosítsd a tagokat olyan halmazokba, amelyek az
\(1/(2^{k-1} + 1)\)-től \(1/2^k\)-ig tartó tagokat tartalmazzák,
\(k = 1, 2, \dots\)) (b) Magyarázd meg, miért ad véges összeget a sor
lebegőpontos aritmetikában való összegzése. (c) Próbáld megjósolni,
mikor szűnik meg változni a részletösszeg IEEE egyszeres, illetve dupla
pontosságú lebegőpontos aritmetikában. A géped lebegőpontos műveletekre
vonatkozó végrehajtási sebességét figyelembe véve próbáld megjósolni,
mennyi ideig tartana minden egyes számítás befejezése. (d) Írj két
programot a sor összegének kiszámítására: egyet egyszeres, egyet dupla
pontossággal. Kövesd az összegzés előrehaladását úgy, hogy az indexet és
a részletösszeget időnként kiírod. Milyen megállási kritériumot
használnál? Milyen eredményt ad ténylegesen a géped? Hasonlítsd össze az
eredményeidet a jóslataiddal, beleértve a szükséges futásidőt is.
(\emph{Figyelem:} Az egyszeres pontosságú változat várhatóan viszonylag
gyorsan befejeződik, a dupla pontosságú azonban sokkal tovább tarthat,
így lehet, hogy nem gyakorlati befejezésig futtatni, még bőkezű
számítási kerettel sem.) - \textbf{1.9.} (a) Írj programot, amely az
\(e^x\) exponenciális függvényt az

\[e^x = 1 + x + \frac{x^2}{2!} + \frac{x^3}{3!} + \cdots\]

végtelen sorral számítja. (b) Természetes sorrendben összegezve milyen
megállási kritériumot használnál? (c) Teszteld a programot
\(x = \pm 1, \pm 5, \pm 10, \pm 15, \pm 20\) értékekre, és hasonlítsd
össze az eredményeket a beépített \texttt{exp(x)} függvénnyel. (d)
Használhatod-e a sort ebben a formában pontos eredmények eléréséhez
\(x < 0\) esetén? (\emph{Útmutatás:} \(e^{-x} = 1/e^x\).) (e) Át tudod-e
rendezni a sort, vagy át tudod-e csoportosítani a tagokat oly módon,
hogy pontosabb eredményeket kapj \(x < 0\) esetén? - \textbf{1.10.} Írj
programot, amely az \(ax^2 + bx + c = 0\) másodfokú egyenletet oldja meg
a szabványos

\[x = \frac{-b \pm \sqrt{b^2 - 4ac}}{2a}\]

képlettel vagy a

\[x = \frac{2c}{-b \mp \sqrt{b^2 - 4ac}}\]

alternatív képlettel. A programnak bemenetként \(a\), \(b\), \(c\)
együtthatókat kell elfogadnia, kimenetként pedig az egyenlet két gyökét
kell előállítania. A programnak észlelnie kell, ha a gyökök nem valósak,
de nem szükséges explicit módon komplex aritmetikát használnia (például
ebben az esetben visszaadhatja a komplex konjugált gyökpár valós és
képzetes részét). Védekezz a felesleges túlcsordulás, alulcsordulás és
kioltás ellen. Mikor kell használnod a két képlet közül melyiket? Tedd
robusztussá a programot a szokatlan bemeneti értékekre is, például
\(a = 0\) vagy \(c = 0\), amelyek egyébként az egyik képletet
megbuktatnák. Bármely gyökhöz, amely a lebegőpontos rendszer tartományán
belül van, pontosan számítsa ki az értékét, még akkor is, ha a másik
gyök kívül esik. Teszteld a programot a következő együtthatóértékekkel:

\begin{longtable}[]{@{}lll@{}}
\toprule\noalign{}
\(a\) & \(b\) & \(c\) \\
\midrule\noalign{}
\endhead
\bottomrule\noalign{}
\endlastfoot
\(6\) & \(5\) & \(-4\) \\
\(6 \times 10^{154}\) & \(5 \times 10^{154}\) &
\(-4 \times 10^{154}\) \\
\(0\) & \(1\) & \(1\) \\
\(1\) & \(-10^5\) & \(1\) \\
\(1\) & \(-4\) & \(3{,}999999\) \\
\(10^{-155}\) & \(-10^{155}\) & \(10^{155}\) \\
\end{longtable}

\begin{itemize}
\tightlist
\item
  \textbf{1.11.} A nem nulla vezető együtthatóval való osztással bármely
  harmadfokú egyenlet az
\end{itemize}

\[x^3 + ax^2 + bx + c = 0\]

alakra hozható. Ha az együtthatók valósak, akkor egy ilyen egyenletnek
legalább egy valós gyöke van, amely zárt alakban a következőképpen
számítható ki. Alkalmazd az \(y = x + a/3\) helyettesítést. Az egyenlet
ekkor a

\[y^3 + py + q = 0\]

alakra redukálódik, ahol

\[p = \frac{3b - a^2}{3}, \qquad q = \frac{2a^3}{27} - \frac{ab}{3} + c.\]

Ha a \(D = (p/3)^3 + (q/2)^2\) diszkrimináns pozitív, akkor csak egy
valós gyök létezik, amelyet az

\[x = -\frac{a}{3} + \sqrt[3]{-\frac{q}{2} + \sqrt{D}} + \sqrt[3]{-\frac{q}{2} - \sqrt{D}}\]

képlet ad meg. Írj valós aritmetikát használó rutint ezzel a módszerrel
(ügyelve arra, hogy az alkalmazott köbgyök valós legyen) egy valós
együtthatós harmadfokú egyenlet valós gyökének kiszámítására abban az
esetben, amikor csak egy gyök valós. Ha különösen ambiciózus vagy,
kiszámíthatod a másik két gyököt is erre az esetre, vagy foglalkozhatsz
a három valós gyök esetével, de ehhez komplex aritmetikára (vagy
bonyolultabb trigonometrikus megoldásra) lesz szükség. Tedd a rutint a
lehető legrobusztusabbá, védekezve a felesleges túlcsordulás,
alulcsordulás és kioltás ellen. Mit kell tennie a rutinnak, ha a vezető
együttható nulla? Teszteld a rutint különféle együtthatóértékekre, az
előző feladathoz hasonlóan. Lásd még az 5.17. számítógépes feladatot. -
\textbf{1.12.} (a) Írj programot, amely egy véges \(x_i\) sorozat
\(\bar{x}\) átlagát és \(\sigma\) szórását számítja. A program
bemenetként egy \(n\) dimenziós \(\boldsymbol{x}\) vektort fogadjon el,
kimenetként pedig a sorozat átlagát és szórását adja meg. A szórásra
próbáld ki mind a kétmenetes

\[\sigma = \left[\frac{1}{n-1} \sum_{i=1}^{n} (x_i - \bar{x})^2\right]^{1/2}\]

képletet, mind az egymenetes

\[\sigma = \left[\frac{1}{n-1} \left(\sum_{i=1}^{n} x_i^2 - n\bar{x}^2\right)\right]^{1/2}\]

képletet, és hasonlítsd össze az eredményeket egy általad választott
bemeneti sorozaton. (b) Ki tudsz találni olyan bemeneti adatsorozatot,
amely drámaian szemlélteti a két matematikailag ekvivalens képlet
numerikus különbségét? (\emph{Figyelem:} Óvakodj egy negatív szám
négyzetgyökének vételétől.) - \textbf{1.13.} Ha egy \(a\) összeget \(r\)
kamatlábbal fektetnek be, és a kamatot évente \(n\)-szer tőkésítik,
akkor az egy év végén kapott végérték

\[f = a(1 + r/n)^n.\]

Ez a kamatos kamat ismerős képlete. Egyszerű kamat esetén \(n = 1\).
Jellemzően negyedévente (\(n = 4\)), vagy akár naponta (\(n = 365\))
tőkésítenek. Nyilvánvaló, hogy minél gyakoribb a tőkésítés, annál
nagyobb a végösszeg, hiszen a korábbi kamatra is kamatot kapunk. De
mekkora különbséget okoz valójában ez a gyakoriság? Írj programot, amely
a kamatos kamat képletet implementálja. Teszteld a programot \(a = 100\)
kezdeti befektetéssel, \(r = 0{,}05\) (5 százalékos) kamatlábbal, és a
következő \(n\) értékekkel: 1, 4, 12, 365. Kísérletezz azzal is, mi
történik, ha \(n\) nagyon nagy lesz. Találsz-e olyan értéket, amelyre a
végösszeg nem nő a tőkésítés gyakoriságával, ahogy kellene? (Ez könnyű,
ha egyszeres pontosságot használsz, de jóval nehezebb, ha duplát.)

Implementáld a kamatos kamat képletet két különböző módon: (a) Ha a
használt programnyelvben nincs hatványozási operátor (pl. C), akkor a
képletet egy olyan ciklussal implementálhatod, amely \(a\)-t ismételten
szorozza \((1 + r/n)\)-nel, összesen \(n\)-szer. Még ha a programnyelved
rendelkezik is hatványozási operátorral (pl. Fortran), implementáld a
képletet egy ilyen ciklussal, és írd ki az eredményeket a bemeneti
értékekre. (b) Az \texttt{exp(x)} és \texttt{log(x)} függvényekkel a
kamatos kamat képlet a

\[f = a \exp\bigl(n\,\log(1 + r/n)\bigr)\]

alakban is írható. Implementáld ezt a képletet a megfelelő beépített
függvények használatával, és hasonlítsd össze az eredményeidet az első
--- ciklust használó --- implementáció eredményeivel ugyanazon bemeneti
értékekre. - \textbf{1.14.} Az \((x - 1)^6\) polinom \(x = 1\)-nél
nulla, máshol pozitív. A polinom kifejtett alakja,
\(x^6 - 6x^5 + 15x^4 - 20x^3 + 15x^2 - 6x + 1\), matematikailag
ekvivalens, de numerikusan nem feltétlenül adja ugyanazt az eredményt.
Számítsd ki és ábrázold e polinom értékeit mindkét forma használatával a
\([0{,}995, 1{,}005]\) intervallum \(101\) egyenletesen elhelyezett
pontjában --- azaz \(0{,}0001\) lépésközzel. A diagramot úgy skálázd,
hogy mind az \(x\), mind a polinom értékei a saját tengelyük teljes
tartományát kihasználják. Meg tudod magyarázni ezt a viselkedést? -
\textbf{1.15.} Egy \(n\)-dimenziós \(\boldsymbol{x}\) vektor euklideszi
normáját az

\[\|\boldsymbol{x}\|_2 = \left(\sum_{i=1}^n x_i^2\right)^{1/2}\]

képlet definiálja. Implementálj robusztus rutint e mennyiség
kiszámítására egy tetszőleges bemeneti \(\boldsymbol{x}\) vektorra. A
rutinnak el kell kerülnie a túlcsordulást és a káros alulcsordulást.
Hasonlítsd össze a robusztus rutin pontosságát és teljesítményét egy
egyszerűbb, naiv implementációval. Ki tudsz találni olyan vektort, amely
jelentősen eltérő eredményeket ad a két rutintól? Mennyi teljesítményt
áldoz fel a robusztus rutin? - \textbf{1.16.} Írj programot, amely \(n\)
véletlen, egyszeres pontosságú lebegőpontos számot --- \(x_i\), a
\([0, 1]\) intervallumon egyenletes eloszlású --- összegez (lásd a 13.1.
táblázatot egy megfelelő véletlenszám-generátorhoz). A számokat az
alábbi módokon összegezd (csak egyszeres pontosságú lebegőpontos
változókat használj, hacsak másképp nem jelezzük): (a) Összegezd a
számokat a generálás sorrendjében, egy dupla pontosságú változót
használva az összeg gyűjtésére. (b) Összegezd a számokat a generálás
sorrendjében, ezúttal egy egyszeres pontosságú gyűjtőt használva. (c)
Használd a következő (Kahan-féle) kompenzált összegzési algoritmust ---
szintén csak egyszeres pontosságot használva ---, a számokat a generálás
sorrendjében összegezve:

\[\begin{aligned}
s &= x_1 \\
c &= 0 \\
\textbf{for } &i = 2 \text{ to } n \\
y &= x_i - c \\
t &= s + y \\
c &= (t - s) - y \\
s &= t \\
\textbf{end}&
\end{aligned}\]

\begin{enumerate}
\def\labelenumi{(\alph{enumi})}
\setcounter{enumi}{3}
\tightlist
\item
  Összegezd a számokat növekvő abszolút érték szerinti sorrendben (ehhez
  a számokat összegzés előtt rendezni kell, amihez használhatsz
  könyvtári rendezőrutint).
\item
  Összegezd a számokat csökkenő abszolút érték szerinti sorrendben (azaz
  az összegzés sorrendjét fordítsd meg a (d)-beli sorrendhez képest).
\end{enumerate}

Futtasd a programot különböző \(n\) értékekre, és hasonlítsd össze a
(a)--(e) módszerek eredményeit. Lényeges eltérés megfigyeléséhez elég
nagy \(n\) értéket érdemes használnod. Hogyan rangsorolhatók a módszerek
pontosság szerint, és miért? Hogyan viszonyulnak egymáshoz a módszerek
költségben? Meg tudod magyarázni, miért működik a (c)-beli algoritmus? -
\textbf{1.17.} Írj programot, amely az

\[x_{k+1} = 2{,}25\,x_k - 0{,}5\,x_{k-1}\]

differenciaegyenlet első \(n\) tagját állítja elő, kezdőértékek

\[x_1 = \frac{1}{3} \quad \text{és} \quad x_2 = \frac{1}{12}.\]

Használd \(n = 25\)-öt egyszeres, vagy \(n = 60\)-at dupla pontosság
esetén. Készíts féllogaritmikus ábrát a kapott értékekről \(k\)
függvényében. A differenciaegyenlet egzakt megoldása

\[x_k = \frac{4^{1-k}}{3},\]

amely \(k\) növekedtével szigorúan monoton csökken. Igazolja-e az ábrád
ezt az elméletileg várt viselkedést? Meg tudod magyarázni az
eredményeidet? (\emph{Útmutatás:} Keresd meg a differenciaegyenlet
általános megoldását.) - \textbf{1.18.} Írj programot, amely az

\[x_{k+1} = 111 - (1130 - 3000/x_{k-1})/x_k\]

differenciaegyenlet első \(n\) tagját állítja elő, kezdőértékek

\[x_1 = \frac{11}{2} \quad \text{és} \quad x_2 = \frac{61}{11}.\]

Használd \(n = 10\)-et egyszeres, vagy \(n = 20\)-at dupla pontosság
esetén. Az egzakt megoldás egy szigorúan monoton növekvő, \(6\)-hoz
konvergáló sorozat. Meg tudod magyarázni az eredményeidet?

\section{Lineáris
egyenletrendszerek}\label{lineuxe1ris-egyenletrendszerek}

\subsection{2.1 Lineáris rendszerek}\label{lineuxe1ris-rendszerek}

A természetben sok összefüggés lineáris, vagyis a hatások az okaikkal
arányosak. A mechanikában például Newton második törvénye, \(F = ma\),
azt mondja ki, hogy az erő arányos a gyorsulással, ahol az arányossági
tényező a tárgy tömege. Ha ismerjük az erőt és a tömeget, akkor ezt a
lineáris egyenletet a gyorsulásra meg tudjuk oldani. Az
elektromosságtanban Ohm törvénye, \(V = iR\), kimondja, hogy egy vezető
két vége közötti feszültség (potenciálkülönbség) arányos a rajta átfolyó
árammal, ahol az arányossági tényező az ellenállás. Ismét, ha ismerjük a
feszültséget és az ellenállást, meg tudjuk oldani az egyenletet az
áramra. A rugalmasságtanban Hooke törvénye szerint a feszültség arányos
a deformációval, ahol az arányossági tényező a Young-modulus.

Magasabb dimenziókban az ilyen lineáris kapcsolatot egy \(\mathcal{L}\)
lineáris transzformáció fejezi ki, amely egy „okokat" tartalmazó
\(\boldsymbol{u}\) vektort egy „hatásokat" tartalmazó \(\boldsymbol{f}\)
vektorhoz rendel:

\[\mathcal{L}\boldsymbol{u} = \boldsymbol{f}.\]

Ohm törvénye és Kirchhoff törvényei például lehetővé teszik, hogy
levezessünk egy egész egyenletrendszert, amely egy sok vezetőből álló
áramkörben a feszültségek és az áramok közötti lineáris kapcsolatot
fejezi ki (lásd a 2.1. példát). Ahogy az elemi lineáris algebrából
megtanultuk, két véges dimenziós vektortér közötti lineáris
transzformációt kényelmesen egy mátrix ábrázol. Mátrix-vektor jelöléssel
egy lineáris algebrai egyenletrendszer alakja

\[A\boldsymbol{x} = \boldsymbol{b},\]

ahol \(A\) egy ismert \(m \times n\)-es mátrix, \(\boldsymbol{b}\) egy
\(m\)-dimenziós vektor, \(\boldsymbol{x}\) pedig egy \(n\)-dimenziós
vektor. Ha ismerjük \(\boldsymbol{x}\)-et, akkor egy ilyen lineáris
kapcsolat lehetővé teszi, hogy mátrix-vektor szorzással előre jelezzük a
\(\boldsymbol{b}\) hatást az \(\boldsymbol{x}\) okból:
\(\boldsymbol{b} = A\boldsymbol{x}\). Érdekesebb módon a lineáris
rendszer lehetőséget adhat a „visszafejtésre": ha ismerjük a hatások
\(\boldsymbol{b}\) vektorát, szeretnénk meghatározni a hozzá tartozó
okok \(\boldsymbol{x}\) vektorát. E fejezet tárgya az ezt lehetővé tevő
numerikus módszerek.

Még amikor az összefüggések nemlineárisak, gyakran jól közelíthetők
lokálisan (azaz egy adott rögzített érték környezetében) egy lineáris
kapcsolattal. Pontosan ezt teszi a kalkulusban a derivált: egy
nemlineáris görbéhez helyi lineáris közelítést --- az érintőegyenest ---
ad. Ez a megfigyelés számos nemlineáris algebrai feladat numerikus
megoldási módszerének alapja. Még a nem algebrai feladatokat is ---
például differenciál- vagy integrálegyenleteket --- végső soron lineáris
algebrai egyenletrendszerekkel közelíthetjük, ahogyan azt a könyv
későbbi részeiben látni fogjuk. Ezért a lineáris egyenletrendszerek
megoldása számos, sokféle gyakorlati számítási feladat megoldására
szolgáló numerikus módszer sarokköve, és ennek megfelelően
elengedhetetlen, hogy az ilyen rendszereket pontosan és hatékonyan
tudjuk megoldani.

Az \(A\boldsymbol{x} = \boldsymbol{b}\) lineáris egyenletrendszer azt a
kérdést teszi fel: „kifejezhető-e a \(\boldsymbol{b}\) vektor az \(A\)
mátrix oszlopainak lineáris kombinációjaként?" (azaz „benne van-e
\(\boldsymbol{b}\) az
\(\operatorname{span}(A) = \{A\boldsymbol{x} : \boldsymbol{x} \in \mathbb{R}^n\}\)
halmazban?"). Ha igen, akkor a rendszert \emph{konzisztensnek} mondjuk,
és ennek a lineáris kombinációnak az együtthatóit az \(\boldsymbol{x}\)
megoldásvektor komponensei adják. Lehet, hogy létezik megoldás, és
lehet, hogy nem; ha van megoldás, az lehet egyértelmű, és lehet, hogy
nem. Ebben a fejezetben csak \emph{négyzetes} rendszereket tekintünk,
azaz \(m = n\) --- a mátrix sorainak és oszlopainak száma megegyezik,
vagyis az egyenletek száma (\(A\) és \(\boldsymbol{b}\) sorai)
megegyezik az ismeretlenek számával (\(\boldsymbol{x}\) komponensei). Az
\(m \neq n\) esetet a 3. fejezetben tárgyaljuk. Az egyszerűség kedvéért
figyelmünket elsősorban valós lineáris rendszerekre fordítjuk; a komplex
rendszereket hasonlóan lehet kezelni.

\textbf{2.1. Példa. Elektromos áramkör.} Tekintsük a 2.1. ábrán látható
elektromos áramkört. Adottak a \(V\) feszültségek és az \(R\)
ellenállások, és azt szeretnénk meghatározni, milyen \(i\) áramok
folynak az áramkör három hurkában. Ehhez az alábbi fizikai törvényeket
alkalmazzuk:

\begin{itemize}
\tightlist
\item
  Ohm törvénye: egy \(R\) ellenálláson egy \(i\) áram irányában eső
  feszültségesés \(iR\).
\item
  Kirchhoff-törvény: zárt hurokban a nettó feszültségesés nulla.
\end{itemize}

Ezeket a törvényeket az áramkör minden hurkára alkalmazva a következő
három lineáris egyenletből álló rendszert kapjuk:

\[\begin{aligned}
i_1 R_1 + (i_1 - i_2) R_2 + (i_1 - i_3) R_3 + V_1 &= 0, \\
(i_2 - i_1) R_2 + (i_2 - i_3) R_5 - V_2 &= 0, \\
(i_3 - i_1) R_3 + i_3 R_4 + (i_3 - i_2) R_5 &= 0,
\end{aligned}\]

amely mátrix alakban így írható:

\[\begin{bmatrix} R_1 + R_2 + R_3 & -R_2 & -R_3 \\ -R_2 & R_2 + R_5 & -R_5 \\ -R_3 & -R_5 & R_3 + R_4 + R_5 \end{bmatrix} \begin{bmatrix} i_1 \\ i_2 \\ i_3 \end{bmatrix} = \begin{bmatrix} -V_1 \\ V_2 \\ 0 \end{bmatrix}.\]

Így ez a feladat \(A\boldsymbol{x} = \boldsymbol{b}\) alakú, és
megoldható az e fejezetben tanulmányozandó módszerekkel.

\pandocbounded{\includegraphics[keepaspectratio]{_page_2_Figure_3.jpeg}}

2.1. ábra: Ellenállásokkal és feszültségforrásokkal felépített
elektromos áramkör.

\subsection{2.2 Létezés és
egyértelműség}\label{luxe9tezuxe9s-uxe9s-egyuxe9rtelmux171suxe9g}

Egy \(n \times n\)-es \(\boldsymbol{A}\) mátrixot
\emph{nemszingulárisnak} nevezünk, ha teljesül rá az alábbi ---
egymással ekvivalens --- feltételek bármelyike:

\begin{enumerate}
\def\labelenumi{\arabic{enumi}.}
\tightlist
\item
  \(\boldsymbol{A}\)-nak van inverze (azaz létezik olyan
  \(n \times n\)-es mátrix, amelyet \(\boldsymbol{A}^{-1}\)-gyel
  jelölünk, amelyre
  \(\boldsymbol{A}\boldsymbol{A}^{-1} = \boldsymbol{A}^{-1}\boldsymbol{A} = \boldsymbol{I}\),
  az egységmátrix).
\item
  \(\det(\boldsymbol{A}) \neq 0\) (azaz \(\boldsymbol{A}\) determinánsa
  nem nulla).
\item
  \(\operatorname{rank}(\boldsymbol{A}) = n\) (egy mátrix rangja a benne
  található lineárisan független sorok vagy oszlopok legnagyobb száma).
\item
  Bármely \(\boldsymbol{z} \neq \boldsymbol{0}\) vektorra
  \(\boldsymbol{A}\boldsymbol{z} \neq \boldsymbol{0}\) (azaz
  \(\boldsymbol{A}\) egyetlen nemtriviális vektort sem annullál).
\end{enumerate}

Egyébként a mátrix \emph{szinguláris}. Az
\(A\boldsymbol{x} = \boldsymbol{b}\) lineáris egyenletrendszer
megoldásának létezése és egyértelműsége attól függ, hogy az
\(\boldsymbol{A}\) mátrix szinguláris vagy nemszinguláris. Ha
\(\boldsymbol{A}\) nemszinguláris, akkor létezik az
\(\boldsymbol{A}^{-1}\) inverze, és az
\(A\boldsymbol{x} = \boldsymbol{b}\) rendszernek mindig létezik egyetlen
\(\boldsymbol{x} = \boldsymbol{A}^{-1}\boldsymbol{b}\) megoldása,
\(\boldsymbol{b}\) értékétől függetlenül. Ha viszont \(\boldsymbol{A}\)
szinguláris, akkor a megoldások száma a \(\boldsymbol{b}\) jobb oldali
vektortól függ: adott \(\boldsymbol{b}\) mellett lehet, hogy nincs
megoldás, de ha létezik olyan \(\boldsymbol{x}\) megoldás, amelyre
\(A\boldsymbol{x} = \boldsymbol{b}\), akkor
\(A(\boldsymbol{x} + \gamma \boldsymbol{z}) = \boldsymbol{b}\) is
teljesül bármely \(\gamma\) skalárra, ahol
\(\boldsymbol{z} \neq \boldsymbol{0}\) olyan vektor, amelyre
\(A\boldsymbol{z} = \boldsymbol{0}\) (ilyen \(\boldsymbol{z}\)-nek
léteznie kell, hiszen ellenkező esetben a definíció 4. feltétele szerint
a mátrix nemszinguláris volna). Ezért egy négyzetes, konzisztens,
szinguláris lineáris rendszer megoldása nem lehet egyértelmű. Adott
négyzetes \(\boldsymbol{A}\) mátrixra és \(\boldsymbol{b}\) jobb oldali
vektorra a lehetőségeket az alábbiak foglalják össze:

\begin{itemize}
\tightlist
\item
  \textbf{Egyértelmű megoldás:} \(\boldsymbol{A}\) nemszinguláris,
  \(\boldsymbol{b}\) tetszőleges.
\item
  \textbf{Végtelen sok megoldás:} \(\boldsymbol{A}\) szinguláris,
  \(\boldsymbol{b} \in \operatorname{span}(\boldsymbol{A})\)
  (konzisztens).
\item
  \textbf{Nincs megoldás:} \(\boldsymbol{A}\) szinguláris,
  \(\boldsymbol{b} \notin \operatorname{span}(\boldsymbol{A})\)
  (inkonzisztens).
\end{itemize}

Két dimenzióban a rendszer minden lineáris egyenlete egy egyenest
határoz meg a síkban. A rendszer megoldása a két egyenes metszéspontja.
Ha a két egyenes nem párhuzamos, akkor van egyetlen metszéspontjuk (a
nemszinguláris eset). Ha a két egyenes párhuzamos, akkor vagy egyáltalán
nem metszik egymást (nincs megoldás), vagy a két egyenes egybeesik (az
egyenes bármely pontja megoldás). Magasabb dimenziókban minden egyenlet
egy hipersíkot határoz meg. A nemszinguláris esetben az egyértelmű
megoldás az összes hipersík metszéspontja.

\textbf{2.2. Példa. Szingularitás és nemszingularitás.} A
\(2 \times 2\)-es

\[\begin{aligned}
2x_1 + 3x_2 &= b_1, \\
5x_1 + 4x_2 &= b_2,
\end{aligned}\]

rendszer, vagy mátrix-vektor jelöléssel

\[\boldsymbol{A}\boldsymbol{x} = \begin{bmatrix} 2 & 3 \\ 5 & 4 \end{bmatrix} \begin{bmatrix} x_1 \\ x_2 \end{bmatrix} = \begin{bmatrix} b_1 \\ b_2 \end{bmatrix} = \boldsymbol{b},\]

\(\boldsymbol{b}\) értékétől függetlenül egyértelmű megoldással
rendelkezik, mert \(\boldsymbol{A}\) nemszinguláris. Például ha
\(\boldsymbol{b} = \begin{bmatrix} 8 & 13 \end{bmatrix}^T\), akkor az
egyértelmű megoldás
\(\boldsymbol{x} = \begin{bmatrix} 1 & 2 \end{bmatrix}^T\).

Viszont a \(2 \times 2\)-es

\[\boldsymbol{A}\boldsymbol{x} = \begin{bmatrix} 2 & 3 \\ 4 & 6 \end{bmatrix} \begin{bmatrix} x_1 \\ x_2 \end{bmatrix} = \begin{bmatrix} b_1 \\ b_2 \end{bmatrix} = \boldsymbol{b}\]

rendszerben az \(\boldsymbol{A}\) mátrix szinguláris, ezért
\(\boldsymbol{b}\) konkrét értékétől függően lehet, hogy van megoldás,
és lehet, hogy nincs, és egyértelmű megoldása semmiképpen nem lehet.
Például ha \(\boldsymbol{b} = \begin{bmatrix} 4 & 7 \end{bmatrix}^T\),
akkor nincs megoldás, ha viszont
\(\boldsymbol{b} = \begin{bmatrix} 4 & 8 \end{bmatrix}^T\), akkor

\[\boldsymbol{x} = \begin{bmatrix} \gamma \\ (4 - 2\gamma)/3 \end{bmatrix}\]

megoldás bármely \(\gamma\) valós számra.

\subsection{2.3 Érzékenység és
kondicionáltság}\label{uxe9rzuxe9kenysuxe9g-uxe9s-kondicionuxe1ltsuxe1g-1}

Miután megfogalmaztuk az \(Ax = b\) lineáris egyenletrendszer
megoldásának létezésére és egyértelműségére vonatkozó feltételeket, most
az \(x\) megoldás érzékenységét vizsgáljuk a bemeneti adatok
perturbációival szemben, amelyek ebben a feladatban az \(A\) mátrix és a
\(b\) jobb oldali vektor. Az ilyen perturbációk méréséhez szükségünk van
a vektorok és mátrixok „méretének" valamilyen fogalmára. A nagyság,
abszolút érték vagy modulus skaláris fogalma általánosítható a vektorok
és mátrixok normáinak fogalmára.

\subsubsection{2.3.1 Vektornormák}\label{vektornormuxe1k}

Bár általánosabb definíció is lehetséges, a továbbiakban használt
vektornormák mind a \(p\)-normák speciális esetei, amelyeket egy
\(p > 0\) egész szám és egy \(n\)-dimenziós \(x\) vektor esetén a
következőképpen definiálunk:

\[\|\boldsymbol{x}\|_p = \left(\sum_{i=1}^n |x_i|^p\right)^{1/p}.\]

Fontos speciális esetek:

\begin{itemize}
\tightlist
\item
  1-norma:
\end{itemize}

\[\|x\|_1 = \sum_{i=1}^n |x_i|,\]

amelyet néha Manhattan-normának is neveznek, mivel két dimenzióban ez
két pont „háztömbökben" mért távolságának felel meg.

\begin{itemize}
\tightlist
\item
  2-norma:
\end{itemize}

\[\|\boldsymbol{x}\|_2 = \left(\sum_{i=1}^n |x_i|^2\right)^{1/2},\]

amely az euklideszi térben megszokott távolságfogalomnak felel meg,
ezért euklideszi normának is nevezzük.

\begin{itemize}
\tightlist
\item
  \(\infty\)-norma:
\end{itemize}

\[\|\boldsymbol{x}\|_{\infty} = \max_{1 \le i \le n} |x_i|,\]

amely \(p \to \infty\) határesetként értelmezhető.

Ezek a normák kvalitatívan mind hasonló eredményeket adnak, de bizonyos
körülmények között egy adott norma a legkönnyebben kezelhető
analitikusan vagy számítástechnikailag. A lineáris egyenletrendszerek
megoldásainak érzékenységi vizsgálatához általában az 1-normát vagy az
\(\infty\)-normát használjuk. A 2-normát később más kontextusokban
fogjuk kiterjedten alkalmazni. E normák közötti különbségeket
\(\mathbb{R}^2\)-ben a 2.2. ábra szemlélteti, amely az egységgömböt,
azaz az \(\{x : \|x\|_p = 1\}\) halmazt mutatja \(p = 1, 2, \infty\)
esetén. (Megjegyzendő, hogy az egységgömb -- amelyet két dimenzióban
pontosabb egységkörnek nevezni -- csak a 2-normában valóban „kerek",
innen kapta a nevét.) Egy vektor normája egyszerűen az a tényező,
amellyel a megfelelő egységgömböt ki kell tágítani vagy össze kell
zsugorítani ahhoz, hogy az a vektort magában foglalja.

\textbf{2.3. Példa. Vektornormák.} A 2.2. ábrán látható
\(x = [-1{,}6, 1{,}2]^T\) vektor esetén

\[\|\boldsymbol{x}\|_1 = 2{,}8, \quad \|\boldsymbol{x}\|_2 = 2{,}0, \quad \|\boldsymbol{x}\|_{\infty} = 1{,}6.\]

\pandocbounded{\includegraphics[keepaspectratio]{_page_5_Figure_2.jpeg}}

2.2. ábra: Egységgömbök különböző vektornormákban.

Általában tetszőleges \(n\)-dimenziós \(\boldsymbol{x}\) vektor esetén

\[\|x\|_1 \ge \|x\|_2 \ge \|x\|_{\infty}.\]

Másrészt fennáll, hogy

\[\|x\|_1 \le \sqrt{n} \|x\|_2\] ,
\(\|x\|_2 \le \sqrt{n} \|x\|_{\infty}\) és
\(\|x\|_1 \le n \|x\|_{\infty}\).

Így adott \(n\) esetén e normák közül bármely kettő legfeljebb egy
konstansban tér el egymástól, tehát abban az értelemben mind
ekvivalensek, hogy ha az egyik kicsi, akkor valamennyinek arányosan
kicsinek kell lennie (valójában minden \(p\)-norma ekvivalens ebben az
értelemben). Ezért mindig azt a normát választhatjuk, amelyik az adott
kontextusban a legkényelmesebb. A könyv hátralévő részében megfelelő
alsó index jelöli majd a konkrét normát, amikor ez szükséges, de az alsó
indexet elhagyjuk, ha közömbös, hogy melyik normát használjuk.

Minden vektor \(p\)-norma rendelkezik a következő fontos
tulajdonságokkal, ahol \(\boldsymbol{x}\) és \(\boldsymbol{y}\)
tetszőleges vektorok:

\begin{itemize}
\tightlist
\item
  \begin{enumerate}
  \def\labelenumi{\arabic{enumi}.}
  \tightlist
  \item
    \(\|\mathbf{x}\| > 0\), ha \(\mathbf{x} \neq \mathbf{0}\).
  \end{enumerate}
\item
  \begin{enumerate}
  \def\labelenumi{\arabic{enumi}.}
  \setcounter{enumi}{1}
  \tightlist
  \item
    \(\|\gamma \boldsymbol{x}\| = |\gamma| \cdot \|\boldsymbol{x}\|\)
    tetszőleges \(\gamma\) skalár esetén.
  \end{enumerate}
\item
  \begin{enumerate}
  \def\labelenumi{\arabic{enumi}.}
  \setcounter{enumi}{2}
  \tightlist
  \item
    \(\|\boldsymbol{x} + \boldsymbol{y}\| \le \|\boldsymbol{x}\| + \|\boldsymbol{y}\|\)
    (háromszög-egyenlőtlenség).
  \end{enumerate}
\end{itemize}

Az általánosabb megközelítésben egy vektornorma \emph{definíciójaként}
bármely olyan valós értékű függvényt tekinthetünk egy vektoron, amely e
három tulajdonságot teljesíti. Megjegyzendő, hogy az első két
tulajdonság együtt magában foglalja, hogy \(\|\boldsymbol{x}\| = 0\)
akkor és csak akkor, ha \(\boldsymbol{x} = \boldsymbol{0}\). A
háromszög-egyenlőtlenség hasznos változata

\[|\|x\| - \|y\|| \le \|x - y\|,\]

amely egyúttal azt is mutatja, hogy a vektornorma folytonos függvény.

\subsubsection{2.3.2 Mátrixnormák}\label{muxe1trixnormuxe1k}

A mátrixok méretének vagy nagyságának mérésére is szükségünk van
valamilyen módszerre. Itt is lehetséges általánosabb definíció, de a
továbbiakban használt mátrixnormákat mind egy alapul szolgáló
vektornorma segítségével definiáljuk. Pontosabban: adott vektornorma
mellett egy \(m \times n\) méretű \(A\) mátrixhoz tartozó mátrixnormát a
következőképpen értelmezzük:

\[\|\boldsymbol{A}\| = \max_{\boldsymbol{x} \neq \boldsymbol{0}} \frac{\|\boldsymbol{A}\boldsymbol{x}\|}{\|\boldsymbol{x}\|}.\]

Az ilyen mátrixnormáról azt mondjuk, hogy a vektornorma által indukált
vagy annak szubordinált normája. Szemléletesen: a mátrix normája azt a
maximális nyújtást méri, amelyet a mátrix az adott vektornormában mérve
bármely vektoron végrehajt.

Egyes mátrixnormák sokkal könnyebben számíthatók ki, mint mások. Például
a vektor 1-normához tartozó mátrixnorma egyszerűen a mátrix maximális
abszolút oszlopösszege,

\[\|A\|_1 = \max_j \sum_{i=1}^m |a_{ij}|,\]

a vektor \(\infty\)-normához tartozó mátrixnorma pedig egyszerűen a
mátrix maximális abszolút sorösszege,

\[\|\mathbf{A}\|_{\infty} = \max_{i} \sum_{j=1}^{n} |a_{ij}|.\]

Ezek megjegyzésének kényelmes módja, hogy ezek a mátrixnormák
\(n \times 1\) méretű mátrix esetén megegyeznek a megfelelő
vektornormákkal. Sajnos a vektor 2-normához tartozó mátrixnormát nem
olyan könnyű kiszámítani (lásd 3.6.1. szakasz).

\textbf{2.4. Példa. Mátrixnormák.} Az

\[\mathbf{A} = \begin{bmatrix} 2 & -1 & 1 \\ 1 & 0 & 1 \\ 3 & -1 & 4 \end{bmatrix}\]

mátrix esetén a maximális abszolút oszlop- és sorösszegekből rendre

\[\|A\|_1 = 6\] és \(\|A\|_{\infty} = 8\) adódik.

E mátrix 2-normáját lásd a 3.17. példában.

A definiált mátrixnormák a következő fontos tulajdonságokkal
rendelkeznek, ahol \(A\) és \(B\) tetszőleges mátrixok:

\begin{itemize}
\tightlist
\item
  \begin{enumerate}
  \def\labelenumi{\arabic{enumi}.}
  \tightlist
  \item
    \(\|A\| > 0\), ha \(A \neq O\).
  \end{enumerate}
\item
  \begin{enumerate}
  \def\labelenumi{\arabic{enumi}.}
  \setcounter{enumi}{1}
  \tightlist
  \item
    \(\|\gamma A\| = |\gamma| \cdot \|A\|\) tetszőleges \(\gamma\)
    skalár esetén.
  \end{enumerate}
\item
  \begin{enumerate}
  \def\labelenumi{\arabic{enumi}.}
  \setcounter{enumi}{2}
  \tightlist
  \item
    \(\|A + B\| \le \|A\| + \|B\|\).
  \end{enumerate}
\item
  \begin{enumerate}
  \def\labelenumi{\arabic{enumi}.}
  \setcounter{enumi}{3}
  \tightlist
  \item
    \(\|AB\| \le \|A\| \cdot \|B\|\).
  \end{enumerate}
\item
  \begin{enumerate}
  \def\labelenumi{\arabic{enumi}.}
  \setcounter{enumi}{4}
  \tightlist
  \item
    \(\|Ax\| \le \|A\| \cdot \|x\|\) tetszőleges \(x\) vektor esetén.
  \end{enumerate}
\end{itemize}

Az általánosabb megközelítésben egy mátrixnorma definíciójaként bármely
olyan valós értékű mátrixfüggvényt tekinthetünk, amely az első három
tulajdonságot teljesíti. A fennmaradó két tulajdonság, amelyeket
szubmultiplikatív vagy konzisztenciafeltételeknek neveznek, az ilyen
általánosabb mátrixnormákra nem feltétlenül teljesül, de a vektor
\(p\)-normák által indukált mátrixnormákra mindig érvényesek. Ismét
megjegyezzük, hogy az első két tulajdonság együtt azt adja, hogy
\(\|A\| = 0\) akkor és csak akkor, ha \(A = O\).

\subsubsection{2.3.3 Mátrix
kondíciószáma}\label{muxe1trix-konduxedciuxf3szuxe1ma}

Egy nemszinguláris \(A\) négyzetes mátrix adott mátrixnormára vonatkozó
kondíciószámát a következőképpen definiáljuk:

\[\operatorname{cond}(\boldsymbol{A}) = \|\boldsymbol{A}\| \cdot \|\boldsymbol{A}^{-1}\|.\]

Megállapodás szerint \(\operatorname{cond}(A) = \infty\), ha \(A\)
szinguláris. A 2.3.4. szakaszban látni fogjuk, hogy ez a fogalom
összhangban van a kondíciószám 1.2.6. szakaszban bevezetett általános
fogalmával annyiban, hogy a mátrix kondíciószáma felső korlátot ad a
lineáris egyenletrendszer megoldásának relatív változásának és a
bemeneti adatok adott relatív változásának arányára.

\textbf{2.5. Példa. Mátrix kondíciószáma.} Mátrixszorzással könnyen
ellenőrizhető, hogy a 2.4. példa mátrixának inverze

\[\mathbf{A}^{-1} = \begin{bmatrix} 0{,}5 & 1{,}5 & -0{,}5 \\ -0{,}5 & 2{,}5 & -0{,}5 \\ -0{,}5 & -0{,}5 & 0{,}5 \end{bmatrix},\]

így

\[\|\mathbf{A}^{-1}\|_1 = 4{,}5\] és
\(\|\mathbf{A}^{-1}\|_{\infty} = 3{,}5\).

Ebből tehát

\[\operatorname{cond}_1(\mathbf{A}) = \|\mathbf{A}\|_1 \cdot \|\mathbf{A}^{-1}\|_1 = 6 \cdot 4{,}5 = 27\]

és

\[\operatorname{cond}_{\infty}(\mathbf{A}) = \|\mathbf{A}\|_{\infty} \cdot \|\mathbf{A}^{-1}\|_{\infty} = 8 \cdot 3{,}5 = 28.\]

E mátrix 2-normára vonatkozó kondíciószámát lásd a 3.17. példában.

A 2.5. példából látható, hogy egy \(n \times n\) méretű mátrix
kondíciószámának számértéke függ az alkalmazott normától (amit a
megfelelő alsó index jelez), de az alapul szolgáló vektornormák
ekvivalenciája miatt ezek az értékek legfeljebb egy rögzített
konstansban térhetnek el (amely \(n\)-től függ), így a kondicionáltság
kvantitatív mértékeként egyaránt hasznosak.

Mivel

\[\|\boldsymbol{A}\| \cdot \|\boldsymbol{A}^{-1}\| = \left(\max_{\boldsymbol{x} \neq \boldsymbol{0}} \frac{\|\boldsymbol{A}\boldsymbol{x}\|}{\|\boldsymbol{x}\|}\right) \cdot \left(\min_{\boldsymbol{x} \neq \boldsymbol{0}} \frac{\|\boldsymbol{A}\boldsymbol{x}\|}{\|\boldsymbol{x}\|}\right)^{-1},\]

egy mátrix kondíciószáma a maximális relatív nyújtás és a maximális
relatív zsugorítás arányát méri, amelyet a mátrix bármely nem nulla
vektoron végrehajt. Ezt úgy is megfogalmazhatjuk, hogy a mátrix
kondíciószáma az egységgömb torzulásának mértékét méri (a megfelelő
vektornormában) a mátrix által való transzformáció során. Minél nagyobb
a kondíciószám, annál torzabb (relatíve nyúlt és vékony) lesz az
egységgömb a mátrix általi transzformáció után. Két dimenzióban például
a 2-normabeli egységkörből egyre inkább szivar alakú ellipszis lesz, az
1-normában vagy az \(\infty\)-normában pedig az egységgömb négyzetből
egyre ferdébb paralelogrammává torzul, ahogy a kondíciószám nő.

\textbf{2.6. Példa. Mátrix kondíciószáma.} A 2.3. ábra négy különböző
mátrix hatását szemlélteti az \(\mathbb{R}^2\)-beli egységkörön a
2-normát használva. \(A_1\) az egységkört 30 fokkal az óramutató
járásával megegyező irányba forgatja el, de egyik vektor euklideszi
hosszát sem változtatja meg, így \(\operatorname{cond}_2(A_1) = 1\).
\(A_2\) az \(e_1\) bázisvektort 2-szeresére nyújtja, az \(e_2\)
bázisvektort pedig 0,5-szeresére zsugorítja, és mindkettő maximális,
tehát \(\operatorname{cond}_2(A_2) = 2/0{,}5 = 4\). \(A_3\) egyszerre
forgat és torzít az egységkörön, de a maximális arány itt is ugyanaz,
mint \(A_2\) esetén, így \(\operatorname{cond}_2(A_3) = 4\). Végül
\(A_4\) egy általánosabb transzformáció, amely esetén a maximális arány
már nem a bázisvektoroknál jelentkezik, de a maximum értéke továbbra is
ugyanaz, tehát \(\operatorname{cond}_2(A_4) = 4\).

\pandocbounded{\includegraphics[keepaspectratio]{_page_8_Figure_3.jpeg}}

2.3. ábra: Az egységkör transzformációja 2-normában különböző
mátrixokkal.

A kondíciószám következő fontos tulajdonságai a definícióból könnyen
levezethetők, és tetszőleges normára érvényesek:

\begin{itemize}
\tightlist
\item
  \begin{enumerate}
  \def\labelenumi{\arabic{enumi}.}
  \tightlist
  \item
    Tetszőleges \(A\) mátrix esetén \(\operatorname{cond}(A) \ge 1\).
  \end{enumerate}
\item
  \begin{enumerate}
  \def\labelenumi{\arabic{enumi}.}
  \setcounter{enumi}{1}
  \tightlist
  \item
    Az egységmátrixra \(\operatorname{cond}(I) = 1\).
  \end{enumerate}
\item
  \begin{enumerate}
  \def\labelenumi{\arabic{enumi}.}
  \setcounter{enumi}{2}
  \tightlist
  \item
    Tetszőleges \(A\) mátrix és nem nulla \(\gamma\) skalár esetén
    \(\operatorname{cond}(\gamma A) = \operatorname{cond}(A)\).
  \end{enumerate}
\item
  \begin{enumerate}
  \def\labelenumi{\arabic{enumi}.}
  \setcounter{enumi}{3}
  \tightlist
  \item
    Tetszőleges \(D = \operatorname{diag}(d_i)\) diagonális mátrix
    esetén \(\operatorname{cond}(D) = (\max |d_i|)/(\min |d_i|)\).
  \end{enumerate}
\end{itemize}

A kondíciószám azt méri, hogy a mátrix mennyire van közel a
szingularitáshoz: egy nagy kondíciószámú mátrix (amit a 2.3.4.
szakaszban fogunk kvantifikálni) közel szinguláris, míg az 1-hez közeli
kondíciószámú mátrix messze van a szingularitástól. A definícióból
nyilvánvaló, hogy egy nemszinguláris mátrix és inverze kondíciószáma
megegyezik. Ez érthető, hiszen ha egy mátrix közel szinguláris, akkor az
inverze is ugyanolyan közel van a szingularitáshoz.

Megjegyzendő, hogy a mátrix determinánsa nem jó indikátora a
szingularitáshoz való közelségnek: bár egy \(A\) mátrix szinguláris, ha
\(\det(A) = 0\), a nem nulla determináns nagysága -- akár nagy, akár
kicsi -- semmilyen információt nem ad arról, hogy a mátrix mennyire
közelít a szingularitáshoz. Például \(\det(\alpha I_n) = \alpha^n\), ami
\(|\alpha| < 1\) esetén tetszőlegesen kicsi lehet, ennek ellenére a
mátrix tetszőleges nem nulla \(\alpha\) esetén tökéletesen jól
kondicionált, kondíciószáma bármely mátrixnormában 1.

Ahogy nemsokára látni fogjuk, a kondíciószám hasznossága a lineáris
egyenletrendszerek megoldásának pontosságának megítélésében rejlik. A
kondíciószám definíciója a mátrix inverzét tartalmazza, így értékének
kiszámítása nyilvánvalóan nem triviális feladat. Valójában a
kondíciószámot közvetlenül a definícióból kiszámítani lényegesen nagyobb
munkát igényelne, mint annak a megoldásnak a kiszámítása, amelynek
pontosságát a kondíciószám segítségével kívánjuk megítélni. A
gyakorlatban ezért a kondíciószámot csupán becsüljük, legfeljebb egy
nagyságrend pontossággal, a megoldási folyamat viszonylag olcsó
melléktermékeként.

Az \(\|A\|\) mátrixnormát könnyen kiszámíthatjuk maximális abszolút
oszlopösszegként (vagy sorösszegként, a használt normától függően). Az
\(\|A^{-1}\|\) alacsony költségű becslése jelenti a kihívást. A normák
tulajdonságaiból tudjuk, hogy ha \(z\) az \(Az = y\) rendszer megoldása,
akkor

\[\|z\| = \|A^{-1}y\| \le \|A^{-1}\| \cdot \|y\|,\]

így

\[\frac{\|\boldsymbol{z}\|}{\|\boldsymbol{y}\|} \le \|\boldsymbol{A}^{-1}\|,\]

és ez a korlát valamely optimálisan választott \(y\) vektor esetén
elérhető. Ezért ha \(y\)-t úgy tudjuk megválasztani, hogy az
\(\|z\|/\|y\|\) hányados a lehető legnagyobb legyen, akkor ésszerű
becslést kapunk \(\|A^{-1}\|\)-re.

\textbf{2.7. Példa. Kondícióbecslés.} Tekintsük az

\[\mathbf{A} = \begin{bmatrix} 0{,}913 & 0{,}659 \\ 0{,}457 & 0{,}330 \end{bmatrix}\]

mátrixot. Ha \(y = [0, 1{,}5]^T\)-t választjuk, akkor
\(z = [-7780, 10780]^T\), tehát

\[\|\boldsymbol{A}^{-1}\|_{1} \approx \frac{\|\boldsymbol{z}\|_{1}}{\|\boldsymbol{y}\|_{1}} \approx 1{,}238 \times 10^{4},\]

és így

\[\operatorname{cond}_1(\mathbf{A}) = \|\mathbf{A}\|_1 \cdot \|\mathbf{A}^{-1}\|_1 \approx 1{,}370 \times 1{,}238 \times 10^4 = 1{,}696 \times 10^4,\]

ami -- mint kiderül -- a megjelenített számjegyek pontosságával egzakt.
E viszonylag nagy kondíciószám következményeit a 2.8. és a 2.17.
példában vizsgáljuk.

A 2.7. példabeli \(y\) vektort gondosan úgy választottuk meg, hogy a
lehető legnagyobb \(\|z\|/\|y\|\) hányadost adja, és így \(\|A^{-1}\|\)
pontos értékét eredményezze. Egy ilyen optimális \(y\) megtalálása
általában megengedhetetlenül költséges lenne, de jóval olcsóbban is
kaphatunk hasznos közelítést. Az egyik heurisztika az, hogy \(y\)-t az
\(A^T y = c\) rendszer megoldásaként választjuk, ahol \(c\) egy olyan
vektor, amelynek komponensei \(\pm 1\) értékűek; az előjeleket egymás
után úgy választjuk meg, hogy az adódó \(y\) a lehető legnagyobb legyen.
Egy másik stratégia egyszerűen az, hogy néhány véletlen \(y\)-t
próbálunk; a közöttük legnagyobb hányados megadása gyakorlati célokra
általában elegendően pontos becslést ad \(\|A^{-1}\|\)-re.

A kondícióbecslés alternatív megközelítése, hogy konvex optimalizálási
feladatként kezeljük, amely heurisztikus algoritmussal a gyakorlatban
nagyon hatékonyan megoldható. További lehetőség -- a 2-norma
használatakor -- a kondíciószámot a szingulárisérték-felbontásból
származtatni (lásd 3.6.1. szakasz), ez azonban megengedhetetlenül
költséges, hacsak az SVD-t egyéb okból amúgy sem számítjuk ki.
Szerencsére a felhasználóknak nem kell aggódniuk e részletek miatt,
mivel a lineáris egyenletrendszerek megoldására szolgáló jó modern
szoftvercsomagok többsége hatékony és megbízható kondícióbecslőt
biztosít, amely az itt vázolt módszerek valamelyikének kifinomult
megvalósításán alapul (lásd 2.1. táblázat).

\subsubsection{2.3.4 Hibakorlátok}\label{hibakorluxe1tok}

A kondíciószám nemcsak a szingularitáshoz való közelség megbízható
indikátora, hanem kvantitatív korlátot is szolgáltat a lineáris
egyenletrendszer számított megoldásának hibájára, ahogy azt most
megmutatjuk. Legyen \(x\) a nemszinguláris \(Ax = b\) lineáris
egyenletrendszer megoldása, és legyen \(\hat{x}\) az
\(A\hat{x} = b + \Delta b\) perturbált jobb oldalú rendszer megoldása.
Ha definiáljuk \(\Delta x = \hat{x} - x\)-et, akkor

\[\boldsymbol{A}\hat{\boldsymbol{x}} = \boldsymbol{A}(\boldsymbol{x} + \Delta\boldsymbol{x}) = \boldsymbol{A}\boldsymbol{x} + \boldsymbol{A}\,\Delta\boldsymbol{x} = \boldsymbol{b} + \Delta\boldsymbol{b}.\]

Mivel \(Ax = b\), ezért \(A\,\Delta x = \Delta b\)-nek kell teljesülnie,
így \(\Delta x = A^{-1}\Delta b\). Normákat véve kapjuk az

\[\|b\| = \|Ax\| \le \|A\| \cdot \|x\|, \quad \text{vagyis} \quad \|x\| \ge \|b\|/\|A\|,\]

és

\[\|\Delta x\| = \|A^{-1}\Delta b\| \le \|A^{-1}\| \cdot \|\Delta b\|\]

egyenlőtlenségeket. E két egyenlőtlenséget összekapcsolva adódik, hogy

\[\frac{\|\Delta \boldsymbol{x}\|}{\|\boldsymbol{x}\|} \leq \|\boldsymbol{A}^{-1}\| \cdot \|\Delta \boldsymbol{b}\| \frac{\|\boldsymbol{A}\|}{\|\boldsymbol{b}\|}.\]

Definíció szerint \(\|A\| \cdot \|A^{-1}\| = \operatorname{cond}(A)\),
így a következő korlátot kapjuk:

\[\frac{\|\Delta \boldsymbol{x}\|}{\|\boldsymbol{x}\|} \leq \operatorname{cond}(\boldsymbol{A}) \frac{\|\Delta \boldsymbol{b}\|}{\|\boldsymbol{b}\|}.\]

Így a mátrix kondíciószáma egy „erősítési tényező", amely a jobb oldali
vektor adott relatív változásából fakadó megoldás maximális relatív
változását korlátozza (vö. a kondíciószám 1.2.6. szakaszban megadott
általános fogalmával).

Hasonló eredmény érvényes az \(A\) mátrix elemeinek relatív változásaira
is. Ha \(Ax = b\) és \((A + E)\hat{x} = b\), akkor

\[\Delta x = \hat{x} - x = A^{-1}(A\hat{x} - b) = -A^{-1}E\hat{x}.\]

Normákat véve az

\[\|\Delta \boldsymbol{x}\| \leq \|\boldsymbol{A}^{-1}\| \cdot \|\boldsymbol{E}\| \cdot \|\hat{\boldsymbol{x}}\|\]

egyenlőtlenséget kapjuk, amelyből \(\operatorname{cond}(A)\)
definíciójával adódik a

\[\frac{\|\Delta \boldsymbol{x}\|}{\|\hat{\boldsymbol{x}}\|} \leq \operatorname{cond}(\boldsymbol{A}) \frac{\|\boldsymbol{E}\|}{\|\boldsymbol{A}\|}\]

korlát.

Az imént ismertetett algebrai levezetések alternatívájaként a lineáris
egyenletrendszerek érzékenységét a differenciálszámítás segítségével is
becsülhetjük. A \(t\) valós paraméter bevezetésével definiáljuk az
\(A(t) = A + tE\) és \(b(t) = b + t\Delta b\) mennyiségeket, és
tekintjük az \(A(t)x(t) = b(t)\) lineáris egyenletrendszer \(x(t)\)
megoldását. Ezt az egyenletet \(t\) szerint differenciálva

\[\boldsymbol{A}'(t)\boldsymbol{x}(t) + \boldsymbol{A}(t)\boldsymbol{x}'(t) = \boldsymbol{b}'(t),\]

vagy ekvivalens módon

\[Ex(t) + A(t)x'(t) = \Delta b\]

adódik. Átrendezve kapjuk, hogy

\[\boldsymbol{x}'(t) = \boldsymbol{A}(t)^{-1} (\Delta \boldsymbol{b} - \boldsymbol{E} \boldsymbol{x}(t)),\]

amiből

\[\boldsymbol{x}'(0) = \boldsymbol{A}^{-1}(\Delta \boldsymbol{b} - \boldsymbol{E}\boldsymbol{x}(0)).\]

A Taylor-tétel szerint \(x(t) = x(0) + tx'(0) + \mathcal{O}(t^2)\),
tehát

\[\boldsymbol{x}(t) - \boldsymbol{x}(0) = t\boldsymbol{x}'(0) + \mathcal{O}(t^2) = t\boldsymbol{A}^{-1}(\Delta\boldsymbol{b} - \boldsymbol{E}\boldsymbol{x}(0)) + \mathcal{O}(t^2).\]

Legyen \(x \equiv x(0)\); normákat véve, majd \(\|x\|\)-szel osztva a

\[\frac{\|\boldsymbol{x}(t) - \boldsymbol{x}\|}{\|\boldsymbol{x}\|} \leq \|\boldsymbol{A}^{-1}\| \left( \frac{\|\Delta \boldsymbol{b}\|}{\|\boldsymbol{x}\|} + \|\boldsymbol{E}\| \right) |t| + \mathcal{O}(t^{2})\]

\[\leq \operatorname{cond}(\boldsymbol{A}) \left( \frac{\|\Delta \boldsymbol{b}\|}{\|\boldsymbol{A}\| \cdot \|\boldsymbol{x}\|} + \frac{\|\boldsymbol{E}\|}{\|\boldsymbol{A}\|} \right) |t| + \mathcal{O}(t^{2})\]

\[\leq \operatorname{cond}(\boldsymbol{A}) \left( \frac{\|\Delta \boldsymbol{b}\|}{\|\boldsymbol{b}\|} + \frac{\|\boldsymbol{E}\|}{\|\boldsymbol{A}\|} \right) |t| + \mathcal{O}(t^{2})\]

korlátot kapjuk.

Így ismét azt látjuk, hogy a megoldás relatív változását a kondíciószám
szorozva a feladat adatainak relatív változásával korlátozza.

Az érzékenységi eredmények geometriai értelmezése két dimenzióban: ha a
két egyenlet által meghatározott két egyenes közel párhuzamos, akkor
metszéspontjuk nem élesen definiált, amennyiben az egyenesek kerekítési
vagy mérési hiba miatt némileg bizonytalanok. Ha viszont az egyenesek
távol vannak a párhuzamostól, mondjuk közel merőlegesek, akkor
metszéspontjuk viszonylag élesen definiált. E két esetet a 2.4. ábra
mutatja be, ahol a szaggatott vonalak jelölik az egyes folytonos
vonalakhoz tartozó bizonytalansági tartományt, úgyhogy a metszéspont
mindkét esetben a satírozott paralelogramma bármelyik pontja lehet. Így
egy nagy kondíciószám a megoldás nagy bizonytalanságával jár.

\pandocbounded{\includegraphics[keepaspectratio]{_page_12_Picture_4.jpeg}}

2.4. ábra: Jól kondicionált és rosszul kondicionált lineáris
egyenletrendszerek.

Összefoglalva: ha a bemeneti adatok gépi pontosságig pontosak, akkor egy
\(Ax = b\) lineáris egyenletrendszer közelítő \(\hat{x}\) megoldásának
relatív hibájára azt várhatjuk, hogy

\[\frac{\|\hat{\boldsymbol{x}}-\boldsymbol{x}\|}{\|\boldsymbol{x}\|}\lessapprox \operatorname{cond}(\boldsymbol{A})\;\epsilon_{\mathrm{mach}}.\]

Ezt az eredményt a hátraható hibaelemzés nyelvén úgy értelmezhetjük,
hogy a számított megoldásról azt várhatjuk: a bemenet pontosságához
képest körülbelül \(\log_{10}(\operatorname{cond}(\mathbf{A}))\) tizedes
értékes jegyet veszít pontosságából. A 2.7. példabeli mátrix
kondíciószáma például nagyobb \(10^4\)-nél, ezért e mátrixot tartalmazó
lineáris egyenletrendszer megoldásában egyetlen helyes számjegyet sem
várhatunk, hacsak a bemeneti adatok négy tizedes értékes jegynél
pontosabbak, és a munkapontosság is négy tizedes értékes jegynél többet
hordoz (lásd a 2.8. és a 2.17. példát).

Az érzékenység kvantitatív mértékeként a mátrix kondíciószáma a lineáris
egyenletrendszerek megoldásának feladatában ugyanazt a szerepet tölti be
-- és ugyanolyan típusú kapcsolatot ad az előreható és a hátraható hiba
között --, mint az 1.2.6. szakaszban definiált általános kondíciószám.
Fontos különbség azonban, hogy a mátrix kondíciószáma soha nem lehet
kisebb 1-nél.

Mielőtt lezárnánk a pontosság kondíciószámok alapján történő
megítélésének témáját, vegyük figyelembe e két megjegyzést:

\begin{itemize}
\item
  A fenti normákon alapuló elemzés a megoldásvektor \emph{legnagyobb}
  komponenseinek relatív hibájára ad korlátot. A kisebb komponensek
  \emph{relatív} hibája ennél jóval nagyobb is lehet, mivel a
  vektornormát a vektor legnagyobb komponensei dominálják.
  Komponensenkénti hibakorlátok is nyerhetők, de ezek kiszámítása
  valamivel összetettebb, és mi ezzel a témával nem foglalkozunk. A
  komponensenkénti korlátok különösen érdekesek, amikor a rendszer
  rosszul skálázott (lásd 2.4.10. szakasz).
\item
  A mátrix kondíciószámát befolyásolja a mátrix skálázása (lásd 2.10.
  példa). Nagy kondíciószám származhat pusztán rossz skálázásból (lásd
  2.20. példa), csakúgy, mint a szingularitáshoz való közelségből. A
  mátrix átskálázása az előbbin segíthet, az utóbbin nem (lásd 2.4.10.
  szakasz).
\end{itemize}

\subsubsection{2.3.5 Maradék}\label{maraduxe9k}

Egy egyenlet megoldásának egyik ellenőrzési módja, hogy behelyettesítjük
az egyenletbe, és megnézzük, mennyire illeszkedik a bal és a jobb oldal.
Az \(Ax = b\) lineáris egyenletrendszer közelítő \(\hat{x}\)
megoldásának maradéka a

\[r = b - A\hat{x}\]

különbség.

Ha \(A\) nemszinguláris, akkor elvben a
\(\|\Delta x\| = \|\hat{x} - x\| = 0\) hiba akkor és csak akkor
teljesül, ha \(\|r\| = 0\). A gyakorlatban azonban ezek a mennyiségek
nem szükségképpen egyszerre kicsik. Először is figyeljük meg, hogy ha az
\(Ax = b\) egyenletet tetszőleges nem nulla konstanssal megszorozzuk,
akkor a megoldás nem változik, de a maradék ugyanezzel a tényezővel
szorzódik. Így a maradék a feladat skálázásától függően tetszőlegesen
naggyá vagy kicsivé tehető, ezért a maradék nagysága önmagában
értelmetlen, hacsak nem a feladat adatainak és a megoldás méretéhez
viszonyítva tekintjük. Ezért a \(\hat{x}\) közelítő megoldás relatív
maradékát \(\|r\|/(\|A\|\cdot\|\hat{x}\|)\)-ként definiáljuk.

A hiba és a maradék kapcsolatának megállapítása érdekében megfigyeljük,
hogy

\[\|\Delta x\| = \|\hat{x} - x\| = \|A^{-1}(A\hat{x} - b)\| = \|-A^{-1}r\| \le \|A^{-1}\| \cdot \|r\|.\]

Mindkét oldalt \(\|\hat{x}\|\)-szel osztva és \(\operatorname{cond}(A)\)
definícióját felhasználva kapjuk, hogy

\[\frac{\|\Delta \boldsymbol{x}\|}{\|\hat{\boldsymbol{x}}\|} \leq \operatorname{cond}(\boldsymbol{A}) \frac{\|\boldsymbol{r}\|}{\|\boldsymbol{A}\| \cdot \|\hat{\boldsymbol{x}}\|}.\]

Így kis relatív maradék akkor és csak akkor jelent kis relatív hibát a
megoldásban, ha \(A\) jól kondicionált.

Ha viszont a nagy maradék következményeit akarjuk látni, tegyük fel,
hogy a számított \(\hat{x}\) megoldás egzaktul kielégíti az

\[(\boldsymbol{A} + \boldsymbol{E})\hat{\boldsymbol{x}} = \boldsymbol{b}\]

egyenletet. Ekkor

\[\|\boldsymbol{r}\| = \|\boldsymbol{b} - \boldsymbol{A}\hat{\boldsymbol{x}}\| = \|\boldsymbol{E}\hat{\boldsymbol{x}}\| \le \|\boldsymbol{E}\| \cdot \|\hat{\boldsymbol{x}}\|,\]

tehát a

\[\frac{\|\boldsymbol{r}\|}{\|\boldsymbol{A}\|\cdot\|\hat{\boldsymbol{x}}\|} \leq \frac{\|\boldsymbol{E}\|}{\|\boldsymbol{A}\|}\]

egyenlőtlenséget kapjuk.

Így a nagy relatív maradék nagy hátraható hibát jelent a mátrixban, azaz
a megoldás kiszámítására használt algoritmus instabil. Másként
fogalmazva: egy stabil algoritmus a feladat kondicionáltságától
függetlenül mindig kis relatív maradékú megoldást ad, ezért egy kis
maradék önmagában keveset mond a közelítő megoldás minőségéről. Erre a
kérdésre a 2.4.5. szakaszban még visszatérünk.

\textbf{2.8. Példa. Kis maradék.} Tekintsük az

\[\boldsymbol{Ax} = \begin{bmatrix} 0{,}913 & 0{,}659 \\ 0{,}457 & 0{,}330 \end{bmatrix} \begin{bmatrix} x_1 \\ x_2 \end{bmatrix} = \begin{bmatrix} 0{,}254 \\ 0{,}127 \end{bmatrix} = \boldsymbol{b}\]

lineáris egyenletrendszert, amelynek mátrixát a 2.7. példában láttuk.
Tekintsük a két közelítő megoldást:

\[\hat{\boldsymbol{x}}_1 = \begin{bmatrix} 0{,}6391 \\ -0{,}5 \end{bmatrix}\]
és
\(\hat{\boldsymbol{x}}_2 = \begin{bmatrix} 0{,}999 \\ -1{,}001 \end{bmatrix}\).

A megfelelő maradékok normái

\[\|\boldsymbol{r}_1\|_1 = 7{,}0 \times 10^{-5}\] és
\(\|\boldsymbol{r}_2\|_1 = 2{,}4 \times 10^{-2}\).

Melyik tehát a jobb megoldás? Hajlamosak lennénk \(\hat{x}_1\)-et
választani a sokkal kisebb maradéka miatt. Azonban e rendszer egzakt
megoldása -- ahogy könnyen ellenőrizhető -- \(x = [1, -1]^T\), tehát
\(\hat{x}_2\) valójában sokkal pontosabb, mint \(\hat{x}_1\). E meglepő
viselkedés oka, hogy az \(A\) mátrix rosszul kondicionált, ahogy a 2.7.
példában láttuk, és nagy kondíciószáma miatt a kis maradék nem jelent
kis hibát a megoldásban. Arra, hogy \(\hat{x}_1\)-et hogyan kaptuk, lásd
a 2.17. példát.

\subsection{2.4 Lineáris egyenletrendszerek
megoldása}\label{lineuxe1ris-egyenletrendszerek-megolduxe1sa}

\subsubsection{2.4.1
Feladatátalakítások}\label{feladatuxe1talakuxedtuxe1sok}

Az \(Ax = b\) lineáris egyenletrendszer megoldásához az 1.1.2.
szakaszban felvázolt általános stratégia azt sugallja, hogy a rendszert
olyan rendszerré kell átalakítanunk, amelynek megoldása megegyezik az
eredetiével, de könnyebben kiszámítható. Milyen típusú átalakítás hagyja
változatlanul egy lineáris rendszer megoldását? A válasz az, hogy az
\(Ax = b\) lineáris rendszer mindkét oldalát balról bármely
nemszinguláris \(M\) mátrixszal megszorozhatjuk anélkül, hogy a
megoldást befolyásolnánk. Hogy lássuk, miért, vegyük észre, hogy az
\(MAz = Mb\) lineáris rendszer megoldása

\[z = (MA)^{-1}Mb = A^{-1}M^{-1}Mb = A^{-1}b = x.\]

\textbf{2.9. Példa. Permutációk.} Ilyen átalakítás fontos példája az a
tény, hogy \(A\) sorai és \(b\) megfelelő elemei átrendezhetők az \(x\)
megoldás megváltoztatása nélkül. Ez intuitíve nyilvánvaló: a rendszer
összes egyenletének egyidejűleg teljesülnie kell, ezért az a sorrend,
amelyben leírjuk őket, közömbös; akár véletlenszerűen is húzhatnánk őket
egy kalapból. Formálisan a sorok ilyen átrendezését úgy valósítjuk meg,
hogy az egyenlet mindkét oldalát balról megszorozzuk egy \(P\)
permutációs mátrixszal, amely olyan négyzetes mátrix, amelynek minden
sorában és oszlopában pontosan egy 1-es szerepel, a többi elem pedig
nulla (azaz egy egységmátrix, amelynek sorai és oszlopai permutálva
vannak). Például

\[\begin{bmatrix} 0 & 0 & 1 \\ 1 & 0 & 0 \\ 0 & 1 & 0 \end{bmatrix} \begin{bmatrix} v_1 \\ v_2 \\ v_3 \end{bmatrix} = \begin{bmatrix} v_3 \\ v_1 \\ v_2 \end{bmatrix}.\]

Egy permutációs mátrix mindig nemszinguláris; valójában az inverze
egyszerűen a transzponáltja, \(P^{-1} = P^T\) (egy \(M\) mátrix
transzponáltja, amelyet \(M^T\) jelöl, olyan mátrix, amelynek oszlopai
\(M\) sorai, azaz ha \(N = M^T\), akkor \(n_{ij} = m_{ji}\)). Így az
átrendezett rendszer \(PAx = Pb\) alakban írható, és az \(x\) megoldás
változatlan marad.

Ha jobbról szorzunk meg egy permutációs mátrixszal, akkor az a mátrix
oszlopait rendezi át a sorok helyett. Az ilyen átalakítás megváltoztatja
a megoldást, de csak úgy, hogy a megoldás komponensei permutálódnak.
Ennek belátásához figyeljük meg, hogy az \(APz = b\) rendszer megoldása

\[z = (AP)^{-1}b = P^{-1}A^{-1}b = P^{T}A^{-1}b = P^{T}x,\]

és ennélfogva az eredeti \(Ax = b\) rendszer megoldása \(x = Pz\).

\textbf{2.10. Példa. Diagonális skálázás.} Egy másik egyszerű, de fontos
átalakítási típus a diagonális skálázás. Emlékezzünk vissza, hogy egy
\(D\) mátrix diagonális, ha \(d_{ij} = 0\) minden \(i \neq j\) esetén,
azaz az egyetlen nemnulla elemei a főátlón lévő \(d_{ii}\),
\(i = 1, \dots, n\) elemek. Ha az \(Ax = b\) lineáris rendszer mindkét
oldalát balról megszorozzuk egy nemszinguláris \(D\) diagonális
mátrixszal, akkor a mátrix minden sorát és a jobb oldalt \(D\) megfelelő
átlóbeli elemével szorozzuk meg, és ezt ezért sorkénti skálázásnak
nevezzük. Elvben a sorkénti skálázás nem változtatja meg a lineáris
rendszer megoldását, a gyakorlatban azonban befolyásolhatja a numerikus
megoldási folyamatot és az adott feladatnál elérhető pontosságot, mint
azt látni fogjuk.

Az oszloponkénti skálázás --- egy lineáris rendszer mátrixának jobbról
való megszorzása egy nemszinguláris \(D\) diagonális mátrixszal --- a
mátrix minden oszlopát \(D\) megfelelő átlóbeli elemével szorozza meg.
Az ilyen átalakítás valóban megváltoztatja a megoldást, gyakorlatilag
azokat a mértékegységeket változtatja meg, amelyekben a megoldás
komponenseit mérjük. A skálázott \(ADz = b\) rendszer megoldása

\[z = (AD)^{-1}b = D^{-1}A^{-1}b = D^{-1}x,\]

és ennélfogva az eredeti \(Ax = b\) rendszer megoldása \(x = Dz\).

\subsubsection{2.4.2 Háromszögű lineáris
egyenletrendszerek}\label{huxe1romszuxf6gux171-lineuxe1ris-egyenletrendszerek}

A következő kérdés az, hogy milyen típusú lineáris rendszer könnyen
megoldható. Tegyük fel, hogy az \(Ax = b\) rendszerben van egy olyan
egyenlet, amely csak egy ismeretlen megoldáskomponenst tartalmaz (azaz
\(A\) ezen sorában csak egy elem nemnulla). Akkor ez az egyenlet könnyen
megoldható (osztással) erre az ismeretlenre. Tegyük fel most, hogy a
rendszerben van egy másik egyenlet, amely csak két ismeretlent
tartalmaz, és ezek egyike a már meghatározott. Ha a már meghatározott
megoldáskomponenst behelyettesítjük ebbe a második egyenletbe, akkor
könnyen megoldhatjuk a másik ismeretlenre. Ha ez a minta folytatódik ---
egyenletenként csak egy új ismeretlen komponens lép fel ---, akkor a
megoldás összes komponense egymás után kiszámítható. Az ilyen speciális
tulajdonságú mátrixot \emph{háromszögnek} nevezzük --- a név okát hamar
látni fogjuk. Mivel a háromszögű lineáris rendszerek ezzel az egymás
utáni helyettesítéssel könnyen megoldhatók, alkalmas cél, amelyre egy
általános lineáris rendszert átalakíthatunk.

Bár a fentebb leírt általános háromszögalak elegendő ahhoz, hogy a
rendszer egymás utáni helyettesítéssel megoldható legyen, számítási
célokra kényelmes két speciális háromszögalakot definiálnunk. Egy \(L\)
mátrix \emph{alsó háromszögmátrix}, ha a főátló feletti összes eleme
nulla (azaz \(\ell_{ij} = 0\) \(i < j\) esetén). Hasonlóan, egy \(U\)
mátrix \emph{felső háromszögmátrix}, ha a főátló alatti összes eleme
nulla (azaz \(u_{ij} = 0\) \(i > j\) esetén). Egy a fent definiált
általánosabb értelemben háromszögalakú mátrix sorainak vagy oszlopainak
alkalmas permutációjával felső vagy alsó háromszögalakba hozható.

Egy \(Lx = b\) alsó háromszögű rendszer esetén az egymás utáni
helyettesítést \emph{előrehelyettesítésnek} nevezzük, és matematikailag
így fejezhetjük ki:

\[x_1 = b_1/\ell_{11}, \quad x_i = \left(b_i - \sum_{j=1}^{i-1} \ell_{ij} x_j\right)/\ell_{ii}, \quad i = 2, \dots, n.\]

Megvalósításának egyik módját a 2.1. Algoritmus mutatja.

\textbf{2.1. Algoritmus. Előrehelyettesítés alsó háromszögű rendszerre.}

\begin{verbatim}
for j = 1 to n                       { ciklus az oszlopokra }
    if ℓjj = 0 then stop             { leállás, ha a mátrix szinguláris }
    xj = bj/ℓjj                      { megoldáskomponens kiszámítása }
    for i = j + 1 to n
        bi = bi − ℓij xj             { jobb oldal frissítése }
    end
end
\end{verbatim}

Hasonlóan, egy \(Ux = b\) felső háromszögű rendszer esetén az egymás
utáni helyettesítést \emph{visszahelyettesítésnek} nevezzük, és
matematikailag így fejezhetjük ki:

\[x_n = b_n/u_{nn}, \quad x_i = \left(b_i - \sum_{j=i+1}^n u_{ij}x_j\right)/u_{ii}, \quad i = n-1, \dots, 1.\]

Megvalósításának egyik módját a 2.2. Algoritmus mutatja.

Mindkét algoritmusban úgy választottuk meg a ciklusváltozók sorrendjét,
hogy a mátrixhoz oszloponként (és nem soronként) férünk hozzá, és a
belső ciklus művelete egy skalár szorozva egy vektorral, majd ehhez
hozzáadva egy vektor --- azaz egy „saxpy" művelet (lásd a 2.7.2.
szakaszt).

\textbf{2.2. Algoritmus. Visszahelyettesítés felső háromszögű
rendszerre.}

\begin{verbatim}
for j = n to 1                       { ciklus az oszlopokra visszafelé }
    if ujj = 0 then stop             { leállás, ha a mátrix szinguláris }
    xj = bj/ujj                      { megoldáskomponens kiszámítása }
    for i = 1 to j − 1
        bi = bi − uij xj             { jobb oldal frissítése }
    end
end
\end{verbatim}

Választhattuk volna a ciklusváltozók fordított sorrendjét is, ekkor a
mátrixhoz soronként férnénk hozzá, és a belső ciklus művelete két vektor
belső szorzata lenne, azaz „sdot". Ezek a megvalósítási választások
jelentős hatással lehetnek a teljesítményre, az alkalmazott programozási
nyelvtől és számítógépes rendszertől függően (lásd a 2.16. számítógépes
feladatot). Azt is érdemes megjegyezni, hogy egy nulla átlóbeli elem
mindkét algoritmust elbukatja, de ez így is várható, mivel egy nulla
átlóbeli elemű háromszögmátrix szükségképpen szinguláris.

\textbf{2.11. Példa. Háromszögű lineáris egyenletrendszer.} Tekintsük a
következő felső háromszögű lineáris rendszert:

\[\begin{bmatrix} 1 & 2 & 2 \\ 0 & -4 & -6 \\ 0 & 0 & -1 \end{bmatrix} \begin{bmatrix} x_1 \\ x_2 \\ x_3 \end{bmatrix} = \begin{bmatrix} 3 \\ -6 \\ 1 \end{bmatrix}.\]

Az utolsó egyenlet, \(-x_3 = 1\), közvetlenül megoldható
\(x_3 = -1\)-re. Ezt az értéket ezután behelyettesíthetjük a második
egyenletbe, hogy \(x_2 = 3\)-at kapjunk, végül pedig \(x_3\)-at és
\(x_2\)-t egyaránt behelyettesítve az első egyenletbe, \(x_1 = -1\)
adódik.

\subsubsection{2.4.3 Elemi eliminációs
mátrixok}\label{elemi-eliminuxe1ciuxf3s-muxe1trixok}

Stratégiánk tehát az, hogy olyan nemszinguláris lineáris transzformációt
tervezzünk, amely egy adott általános lineáris rendszert olyan
háromszögű lineáris rendszerré alakít, amely utána az egymás utáni
helyettesítéssel könnyen megoldható. Ehhez tehát olyan transzformációra
van szükségünk, amely az adott mátrix kiválasztott nemnulla elemeit
nullákkal helyettesíti. Ez a mátrix sorainak alkalmas lineáris
kombinációival érhető el, ahogyan azt most megmutatjuk.

Tekintsük az \(a = \begin{bmatrix} a_1 & a_2 \end{bmatrix}^T\)
2-vektort. Ha \(a_1 \neq 0\), akkor

\[\begin{bmatrix} 1 & 0 \\ -a_2/a_1 & 1 \end{bmatrix} \begin{bmatrix} a_1 \\ a_2 \end{bmatrix} = \begin{bmatrix} a_1 \\ 0 \end{bmatrix}.\]

Általánosabban: egy \(n\)-dimenziós \(a\) vektor \(k\)-adik pozíciója
alatti minden elemet annullálhatunk --- feltéve, hogy \(a_k \neq 0\) ---
a következő transzformációval:

\[\boldsymbol{M}_{k}\,\boldsymbol{a} = \begin{bmatrix} 1 & \cdots & 0 & 0 & \cdots & 0 \\ \vdots & \ddots & \vdots & \vdots & \ddots & \vdots \\ 0 & \cdots & 1 & 0 & \cdots & 0 \\ 0 & \cdots & -m_{k+1} & 1 & \cdots & 0 \\ \vdots & \ddots & \vdots & \vdots & \ddots & \vdots \\ 0 & \cdots & -m_{n} & 0 & \cdots & 1 \end{bmatrix} \begin{bmatrix} a_{1} \\ \vdots \\ a_{k} \\ a_{k+1} \\ \vdots \\ a_{n} \end{bmatrix} = \begin{bmatrix} a_{1} \\ \vdots \\ a_{k} \\ 0 \\ \vdots \\ 0 \end{bmatrix},\]

ahol \(m_i = a_i/a_k\), \(i = k+1, \dots, n\). Az \(a_k\) osztót
\emph{főelemnek} (pivotnak) nevezzük. Az ilyen alakú mátrixot néha
\emph{elemi eliminációs mátrixnak} vagy \emph{Gauss-transzformációnak}
nevezzük, és egy vektorra gyakorolt hatása az, hogy a \(k\)-adik sor
többszörösét hozzáadja minden további sorhoz, úgy megválasztva az
\(m_i\) szorzótényezőket, hogy az eredmény minden esetben nulla legyen.
Vegyük észre az elemi eliminációs mátrixok alábbi tulajdonságait:

\begin{enumerate}
\def\labelenumi{\arabic{enumi}.}
\tightlist
\item
  \(M_k\) alsó háromszögmátrix, amelynek főátlójában csupa egyes áll,
  ezért szükségképpen nemszinguláris.
\item
  \(M_k = I - m_k e_k^T\), ahol
  \(m_k = [0, \dots, 0, m_{k+1}, \dots, m_n]^T\), és \(e_k\) az
  egységmátrix \(k\)-adik oszlopa.
\item
  \(M_k^{-1} = I + m_k e_k^T\), ami azt jelenti, hogy \(M_k^{-1}\) ---
  amelyet \(L_k\)-val fogunk jelölni --- ugyanaz, mint \(M_k\), csak a
  szorzótényezők előjele megváltozik.
\item
  Ha \(M_j\), \(j > k\), egy másik elemi eliminációs mátrix \(m_j\)
  szorzótényező-vektorral, akkor
\end{enumerate}

\[M_k M_j = I - m_k e_k^T - m_j e_j^T + m_k e_k^T m_j e_j^T = I - m_k e_k^T - m_j e_j^T,\]

mivel \(e_k^T m_j = 0\). Így szorzatuk lényegében „uniójuk". Mivel
azonos alakúak, hasonló eredmény érvényes inverzeik \(L_k L_j\)
szorzatára is. Figyeljük meg, hogy a szorzás sorrendje lényeges; ezek az
eredmények a fordított szorzatra nem teljesülnek.

\textbf{2.12. Példa. Elemi eliminációs mátrixok.} Ha
\(a = \begin{bmatrix} 2 & 4 & -2 \end{bmatrix}^T\), akkor

\[\boldsymbol{M}_1 \boldsymbol{a} = \begin{bmatrix} 1 & 0 & 0 \\ -2 & 1 & 0 \\ 1 & 0 & 1 \end{bmatrix} \begin{bmatrix} 2 \\ 4 \\ -2 \end{bmatrix} = \begin{bmatrix} 2 \\ 0 \\ 0 \end{bmatrix}, \quad \boldsymbol{M}_2 \boldsymbol{a} = \begin{bmatrix} 1 & 0 & 0 \\ 0 & 1 & 0 \\ 0 & 0{,}5 & 1 \end{bmatrix} \begin{bmatrix} 2 \\ 4 \\ -2 \end{bmatrix} = \begin{bmatrix} 2 \\ 4 \\ 0 \end{bmatrix}.\]

Megjegyezzük továbbá, hogy

\[\boldsymbol{L}_1 = \boldsymbol{M}_1^{-1} = \begin{bmatrix} 1 & 0 & 0 \\ 2 & 1 & 0 \\ -1 & 0 & 1 \end{bmatrix}, \quad \boldsymbol{L}_2 = \boldsymbol{M}_2^{-1} = \begin{bmatrix} 1 & 0 & 0 \\ 0 & 1 & 0 \\ 0 & -0{,}5 & 1 \end{bmatrix},\]

és

\[\boldsymbol{M}_1 \boldsymbol{M}_2 = \begin{bmatrix} 1 & 0 & 0 \\ -2 & 1 & 0 \\ 1 & 0{,}5 & 1 \end{bmatrix}, \quad \boldsymbol{L}_1 \boldsymbol{L}_2 = \begin{bmatrix} 1 & 0 & 0 \\ 2 & 1 & 0 \\ -1 & -0{,}5 & 1 \end{bmatrix}.\]

\subsubsection{2.4.4 Gauss-kiküszöbölés és
LU-felbontás}\label{gauss-kikuxfcszuxf6buxf6luxe9s-uxe9s-lu-felbontuxe1s}

Elemi eliminációs mátrixokat használva viszonylag egyszerűen
redukálhatunk egy általános \(Ax = b\) lineáris rendszert felső
háromszögalakra. Először a 2.4.3. szakaszban megadott recept szerint
választunk egy \(M_1\) elemi eliminációs mátrixot, úgy, hogy a főelem az
első átlóbeli elem, \(a_{11}\) legyen, hogy \(A\) első oszlopa nullává
váljon az első sor alatt, amikor \(M_1\)-gyel balról szorozzuk.
Természetesen \(A\) többi oszlopát, valamint a \(b\) jobb oldali vektort
is meg kell szorozni \(M_1\)-gyel, így az új rendszer
\(M_1 A x = M_1 b\) alakú lesz, de korábbi tárgyalásunk szerint a
megoldás változatlan.

Ezután a második átlóbeli elemet használjuk főelemként egy olyan \(M_2\)
elemi eliminációs mátrix meghatározásához, amely az új \(M_1 A\) mátrix
második oszlopának minden elemét a második sor alatt annullálja.
\(M_2\)-t ismét az egész mátrixra és a jobb oldali vektorra kell
alkalmazni, így a \(M_2 M_1 A x = M_2 M_1 b\) módosított lineáris
rendszert kapjuk. Megjegyzendő, hogy \(M_1 A\) első oszlopát \(M_2\) nem
befolyásolja, mivel az érintett sorokban minden eleme nulla. Ez a
folyamat minden egyes egymás utáni oszlopra folytatódik, amíg a mátrix
összes főátló alatti eleme nem annullálódik. Ha bevezetjük az
\(M = M_{n-1} \cdots M_1\) mátrixot, akkor a transzformált lineáris
rendszer

\[MAx = M_{n-1} \cdots M_1 Ax = M_{n-1} \cdots M_1 b = Mb\]

felső háromszögű, és visszahelyettesítéssel megoldható, hogy megkapjuk
az eredeti \(Ax = b\) lineáris rendszer megoldását.

Az imént leírt eljárást \emph{Gauss-kiküszöbölésnek} nevezzük.
\emph{LU-felbontásnak} is nevezik, mert az \(A\) mátrixot egy egységnyi
átlójú alsó háromszögmátrix, \(L\), és egy felső háromszögmátrix, \(U\),
szorzatára bontja. Ennek belátásához emlékezzünk vissza, hogy az
\(L_k L_j\) szorzat egységnyi átlójú alsó háromszögű, ha \(k < j\), így

\[L = M^{-1} = (M_{n-1} \cdots M_1)^{-1} = M_1^{-1} \cdots M_{n-1}^{-1} = L_1 \cdots L_{n-1}\]

egységnyi átlójú alsó háromszögmátrix. Már láttuk, hogy konstrukció
szerint az \(U = MA\) mátrix felső háromszögű. Ennek megfelelően \(A\)-t

\[A = LU\]

szorzatként állítottuk elő, ahol \(L\) egységnyi átlójú alsó
háromszögmátrix és \(U\) felső háromszögmátrix. Egy ilyen felbontás
birtokában az \(Ax = b\) lineáris rendszer \(LUx = b\) alakban írható,
és így előbb az \(Ly = b\) alsó háromszögű rendszert
előrehelyettesítéssel, majd az \(Ux = y\) felső háromszögű rendszert
visszahelyettesítéssel megoldva megoldható. Megjegyzendő, hogy az \(y\)
köztes megoldás megegyezik a korábbi megfogalmazásban szereplő
transzformált jobb oldali \(Mb\) vektorral. Így a Gauss-kiküszöbölés és
az LU-felbontás ugyanannak a megoldási folyamatnak két kifejezésmódja.
Azzal azonban, hogy az előre- és visszahelyettesítő háromszöges
megoldási fázisokat a felbontástól elkülönítve hangsúlyozza, az LU-alak
talán világosabbá teszi, hogy a felbontási fázist nem kell megismételni,
amikor ugyanazzal az \(A\) mátrixszal, de más jobb oldali vektorokkal
további rendszereket oldunk meg, mivel az \(L\) és \(U\) tényezők újra
felhasználhatók.

A Gauss-kiküszöbölési eljárást formálisan a 2.3. Algoritmus foglalja
össze. Ez az algoritmus egyúttal \(A\) LU-felbontását is kiszámítja:
\(L\) főátló alatti elemeit \(\ell_{ik} = m_{ik}\) adja meg, \(U\)
átlóbeli és főátló feletti elemei pedig \(A\) megfelelő elemeit írják
felül. A gyakorlatban nem szükséges \(A\) főátló alatti elemeit explicit
módon nullára állítani, mivel csak a \(U\) és \(L\) eredményül kapott
háromszögmátrixok érdekelnek bennünket. Sőt: \(A\) főátló alatti elemei,
amelyeket kinulláztunk, tökéletes helyet nyújtanak \(L\) főátló alatti
elemeinek tárolására; a 2.3. Algoritmus lényegében „helyben" végzi el a
felbontást, ha \(m_{ik}\) helyett \(a_{ik}\)-t írunk minden
előforduláskor (lásd a 2.4.6. szakaszt ezekről és más megvalósítási
részletekről). Egy \(Ax = b\) lineáris rendszer megoldása során a \(b\)
jobb oldali vektor szükséges transzformációját beépíthetjük az
LU-felbontási folyamatba, vagy különálló lépésként végezhetjük el, a
2.1. Algoritmust használva az \(Ly = b\) alsó háromszögű rendszer
megoldására, miután \(L\)-et a 2.3. Algoritmussal meghatároztuk.
Bármelyik esetben a 2.2. Algoritmust használjuk ezután az \(Ux = y\)
felső háromszögű rendszer megoldására, hogy megkapjuk az \(x\)
megoldást.

\textbf{2.3. Algoritmus. LU-felbontás Gauss-kiküszöböléssel.}

\begin{verbatim}
for k = 1 to n − 1                       { ciklus az oszlopokra }
    if akk = 0 then stop                 { leállás, ha a főelem nulla }
    for i = k + 1 to n                   { szorzótényezők kiszámítása
        mik = aik/akk                        az aktuális oszlopra }
    end
    for j = k + 1 to n                   { transzformáció alkalmazása
        for i = k + 1 to n                   a maradék részmátrixra }
            aij = aij − mik akj
        end
    end
end
\end{verbatim}

\textbf{2.13. Példa. Gauss-kiküszöbölés.} A Gauss-kiküszöbölést a
következő lineáris egyenletrendszer megoldásán illusztráljuk:

\[x_1 + 2x_2 + 2x_3 = 3,\] \[4x_1 + 4x_2 + 2x_3 = 6,\]
\[4x_1 + 6x_2 + 4x_3 = 10,\]

vagy mátrix jelöléssel

\[\boldsymbol{A}\boldsymbol{x} = \begin{bmatrix} 1 & 2 & 2 \\ 4 & 4 & 2 \\ 4 & 6 & 4 \end{bmatrix} \begin{bmatrix} x_1 \\ x_2 \\ x_3 \end{bmatrix} = \boldsymbol{b}.\]

\(A\) első oszlopának főátló alatti elemeinek annullálásához az első sor
négyszeresét vonjuk ki a második és a harmadik sorból:

\[M_1 A = \begin{bmatrix} 1 & 0 & 0 \\ -4 & 1 & 0 \\ -4 & 0 & 1 \end{bmatrix} \begin{bmatrix} 1 & 2 & 2 \\ 4 & 4 & 2 \\ 4 & 6 & 4 \end{bmatrix} = \begin{bmatrix} 1 & 2 & 2 \\ 0 & -4 & -6 \\ 0 & -2 & -4 \end{bmatrix},\]

\[M_1 b = \begin{bmatrix} 1 & 0 & 0 \\ -4 & 1 & 0 \\ -4 & 0 & 1 \end{bmatrix} \begin{bmatrix} 3 \\ 6 \\ 10 \end{bmatrix} = \begin{bmatrix} 3 \\ -6 \\ -2 \end{bmatrix}.\]

Most az \(M_1 A\) második oszlopa főátló alatti elemének annullálásához
a második sor \(0{,}5\)-szörösét vonjuk ki a harmadik sorból:

\[\mathbf{M}_{2}\mathbf{M}_{1}\mathbf{A} = \begin{bmatrix} 1 & 0 & 0 \\ 0 & 1 & 0 \\ 0 & -0{,}5 & 1 \end{bmatrix} \begin{bmatrix} 1 & 2 & 2 \\ 0 & -4 & -6 \\ 0 & -2 & -4 \end{bmatrix} = \begin{bmatrix} 1 & 2 & 2 \\ 0 & -4 & -6 \\ 0 & 0 & -1 \end{bmatrix},\]

\[\mathbf{M}_{2}\mathbf{M}_{1}\mathbf{b} = \begin{bmatrix} 1 & 0 & 0 \\ 0 & 1 & 0 \\ 0 & -0{,}5 & 1 \end{bmatrix} \begin{bmatrix} 3 \\ -6 \\ -2 \end{bmatrix} = \begin{bmatrix} 3 \\ -6 \\ 1 \end{bmatrix}.\]

Ezzel az eredeti rendszert a következő ekvivalens felső háromszögű
rendszerre redukáltuk:

\[\boldsymbol{U}\boldsymbol{x} = \begin{bmatrix} 1 & 2 & 2 \\ 0 & -4 & -6 \\ 0 & 0 & -1 \end{bmatrix} \begin{bmatrix} x_1 \\ x_2 \\ x_3 \end{bmatrix} = \begin{bmatrix} 3 \\ -6 \\ 1 \end{bmatrix} = \boldsymbol{M}\boldsymbol{b} = \boldsymbol{y},\]

amelyet most visszahelyettesítéssel oldhatunk meg (mint a 2.11.
Példában), és így \(x = \begin{bmatrix} -1 & 3 & -1 \end{bmatrix}^T\)
adódik. Az LU-felbontás explicit felírásához

\[\boldsymbol{L}_1\boldsymbol{L}_2 = \begin{bmatrix} 1 & 0 & 0 \\ 4 & 1 & 0 \\ 4 & 0 & 1 \end{bmatrix} \begin{bmatrix} 1 & 0 & 0 \\ 0 & 1 & 0 \\ 0 & 0{,}5 & 1 \end{bmatrix} = \begin{bmatrix} 1 & 0 & 0 \\ 4 & 1 & 0 \\ 4 & 0{,}5 & 1 \end{bmatrix} = \boldsymbol{L},\]

így

\[\mathbf{A} = \begin{bmatrix} 1 & 2 & 2 \\ 4 & 4 & 2 \\ 4 & 6 & 4 \end{bmatrix} = \begin{bmatrix} 1 & 0 & 0 \\ 4 & 1 & 0 \\ 4 & 0{,}5 & 1 \end{bmatrix} \begin{bmatrix} 1 & 2 & 2 \\ 0 & -4 & -6 \\ 0 & 0 & -1 \end{bmatrix} = \mathbf{L}\mathbf{U}.\]

\subsubsection{2.4.5
Főelemkiválasztás}\label{fux151elemkivuxe1lasztuxe1s}

A Gauss-kiküszöbölési eljárással, ahogy eddig leírtuk, van egy
kézenfekvő probléma, és van egy másik, némileg finomabb is. A kézenfekvő
potenciális nehézség az, hogy az eljárás összeomlik, ha a mátrix még
redukálatlan részének vezető átlóbeli eleme valamelyik lépésben nulla,
mivel egy adott oszlop \(m_i\) szorzótényezőinek kiszámításához az
oszlop átlóbeli elemével való osztásra van szükség. A probléma megoldása
majdnem ugyanilyen kézenfekvő: ha a \(k\)-adik lépésben az átlóbeli elem
nulla, akkor cseréljük ki a rendszer \(k\)-adik sorát (mind a mátrixét,
mind a jobb oldali vektorét) egy olyan későbbi sorral, amelynek
\(k\)-adik oszlopbeli eleme nemnulla. A 2.9. Példából tudjuk, hogy egy
ilyen csere nem változtatja meg a rendszer megoldását. Nemnulla átlóbeli
elemet főelemként véve az eljárás a szokásos módon folytatódhat. A sorok
ilyen cseréjét \emph{főelemkiválasztásnak} nevezzük.

\textbf{2.14. Példa. Főelemkiválasztás és szingularitás.} A
főelemkiválasztás esetleges szüksége semmilyen kapcsolatban nem áll
azzal, hogy a mátrix szinguláris-e. Például az

\[\boldsymbol{A} = \begin{bmatrix} 0 & 1 \\ 1 & 0 \end{bmatrix}\]

mátrix nemszinguláris, mégsincs LU-felbontása, hacsak nem cseréljük fel
a sorokat, míg a szinguláris

\[\boldsymbol{A} = \begin{bmatrix} 1 & 1 \\ 1 & 1 \end{bmatrix}\]

mátrixnak van LU-felbontása:

\[A = \begin{bmatrix} 1 & 1 \\ 1 & 1 \end{bmatrix} = \begin{bmatrix} 1 & 0 \\ 1 & 1 \end{bmatrix} \begin{bmatrix} 1 & 1 \\ 0 & 0 \end{bmatrix} = LU.\]

De mi van akkor, ha a \(k\)-adik oszlop átlóján vagy az alatt nincs
nemnulla elem? Akkor ebben a lépésben nincs mit tenni, mivel az
annullálandó elemek már eleve nullák, így egyszerűen továbbléphetünk a
következő oszlopra (azaz \(M_k = I\)). Jegyezzük meg, hogy ez a lépés
nullát hagy az átlón, és így az eredményül kapott \(U\) felső
háromszögmátrix szinguláris lesz, de az LU-felbontás még így is
befejezhető. Ez viszont azt jelenti, hogy az ezt követő
visszahelyettesítés meghiúsul, mert \(U\) minden átlóbeli elemével
osztani kellene, de ez nem meglepő, hiszen az eredeti mátrixnak
mindenképpen szingulárisnak kellett lennie. Sokkal alattomosabb probléma
az, hogy lebegőpontos aritmetikában előfordulhat, hogy nem pontosan
nullát kapunk, hanem csak egy nagyon kicsi átlóbeli elemet --- és ezzel
érkezünk meg a finomabb ponthoz.

Elvben bármely nemnulla érték megfelelhet főelemnek a szorzótényezők
kiszámítására, de véges pontosságú aritmetikában a választást
körültekintően kell meghozni, hogy minimalizáljuk a numerikus hiba
terjedését. Konkrétan korlátozni szeretnénk a szorzótényezők nagyságát,
hogy a korábbi kerekítési hibák ne erősödjenek fel, amikor a mátrix és a
jobb oldal maradék részét megszorozzuk az egyes elemi eliminációs
mátrixokkal. A szorzótényezők nagysága soha nem fogja meghaladni az
1-et, ha minden oszlopban az átlón vagy az alatt található legnagyobb
abszolút értékű elemet választjuk főelemnek. Az ilyen eljárást
\emph{részleges főelemkiválasztásnak} nevezzük, és a gyakorlatban
elengedhetetlen az általános lineáris rendszereken végzett
Gauss-kiküszöbölés numerikusan stabil megvalósításához.

\textbf{2.15. Példa. Kis főelemek.} Véges pontosságú aritmetikában
nemcsak a nulla főelemeket, hanem a \emph{kis} főelemeket is kerülnünk
kell, hogy megelőzzük az elfogadhatatlan hibanövekedést, mint az alábbi
példa mutatja. Legyen

\[A = \begin{bmatrix} \epsilon & 1 \\ 1 & 1 \end{bmatrix},\]

ahol \(\epsilon\) egy olyan pozitív szám, amely az adott lebegőpontos
rendszer \(\epsilon_{\text{mach}}\) egységnyi kerekítési hibájánál
kisebb. Ha nem cseréljük fel a sorokat, akkor a főelem \(\epsilon\), és
a kapott szorzótényező \(-1/\epsilon\), így az eliminációs mátrix

\[M = \begin{bmatrix} 1 & 0 \\ -1/\epsilon & 1 \end{bmatrix},\]

és ennélfogva

\[\boldsymbol{L} = \begin{bmatrix} 1 & 0 \\ 1/\epsilon & 1 \end{bmatrix} \quad \text{és} \quad \boldsymbol{U} = \begin{bmatrix} \epsilon & 1 \\ 0 & 1 - 1/\epsilon \end{bmatrix} = \begin{bmatrix} \epsilon & 1 \\ 0 & -1/\epsilon \end{bmatrix}\]

lebegőpontos aritmetikában. De akkor

\[\boldsymbol{L}\boldsymbol{U} = \begin{bmatrix} 1 & 0 \\ 1/\epsilon & 1 \end{bmatrix} \begin{bmatrix} \epsilon & 1 \\ 0 & -1/\epsilon \end{bmatrix} = \begin{bmatrix} \epsilon & 1 \\ 1 & 0 \end{bmatrix} \neq \boldsymbol{A}.\]

Egy kis főelem és az ennek megfelelő nagy szorzótényező használata a
transzformált mátrixban helyrehozhatatlan információvesztést okozott. Ha
viszont felcseréljük a sorokat, akkor a főelem 1, és a kapott
szorzótényező \(-\epsilon\), így az eliminációs mátrix

\[\boldsymbol{M} = \begin{bmatrix} 1 & 0 \\ -\epsilon & 1 \end{bmatrix},\]

és ennélfogva

\[\boldsymbol{L} = \begin{bmatrix} 1 & 0 \\ \epsilon & 1 \end{bmatrix} \quad \text{és} \quad \boldsymbol{U} = \begin{bmatrix} 1 & 1 \\ 0 & 1 - \epsilon \end{bmatrix} = \begin{bmatrix} 1 & 1 \\ 0 & 1 \end{bmatrix}\]

lebegőpontos aritmetikában. Így

\[\boldsymbol{L}\boldsymbol{U} = \begin{bmatrix} 1 & 0 \\ \epsilon & 1 \end{bmatrix} \begin{bmatrix} 1 & 1 \\ 0 & 1 \end{bmatrix} = \begin{bmatrix} 1 & 1 \\ \epsilon & 1 \end{bmatrix},\]

ami a permutálás utáni helyes eredmény.

Bár az előbbi példa meglehetősen szélsőséges, az elv általánosan
érvényes: a nagyobb főelemek kisebb szorzótényezőkhöz, és így kisebb
hibákhoz vezetnek. Különösen, ha minden oszlopban az átlón vagy az alatt
található legnagyobb elemet használjuk főelemnek --- mint a 2.4.
Algoritmusban ---, akkor a szorzótényezők nagysága 1-gyel felülről
korlátozott. A részleges főelemkiválasztás által megkövetelt sorcserék
némileg bonyolultabbá teszik az LU-felbontás korábban megadott formális
leírását. Konkrétan minden \(M_k\) elemi eliminációs mátrix előtt egy
\(P_k\) permutációs mátrix áll, amely sorcserékkel a \(k\)-adik oszlop
átlón vagy alatt található legnagyobb abszolút értékű elemét az átlós
főelempozícióba hozza. Továbbra is \(MA = U\), ahol \(U\) felső
háromszögű, de most

\[\boldsymbol{M} = \boldsymbol{M}_{n-1} \boldsymbol{P}_{n-1} \cdots \boldsymbol{M}_1 \boldsymbol{P}_1.\]

\(M^{-1}\) továbbra is háromszögalakú a korábban definiált általános
értelemben, de a permutációk miatt \(M^{-1}\) nem szükségképpen alsó
háromszögű, bár továbbra is \(L\)-lel jelöljük. Így az „LU"-felbontás
szó szerint már nem „alsó szorozva felső" háromszögmátrixot jelent, de a
lineáris rendszerek egymás utáni helyettesítéssel való megoldására
ugyanolyan hasznos.

\textbf{2.4. Algoritmus. LU-felbontás Gauss-kiküszöböléssel részleges
főelemkiválasztással.}

\begin{verbatim}
for k = 1 to n − 1                           { ciklus az oszlopokra }
    Find index p such that                   { főelem keresése az
        |apk| ≥ |aik| for k ≤ i ≤ n              aktuális oszlopban }
    if p ≠ k then                            { sorok cseréje,
        interchange rows k and p                 ha szükséges }
    if akk = 0 then                          { az aktuális oszlop átugrása,
        continue with next k                     ha már eleve nulla }
    for i = k + 1 to n                       { szorzótényezők kiszámítása
        mik = aik/akk                            az aktuális oszlopra }
    end
    for j = k + 1 to n                       { transzformáció alkalmazása
        for i = k + 1 to n                       a maradék részmátrixra }
            aij = aij − mik akj
        end
    end
end
\end{verbatim}

Megjegyezzük, hogy a

\[P = P_{n-1} \cdots P_1\]

permutációs mátrix a részleges főelemkiválasztás által meghatározott
sorrendbe rendezi \(A\) sorait. Egy másik értelmezés szerint tehát a
részleges főelemkiválasztásra úgy is tekinthetünk, mint egy olyan
sorrend meghatározásának módjára, amelyben a numerikus stabilitás
érdekében nem volna szükség semmilyen cserére (persze egy ilyen
sorrendet előre meghatározni nem lehet). Így a

\[PA = LU\]

felbontást kapjuk, ahol most \(L\) valóban alsó háromszögmátrix. Az
\(Ax = b\) lineáris rendszer megoldásához először az \(Ly = Pb\) alsó
háromszögű rendszert előrehelyettesítéssel, majd az \(Ux = y\) felső
háromszögű rendszert visszahelyettesítéssel oldjuk meg.

\textbf{2.16. Példa. Gauss-kiküszöbölés részleges főelemkiválasztással.}
A 2.13. Példában nem használtunk sorcseréket, és néhány szorzótényező
nagysága meghaladta az 1-et. Most illusztráció kedvéért megismételjük
ezt a példát, ezúttal részleges főelemkiválasztást alkalmazva. A 2.13.
Példa rendszere

\[\boldsymbol{A}\boldsymbol{x} = \begin{bmatrix} 1 & 2 & 2 \\ 4 & 4 & 2 \\ 4 & 6 & 4 \end{bmatrix} \begin{bmatrix} x_1 \\ x_2 \\ x_3 \end{bmatrix} = \begin{bmatrix} 3 \\ 6 \\ 10 \end{bmatrix} = \boldsymbol{b}.\]

Az első oszlop legnagyobb eleme 4, ezért felcseréljük az első két sort a

\[\mathbf{P}_1 = \begin{bmatrix} 0 & 1 & 0 \\ 1 & 0 & 0 \\ 0 & 0 & 1 \end{bmatrix}\]

permutációs mátrixszal, így a permutált rendszer

\[\boldsymbol{P}_{1}\boldsymbol{A}\boldsymbol{x} = \begin{bmatrix} 4 & 4 & 2 \\ 1 & 2 & 2 \\ 4 & 6 & 4 \end{bmatrix} \begin{bmatrix} x_{1} \\ x_{2} \\ x_{3} \end{bmatrix} = \begin{bmatrix} 6 \\ 3 \\ 10 \end{bmatrix} = \boldsymbol{P}_{1}\boldsymbol{b}.\]

Az első oszlop főátló alatti elemeinek annullálásához a

\[\mathbf{M}_1 = \begin{bmatrix} 1 & 0 & 0 \\ -0{,}25 & 1 & 0 \\ -1 & 0 & 1 \end{bmatrix}\]

eliminációs mátrixot használjuk, így a transzformált rendszer

\[\boldsymbol{M}_1\boldsymbol{P}_1\boldsymbol{A}\boldsymbol{x} = \begin{bmatrix} 4 & 4 & 2 \\ 0 & 1 & 1{,}5 \\ 0 & 2 & 2 \end{bmatrix} \begin{bmatrix} x_1 \\ x_2 \\ x_3 \end{bmatrix} = \begin{bmatrix} 6 \\ 1{,}5 \\ 4 \end{bmatrix} = \boldsymbol{M}_1\boldsymbol{P}_1\boldsymbol{b}.\]

A második oszlopban az átlón vagy az alatt található legnagyobb elem 2,
ezért felcseréljük az utolsó két sort a

\[P_2 = \begin{bmatrix} 1 & 0 & 0 \\ 0 & 0 & 1 \\ 0 & 1 & 0 \end{bmatrix}\]

permutációs mátrixszal, így a permutált rendszer

\[P_2 M_1 P_1 A x = \begin{bmatrix} 4 & 4 & 2 \\ 0 & 2 & 2 \\ 0 & 1 & 1{,}5 \end{bmatrix} \begin{bmatrix} x_1 \\ x_2 \\ x_3 \end{bmatrix} = \begin{bmatrix} 6 \\ 4 \\ 1{,}5 \end{bmatrix} = P_2 M_1 P_1 b.\]

A második oszlop főátló alatti elemének annullálásához a

\[\mathbf{M}_2 = \begin{bmatrix} 1 & 0 & 0 \\ 0 & 1 & 0 \\ 0 & -0{,}5 & 1 \end{bmatrix}\]

eliminációs mátrixot használjuk, így a transzformált rendszer

\[\boldsymbol{M}_2\boldsymbol{P}_2\boldsymbol{M}_1\boldsymbol{P}_1\boldsymbol{A}\boldsymbol{x} = \begin{bmatrix} 4 & 4 & 2 \\ 0 & 2 & 2 \\ 0 & 0 & 0{,}5 \end{bmatrix} \begin{bmatrix} x_1 \\ x_2 \\ x_3 \end{bmatrix} = \begin{bmatrix} 6 \\ 4 \\ -0{,}5 \end{bmatrix} = \boldsymbol{M}_2\boldsymbol{P}_2\boldsymbol{M}_1\boldsymbol{P}_1\boldsymbol{b}.\]

Ezzel az eredeti rendszert egy ekvivalens felső háromszögű rendszerre
redukáltuk, amelyet most visszahelyettesítéssel oldhatunk meg, és
ugyanazt a megoldást kapjuk, mint korábban:
\(x = \begin{bmatrix} -1 & 3 & -1 \end{bmatrix}^T\).

Az LU-felbontás explicit felírásához

\[L = M^{-1} = (M_2 P_2 M_1 P_1)^{-1} = P_1^T L_1 P_2^T L_2 =\]

\[\begin{bmatrix} 0 & 1 & 0 \\ 1 & 0 & 0 \\ 0 & 0 & 1 \end{bmatrix} \begin{bmatrix} 1 & 0 & 0 \\ 0{,}25 & 1 & 0 \\ 1 & 0 & 1 \end{bmatrix} \begin{bmatrix} 1 & 0 & 0 \\ 0 & 0 & 1 \\ 0 & 1 & 0 \end{bmatrix} \begin{bmatrix} 1 & 0 & 0 \\ 0 & 1 & 0 \\ 0 & 0{,}5 & 1 \end{bmatrix} = \begin{bmatrix} 0{,}25 & 0{,}5 & 1 \\ 1 & 0 & 0 \\ 1 & 1 & 0 \end{bmatrix},\]

és ennélfogva

\[\mathbf{A} = \begin{bmatrix} 1 & 2 & 2 \\ 4 & 4 & 2 \\ 4 & 6 & 4 \end{bmatrix} = \begin{bmatrix} 0{,}25 & 0{,}5 & 1 \\ 1 & 0 & 0 \\ 1 & 1 & 0 \end{bmatrix} \begin{bmatrix} 4 & 4 & 2 \\ 0 & 2 & 2 \\ 0 & 0 & 0{,}5 \end{bmatrix} = \mathbf{L}\mathbf{U}.\]

Vegyük észre, hogy \(L\) nem alsó háromszögű, de az általánosabb
értelemben háromszögű (egy alsó háromszögmátrix permutációja).
Alternatívaként választhatjuk a

\[\boldsymbol{P} = \boldsymbol{P}_2 \boldsymbol{P}_1 = \begin{bmatrix} 1 & 0 & 0 \\ 0 & 0 & 1 \\ 0 & 1 & 0 \end{bmatrix} \begin{bmatrix} 0 & 1 & 0 \\ 1 & 0 & 0 \\ 0 & 0 & 1 \end{bmatrix} = \begin{bmatrix} 0 & 1 & 0 \\ 0 & 0 & 1 \\ 1 & 0 & 0 \end{bmatrix},\]

és

\[\boldsymbol{L} = \begin{bmatrix} 1 & 0 & 0 \\ 1 & 1 & 0 \\ 0{,}25 & 0{,}5 & 1 \end{bmatrix}\]

megválasztást, így

\[\boldsymbol{P}\boldsymbol{A} = \begin{bmatrix} 0 & 1 & 0 \\ 0 & 0 & 1 \\ 1 & 0 & 0 \end{bmatrix} \begin{bmatrix} 1 & 2 & 2 \\ 4 & 4 & 2 \\ 4 & 6 & 4 \end{bmatrix} = \begin{bmatrix} 1 & 0 & 0 \\ 1 & 1 & 0 \\ 0{,}25 & 0{,}5 & 1 \end{bmatrix} \begin{bmatrix} 4 & 4 & 2 \\ 0 & 2 & 2 \\ 0 & 0 & 0{,}5 \end{bmatrix} = \boldsymbol{L}\boldsymbol{U},\]

ahol most \(L\) valóban alsó háromszögmátrix, de \(A\) permutálva van.

A „részleges" főelemkiválasztás elnevezés abból a tényből ered, hogy
csak az aktuális oszlopban keresünk alkalmas főelemet. Egy kimerítőbb
főelemkiválasztási stratégia a \emph{teljes főelemkiválasztás}, amelyben
a teljes maradék, redukálatlan részmátrixban keressük a legnagyobb
elemet, majd azt permutáljuk az átlós főelempozícióba. Vegyük észre,
hogy ez a sorok cseréje mellett az oszlopok cseréjét is igényli, és
ezért a

\[PAQ = LU\]

alakú felbontáshoz vezet, ahol \(L\) egységnyi átlójú alsó
háromszögmátrix, \(U\) felső háromszögmátrix, \(P\) és \(Q\) pedig
permutációs mátrixok, amelyek rendre \(A\) sorait és oszlopait rendezik
át. Az \(Ax = b\) lineáris rendszer megoldásához először az \(Ly = Pb\)
alsó háromszögű rendszert előrehelyettesítéssel, majd az \(Uz = y\)
felső háromszögű rendszert visszahelyettesítéssel oldjuk meg, végül
pedig a megoldás komponenseit permutáljuk, hogy megkapjuk \(x = Qz\)-t.
Bár a teljes főelemkiválasztás numerikus stabilitása elméletben jobb,
sokkal költségesebb főelem-keresést igényel, mint a részleges
főelemkiválasztás. Mivel a részleges főelemkiválasztás numerikus
stabilitása a gyakorlatban bőven elegendő, általános lineáris rendszerek
Gauss-kiküszöböléssel történő megoldásában szinte univerzálisan ezt
használják.

A főelem kiválasztása az egyes mátrixelemek nagyságától függ, így a
konkrét választás nyilvánvalóan függ a mátrix skálázásától. A mátrix
diagonális skálázása (emlékezzünk a 2.10. Példára) eltérő
főelem-sorrendhez vezethet. Például egy adott oszlop bármelyik nemnulla
eleme a legnagyobbá tehető abszolút értékben, ha az adott sornak kellően
nagy súlyt adunk. Ez azonban nem jelenti azt, hogy tetszőleges
főelem-sorrend elfogadható: egy rosszul eltorzított skálázás rosszul
kondicionált rendszerhez és ennek megfelelően pontatlan megoldáshoz
vezethet (lásd a 2.20. Példát). Egy jól megfogalmazott feladatban az
ismeretlen változók mérésére megfelelően arányos mértékegységeket kell
használni (oszloponkénti skálázás), és az egyes egyenleteknek a
viszonylagos fontosságukat helyesen tükröző súlyozást kell kapniuk
(sorkénti skálázás). Figyelembe kell venni a bemeneti adatok
viszonylagos pontosságát is. Ezek között a körülmények között a
főelemkiválasztási eljárás általában olyan megoldást produkál, amely
annyira pontos, amennyire a feladat indokolja (lásd a 2.3.4. szakaszt).

A 2.3.5. szakaszban láttuk, hogy egy kiszámított megoldás relatív
maradéka kielégíti a

\[\frac{\|\boldsymbol{r}\|}{\|\boldsymbol{A}\|\cdot\|\hat{\boldsymbol{x}}\|} \leq \frac{\|\boldsymbol{E}\|}{\|\boldsymbol{A}\|}\]

egyenlőtlenséget, ahol \(E\) az \(A\) mátrixhoz tartozó hátraható hiba.
De mekkora is \(\|E\|\) valójában a gyakorlatban? Wilkinson {[}499{]}
megmutatta, hogy a Gauss-kiküszöböléssel végzett LU-felbontás esetén

\[\frac{\|\boldsymbol{E}\|}{\|\boldsymbol{A}\|} \le \rho \; n^2 \; \epsilon_{\text{mach}}\]

alakú korlát érvényes, ahol \(\rho\) --- amelyet \emph{növekedési
tényezőnek} nevezünk --- \(U\) legnagyobb abszolút értékű elemének és
\(A\) legnagyobb abszolút értékű elemének hányadosa (technikailag a
növekedési tényező a felbontási folyamat \emph{bármely} lépésében
előállított legnagyobb elemtől függ, de ez rendszerint az utolsó, azaz
\(U\)). Főelemkiválasztás nélkül \(\rho\) tetszőlegesen nagy lehet, és
ezért a főelemkiválasztás nélküli Gauss-kiküszöbölés instabil, amint azt
már láttuk. Részleges főelemkiválasztás mellett a növekedési tényező
legrosszabb esetben elérheti a \(2^{n-1}\)-t (mivel legrosszabb esetben
az elemek mérete az elimináció minden lépésében megduplázódhat), de az
ilyen viselkedés rendkívül ritka. A gyakorlatban kicsi vagy egyáltalán
nincs növekedés, és egy realisztikus korlátot a

\[\frac{\|\boldsymbol{E}\|}{\|\boldsymbol{A}\|} \lessapprox n \ \epsilon_{\mathrm{mach}}\]

reláció ad. Ez a reláció azt jelenti, hogy egy lineáris rendszernek
részleges főelemkiválasztásos Gauss-kiküszöböléssel, majd
visszahelyettesítéssel való megoldása szinte mindig nagyon kicsi relatív
maradékot eredményez, függetlenül attól, mennyire rosszul kondicionált a
rendszer. Így egy kicsi relatív maradék önmagában nem feltétlenül
jelenti, hogy a kiszámított megoldás pontos, hacsak a rendszer nem jól
kondicionált. A teljes főelemkiválasztás még kisebb növekedési tényezőt
eredményez, mind elméletben, mind a gyakorlatban, de az általa nyújtott
további stabilitási tartalék rendszerint nem éri meg az extra költséget.

\textbf{2.17. Példa. Kis maradék.} Tekintsük a

\[\boldsymbol{A}\boldsymbol{x} = \begin{bmatrix} 0{,}913 & 0{,}659 \\ 0{,}457 & 0{,}330 \end{bmatrix} \begin{bmatrix} x_1 \\ x_2 \end{bmatrix} = \begin{bmatrix} 0{,}254 \\ 0{,}127 \end{bmatrix} = \boldsymbol{b}\]

lineáris rendszert a 2.8. Példából. Négyjegyű decimális aritmetika
használatával a Gauss-kiküszöbölés a

\[\begin{bmatrix} 0{,}9130 & 0{,}6590 \\ 0 & 0{,}0002 \end{bmatrix} \begin{bmatrix} x_1 \\ x_2 \end{bmatrix} = \begin{bmatrix} 0{,}2540 \\ -0{,}0001 \end{bmatrix}\]

háromszögű rendszert adja, majd a visszahelyettesítés az

\[\hat{\boldsymbol{x}} = \begin{bmatrix} 0{,}6391 \\ -0{,}5 \end{bmatrix}\]

kiszámított megoldást szolgáltatja. Amint a 2.8. Példában láttuk, ennek
a megoldásnak a maradéknormája
\(\|\boldsymbol{r}\|_1 = 7{,}04 \times 10^{-5}\), ami olyan kicsi,
amilyenre négyjegyű aritmetikával számíthatunk. Mégis a rendszer pontos
megoldása könnyen ellenőrizhetően \(\boldsymbol{x} = [1, -1]^T\), így a
hiba ugyanakkora, mint a megoldás. Ennek a jelenségnek az oka az, hogy
az \(\boldsymbol{A}\) mátrix közel szinguláris: amint a 2.7. Példában
láttuk, kondíciószáma nagyobb, mint \(10^4\). Az \(x_2\)-t meghatározó
osztás két olyan mennyiség között történik, amelyek mindegyike a
kerekítési hiba nagyságrendjébe esik (négyjegyű aritmetikában), és ezért
az eredmény lényegében tetszőleges. Mégis, konstrukció szerint, amikor
ezt az \(x_2\)-re vonatkozó tetszőleges értéket ezután behelyettesítjük
az első egyenletbe, egy olyan \(x_1\) érték adódik, amelyre az első
egyenlet teljesül. Így kicsi maradékot, de rossz megoldást kapunk.

Amint az imént láttuk, a Gauss-kiküszöbölés stabilitásához általában
szükség van a főelemkiválasztásra. Vannak azonban olyan mátrixosztályok,
amelyekre a Gauss-kiküszöbölés főelemkiválasztás nélkül is stabil.
Például ha az \(\boldsymbol{A}\) mátrix \emph{oszloponként diagonálisan
domináns}, ami azt jelenti, hogy minden átlóbeli elem abszolút értéke
nagyobb, mint az oszlopa többi eleme abszolút értékeinek összege,

\[\sum_{i=1, i \neq j}^{n} |a_{ij}| < |a_{jj}|, \quad j = 1, \dots, n,\]

akkor az LU-felbontás Gauss-kiküszöböléssel történő kiszámításához nincs
szükség főelemkiválasztásra. Ha ilyen mátrixra részleges
főelemkiválasztást alkalmazunk, akkor ténylegesen nem történik sorcsere.
Egy másik fontos osztály, amelyre nincs szükség főelemkiválasztásra, a
szimmetrikus és pozitív definit mátrixok, amelyeket a 2.5. szakaszban
fogunk definiálni. A felesleges főelem-keresések elkerülése jelentős
időmegtakarítást jelenthet a felbontás kiszámításában.

\subsubsection{2.4.6 A Gauss-kiküszöbölés
megvalósítása}\label{a-gauss-kikuxfcszuxf6buxf6luxe9s-megvaluxf3suxedtuxe1sa}

A Gauss-kiküszöbölés, azaz az LU-felbontás, általános alakja egy
háromszorosan egymásba ágyazott ciklus, amint azt a 2.5. Algoritmus
sematikusan mutatja. A \texttt{for} ciklusok \(i\), \(j\) és \(k\)
indexeit tetszőleges sorrendben használhatjuk, így a ciklusok összesen
\(3! = 6\)-féleképpen rendezhetők el. Néhány jelzett aritmetikai művelet
nagyobb hatékonyság érdekében kivihető a legbelső ciklusból (mint
például a 2.3. Algoritmusban), és a műveletek olyan további átrendezései
is lehetségesek, amelyek esetleg nem szigorúan egymásba ágyazott
ciklusokat eredményeznek. Az alapalgoritmus ezen változatainak eltérő
hozzáférési mintáik (például sorkénti vagy oszloponkénti) vannak, és
eltérőek abban, mennyire képesek kihasználni egy adott számítógép
architekturális jellemzőit (például gyorsítótár, lapozás, vektorizáció,
több processzor). Így teljesítményük adott számítógépen vagy különböző
számítógépeken nagyon eltérő lehet, és nincs olyan egyetlen elrendezés,
amely egyöntetűen felsőbbrendű volna.

\textbf{2.5. Algoritmus. Általános Gauss-kiküszöbölés.}

\begin{verbatim}
for ______

    for ______

        for ______

            aij = aij − (aik/akk) akj
        end
    end
end
\end{verbatim}

Számos megvalósítási részlet ehhez hasonló módon változó lehet. Például
az általunk leírt részleges főelemkiválasztási eljárás oszlopok mentén
keres és sorokat cserél, de alternatívaként keresni lehetne sorok mentén
és oszlopokat cserélni. Úgy is vettük, hogy \(\boldsymbol{L}\)-nek
egységnyi átlója van, de intézhetnénk úgy is, hogy
\(\boldsymbol{U}\)-nak legyen egységnyi átlója. A Gauss-kiküszöbölés
ezen változatai közül némelyiket olyan fontosnak tartják, hogy nevet is
kaptak, mint például a Crout- és Doolittle-módszer.

Bár a Gauss-kiküszöbölés sokféle lehetséges változata drámai hatással
lehet a teljesítményre, nemszinguláris \(A\) mátrix esetén mind
lényegében ugyanazt a felbontást produkálják. Ha a sorok
főelem-kiválasztási sorrendje ugyanaz, akkor két LU-felbontás,
\(PA = LU = \hat{L}\hat{U}\), esetén ez a kifejezés azt jelenti, hogy
\(\hat{L}^{-1}L = \hat{U}U^{-1} = D\) egyszerre alsó és felső
háromszögű, és ennélfogva diagonális. Ha mind \(L\)-ről, mind
\(\hat{L}\)-ről feltesszük, hogy egységnyi átlójú alsó háromszögmátrix,
akkor \(D\) valójában az \(I\) egységmátrix kell legyen, és így
\(L = \hat{L}\) és \(U = \hat{U}\), vagyis a felbontás egyértelmű. E
feltevés nélkül azonban még mindig következtethetünk arra, hogy az
LU-felbontás a tényezők diagonális skálázása erejéig egyértelmű. Ezt az
egyértelműséget az LDU-felbontás, \(PA = LDU\), teszi explicitté, ahol
\(L\) egységnyi átlójú alsó háromszögmátrix, \(U\) egységnyi átlójú
felső háromszögmátrix, \(D\) pedig diagonális.

A tárkezelés egy másik fontos megvalósítási kérdés. A tárgyalt számos
mátrix --- az \(M_k\) elemi eliminációs mátrixok, azok \(L_k\) inverzei
és a \(P_k\) permutációs mátrixok --- csupán formálisan írja le a
felbontási folyamatot. A gyakorlati megvalósításban nem képezzük őket
explicit módon. A tárhely megtakarítása érdekében az \(L\) és \(U\)
tényezők felülírják a bemeneti \(A\) mátrix kezdeti tárterületét: a
transzformált \(U\) mátrix \(A\) felső háromszögét (az átlót is
beleértve) foglalja el, a szigorú alsó háromszögű \(L\)-t alkotó
szorzótényezők pedig \(A\) (most már nulla) szigorú alsó háromszögét.
\(L\) egységnyi átlóját nem kell tárolni.

Az adatmozgatás minimalizálása érdekében a főelemkiválasztás által
megkövetelt sorcseréket rendszerint nem végezzük el ténylegesen.
Ehelyett a sorok eredeti helyükön maradnak, és egy segéd egészértékű
vektor nyilvántartja az új sorsorrendet. Megjegyzendő, hogy egyetlen
ilyen vektor elegendő, mivel az összes csere nettó hatása még mindig
csak az \(1, \ldots, n\) egészek egy permutációja.

\subsubsection{2.4.7 A lineáris egyenletrendszerek megoldásának
komplexitása}\label{a-lineuxe1ris-egyenletrendszerek-megolduxe1suxe1nak-komplexituxe1sa}

Egy \(n \times n\)-es mátrix LU-felbontásának Gauss-kiküszöböléssel való
kiszámítása körülbelül \(n^3/3\) lebegőpontos szorzást és hasonló számú
összeadást igényel. Az eredményül kapott háromszögű rendszer egyetlen
jobb oldali vektorra való megoldása előre- és visszahelyettesítéssel
körülbelül \(n^2\) szorzást és hasonló számú összeadást igényel. Így
ahogy a mátrix \(n\) mérete nő, a lineáris rendszerek megoldásának
költségében az LU-felbontás fázisa egyre dominánsabbá válik.

Egy lineáris rendszert a mátrix explicit invertálásával is megoldhatunk,
úgy, hogy a megoldást \(x = A^{-1}b\) adja. De \(A^{-1}\) kiszámítása
egyenértékű \(n\) darab lineáris rendszer megoldásával: \(A\) egy
LU-felbontását, majd \(n\) darab előre- és visszahelyettesítést igényel,
egyet az egységmátrix minden oszlopára. A teljes műveletszám körülbelül
\(n^3\) szorzás és hasonló számú összeadás (kihasználva a jobb oldali
vektorok nullait az előrehelyettesítésnél). Az explicit invertálás tehát
háromszor költségesebb, mint az LU-felbontás.

Az azt követő \(x = A^{-1}b\) mátrix-vektor szorzás, amellyel a lineáris
rendszert megoldjuk, körülbelül \(n^2\) szorzást és hasonló számú
összeadást igényel, ami hasonló az előre- és visszahelyettesítés
összköltségéhez. Így, még több jobb oldali vektor esetén is, a
mátrixinverzió költségesebb, mint az LU-felbontás, ha lineáris
rendszereket szeretnénk megoldani. Ezen felül az explicit invertálás
pontatlanabb választ ad. Egyszerű példaként, ha a \(3x = 18\)
\(1 \times 1\)-es lineáris rendszert osztással oldjuk meg, akkor
\(x = 18/3 = 6\)-ot kapunk, de az explicit invertálás
\(x = 3^{-1} \times 18 = 0{,}333 \times 18 = 5{,}99\)-et adna háromjegyű
aritmetikát használva. Ebben a kis példában az invertáláshoz egy plusz
aritmetikai műveletre van szükség, és kevésbé pontos eredményt ad. Az
invertálás ezen hátrányai csak súlyosbodnak, ahogy a rendszer mérete nő.

Az explicit mátrixinverzek kényelmes jelölésként gyakran előfordulnak
különféle képletekben, de ez a gyakorlat nem jelenti azt, hogy egy ilyen
képlet megvalósításához explicit inverzre volna szükség. Csupán
megfelelő jobb oldallal --- amely maga is lehet mátrix --- egy lineáris
rendszert kell megoldani. Így például egy \(A^{-1}B\) alakú szorzatot
\(A\) LU-felbontásával, majd \(B\) minden oszlopát használva elvégzett
előre- és visszahelyettesítésekkel kell kiszámítani. A gyakorlatban
rendkívül ritka, hogy tényleg explicit mátrixinverzre volna szükség, így
amikor csak egy mátrixinverzet látsz egy képletben, mindig arra gondolj:
„oldj meg egy rendszert", ne arra, hogy „invertálj egy mátrixot".

Egy másik lineáris rendszer megoldási módszer, amelyet kerülni kell, a
Cramer-szabály, amelyben a megoldás minden komponensét determinánsok
hányadosaként számítjuk ki. Bár elemi lineáris algebrai kurzusokon
gyakran tanítják, ez a módszer nem triviális méretű teli mátrixok esetén
csillagászatilag drága. A Cramer-szabály főként elméleti eszközként
hasznos.

\subsubsection{2.4.8
Gauss--Jordan-kiküszöbölés}\label{gaussjordan-kikuxfcszuxf6buxf6luxe9s}

A Gauss-kiküszöbölés mögötti motiváció az, hogy egy általános mátrixot
háromszögalakra redukáljunk, mivel az így kapott lineáris rendszer
könnyen megoldható. A diagonális lineáris rendszerek azonban még
könnyebben megoldhatók, így a diagonális alak még kívánatosabb célnak
tűnik. A \emph{Gauss--Jordan-kiküszöbölés} a standard Gauss-kiküszöbölés
olyan változata, amelyben a mátrixot nemcsak háromszög-, hanem
diagonális alakra redukáljuk. A mátrixelemek kiküszöbölésére ugyanolyan
jellegű sorkombinációkat használunk, mint a standard
Gauss-kiküszöbölésben, de ezeket mind az átló alatti, mind az átló
feletti elemek annullálására alkalmazzuk. Így egy adott
\(\boldsymbol{a}\) oszlopvektorra alkalmazott eliminációs mátrix alakja

\[\begin{bmatrix} 1 & \cdots & 0 & -m_1 & 0 & \cdots & 0 \\ \vdots & \ddots & \vdots & \vdots & \vdots & \ddots & \vdots \\ 0 & \cdots & 1 & -m_{k-1} & 0 & \cdots & 0 \\ 0 & \cdots & 0 & 1 & 0 & \cdots & 0 \\ 0 & \cdots & 0 & -m_{k+1} & 1 & \cdots & 0 \\ \vdots & \ddots & \vdots & \vdots & \vdots & \ddots & \vdots \\ 0 & \cdots & 0 & -m_n & 0 & \cdots & 1 \end{bmatrix} \begin{bmatrix} a_1 \\ \vdots \\ a_{k-1} \\ a_k \\ a_{k+1} \\ \vdots \\ a_n \end{bmatrix} = \begin{bmatrix} 0 \\ \vdots \\ 0 \\ a_k \\ 0 \\ \vdots \\ 0 \end{bmatrix},\]

ahol \(m_i = a_i/a_k\), \(i = 1, \dots, n\). Ez az eljárás körülbelül
\(n^3/2\) szorzást és hasonló számú összeadást igényel, ami 50
százalékkal drágább, mint a standard Gauss-kiküszöbölés.

A kiküszöbölési fázis során ugyanezeket a sorműveleteket alkalmazzuk a
lineáris egyenletrendszer jobb oldali vektorára (vagy vektoraira) is. Ha
egyszer a kiküszöbölési fázis befejeződött, és a mátrix diagonális
alakba került, akkor a lineáris rendszer megoldásának komponensei
egyszerűen úgy számíthatók ki, hogy a transzformált jobb oldal minden
elemét a mátrix megfelelő átlóbeli elemével elosztjuk. Ez a számítás
összesen csak \(n\) osztást igényel, ami jelentősen olcsóbb, mint egy
háromszögű rendszer megoldása, de nem elég a költségesebb kiküszöbölési
fázis kompenzálására. A Gauss--Jordan-kiküszöbölésnek az a numerikus
hátránya is megvan, hogy a szorzótényezők nagysága még főelemkiválasztás
esetén is meghaladhatja az 1-et.

Nagyobb összköltsége ellenére bizonyos helyzetekben a
Gauss--Jordan-kiküszöbölés előnyben részesíthető, mert végső megoldási
fázisa rendkívül egyszerű. Például néha párhuzamos számítógépekre való
megvalósításhoz ajánlják, mivel a felbontási fázis alatt egyenletes a
munkaterhelés, és a megoldás minden komponense egyszerre számítható ki
--- nem pedig egyesével, ahogyan a közönséges visszahelyettesítésnél.

A Gauss--Jordan-kiküszöbölést néha egy mátrix inverzének explicit
kiszámítására is használják, ha arra szükség van. Ha a jobb oldali
mátrixot az \(I\) egységmátrixra inicializáljuk, és az adott \(A\)
mátrixot Gauss--Jordan-kiküszöböléssel az egységmátrixra redukáljuk,
akkor a transzformált jobb oldali mátrix \(A\) inverze lesz. Az inverz
kiszámításához a Gauss--Jordan-kiküszöbölés körülbelül ugyanannyi
műveletet igényel, mint a Gauss-kiküszöböléssel végzett explicit
invertálás, amelyet előre- és visszahelyettesítés követ.

\textbf{2.18. Példa. Gauss--Jordan-kiküszöbölés.} A
Gauss--Jordan-kiküszöbölést annak használatával illusztráljuk, hogy
kiszámítjuk a 2.13. Példa mátrixának inverzét. Az egyszerűség kedvéért a
főelemkiválasztást kihagyjuk. Az \(\boldsymbol{A}\) mátrixszal kezdünk,
jobb oldalon az \(\boldsymbol{I}\) egységmátrixszal bővítve, és
ismételten eliminációs mátrixokat alkalmazunk \(\boldsymbol{A}\) átlón
kívüli elemeinek annullálására, amíg el nem érjük a diagonális alakot;
majd a megmaradt átlóbeli elemekkel skálázva a bal oldalon az
egységmátrixot, így a jobb oldalon az \(A^{-1}\) inverz mátrixot
állítjuk elő.

\[\begin{bmatrix} 1 & 0 & 0 \\ -4 & 1 & 0 \\ -4 & 0 & 1 \end{bmatrix} \begin{bmatrix} 1 & 2 & 2 & 1 & 0 & 0 \\ 4 & 4 & 2 & 0 & 1 & 0 \\ 4 & 6 & 4 & 0 & 0 & 1 \end{bmatrix} = \begin{bmatrix} 1 & 2 & 2 & 1 & 0 & 0 \\ 0 & -4 & -6 & -4 & 1 & 0 \\ 0 & -2 & -4 & -4 & 0 & 1 \end{bmatrix},\]

\[\begin{bmatrix} 1 & 0{,}5 & 0 \\ 0 & 1 & 0 \\ 0 & -0{,}5 & 1 \end{bmatrix} \begin{bmatrix} 1 & 2 & 2 & 1 & 0 & 0 \\ 0 & -4 & -6 & -4 & 1 & 0 \\ 0 & -2 & -4 & -4 & 0 & 1 \end{bmatrix} = \begin{bmatrix} 1 & 0 & -1 & -1 & 0{,}5 & 0 \\ 0 & -4 & -6 & -4 & 1 & 0 \\ 0 & 0 & -1 & -2 & -0{,}5 & 1 \end{bmatrix},\]

\[\begin{bmatrix} 1 & 0 & -1 \\ 0 & 1 & -6 \\ 0 & 0 & 1 \end{bmatrix} \begin{bmatrix} 1 & 0 & -1 & -1 & 0{,}5 & 0 \\ 0 & -4 & -6 & -4 & 1 & 0 \\ 0 & 0 & -1 & -2 & -0{,}5 & 1 \end{bmatrix} = \begin{bmatrix} 1 & 0 & 0 & 1 & 1 & -1 \\ 0 & -4 & 0 & 8 & 4 & -6 \\ 0 & 0 & -1 & -2 & -0{,}5 & 1 \end{bmatrix},\]

\[\begin{bmatrix} 1 & 0 & 0 \\ 0 & -0{,}25 & 0 \\ 0 & 0 & -1 \end{bmatrix} \begin{bmatrix} 1 & 0 & 0 & 1 & 1 & -1 \\ 0 & -4 & 0 & 8 & 4 & -6 \\ 0 & 0 & -1 & -2 & -0{,}5 & 1 \end{bmatrix} = \begin{bmatrix} 1 & 0 & 0 & 1 & 1 & -1 \\ 0 & 1 & 0 & -2 & -1 & 1{,}5 \\ 0 & 0 & 1 & 2 & 0{,}5 & -1 \end{bmatrix},\]

így

\[A^{-1} = \begin{bmatrix} 1 & 1 & -1 \\ -2 & -1 & 1{,}5 \\ 2 & 0{,}5 & -1 \end{bmatrix}.\]

\subsubsection{2.4.9 Módosított feladatok
megoldása}\label{muxf3dosuxedtott-feladatok-megolduxe1sa}

Számos gyakorlati helyzetben a lineáris rendszerek nem elszigetelten
fordulnak elő, hanem rokon feladatok olyan sorozatának részeként,
amelyek valamilyen szisztematikus módon változnak. Előfordulhat például,
hogy olyan \(Ax = b\) lineáris rendszerek sorozatát kell megoldani,
amelyeknek ugyanaz az \(A\) mátrixuk, de eltérőek a \(b\) jobb oldalaik.
Az első rendszer Gauss-kiküszöböléssel való megoldása után a már
kiszámított \(L\) és \(U\) tényezőket fel lehet használni a további
rendszereknek előre- és visszahelyettesítéssel való megoldására. A
felbontási fázist a későbbi lineáris rendszerek megoldása során nem kell
megismételni, hacsak a mátrix nem változik. Ez az eljárás jelentős
munkamegtakarítást jelent, mivel a további háromszögű megoldások
költsége csupán \(\mathcal{O}(n^2)\), ellentétben egy felbontás
\(\mathcal{O}(n^3)\) költségével.

Sőt: bizonyos fontos különleges esetekben még akkor is elkerülhető egy
új felbontás, ha a mátrix változik. Ilyen, gyakran előforduló eset egy
olyan \(n \times n\)-es mátrix hozzáadása vagy kivonása, amely két
nemnulla \(n\)-dimenziós vektor, \(\boldsymbol{u}\) és
\(\boldsymbol{v}\), \(\boldsymbol{u}\boldsymbol{v}^T\) külső szorzata.
Ezt \emph{egyrangú módosításnak} nevezzük, mert az
\(\boldsymbol{u}\boldsymbol{v}^T\) külső szorzatmátrix rangja egy (azaz
csak egy lineárisan független sora vagy oszlopa van), és bármely
egyrangú mátrix kifejezhető két vektor ilyen külső szorzataként (lásd a
2.25. Feladatot). Például ha az \(\boldsymbol{A}\) mátrix egyetlen eleme
változik meg --- például a \((j, k)\) elem \(a_{jk}\)-ról
\(\tilde{a}_{jk}\)-ra ---, akkor az új mátrix
\(\boldsymbol{A} - \alpha \boldsymbol{e}_j \boldsymbol{e}_k^T\), ahol
\(\boldsymbol{e}_j\) és \(\boldsymbol{e}_k\) az egységmátrix megfelelő
oszlopai, és \(\alpha = a_{jk} - \tilde{a}_{jk}\).

A Sherman--Morrison-formula,

\[(A - uv^T)^{-1} = A^{-1} + A^{-1}u (1 - v^T A^{-1}u)^{-1} v^T A^{-1},\]

ahol \(u\) és \(v\) \(n\)-dimenziós vektorok, megadja egy olyan mátrix
inverzét, amely egy olyan mátrix egyrangú módosításából származik,
amelynek inverze már ismert; ez közvetlen szorzással könnyen igazolható
(lásd a 2.27. Feladatot). Ennek a képletnek a kiértékelése csak
\(\mathcal{O}(n^2)\) munkát (mátrix-vektor szorzásokat) igényel, szemben
a módosított mátrix nulláról való invertálásához szükséges
\(\mathcal{O}(n^3)\) munkával.

Az \((A - uv^T)x = b\) új mátrixú lineáris rendszerre a
Sherman--Morrison-formula az

\[x = (A - uv^T)^{-1}b = A^{-1}b + A^{-1}u(1 - v^T A^{-1}u)^{-1} v^T A^{-1}b\]

megoldást adja, de kerüljük az explicit inverzeket. Ha van
LU-felbontásunk az eredeti \(A\) mátrixra, akkor a módosított rendszer
megoldása a 2.6. Algoritmussal számítható, amely háromszögű rendszerek
megoldását és vektorok belső szorzatának kiszámítását tartalmazza, így
nem igényel explicit inverzet, és csak \(\mathcal{O}(n^2)\) munkát
igényel. Megjegyzendő, hogy az első lépés \(b\)-től független, és ezért
nem kell megismételni, ha több \(b\) jobb oldali vektor van.

\textbf{2.6. Algoritmus. Megoldás egyrangú frissítése.}

\begin{verbatim}
Oldjuk meg Az = u-t z-re, így z = A^{-1} u
Oldjuk meg Ay = b-t y-ra, így y = A^{-1} b
x = y + ((v^T y) / (1 - v^T z)) z
\end{verbatim}

Hasonló technikákkal az inverz vagy a megoldás helyett a felbontást is
frissíthetjük. Ezeknek a frissítési képleteknek a használatakor azonban
óvatosnak kell lennünk, mert a mátrix változása során az egymást követő
frissítésekben általánosságban nincs garancia a numerikus stabilitásra.
A Woodbury-formula,

\[(A - UV^T)^{-1} = A^{-1} + A^{-1}U(I - V^T A^{-1} U)^{-1} V^T A^{-1},\]

ahol \(U\) és \(V\) \(n \times k\)-as mátrixok, általánosítja a
Sherman--Morrison-formulát a mátrix egy \(k\)-rangú módosítására (lásd a
2.28. Feladatot).

\textbf{2.19. Példa. Megoldás egyrangú frissítése.} Az egyrangú
frissítést a

\[\begin{bmatrix} 1 & 2 & 2 \\ 4 & 4 & 2 \\ 4 & 4 & 4 \end{bmatrix} \begin{bmatrix} x_1 \\ x_2 \\ x_3 \end{bmatrix} = \begin{bmatrix} 3 \\ 6 \\ 10 \end{bmatrix}\]

lineáris rendszer megoldásán illusztráljuk, amely a 2.13. Példa
rendszerének egyrangú módosítása, mivel csak az \(A\) mátrix \((3, 2)\)
eleme változott 6-ról 4-re. A frissítővektorok egyik választási
lehetősége

\[\boldsymbol{u} = \begin{bmatrix} 0 \\ 0 \\ 1 \end{bmatrix} \quad \text{és} \quad \boldsymbol{v} = \begin{bmatrix} 0 \\ 2 \\ 0 \end{bmatrix},\]

így a módosított rendszer mátrixa \(A - uv^T =\)

\[\begin{bmatrix} 1 & 2 & 2 \\ 4 & 4 & 2 \\ 4 & 6 & 4 \end{bmatrix} - \begin{bmatrix} 0 \\ 0 \\ 1 \end{bmatrix} \begin{bmatrix} 0 & 2 & 0 \end{bmatrix} = \begin{bmatrix} 1 & 2 & 2 \\ 4 & 4 & 2 \\ 4 & 6 & 4 \end{bmatrix} - \begin{bmatrix} 0 & 0 & 0 \\ 0 & 0 & 0 \\ 0 & 2 & 0 \end{bmatrix} = \begin{bmatrix} 1 & 2 & 2 \\ 4 & 4 & 2 \\ 4 & 4 & 4 \end{bmatrix},\]

a \(b\) jobb oldali vektor pedig nem változott.

A 2.13. Példában korábban kiszámított LU-felbontást felhasználva
megoldhatjuk \(Az = u\)-t, így
\(z = \begin{bmatrix} -1 & 1{,}5 & -1 \end{bmatrix}^T\), és \(Ay = b\)-t
már megoldottuk, így
\(y = \begin{bmatrix} -1 & 3 & -1 \end{bmatrix}^T\). Az utolsó lépés
ezek után a frissített megoldás kiszámítása:

\[\boldsymbol{x} = \boldsymbol{y} + \frac{\boldsymbol{v}^T \boldsymbol{y}}{1 - \boldsymbol{v}^T \boldsymbol{z}} \, \boldsymbol{z} = \begin{bmatrix} -1 \\ 3 \\ -1 \end{bmatrix} + \frac{6}{1 - 3} \begin{bmatrix} -1 \\ 1{,}5 \\ -1 \end{bmatrix} = \begin{bmatrix} 2 \\ -1{,}5 \\ 2 \end{bmatrix}.\]

Így a módosított rendszer megoldását anélkül számítottuk ki, hogy a
módosított mátrixot újra felbontottuk volna.

\subsubsection{2.4.10 A pontosság
javítása}\label{a-pontossuxe1g-javuxedtuxe1sa}

Bár egy lineáris rendszer megoldásában várható pontosság látszólag kőbe
vésett, bizonyos esetekben a pontosság javítható a rendszer
átskálázásával vagy a kezdetben kiszámított megoldás iteratív
javításával. Ezek az eszközök nem mindig vihetők véghez, de kipróbálásuk
megérheti, különösen erősen rosszul kondicionált rendszerek esetén.

A 2.10. Példából emlékezhetünk rá, hogy egy lineáris rendszer diagonális
skálázása vagy változatlanul hagyja a megoldást (sorkénti skálázás),
vagy úgy változtatja meg, hogy a megoldás könnyen visszanyerhető
(oszloponkénti skálázás). A gyakorlatban azonban a skálázás befolyásolja
a rendszer kondicionáltságát és a Gauss-kiküszöbölésben a főelemek
kiválasztását, amelyek mindegyike viszont a kiszámított megoldás
pontosságát befolyásolja. Így egy lineáris rendszer sorkénti és
oszloponkénti skálázása lehetőséget adhat a numerikus stabilitás és
pontosság javítására (vagy rontására).

A pontosság rendszerint javul, ha a mátrix minden eleme körülbelül
azonos nagyságrendű --- vagy ami még jobb, ha a mátrixelemekben lévő
bizonytalanságok is körülbelül azonos méretűek. Néha puszta ránézésre is
nyilvánvaló, hogyan skálázzuk a mátrixot úgy, hogy ilyen egyensúly
álljon elő: az egyes változókhoz tartozó mértékegységek megválasztásával
és az egyes egyenleteknek a viszonylagos fontosságuk és bizonytalanságuk
szerinti súlyozásával. Olyan általános, automatikus technikát azonban
még soha nem fejlesztettek ki, amely hatékony és megbízható módon
optimális skálázást eredményezne. Ráadásul maga a skálázási folyamat is
bevezethet kerekítési hibákat, hacsak nem vigyázunk (például ha csak az
aritmetikai alap hatványait használjuk skálázási tényezőként).

\textbf{2.20. Példa. Rossz skálázás.} Egyszerű példaként az

\[\begin{bmatrix} 1 & 0 \\ 0 & \epsilon \end{bmatrix} \begin{bmatrix} x_1 \\ x_2 \end{bmatrix} = \begin{bmatrix} 1 \\ \epsilon \end{bmatrix}\]

lineáris rendszer kondíciószáma \(1/\epsilon\), és ezért nagyon rosszul
kondicionált, ha \(\epsilon\) nagyon kicsi. Ez a rossz kondicionáltság
azt jelenti, hogy a bemeneti adatok kis perturbációi viszonylag nagy
változásokat okozhatnak a megoldásban. Például ha a jobb oldalt a
\(\begin{bmatrix} 0 & -\epsilon \end{bmatrix}^T\) vektorral
perturbáljuk, a megoldás \(\begin{bmatrix} 1 & 1 \end{bmatrix}^T\)-ról
\(\begin{bmatrix} 1 & 0 \end{bmatrix}^T\)-ra változik. Ha viszont a
második sort előbb \(1/\epsilon\)-nal szorozzuk, akkor a rendszer
tökéletesen jól kondicionálttá válik, és ugyanaz a perturbáció most a
megoldásban ezzel arányos kis változást eredményez. Így a látszólagos
rossz kondicionáltság pusztán a rossz skálázásnak volt köszönhető.
Sajnos az, hogy általános mátrixok esetén hogyan korrigáljuk a rossz
skálázást, sokkal kevésbé kézenfekvő.

Az \emph{iteratív finomítás} egy másik eszköz, amellyel egy kiszámított
megoldás pontossága esetleg javítható. Tegyük fel, hogy az \(Ax = b\)
lineáris rendszerre kiszámítottunk egy közelítő \(x_0\) megoldást ---
például valamilyen LU-felbontási alakkal ---, és kiszámítjuk a

\[\boldsymbol{r}_0 = \boldsymbol{b} - \boldsymbol{A}\boldsymbol{x}_0\]

maradékot. Természetesen azt szeretnénk, ha a maradék megfelelően kicsi
volna, ha viszont nem az, akkor a korábban kiszámított LU-tényezőkkel
megoldhatjuk az

\[As_0 = r_0\]

lineáris rendszert \(s_0\)-ra, és az

\[\boldsymbol{x}_1 = \boldsymbol{x}_0 + \boldsymbol{s}_0\]

értéket vesszük új, „jobb" közelítő megoldásnak, mivel

\[Ax_1 = A(x_0 + s_0) = Ax_0 + As_0 = (b - r_0) + r_0 = b.\]

Ez az eljárás ismételhető a megoldás egymás utáni finomítására,
konvergenciáig, esetleg olyan megoldást előállítva, amelynek maradéka
olyan kicsi, amennyire az alkalmazott aritmetikai pontosságon csak
lehet.

Sajnos az iteratív finomítás dupla tárhelyet igényel, mivel mind az
eredeti mátrixra, mind annak LU-felbontására szükség van (rendre a
maradék kiszámításához és a későbbi rendszerek megoldásához). Ráadásul
ahhoz, hogy az iteratív finomítás maximális haszonnal járjon, a
maradékot rendszerint a kezdeti megoldás kiszámításához használt
pontosságnál nagyobb pontossággal kell kiszámítani (emlékezzünk az 1.17.
Példára).

Ezen okok miatt az iteratív finomítás gyakran nem praktikus rutinszerű
használatra, bizonyos körülmények között azonban mégis hasznos lehet. Az
iteratív finomítással például visszanyerhető a teljes pontosság a
rosszul skálázott rendszerek esetén, és olykor stabilizálhatók vele az
egyébként potenciálisan instabil megoldási módszerek. Ironikus módon, ha
a kezdeti megoldás viszonylag rossz, akkor a maradék elegendően nagy
lehet ahhoz, hogy extra pontosság nélkül is kielégítő pontossággal
legyen kiszámítható. Az iteratív finomításra a 11.7. Példában térünk
vissza.

\subsection{2.5 Lineáris egyenletrendszerek speciális
típusai}\label{lineuxe1ris-egyenletrendszerek-speciuxe1lis-tuxedpusai}

Eddig azt feltételeztük, hogy a lineáris rendszer mátrixa általános és
\emph{sűrű}, azaz a mátrix bejegyzéseinek lényegében mindegyike
nemnulla. Ha a mátrix rendelkezik valamilyen speciális tulajdonsággal,
akkor a lineáris rendszer megoldásakor gyakran műveletet és tárhelyet
lehet megtakarítani. A kihasználható speciális tulajdonságokra néhány
példa a következő:

\begin{itemize}
\tightlist
\item
  Szimmetrikus: \(\mathbf{A} = \mathbf{A}^T\), azaz \(a_{ij} = a_{ji}\)
  minden \(i, j\) esetén.
\item
  Pozitív definit: \(\mathbf{x}^T \mathbf{A} \mathbf{x} > 0\) minden
  \(\mathbf{x} \neq \mathbf{0}\) vektorra.
\item
  Sávos: \(a_{ij} = 0\) minden \(|i - j| > \beta\) esetén, ahol
  \(\beta\) az \(\boldsymbol{A}\) mátrix sávszélessége. Fontos speciális
  eset a tridiagonális mátrix, amelyre \(\beta = 1\).
\item
  Ritka: az \(A\) bejegyzéseinek többsége nulla.
\end{itemize}

A szimmetrikus és sávos rendszerek kezelésének technikái a sűrű
rendszerekre alkalmazott Gauss-kiküszöbölés viszonylag egyszerű
változatai. Az általánosabb nemnulla-mintázatú ritka lineáris rendszerek
ezzel szemben olyan kifinomultabb algoritmusokat és adatszerkezeteket
igényelnek, amelyek elkerülik a mátrixban lévő nullák tárolását és a
velük végzett műveleteket (lásd a 11.4.1. szakaszt).

A valós mátrixokra az imént definiált tulajdonságoknak komplex mátrixok
esetén is megvannak a megfelelőik, de a komplex esetben a szokásos
mátrixtranszponált helyébe a konjugált transzponált lép, amelyet \(H\)
felső indexszel jelölünk. Ha \(z = \alpha + i\beta\) komplex szám, ahol
\(\alpha\) és \(\beta\) valós számok és \(i = \sqrt{-1}\), akkor a
komplex konjugáltját a \(\bar{z} = \alpha - i\beta\) definíció adja
(lásd az 1.3.11. szakaszt). Egy \(\boldsymbol{A}\) mátrix konjugált
transzponáltját ezek után az
\(\{\boldsymbol{A}^H\}_{ij} = \bar{a}_{ji}\) képlet adja meg.
Természetesen egy valós \(\boldsymbol{A}\) mátrixra
\(\boldsymbol{A}^H = \boldsymbol{A}^T\). Egy komplex mátrix
\emph{hermitikus}, ha \(\boldsymbol{A} = \boldsymbol{A}^H\), és pozitív
definit, ha \(\boldsymbol{x}^H \boldsymbol{A} \boldsymbol{x} > 0\)
minden \(\boldsymbol{x} \neq \boldsymbol{0}\) komplex vektorra.

\subsubsection{2.5.1 Szimmetrikus pozitív definit
rendszerek}\label{szimmetrikus-pozituxedv-definit-rendszerek}

Ha az \(\boldsymbol{A}\) mátrix szimmetrikus és pozitív definit, akkor
az LU-felbontás úgy is elrendezhető, hogy
\(\boldsymbol{U} = \boldsymbol{L}^T\) legyen, azaz
\(\boldsymbol{A} = \boldsymbol{L}\boldsymbol{L}^T\), ahol
\(\boldsymbol{L}\) alsó háromszögmátrix pozitív diagonális
bejegyzésekkel (de általában nem egységdiagonálissal). Ezt nevezik az
\(\boldsymbol{A}\) Cholesky-felbontásának, és egy algoritmus a
kiszámítására abból kapható, ha az \(\boldsymbol{A}\) és a
\(\boldsymbol{L}\boldsymbol{L}^T\) megfelelő bejegyzéseit egyenlővé
tesszük, majd az \(\boldsymbol{L}\) bejegyzéseit a helyes sorrendben
állítjuk elő. Például a \(2 \times 2\)-es esetben

\[\begin{bmatrix} a_{11} & a_{21} \\ a_{21} & a_{22} \end{bmatrix} = \begin{bmatrix} \ell_{11} & 0 \\ \ell_{21} & \ell_{22} \end{bmatrix} \begin{bmatrix} \ell_{11} & \ell_{21} \\ 0 & \ell_{22} \end{bmatrix},\]

amiből adódik, hogy

\[\ell_{11} = \sqrt{a_{11}}, \quad \ell_{21} = a_{21}/\ell_{11}, \quad \ell_{22} = \sqrt{a_{22} - \ell_{21}^2}.\]

Az ezáltal nyert általános eljárás egyik lehetséges megszervezése a 2.7.
algoritmus, amelyben a Cholesky-faktor \(\boldsymbol{L}\) felülírja az
eredeti \(\boldsymbol{A}\) mátrixot.

A Cholesky-felbontási algoritmusnak számos olyan tulajdonsága van, amely
nagyon vonzóvá és népszerűvé teszi szimmetrikus pozitív definit mátrixok
esetén:

\begin{itemize}
\tightlist
\item
  A szükséges \(n\) négyzetgyökvonás mindegyike pozitív számra történik,
  így az algoritmus jól definiált.
\item
  Numerikus stabilitás céljából nincs szükség főelemkiválasztásra.
\end{itemize}

\textbf{2.7. Algoritmus. Cholesky-felbontás.}

\begin{verbatim}
for k = 1 to n
    akk =
          √
            akk
    for i = k + 1 to n
        aik = aik/akk
    end
    for j = k + 1 to n
        for i = j to n
            aij = aij − aik · ajk
        end
    end
end
                                        { ciklus az oszlopokon }
                                        { az aktuális oszlop skálázása }
                                        { minden hátralévő oszlopból
                                             kivonjuk az aktuális oszlop
                                             egy megfelelő többszörösét }
\end{verbatim}

\begin{itemize}
\tightlist
\item
  Csak az \(A\) alsó háromszöge kerül elérésre, ezért a szigorú felső
  háromszög-részét nem kell tárolni.
\item
  Csak körülbelül \(n^3/6\) szorzásra és hozzávetőleg ugyanennyi
  összeadásra van szükség.
\end{itemize}

Így a Cholesky-felbontás csak körülbelül feleannyi műveletet és
feleannyi tárhelyet igényel, mint egy általános mátrix
Gauss-kiküszöböléssel való LU-felbontása. Sajnos ahhoz, hogy a
tárhelynek ezt a megtakarítását ki is használjuk, rendszerint szükség
van arra, hogy a szimmetrikus mátrix egyik háromszögét egy egydimenziós
tömbbe csomagoljuk, ami kényelmetlenebb, mint a mátrixok szokásos
kétdimenziós tárolása. Ezért a lineáris algebrai szoftvercsomagok
szimmetrikus mátrixokra általában mind a csomagolt tárolású, mind a
szabványos kétdimenziós tömbtárolású változatot felkínálják, hogy a
felhasználó a kényelem és a tárhelytakarékosság között választhasson.

Bizonyos körülmények között előnyös lehet a Cholesky-felbontást
\(\boldsymbol{A} = \boldsymbol{L}\boldsymbol{D}\boldsymbol{L}^T\)
alakban felírni, ahol \(\boldsymbol{L}\) egységdiagonálisú alsó
háromszögmátrix, \(\boldsymbol{D}\) pedig pozitív diagonális bejegyzésű
diagonális mátrix. Egy ilyen felbontás a szabványos Cholesky-algoritmus
egy egyszerű változatával számítható ki, és megvan az az előnye, hogy
nem igényel négyzetgyökvonást. A
\(\boldsymbol{L}\boldsymbol{D}\boldsymbol{L}^T\)-felbontásban a
\(\boldsymbol{D}\) diagonális bejegyzései egyszerűen a
\(\boldsymbol{L}\boldsymbol{L}^T\)-felbontásban szereplő
\(\boldsymbol{L}\) diagonális bejegyzéseinek négyzetei.

\textbf{2.21. Példa. Cholesky-felbontás.} Az algoritmus szemléltetésére
kiszámítjuk a szimmetrikus pozitív definit

\[\mathbf{A} = \begin{bmatrix} 3 & -1 & -1 \\ -1 & 3 & -1 \\ -1 & -1 & 3 \end{bmatrix}\]

mátrix Cholesky-felbontását. Mivel az algoritmus csak a mátrix alsó
háromszögét érinti, a mátrix alsó háromszögének egymást követő
transzformációit mutatjuk be. Az első oszlopot a diagonális
bejegyzésének négyzetgyökével, \(\sqrt{3} \approx 1{,}7321\)-gyel
elosztva

\[\begin{bmatrix} 1{,}7321 \\ -0{,}5774 & 3 \\ -0{,}5774 & -1 & 3 \end{bmatrix}\]

adódik. A második oszlopot úgy módosítjuk, hogy kivonjuk belőle a
\((2, 1)\) bejegyzésnek, \(-0{,}5774\)-nek, az első oszlop megfelelő
részével vett szorzatát, a harmadik oszlopot pedig úgy, hogy kivonjuk
belőle a \((3, 1)\) bejegyzésnek, szintén \(-0{,}5774\)-nek, az első
oszlop megfelelő részével vett szorzatát, ami

\[\begin{bmatrix} 1{,}7321 \\ -0{,}5774 & 2{,}6667 \\ -0{,}5774 & -1{,}3333 & 2{,}6667 \end{bmatrix}\]

eredményre vezet. A második oszlopot ezután elosztjuk a diagonális
bejegyzésének négyzetgyökével, \(\sqrt{2{,}6667} \approx 1{,}6330\)-cal,
amivel

\[\begin{bmatrix} 1{,}7321 \\ -0{,}5774 & 1{,}6330 \\ -0{,}5774 & -0{,}8165 & 2{,}6667 \end{bmatrix}\]

adódik. A harmadik oszlopot úgy módosítjuk, hogy kivonjuk belőle a
\((3, 2)\) bejegyzésnek, \(-0{,}8165\)-nek, a második oszlop megfelelő
részével vett szorzatát, ami

\[\begin{bmatrix} 1{,}7321 \\ -0{,}5774 & 1{,}6330 \\ -0{,}5774 & -0{,}8165 & 2{,}0000 \end{bmatrix}\]

eredményre vezet. A harmadik diagonális bejegyzés négyzetgyökét véve a
végeredmény

\[\mathbf{L} = \begin{bmatrix} 1{,}7321 \\ -0{,}5774 & 1{,}6330 \\ -0{,}5774 & -0{,}8165 & 1{,}4142 \end{bmatrix}.\]

\subsubsection{2.5.2 Szimmetrikus indefinit
rendszerek}\label{szimmetrikus-indefinit-rendszerek}

Ha az \(A\) mátrix szimmetrikus, de indefinit (azaz \(x^T A x\)
\(x\)-től függően pozitív és negatív értékeket is felvehet), akkor a
Cholesky-felbontás nem alkalmazható, és a numerikus stabilitáshoz
általában valamilyen főelemkiválasztásra van szükség. Nyilvánvaló, hogy
ha meg akarjuk őrizni a mátrix szimmetriáját, a főelemkiválasztásnak
szimmetrikusnak kell lennie, azaz \(P A P^T\) alakúnak, ahol \(P\)
permutációs mátrix.

Szeretnénk olyan \(P A P^T = L D L^T\) alakú felbontást kapni, ahol
\(L\) egységdiagonálisú alsó háromszögmátrix és \(D\) diagonális. Sajnos
diagonális \(D\)-vel ilyen felbontás nem feltétlenül létezik, és
létezése esetén is általában nem számítható ki stabilan pusztán
szimmetrikus főelemkiválasztással. A legjobb, amit tehetünk, hogy
\(D\)-t vagy tridiagonálisnak, vagy \(1 \times 1\)-es és
\(2 \times 2\)-es diagonális blokkokkal rendelkező blokkdiagonálisnak
választjuk. (Blokkmátrixnak nevezünk egy olyan mátrixot, amelynek
bejegyzéseit kompatibilis méretű részmátrixokra, azaz „blokkokra''
osztjuk. Egy blokkdiagonális mátrixban ezek közül a részmátrixok közül
mindegyik nulla, kivéve a fő blokkdiagonálisban lévőket.)

Hatékony algoritmusokat dolgozott ki Aasen a tridiagonális felbontásra,
valamint Bunch és Parlett (Bunch és Kaufman, illetve mások későbbi
továbbfejlesztéseivel a főelemkiválasztási eljárásban) a blokkdiagonális
felbontásra (lásd {[}198{]}). Mindkét esetben a főelemkiválasztási
eljárás olyan stabil felbontást ad, amely csak körülbelül \(n^3/6\)
szorzást és hozzávetőleg ugyanennyi összeadást igényel. Mindkét esetben
a rákövetkező megoldási szakasznak is csak \(\mathcal{O}(n^2)\) munkára
van szüksége. Így a szimmetrikus indefinit rendszerek megoldásának
költsége hasonló a pozitív definit rendszerek Cholesky-felbontással való
megoldásának költségéhez, és csak körülbelül a fele a nemszimmetrikus
rendszerek Gauss-kiküszöböléssel való megoldásának költségének.

\subsubsection{2.5.3 Sávos rendszerek}\label{suxe1vos-rendszerek}

A Gauss-kiküszöbölés sávmátrixokra alig különbözik az általános esettől:
az egyedüli algoritmikus változtatások a ciklusok tartományaiban vannak.
Természetesen olyan adatszerkezetet érdemes használni a sávmátrixhoz,
amely elkerüli a sávon kívüli nulla bejegyzések tárolását. Gyakori
választás, amikor a sáv sűrű, hogy a mátrixot kétdimenziós tömbben,
diagonálisonként tároljuk. Ha numerikus stabilitáshoz szükség van
főelemkiválasztásra, akkor az algoritmus kissé bonyolultabbá válik
annyiban, hogy a sávszélesség növekedhet (de legfeljebb a kétszeresére).
Így egy tetszőleges sávszélességű sávos rendszereket kezelő általános
célú megoldó igen hasonló egy általános mátrixokra szolgáló
Gauss-kiküszöbölési kódhoz.

Rögzített kis sávszélesség esetén azonban egy sávos rendszermegoldó
rendkívül egyszerű lehet, különösen ha a stabilitáshoz nincs szükség
főelemkiválasztásra. Tekintsük például a tridiagonális

\[\mathbf{A} = \begin{bmatrix} b_1 & c_1 & 0 & \cdots & 0 \\ a_2 & b_2 & c_2 & \ddots & \vdots \\ 0 & \ddots & \ddots & \ddots & 0 \\ \vdots & \ddots & a_{n-1} & b_{n-1} & c_{n-1} \\ 0 & \cdots & 0 & a_n & b_n \end{bmatrix}\]

mátrixot. Ha a stabilitáshoz nincs szükség főelemkiválasztásra --- ami a
gyakorlatban előforduló tridiagonális rendszerek esetén gyakran teljesül
(például ha a mátrix diagonálisan domináns vagy pozitív definit) ---,
akkor a Gauss-kiküszöbölés a 2.8. algoritmusra redukálódik, és az
\(\boldsymbol{A}\) mátrix keletkező háromszögű faktorait a következő
képletek adják:

\[\boldsymbol{L} = \begin{bmatrix} 1 & 0 & \cdots & \cdots & 0 \\ m_2 & 1 & \ddots & \ddots & \vdots \\ 0 & \ddots & \ddots & \ddots & \vdots \\ \vdots & \ddots & m_{n-1} & 1 & 0 \\ 0 & \cdots & 0 & m_n & 1 \end{bmatrix}, \quad \boldsymbol{U} = \begin{bmatrix} d_1 & c_1 & 0 & \cdots & 0 \\ 0 & d_2 & c_2 & \ddots & \vdots \\ \vdots & \ddots & \ddots & \ddots & 0 \\ \vdots & \ddots & \ddots & \ddots & 0 \\ \vdots & \ddots & \ddots & d_{n-1} & c_{n-1} \\ 0 & \cdots & \cdots & 0 & d_n \end{bmatrix}.\]

Általában egy \(\beta\) sávszélességű sávos rendszer csak
\(\mathcal{O}(\beta n)\) tárhelyet igényel, a felbontás pedig csak
\(\mathcal{O}(\beta^2 n)\) munkát; mindkettő jelentős megtakarítást
jelent a telt rendszerekhez képest, ha \(\beta \ll n\).

\textbf{2.8. Algoritmus. Tridiagonális LU-felbontás főelemkiválasztás
nélkül.}

\begin{verbatim}
d1 = b1
for i = 2 to n
    mi = ai/di−1
    di = bi − mici−1
end
                                 { ciklus az oszlopokon }
                                 { szorzó kiszámítása }
                                 { transzformáció alkalmazása }
\end{verbatim}

\subsection{2.6 Iterációs módszerek lineáris
egyenletrendszerekhez}\label{iteruxe1ciuxf3s-muxf3dszerek-lineuxe1ris-egyenletrendszerekhez}

A Gauss-kiküszöbölés egy példa a lineáris rendszerek megoldására
szolgáló direkt módszerre, azaz olyan módszerre, amely (pontos
aritmetikát feltételezve) egy lineáris rendszer pontos megoldását véges
számú lépésben állítja elő. Az iterációs módszerek ezzel szemben a
megoldás egy kezdeti becsléséből indulnak ki, és azt egymást követően
javítják, amíg a megoldás olyan pontos nem lesz, amilyet kívánunk.
Elméletben végtelen sok iterációra lehet szükség ahhoz, hogy a pontos
megoldáshoz konvergáljon, a gyakorlatban azonban az iterációkat akkor
állítjuk le, amikor a \(\|b - Ax\|\) maradék --- vagy a hiba valamilyen
más mérőszáma --- olyan kicsi, amilyet kívánunk. Egyes feladattípusok
esetén az iterációs módszereknek jelentős előnyeik lehetnek a direkt
módszerekkel szemben. A lineáris egyenletrendszerek iterációs
módszereinek tárgyalását a 11. fejezetre halasztjuk, ahol a parciális
differenciálegyenletek numerikus megoldását tekintjük, ami olyan ritka
lineáris rendszerekre vezet, amelyeket gyakran iterációs módszerekkel
lehet a legjobban megoldani.

\subsection{2.7 Szoftverek lineáris
egyenletrendszerekhez}\label{szoftverek-lineuxe1ris-egyenletrendszerekhez}

Szinte minden tudományos számítási szoftverkönyvtár tartalmaz különféle
típusú lineáris rendszerek megoldására szolgáló rutinokat. A 2.1.
táblázat valós, általános, sűrű lineáris rendszerek megoldására,
valamint a kondíciószám becslésére szolgáló megfelelő rutinok listáját
tartalmazza néhány széles körben elérhető szoftvergyűjteményből. Egyes
csomagok a rutinnevekben különböző elő- vagy utótagokat használnak az
adattípus jelölésére, tipikusan az \texttt{s} az egyszeres pontosságú
valós, \texttt{d} a dupla pontosságú valós, \texttt{c} az egyszeres
pontosságú komplex és \texttt{z} a dupla pontosságú komplex változatra
utal; itt csak az egyszeres pontosságú valós változatokat soroljuk fel.
A legtöbb ilyen szubrutinkönyvtárban speciálisabb rutinok is
rendelkezésre állnak a lineáris rendszerek bizonyos típusaira, mint a
szimmetrikus, pozitív definit, sávos rendszerek vagy ezek kombinációi.
Ezek közül a rutinok közül néhányat a 2.2. táblázat sorol fel; más
rutinok is rendelkezésre állhatnak, amelyek tárhelytakarékosabbak, vagy
egyéb speciális feladatokra készültek.

Az \(A\boldsymbol{x} = \boldsymbol{b}\) lineáris rendszerek megoldására
szolgáló hagyományos szoftver olykor egyetlen rutinként van
megvalósítva, de felosztható két rutinra is: az egyik a felbontás
kiszámítására, a másik a keletkező háromszögű rendszer megoldására.
Mindkét esetben a felbontás megismétlésére nincs szükség, ha ugyanazzal
a mátrixszal, de más jobb oldalakkal szeretnénk további megoldásokat. A
tipikusan megkövetelt bemenet a következőket tartalmazza: egy
kétdimenziós tömb az \(A\) mátrixszal, egy egydimenziós tömb a
\(\boldsymbol{b}\) jobb oldali vektorral (vagy egy kétdimenziós tömb
több jobb oldali vektorra), a rendszer egész \(n\) rendje, az \(A\)-t
tartalmazó tömb vezető dimenziója (hogy a szubrutin helyesen tudja
értelmezni a tömbindexeket), valamint esetenként némi munkaterület és
egy zászló az elvégzendő konkrét feladat jelzésére. Visszatéréskor az
\(\boldsymbol{x}\) megoldás rendszerint felülírja a \(\boldsymbol{b}\)
tárhelyét, a felbontás pedig az \(A\) tárhelyét. További kimenetként
szerepelhet egy állapotjelző zászló a hibák vagy figyelmeztetések
jelzésére, valamint a mátrix kondíciószámának egy becslése (vagy olykor
a kondíciószám reciproka). A kondíciószám becslésének többletköltsége
miatt ez a szolgáltatás általában opcionális.

\begin{longtable}[]{@{}
  >{\raggedright\arraybackslash}p{(\linewidth - 6\tabcolsep) * \real{0.2571}}
  >{\raggedright\arraybackslash}p{(\linewidth - 6\tabcolsep) * \real{0.2286}}
  >{\raggedright\arraybackslash}p{(\linewidth - 6\tabcolsep) * \real{0.2286}}
  >{\raggedright\arraybackslash}p{(\linewidth - 6\tabcolsep) * \real{0.2857}}@{}}
\toprule\noalign{}
\begin{minipage}[b]{\linewidth}\raggedright
Forrás
\end{minipage} & \begin{minipage}[b]{\linewidth}\raggedright
Felbontás
\end{minipage} & \begin{minipage}[b]{\linewidth}\raggedright
Megoldás
\end{minipage} & \begin{minipage}[b]{\linewidth}\raggedright
Kondíció\-becslés
\end{minipage} \\
\midrule\noalign{}
\endhead
\bottomrule\noalign{}
\endlastfoot
{[}152{]}FMM & decomp & solve & \\
GSL & gsllinalg & gsllinalg & \\
& LUdecomp & LUsolve & \\
HSL & ma21 & ma21 & \\
IMSL & lftrg & lfsrg & lfcrg \\
{[}262{]}KMN & sgefs & sgefs & sgefs \\
{[}9{]}LAPACK & sgetrf & sgetrs & sgecon \\
{[}116{]}LINPACK & sgefa & sgesl & sgeco \\
MATLAB & lu & \textbackslash{} & rcond/condest \\
NAG & f07adf & f07aef & f07agf \\
{[}220{]}NAPACK & fact & solve & con \\
{[}377{]}NR & ludcmp & lubksb & \\
{[}297{]}NUMAL & dec & sol & \\
SciPy & linalg. & linalg. & numpy.linalg. \\
& lufactor & lusolve & cond \\
SLATEC & sgefa & sgesl & sgeco \\
\end{longtable}

2.1. táblázat: Szoftverek általános lineáris rendszerek megoldására

Lineáris rendszerek megoldása interaktív környezetben, például
MATLAB-ban egyszerűbb, mint hagyományos szoftverrel, mert a csomag
belsőleg nyilvántartja az olyan részleteket, mint a vektorok és mátrixok
méretei, és sok mátrixművelet be van építve a nyelv szintaxisába és
szemantikájába. Például az \(A\boldsymbol{x} = \boldsymbol{b}\) lineáris
rendszer megoldását MATLAB-ban a „bal oldali osztás'' operátor adja,
amelyet fordított törtvonal jelöl, tehát
\texttt{x\ =\ A\ \textbackslash{}\ b}. A megoldást belsőleg
LU-felbontással, majd előre- és visszahelyettesítéssel számítja a
rendszer, de a felhasználónak erről nem kell tudnia. Az LU-felbontás
kifejezetten is kiszámítható, ha szükséges, a MATLAB \texttt{lu}
függvényével: \texttt{{[}L,\ U{]}\ =\ lu(A)}, illetve ha a mátrix
szimmetrikus és pozitív definit, akkor a Cholesky-felbontását az
\texttt{L\ =\ chol(A)} adja.

\subsubsection{2.7.1 LINPACK és LAPACK}\label{linpack-uxe9s-lapack}

A LINPACK átfogó szoftvercsomag sokféle lineáris egyenletrendszer
megoldására, az általános sűrű rendszereken kívül olyanokat is
beleértve, amelyeknek különféle speciális tulajdonságaik vannak, például
szimmetrikusak vagy sávosak. A lineáris rendszerek megoldása annyira
alapvető fontosságú a tudományos számításokban, hogy a LINPACK a
számítógépek teljesítményének összehasonlítására szolgáló szabványos
benchmarkká vált. A LINPACK kézikönyv {[}116{]} hasznos forrás
gyakorlati tanácsokhoz a lineáris egyenletrendszerek megoldásához.

\begin{longtable}[]{@{}
  >{\raggedright\arraybackslash}p{(\linewidth - 6\tabcolsep) * \real{0.2368}}
  >{\raggedright\arraybackslash}p{(\linewidth - 6\tabcolsep) * \real{0.2632}}
  >{\raggedright\arraybackslash}p{(\linewidth - 6\tabcolsep) * \real{0.2763}}
  >{\raggedright\arraybackslash}p{(\linewidth - 6\tabcolsep) * \real{0.2237}}@{}}
\toprule\noalign{}
\begin{minipage}[b]{\linewidth}\raggedright
\end{minipage} & \begin{minipage}[b]{\linewidth}\raggedright
Szimmetrikus
\end{minipage} & \begin{minipage}[b]{\linewidth}\raggedright
Szimmetrikus
\end{minipage} & \begin{minipage}[b]{\linewidth}\raggedright
Általános
\end{minipage} \\
\midrule\noalign{}
\endhead
\bottomrule\noalign{}
\endlastfoot
Forrás & pozitív definit & indefinit & sávos \\
GSL & gsllinalg & gsllinalg & \\
& choleskydecomp & mcholeskydecomp & \\
HSL & ma22 & ma29 & ma35 \\
IMSL & lftds/lfsds & lftsf/lfssf & lftrb/lfsrb \\
{[}9{]}LAPACK & spotrf/spotrs & ssytrf/ssytrs & sgbtrf/sgbtrs \\
{[}116{]}LINPACK & spofa/sposl & ssifa/ssisl & sgbfa/sgbsl \\
NAG & f07fdf/f07fef & f07mdf/f07mef & f07bdf/f07bef \\
{[}220{]}NAPACK & sfact/ssolve & ifact/isolve & bfact/bsolve \\
{[}377{]}NR & choldc/cholsl & & bandec/banbks \\
{[}297{]}NUMAL & chldec2/chlsol2 & decsym2/solsym2 & decbnd/solbnd \\
SciPy & linalg. & linalg. & linalg. \\
& cholesky & ldl & solvebanded \\
SLATEC & spofa/sposl & ssifa/ssisl & sgbfa/sgbsl \\
\end{longtable}

2.2. táblázat: Szoftverek speciális lineáris rendszerek megoldására

A LAPACK nevű utódcsomag a teljes LINPACK-gyűjteményt frissíti, nagyobb
teljesítményt nyújtva a modern számítógép-architektúrákon, egyes
párhuzamos számítógépeket is beleértve. A LAPACK újabb algoritmusai sok
esetben nagyobb pontosságot, robusztusságot és funkcionalitást is
elérnek, mint a LINPACK-beli elődeik. A LAPACK a lineáris algebra összes
jelentős számítási feladatához egyaránt tartalmaz egyszerű és szakértői
meghajtókat, valamint a különféle felbontásokhoz, a háromszögű
rendszerek megoldásához, a normabecsléshez, a skálázáshoz és az iteratív
finomításhoz szükséges sokféle számítási és kiegészítő rutint. Mind a
LINPACK, mind a LAPACK elérhető a Netlibről, és számos egyéb könyvtár és
csomag lineáris rendszermegoldóit közvetlenül ezekre építik.

\subsubsection{2.7.2 Basic Linear Algebra Subprograms
(BLAS)}\label{basic-linear-algebra-subprograms-blas}

A LINPACK és a LAPACK magas szintű rutinjai az alacsonyabb szintű Basic
Linear Algebra Subprograms (BLAS) csomagon alapulnak. A BLAS-t
eredetileg arra tervezték, hogy egységbe zárja a vektorokon végzett
alapvető műveleteket, hogy azok egy adott számítógép-architektúrára
optimalizálhatók legyenek, miközben a magasabb szintű, őket hívó rutinok
hordozhatók maradnak. Az újabb számítógép-architektúrák arra
ösztönöztek, hogy magasabb szintű BLAS-t fejlesszenek ki, amely a
mátrix-vektor és a mátrix-mátrix műveleteket zárja egységbe a
hierarchikus memóriának --- például a gyorsítótárnak, a
vektorregisztereknek és a lapozásos virtuális memóriának --- a jobb
kihasználása érdekében. A 2.3. táblázat felsorolja az egyes szintek
néhány legfontosabb BLAS-rutinját.

A jó teljesítmény kulcsa az adat-újrafelhasználás, azaz az, hogy egy
adott adaton minél több aritmetikai műveletet végezzünk, miközben az a
memóriahierarchia leggyorsabban elérhető részében van tárolva. A 3.
szintű BLAS-nak nagyobb lehetősége van az adat-újrafelhasználásra, mert
\(\mathcal{O}(n^3)\) műveletet végez \(\mathcal{O}(n^2)\) adaton, míg az
alacsonyabb szintű BLAS esetén a műveletek száma az adatok számával
arányos. A BLAS általános változatai elérhetők a Netlibről, és számos
számítógépgyártó kínál egyedi változatokat, amelyek a saját rendszerükön
a legjobb teljesítményre vannak optimalizálva.

\begin{longtable}[]{@{}
  >{\raggedright\arraybackslash}p{(\linewidth - 10\tabcolsep) * \real{0.0897}}
  >{\raggedright\arraybackslash}p{(\linewidth - 10\tabcolsep) * \real{0.1410}}
  >{\raggedright\arraybackslash}p{(\linewidth - 10\tabcolsep) * \real{0.1923}}
  >{\raggedright\arraybackslash}p{(\linewidth - 10\tabcolsep) * \real{0.1282}}
  >{\raggedright\arraybackslash}p{(\linewidth - 10\tabcolsep) * \real{0.4231}}
  >{\raggedright\arraybackslash}p{(\linewidth - 10\tabcolsep) * \real{0.0256}}@{}}
\toprule\noalign{}
\begin{minipage}[b]{\linewidth}\raggedright
Szint
\end{minipage} & \begin{minipage}[b]{\linewidth}\raggedright
TOMS\#
\end{minipage} & \begin{minipage}[b]{\linewidth}\raggedright
Munka
\end{minipage} & \begin{minipage}[b]{\linewidth}\raggedright
Példák
\end{minipage} & \begin{minipage}[b]{\linewidth}\raggedright
Funkció
\end{minipage} & \begin{minipage}[b]{\linewidth}\raggedright
\end{minipage} \\
\midrule\noalign{}
\endhead
\bottomrule\noalign{}
\endlastfoot
1 & 539 & O(n) & saxpy & Skalár · vektor plusz vektor & \\
& & & sdot & Két vektor belső szorzata & \\
& & & snrm2 & Vektor euklideszi normája & \\
2 & 656 & 2O(n) & sgemv & Mátrix-vektor szorzás & \\
& & & strsv & Háromszögű megoldás & \\
& & & sger & Elsőrendű frissítés & \\
3 & 679 & 3O(n) & sgemm & Mátrix-mátrix szorzás & \\
& & & strsm & Többszörös háromszögű megoldás & \\
& & & ssyrk & \(k\)-adrendű frissítés & \\
\end{longtable}

2.3. táblázat: Példák az alapvető lineáris algebrai szubrutinokra (BLAS)

\subsection{2.8 Történeti jegyzetek és további
olvasnivaló}\label{tuxf6rtuxe9neti-jegyzetek-uxe9s-tovuxe1bbi-olvasnivaluxf3-1}

Gauss 1810-ben fogalmazta meg a kiküszöbölési módszerét, de nem egy
mátrix felbontásaként, hanem egy kvadratikus alak egyszerűsítéseként
írta le. A mátrixokat ugyanis csak 1855-ben vezette be Cayley, és a
Gauss-kiküszöbölés mátrixfelbontásként való értelmezését csak az 1940-es
években ismerte fel Dwyer, von Neumann és mások (lásd {[}431{]}). A
digitális számítások korai úttörőit --- például von Neumannt {[}487{]}
és Turingot {[}471{]} (a mátrix kondíciószámának fogalma Turingtól
származik) --- súlyosan aggasztotta, hogy a nagy lineáris rendszerek
Gauss-kiküszöböléssel történő megoldásakor felhalmozódó kerekítési hiba
használhatatlanná teszi-e az eredményeket, és kezdetben jelentős
pesszimizmus övezte ezt a kérdést. A számítási tapasztalatok hamarosan
megmutatták azonban, hogy a módszer a gyakorlatban meglepően stabil és
pontos, és a hibaelemzések végül követték a tapasztalatot, hogy
magyarázatot adjanak erre a szerencsére (lásd különösen Wilkinson
munkáit {[}499, 500, 501{]}).

Mint kiderült, a részleges főelemkiválasztással végzett
Gauss-kiküszöbölés műveletszáma rosszabb az optimálisnál {[}444{]}, a
legrosszabb esetben instabil {[}499{]}, és elméleti értelemben nem
valósítható meg hatékonyan párhuzamosan {[}485{]}. Mégis a gyakorlatban
--- még párhuzamos számítógépeken is --- következetesen eredményes, és a
tudományos számítások egyik fő igáslova. A legtöbb numerikus algoritmus
Murphy törvényének engedelmeskedik --- „ha valami elromolhat, el is
romlik" ---, ám úgy tűnik, a Gauss-kiküszöbölés boldog kivétel. E
figyelemre méltó algoritmus egyes „rejtélyeinek" további tárgyalásához
lásd {[}461, 467{]}-et; „számítási ellenpéldákhoz" lásd {[}156,
510{]}-et.

A lineáris algebra hátteréhez az olvasó Strang kiváló tankönyveit
{[}436, 441{]} veheti kézbe. Haladóbb mátrixelmélethez lásd {[}241,
242{]}-t. További példák, feladatok és a számítási lineáris algebra
gyakorlati alkalmazásai megtalálhatók {[}238, 276, 314{]}-ben. A
mátrixszámítások enciklopédikus referenciája {[}198{]}; tankönyvi
tárgyalások közé tartoznak {[}40, 96, 104, 185, 220, 257, 426, 430, 466,
493, 496{]}. A lineáris rendszerek megoldásának egyik nagy hatású korai
műve --- és az elsők egyike, amely magas színvonalú szoftvert is
tartalmaz --- a {[}153{]}. A mátrixszámítások hasznos, Fortran és MATLAB
nyelvű gyakorlati kézikönyve a {[}82{]}.

A lineáris rendszerek és a lineáris algebra számos egyéb feladatának
hibaelemzéséről és perturbációelméletéről szóló átfogó tárgyaláshoz lásd
{[}237, 433{]}-at. A kondíciószám-becslésről áttekintést ad {[}235{]}. A
lineáris algebrában a (normaszerű helyett) komponensenkénti
perturbációelmélet részletes áttekintését adja {[}236{]}. A LINPACK és a
LAPACK dokumentációja {[}116{]}, illetve {[}9{]}. A BLAS (Basic Linear
Algebra Subprograms) könyvtárról lásd {[}114, 115, 300{]}-at. A
számítási környezetnek a Gauss-kiküszöbölés és más mátrixszámítások
teljesítményére gyakorolt hatását vizsgáló egyik legkorábbi cikk
{[}330{]} volt. Az erről a mára igen terjedelmessé vált témáról egy
ízelítőhöz lásd {[}105, 117, 118, 357{]}-et.

\subsubsection{Áttekintő
kérdések}\label{uxe1ttekintux151-kuxe9rduxe9sek-1}

\begin{itemize}
\tightlist
\item
  \textbf{2.1.} Igaz vagy hamis: Ha egy \(A\) mátrix nemszinguláris,
  akkor az \(Ax = b\) lineáris egyenletrendszer megoldásainak száma a
  \(b\) jobb oldali vektor konkrét választásától függ.
\item
  \textbf{2.2.} Igaz vagy hamis: Ha egy mátrix determinánsa nagyon
  kicsi, akkor a mátrix majdnem szinguláris.
\item
  \textbf{2.3.} Igaz vagy hamis: Szimmetrikus \(A\) mátrixra mindig
  \(\|A\|_1 = \|A\|_\infty\).
\item
  \textbf{2.4.} Igaz vagy hamis: Ha egy háromszögmátrixnak van nullával
  egyenlő eleme a főátlóban, akkor a mátrix szükségképpen szinguláris.
\item
  \textbf{2.5.} Igaz vagy hamis: Ha egy tetszőleges mátrixnak van
  nullával egyenlő eleme a főátlóban, akkor a mátrix szükségképpen
  szinguláris.
\item
  \textbf{2.6.} Igaz vagy hamis: Egy alulhatározott \(Ax = b\) lineáris
  egyenletrendszernek, ahol \(A\) egy \(m \times n\)-es mátrix
  \(m < n\)-nel, mindig van megoldása.
\item
  \textbf{2.7.} Igaz vagy hamis: Két felső háromszögmátrix szorzata
  felső háromszögmátrix.
\item
  \textbf{2.8.} Igaz vagy hamis: Két szimmetrikus mátrix szorzata
  szimmetrikus.
\item
  \textbf{2.9.} Igaz vagy hamis: Egy nemszinguláris felső
  háromszögmátrix inverze felső háromszögmátrix.
\item
  \textbf{2.10.} Igaz vagy hamis: Ha egy \(n \times n\)-es \(A\) mátrix
  sorai lineárisan összefüggők, akkor az oszlopai is lineárisan
  összefüggők.
\item
  \textbf{2.11.} Igaz vagy hamis: Az \(Ax = b\) lineáris
  egyenletrendszernek akkor és csak akkor van megoldása, ha az
  \(m \times n\)-es \(A\) mátrix és az \(m \times (n+1)\)-es
  \([\,A\ b\,]\) kibővített mátrix rangja azonos.
\item
  \textbf{2.12.} Igaz vagy hamis: Ha \(A\) tetszőleges \(n \times n\)-es
  mátrix és \(P\) tetszőleges \(n \times n\)-es permutációs mátrix,
  akkor \(PA = AP\).
\item
  \textbf{2.13.} Igaz vagy hamis: Ha sorcserék megengedettek, akkor az
  LU-felbontás mindig létezik, még szinguláris \(A\) mátrixra is.
\item
  \textbf{2.14.} Igaz vagy hamis: Ha egy lineáris rendszer jól
  kondicionált, akkor a Gauss-kiküszöbölésben nincs szükség
  főelemkiválasztásra.
\item
  \textbf{2.15.} Igaz vagy hamis: Ha egy mátrix szinguláris, akkor nem
  lehet LU-felbontása.
\item
  \textbf{2.16.} Igaz vagy hamis: Ha egy nemszinguláris szimmetrikus
  mátrix nem pozitív definit, akkor nem lehet Cholesky-felbontása.
\item
  \textbf{2.17.} Igaz vagy hamis: Egy szimmetrikus pozitív definit
  mátrix mindig jól kondicionált.
\item
  \textbf{2.18.} Igaz vagy hamis: A főelemkiválasztás nélküli
  Gauss-kiküszöbölés csak akkor hiúsul meg, ha a mátrix rosszul
  kondicionált vagy szinguláris.
\item
  \textbf{2.19.} Igaz vagy hamis: Ha egy mátrix LU-felbontását már
  kiszámítottuk egy lineáris rendszer megoldásához, akkor az ugyanazzal
  a mátrixszal, de eltérő jobb oldali vektorokkal adott további lineáris
  rendszerek a mátrix újrafelbontása nélkül megoldhatók.
\item
  \textbf{2.20.} Igaz vagy hamis: LU-felbontáson és háromszögű
  megoldáson alapuló explicit mátrixinvertálásban a munka többségét a
  felbontás adja.
\item
  \textbf{2.21.} Igaz vagy hamis: Ha \(x\) tetszőleges \(n\)-dimenziós
  vektor, akkor \(\|x\|_1 \ge \|x\|_\infty\).
\item
  \textbf{2.22.} Igaz vagy hamis: Egy szinguláris mátrix normája nulla.
\item
  \textbf{2.23.} Igaz vagy hamis: Ha \(\|A\| = 0\), akkor \(A = 0\).
\item
  \textbf{2.24.} Igaz vagy hamis: \(\|A\|_1 = \|A^T\|_\infty\).
\item
  \textbf{2.25.} Igaz vagy hamis: Ha \(A\) tetszőleges \(n \times n\)-es
  nemszinguláris mátrix, akkor
  \(\operatorname{cond}(A) = \operatorname{cond}(A^{-1})\).
\item
  \textbf{2.26.} Igaz vagy hamis: Nemszinguláris lineáris
  egyenletrendszer megoldásakor a részleges főelemkiválasztással végzett
  Gauss-kiküszöbölés általában kis maradékot ad, még ha a mátrix rosszul
  kondicionált is.
\item
  \textbf{2.27.} Igaz vagy hamis: A részleges főelemkiválasztással
  végzett Gauss-kiküszöbölésben a szorzótényezők abszolút értékben
  \(1\)-gyel korlátosak, ezért az egymást követő redukált mátrixok
  elemei abszolút értékben nem növekedhetnek.
\item
  \textbf{2.28.} Lehet-e egy \(Ax = b\) lineáris egyenletrendszernek
  pontosan két különböző megoldása?
\item
  \textbf{2.29.} Meghatározható-e egy \(Ax = b\) lineáris rendszer
  megoldásainak száma valaha is pusztán az \(A\) mátrixból, a \(b\) jobb
  oldali vektor ismerete nélkül?
\item
  \textbf{2.30.} Egy \(Ax = b\) négyzetes lineáris egyenletrendszer
  megoldásakor melyik volna a súlyosabb nehézség: ha \(A\) sorai
  lineárisan összefüggők, vagy ha \(A\) oszlopai lineárisan összefüggők?
  Magyarázd meg.
\item
  \textbf{2.31.} (a) Mondj ki egy definiáló tulajdonságot egy \(A\)
  szinguláris mátrixra. (b) Tegyük fel, hogy az \(Ax = b\) lineáris
  rendszernek két különböző \(x\) és \(y\) megoldása van. Az (a)-ban
  megadott tulajdonsággal bizonyítsd be, hogy \(A\) szükségképpen
  szinguláris.
\item
  \textbf{2.32.} Adott egy nemszinguláris \(Ax = b\) lineáris
  egyenletrendszer. Az alábbi műveletek mindegyikének milyen hatása van
  az \(x\) megoldásvektorra? (a) \([\,A\ b\,]\) sorainak permutálása;
  (b) \(A\) oszlopainak permutálása; (c) az egyenlet mindkét oldalának
  balról való szorzása egy \(M\) nemszinguláris mátrixszal.
\item
  \textbf{2.33.} Tegyük fel, hogy egy \(Ax = b\) lineáris
  egyenletrendszer mindkét oldalát megszorozzuk egy nem nulla \(\alpha\)
  skalárral. (a) Megváltozik-e ettől a valódi \(x\) megoldás? (b)
  Megváltozik-e ettől egy adott \(x\)-re az \(r = b - Ax\)
  maradékvektor? (c) Milyen következtetés vonható le egy számított
  megoldás minőségének megítéléséről?
\item
  \textbf{2.34.} Tegyük fel, hogy egy \(Ax = b\) lineáris
  egyenletrendszer mindkét oldalát balról megszorozzuk egy
  nemszinguláris diagonális mátrixszal. (a) Megváltozik-e ettől a valódi
  \(x\) megoldás? (b) Befolyásolhatja-e ez a rendszer kondicionáltságát?
  (c) Befolyásolhatja-e ez a Gauss-kiküszöbölésben a főelemek
  választását?
\item
  \textbf{2.35.} Adj meg egy elemi eliminációs mátrixot, amely az alábbi
  vektor utolsó két komponensét nullává teszi:
\end{itemize}

\[\begin{bmatrix} 3 \\ 2 \\ -1 \\ 4 \end{bmatrix}.\]

\begin{itemize}
\tightlist
\item
  \textbf{2.36.} (a) Adj meg egy \(4 \times 4\)-es permutációs mátrixot,
  amely bármely \(4\)-dimenziós vektor 2. és 4. komponensét felcseréli.
  (b) Adj meg egy \(4 \times 4\)-es permutációs mátrixot, amely bármely
  \(4\)-dimenziós vektor komponenseinek sorrendjét megfordítja.
\item
  \textbf{2.37.} Szinguláris mátrix és egzakt aritmetika használata
  esetén egy lineáris rendszer Gauss-kiküszöböléssel történő
  megoldásában hol fog a megoldási folyamat megakadni (a) részleges
  főelemkiválasztással? (b) főelemkiválasztás nélkül?
\item
  \textbf{2.38.} (a) Mi a különbség a Gauss-kiküszöbölésben a részleges
  és a teljes főelemkiválasztás között? (b) Mondj a másikhoz képest
  egy-egy előnyt mindkét típusú főelemkiválasztásra.
\item
  \textbf{2.39.} Tekintsük az alábbi \(A\) mátrixot, amelynek
  LU-felbontását Gauss-kiküszöböléssel szeretnénk kiszámítani:
\end{itemize}

\[\boldsymbol{A} = \begin{bmatrix} 4 & -8 & 1 \\ 6 & 5 & 7 \\ 0 & -10 & -3 \end{bmatrix}.\]

Mi lesz a kezdeti főelem, ha (a) nem használunk főelemkiválasztást? (b)
részleges főelemkiválasztást használunk? (c) teljes főelemkiválasztást
használunk?

\begin{itemize}
\tightlist
\item
  \textbf{2.40.} Adj két okot arra, hogy a Gauss-kiküszöbölés
  numerikusan stabil megvalósításához miért lényeges a
  főelemkiválasztás.
\item
  \textbf{2.41.} Ha \(A\) egy rosszul kondicionált mátrix, és az
  LU-felbontását részleges főelemkiválasztással végzett
  Gauss-kiküszöböléssel számítjuk, azt várnád-e, hogy a rossz
  kondicionáltság \(L\)-ben, \(U\)-ban vagy mindkettőben tükröződik?
  Miért?
\item
  \textbf{2.42.} (a) Mi az alábbi mátrix inverze?
\end{itemize}

\[\begin{bmatrix} 1 & 0 & 0 & 0 \\ 0 & 1 & 0 & 0 \\ 0 & m_1 & 1 & 0 \\ 0 & m_2 & 0 & 1 \end{bmatrix}\]

\begin{enumerate}
\def\labelenumi{(\alph{enumi})}
\setcounter{enumi}{1}
\tightlist
\item
  Hogyan merülhet fel egy ilyen mátrix a számítási gyakorlatban?
\end{enumerate}

\begin{itemize}
\tightlist
\item
  \textbf{2.43.} (a) Felírható-e minden \(n \times n\)-es nemszinguláris
  \(\boldsymbol{A}\) mátrix
  \(\boldsymbol{A} = \boldsymbol{L}\boldsymbol{U}\) szorzat alakban,
  ahol \(\boldsymbol{L}\) alsó háromszögmátrix és \(\boldsymbol{U}\)
  felső háromszögmátrix? (b) Ha igen, milyen algoritmus valósítja ezt
  meg? Ha nem, adj ellenpéldát a szemléltetésre.
\item
  \textbf{2.44.} Adott egy \(n \times n\)-es nemszinguláris
  \(\boldsymbol{A}\) mátrix és egy további \(n \times n\)-es
  \(\boldsymbol{B}\) mátrix. Mi a legjobb módja az \(n \times n\)-es
  \(\boldsymbol{A}^{-1}\boldsymbol{B}\) mátrix kiszámításának?
\item
  \textbf{2.45.} Ha \(A\) és \(B\) \(n \times n\)-es mátrixok, \(A\)
  nemszinguláris, és \(c\) egy \(n\)-dimenziós vektor, hogyan számítanád
  ki hatékonyan az \(A^{-1}Bc\) szorzatot?
\item
  \textbf{2.46.} Ha \(A\) egy \(n \times n\)-es mátrix és \(x\) egy
  \(n\)-dimenziós vektor, az alábbi számítások közül melyik igényel
  kevesebb munkát? Magyarázd meg. (a)
  \(\boldsymbol{y} = (\boldsymbol{x}\,\boldsymbol{x}^T)\,\boldsymbol{A}\);
  (b)
  \(\boldsymbol{y} = \boldsymbol{x}\,(\boldsymbol{x}^T\,\boldsymbol{A})\).
\item
  \textbf{2.47.} Hogyan viszonyul egy \(n \times n\)-es háromszögű
  lineáris egyenletrendszer megoldásához szükséges számítási munka egy
  általános \(n \times n\)-es rendszer megoldásához szükségesével?
\item
  \textbf{2.48.} Tegyük fel, hogy egy nemszinguláris \(A\) mátrix
  \(A = LU\) LU-felbontását már kiszámítottad. Hogyan használnád fel az
  \(A^T x = b\) lineáris rendszer megoldásához?
\item
  \textbf{2.49.} Ha \(L\) egy nemszinguláris alsó háromszögmátrix, \(P\)
  egy permutációs mátrix, és \(b\) egy adott vektor, hogyan oldanád meg
  az alábbi lineáris rendszerek mindegyikét? (a) \(LPx = b\); (b)
  \(PLx = b\).
\item
  \textbf{2.50.} Az \(\mathbb{R}^2\) síkban lehetséges-e olyan
  \(\boldsymbol{x} \neq \boldsymbol{0}\) vektor, amelyre
  \(\|\boldsymbol{x}\|_1 = \|\boldsymbol{x}\|_\infty\)? Ha igen, adj
  példát.
\item
  \textbf{2.51.} Az \(\mathbb{R}^2\) síkban lehetséges-e olyan
  \(\boldsymbol{x}\) és \(\boldsymbol{y}\) vektor, amelyekre
  \(\|\boldsymbol{x}\|_1 > \|\boldsymbol{y}\|_1\), de
  \(\|\boldsymbol{x}\|_\infty < \|\boldsymbol{y}\|_\infty\)? Ha igen,
  adj példát.
\item
  \textbf{2.52.} Általában melyik mátrixnormát könnyebb kiszámítani,
  \(\|A\|_1\)-et vagy \(\|A\|_2\)-t? Miért?
\item
  \textbf{2.53.} (a) Egy mátrix determinánsának nagysága jó mutatója-e
  annak, hogy a mátrix majdnem szinguláris-e? (b) Ha igen, miért? Ha
  nem, mi a közelszingularitás jobb mutatója?
\item
  \textbf{2.54.} (a) Hogyan definiáljuk egy \(A\) mátrix kondíciószámát
  egy adott mátrixnormára? (b) Hogyan használjuk a kondíciószámot az
  \(Ax = b\) lineáris rendszer számított megoldásának
  pontosságbecslésére?
\item
  \textbf{2.55.} Miért nem triviális feladat egy általános mátrix
  kondíciószámának kiszámítása?
\item
  \textbf{2.56.} Adj példát egy \(3 \times 3\)-as \(A\) mátrixra --- az
  egységmátrixtól különbözőre ---, amelyre
  \(\operatorname{cond}(A) = 1\).
\item
  \textbf{2.57.} (a) Mi az alábbi mátrix kondíciószáma az \(1\)-normát
  használva?
\end{itemize}

\[\begin{bmatrix} 4 & 0 & 0 \\ 0 & -6 & 0 \\ 0 & 0 & 2 \end{bmatrix}\]

\begin{enumerate}
\def\labelenumi{(\alph{enumi})}
\setcounter{enumi}{1}
\tightlist
\item
  Eltér-e a válaszod, ha a \(\infty\)-normát használjuk?
\end{enumerate}

\begin{itemize}
\tightlist
\item
  \textbf{2.58.} Tegyük fel, hogy az \(n \times n\)-es \(A\) mátrix
  tökéletesen jól kondicionált, azaz \(\operatorname{cond}(A) = 1\). Az
  alábbi mátrixok közül melyek szükségképpen szintén ezzel a
  tulajdonsággal rendelkeznek? (a) \(c\boldsymbol{A}\), ahol \(c\)
  tetszőleges nem nulla skalár; (b) \(\boldsymbol{D}\boldsymbol{A}\),
  ahol \(\boldsymbol{D}\) egy nemszinguláris diagonális mátrix; (c)
  \(PA\), ahol \(P\) tetszőleges permutációs mátrix; (d)
  \(\boldsymbol{B}\boldsymbol{A}\), ahol \(\boldsymbol{B}\) tetszőleges
  nemszinguláris mátrix; (e) \(A^{-1}\), azaz \(A\) inverze; (f)
  \(\boldsymbol{A}^T\), azaz \(\boldsymbol{A}\) transzponáltja.
\item
  \textbf{2.59.} Legyen \(\boldsymbol{A} = \operatorname{diag}(1/2)\)
  egy \(n \times n\)-es diagonális mátrix, amelynek minden diagonális
  eleme \(1/2\). (a) Mi \(\det(\boldsymbol{A})\) értéke? (b) Mi
  \(\operatorname{cond}(A)\) értéke? (c) Milyen következtetést lehet
  levonni ezekből az eredményekből?
\item
  \textbf{2.60.} Tegyük fel, hogy az \(n \times n\)-es
  \(\boldsymbol{A}\) mátrix egzaktul szinguláris, de a lebegőpontos
  ábrázolása, \(fl(\boldsymbol{A})\), nemszinguláris. Ebben az esetben
  mit várnál a \(\operatorname{cond}(fl(\boldsymbol{A}))\) kondíciószám
  nagyságrendjének?
\item
  \textbf{2.61.} Sorold be az alábbi mátrixok mindegyikét, hogy jól vagy
  rosszul kondicionált-e:
\end{itemize}

\[\begin{array}{lll}
(a) & \begin{bmatrix} 10^{10} & 0 \\ 0 & 10^{-10} \end{bmatrix} \\
(b) & \begin{bmatrix} 10^{10} & 0 \\ 0 & 10^{10} \end{bmatrix} \\
(c) & \begin{bmatrix} 10^{-10} & 0 \\ 0 & 10^{-10} \end{bmatrix} \\
(d) & \begin{bmatrix} 1 & 2 \\ 2 & 4 \end{bmatrix}
\end{array}\]

\begin{itemize}
\tightlist
\item
  \textbf{2.62.} Az alábbiak közül melyek jó mutatói annak, hogy egy
  mátrix majdnem szinguláris? (a) Determinánsa kicsi. (b) Normája kicsi.
  (c) Normája nagy. (d) Kondíciószáma nagy.
\item
  \textbf{2.63.} (a) Egy \(Ax = b\) lineáris rendszer megoldásakor mit
  értünk egy \(\hat{x}\) közelítő megoldás maradékán? (b) Kis relatív
  maradék mindig azt jelenti-e, hogy a megoldás pontos? Miért? (c) Nagy
  relatív maradék mindig azt jelenti-e, hogy a megoldás pontatlan?
  Miért?
\item
  \textbf{2.64.} Egy \(10\) decimális jegy precíziójú lebegőpontos
  rendszerben, ha egy olyan lineáris rendszert oldunk meg részleges
  főelemkiválasztással végzett Gauss-kiküszöböléssel, amelynek mátrixa
  \(10^3\) kondíciószámú, és amelynek bemeneti adatai teljes gépi
  precízióval pontosak, körülbelül hány értékes jegyet várnál a
  megoldásban?
\item
  \textbf{2.65.} Tegyük fel, hogy egy \(Ax = b\) lineáris
  egyenletrendszert egy olyan számítógépen oldasz meg, amelynek
  lebegőpontos számrendszere \(12\) decimális jegy precíziójú, és hogy a
  feladat adatai teljes gépi precízióval pontosak. Körülbelül milyen
  nagy lehet az \(A\) mátrix kondíciószáma, mielőtt a számított \(x\)
  megoldás egyetlen értékes jegyet sem tartalmaz?
\item
  \textbf{2.66.} Milyen körülmények között jelenti-e egy kis
  \(r = b - Ax\) maradékvektor, hogy \(x\) az \(Ax = b\) lineáris
  rendszer pontos megoldása?
\item
  \textbf{2.67.} Legyen \(A\) egy tetszőleges négyzetes mátrix, és \(c\)
  egy tetszőleges skalár. Az alábbi állítások közül melyek kell
  szükségképpen teljesüljenek? (a) \(\|cA\| = |c| \cdot \|A\|\). (b)
  \(\operatorname{cond}(cA) = |c| \cdot \operatorname{cond}(A)\).
\item
  \textbf{2.68.} (a) Mi a fő különbség a Gauss-kiküszöbölés és a
  Gauss--Jordan-kiküszöbölés között? (b) Mondj a másikhoz képest egy-egy
  előnyt mindkét típusú kiküszöbölésre.
\item
  \textbf{2.69.} Rangsorold az alábbi módszereket az \(n\)-ed rendű
  általános lineáris egyenletrendszer megoldásához szükséges munka
  mennyisége szerint: (a) Gauss--Jordan-kiküszöbölés; (b) részleges
  főelemkiválasztással végzett Gauss-kiküszöbölés; (c) Cramer-szabály;
  (d) explicit mátrixinvertálás, majd mátrix-vektor szorzás.
\item
  \textbf{2.70.} (a) Mennyi tárhely szükséges egy egyrangú
  \(n \times n\)-es mátrix hatékony tárolásához? (b) Hány aritmetikai
  művelet szükséges egy \(n\)-dimenziós vektor és egy egyrangú
  \(n \times n\)-es mátrix szorzatának hatékony kiszámításához?
\item
  \textbf{2.71.} Egy \(Ax = b\) lineáris rendszer megoldásához a
  közönséges Gauss-kiküszöbölés és a Gauss--Jordan-kiküszöbölés
  összehasonlításakor (a) melyiknek költségesebb a felbontása? (b)
  melyiknek költségesebb a visszahelyettesítése? (c) melyiknek nagyobb
  az összköltsége?
\item
  \textbf{2.72.} A lineáris rendszerek megoldására szolgáló alábbi
  kiküszöbölési algoritmusok mindegyikére: van-e olyan főelemválasztási
  stratégia, amely biztosítani tudja, hogy minden szorzótényező abszolút
  értéke legfeljebb \(1\)? (a) Gauss-kiküszöbölés; (b)
  Gauss--Jordan-kiküszöbölés.
\item
  \textbf{2.73.} Egy \(A\) mátrix mely két tulajdonsága együtt
  implikálja, hogy \(A\)-nak van Cholesky-felbontása?
\item
  \textbf{2.74.} Sorolj fel három előnyt, amelyet a Cholesky-felbontás
  az LU-felbontáshoz képest nyújt.
\item
  \textbf{2.75.} Egy \(n \times n\)-es szimmetrikus pozitív definit
  mátrix Cholesky-felbontásának kiszámításához hány négyzetgyökre van
  szükség?
\item
  \textbf{2.76.} Legyen \(A = \{a_{ij}\}\) egy \(n \times n\)-es
  szimmetrikus pozitív definit mátrix. (a) Mi a Cholesky-tényezőjének
  \((1,1)\) eleme? (b) Mi a Cholesky-tényezőjének \((n,1)\) eleme?
\item
  \textbf{2.77.} Mi az alábbi mátrix Cholesky-felbontása?
\end{itemize}

\[\begin{bmatrix} 4 & 2 \\ 2 & 2 \end{bmatrix}\]

\begin{itemize}
\tightlist
\item
  \textbf{2.78.} (a) Lehetséges-e általában egy szimmetrikus indefinit
  lineáris rendszert hasonló költséggel megoldani, mint amennyibe egy
  szimmetrikus pozitív definit lineáris rendszer Cholesky-felbontással
  történő megoldása kerül? (b) Ha igen, milyen algoritmus valósítja ezt
  meg? Ha nem, miért?
\item
  \textbf{2.79.} Adj két okot arra, hogy a lineáris rendszerek
  megoldására szolgáló iteratív finomítás miért gyakran nem valósítható
  meg a gyakorlatban.
\item
  \textbf{2.80.} Tegyük fel, hogy egy \(n \times n\)-es \(Ax = b\)
  lineáris rendszert LU-felbontással és visszahelyettesítéssel már
  megoldottál. Mi a további költsége (a nagyságrendi becslés elegendő)
  egy új rendszer megoldásának (a) ugyanazzal a mátrixszal, de más jobb
  oldali vektorral? (b) ha a mátrixot egy egyrangú mátrix hozzáadásával
  módosítjuk? (c) ha a mátrix teljesen megváltozik?
\end{itemize}

\subsubsection{Feladatok}\label{feladatok-1}

\begin{itemize}
\tightlist
\item
  \textbf{2.1.} A 2.2. szakaszban négy definiáló tulajdonságot adtunk
  meg egy szinguláris mátrixra. Mutasd meg, hogy ez a négy tulajdonság
  valóban ekvivalens.
\item
  \textbf{2.2.} Tegyük fel, hogy egy \(n \times n\)-es \(A\) mátrix
  minden sorösszege nulla. Mutasd meg, hogy \(A\) szükségképpen
  szinguláris.
\item
  \textbf{2.3.} Tegyük fel, hogy \(A\) egy szinguláris \(n \times n\)-es
  mátrix. Bizonyítsd be, hogy ha az \(Ax = b\) lineáris rendszernek van
  legalább egy \(x\) megoldása, akkor végtelen sok megoldása van.
\item
  \textbf{2.4.} (a) Mutasd meg, hogy az alábbi mátrix szinguláris.
\end{itemize}

\[\boldsymbol{A} = \begin{bmatrix} 1 & 1 & 0 \\ 1 & 2 & 1 \\ 1 & 3 & 2 \end{bmatrix}\]

\begin{enumerate}
\def\labelenumi{(\alph{enumi})}
\setcounter{enumi}{1}
\tightlist
\item
  Ha \(b = \begin{bmatrix} 2 & 4 & 6 \end{bmatrix}^T\), hány megoldása
  van az \(Ax = b\) rendszernek?
\end{enumerate}

\begin{itemize}
\tightlist
\item
  \textbf{2.5.} Mi az alábbi mátrix inverze?
\end{itemize}

\[\mathbf{A} = \begin{bmatrix} 1 & 0 & 0 \\ 1 & -1 & 0 \\ 1 & -2 & 1 \end{bmatrix}\]

\begin{itemize}
\tightlist
\item
  \textbf{2.6.} Legyen \(A\) olyan \(n \times n\)-es mátrix, amelyre
  \(A^2 = O\), azaz a nullmátrix. Mutasd meg, hogy \(A\) szükségképpen
  szinguláris.
\item
  \textbf{2.7.} Legyen
\end{itemize}

\[\boldsymbol{A} = \begin{bmatrix} 1 & 1+\epsilon \\ 1-\epsilon & 1 \end{bmatrix}.\]

\begin{enumerate}
\def\labelenumi{(\alph{enumi})}
\tightlist
\item
  Mi \(A\) determinánsa? (b) Lebegőpontos aritmetikában \(\epsilon\)
  mely értéktartományára lesz a determináns számított értéke nulla? (c)
  Mi \(A\) LU-felbontása? (d) Lebegőpontos aritmetikában \(\epsilon\)
  mely értéktartományára lesz \(U\) számított értéke szinguláris?
\end{enumerate}

\begin{itemize}
\tightlist
\item
  \textbf{2.8.} Legyen \(A\) és \(B\) tetszőleges két \(n \times n\)-es
  mátrix. (a) Bizonyítsd be, hogy \((AB)^T = B^T A^T\). (b) Ha \(A\) és
  \(B\) mindkettő nemszinguláris, bizonyítsd be, hogy
  \((AB)^{-1} = B^{-1} A^{-1}\).
\item
  \textbf{2.9.} Ha \(A\) tetszőleges nemszinguláris valós mátrix, mutasd
  meg, hogy \((A^{-1})^T = (A^T)^{-1}\). Ennek következtében az
  \(A^{-T}\) jelölés egyértelműen használható e mátrix jelölésére.
  Hasonlóan, ha \(A\) tetszőleges nemszinguláris komplex mátrix, akkor
  \((A^{-1})^H = (A^H)^{-1}\), és az \(A^{-H}\) jelölés egyértelműen
  használható e mátrix jelölésére.
\item
  \textbf{2.10.} Legyen \(P\) egy tetszőleges permutációs mátrix. (a)
  Bizonyítsd be, hogy \(P^{-1} = P^T\). (b) Bizonyítsd be, hogy \(P\)
  kifejezhető páronkénti cserék szorzataként.
\item
  \textbf{2.11.} Írj részletes algoritmust egy \(Lx = b\) alsó
  háromszögű lineáris rendszer előrehelyettesítéssel történő
  megoldására.
\item
  \textbf{2.12.} Ellenőrizd, hogy egy \(n\)-ed rendű alsó háromszögű
  rendszer előrehelyettesítéssel történő megoldásának műveletszámában
  (szorzások vagy összeadások számában) a domináns tag \(n^2/2\).
\item
  \textbf{2.13.} Hogyan oldanád meg a következő alakú particionált
  lineáris rendszert?
\end{itemize}

\[\begin{bmatrix} \boldsymbol{L}_1 & \boldsymbol{O} \\ \boldsymbol{B} & \boldsymbol{L}_2 \end{bmatrix} \begin{bmatrix} \boldsymbol{x} \\ \boldsymbol{y} \end{bmatrix} = \begin{bmatrix} \boldsymbol{b} \\ \boldsymbol{c} \end{bmatrix},\]

ahol \(\boldsymbol{L}_1\) és \(\boldsymbol{L}_2\) nemszinguláris alsó
háromszögmátrixok, és a megoldás- és jobb oldali vektorokat ennek
megfelelően particionáljuk. Mutasd be a konkrét lépéseket a megadott
almátrixok és vektorok segítségével.

\begin{itemize}
\tightlist
\item
  \textbf{2.14.} Bizonyítsd be az elemi eliminációs mátrixok 2.4.3.
  szakaszban felsorolt négy tulajdonságát.
\item
  \textbf{2.15.} (a) Bizonyítsd be, hogy két alsó háromszögmátrix
  szorzata alsó háromszögmátrix. (b) Bizonyítsd be, hogy egy
  nemszinguláris alsó háromszögmátrix inverze alsó háromszögmátrix.
\item
  \textbf{2.16.} (a) Mi az alábbi mátrix LU-felbontása?
\end{itemize}

\[\begin{bmatrix} 1 & a \\ c & b \end{bmatrix}\]

\begin{enumerate}
\def\labelenumi{(\alph{enumi})}
\setcounter{enumi}{1}
\tightlist
\item
  Milyen feltétel mellett szinguláris ez a mátrix?
\end{enumerate}

\begin{itemize}
\tightlist
\item
  \textbf{2.17.} Írd fel az alábbi mátrix LU-felbontását (mutasd be mind
  az \(L\)-t, mind az \(U\)-t explicit módon):
\end{itemize}

\[\begin{bmatrix} 1 & -1 & 0 \\ -1 & 2 & -1 \\ 0 & -1 & 1 \end{bmatrix}.\]

\begin{itemize}
\tightlist
\item
  \textbf{2.18.} Bizonyítsd be, hogy az
\end{itemize}

\[\mathbf{A} = \begin{bmatrix} 0 & 1 \\ 1 & 0 \end{bmatrix}\]

mátrixnak nincs LU-felbontása, azaz nem létezik olyan \(L\) alsó
háromszögmátrix és \(U\) felső háromszögmátrix, amelyekre \(A = LU\).

\begin{itemize}
\tightlist
\item
  \textbf{2.19.} Legyen \(A\) egy \(n \times n\)-es nemszinguláris
  mátrix. Tekintsük az alábbi algoritmust: 1. Végigfutva \(A\) \(1\)-től
  \(n\)-ig terjedő oszlopain, permutáljuk a sorokat szükség szerint úgy,
  hogy minden oszlopban a főátlón vagy az alatta lévő elemek közül a
  főátlóbeli elem legyen az abszolút értékben legnagyobb. Az eredmény
  \(PA\) valamilyen \(P\) permutációs mátrixszal. 2. Ezután
  főelemkiválasztás nélküli Gauss-kiküszöböléssel számítsuk ki \(PA\)
  LU-felbontását. (a) Numerikusan stabil-e ez az algoritmus? (b) Ha
  igen, magyarázd meg miért. Ha nem, adj ellenpéldát a szemléltetésre.
\item
  \textbf{2.20.} Bizonyítsd be, hogy ha részleges főelemkiválasztással
  végzett Gauss-kiküszöbölést alkalmazunk egy oszloponként diagonálisan
  domináns \(A\) mátrixra, akkor sorcserék nem fognak bekövetkezni.
\item
  \textbf{2.21.} Ha \(A\), \(B\) és \(C\) \(n \times n\)-es mátrixok,
  \(B\) és \(C\) nemszinguláris, és \(b\) egy \(n\)-dimenziós vektor,
  hogyan implementálnád az
\end{itemize}

\[x = B^{-1}(2A + I)(C^{-1} + A)b\]

képletet anélkül, hogy bármelyik mátrixinverzet kiszámítanád?

\begin{itemize}
\tightlist
\item
  \textbf{2.22.} Ellenőrizd, hogy egy \(n\)-ed rendű mátrix
  Gauss-kiküszöböléssel történő LU-felbontásának műveletszámában
  (szorzások vagy összeadások számában) a domináns tag \(n^3/3\).
\item
  \textbf{2.23.} Ellenőrizd, hogy egy \(n\)-ed rendű mátrix
  Gauss-kiküszöböléssel történő inverzének kiszámításához tartozó
  műveletszámban (szorzások vagy összeadások számában) a domináns tag
  \(n^3\).
\item
  \textbf{2.24.} Ellenőrizd, hogy egy \(n\)-ed rendű mátrixon
  végrehajtott Gauss--Jordan-kiküszöbölés műveletszámában (szorzások
  vagy összeadások számában) a domináns tag \(n^3/2\).
\item
  \textbf{2.25.} (a) Ha \(u\) és \(v\) nem nulla \(n\)-dimenziós
  vektorok, bizonyítsd be, hogy a \(uv^T\) \(n \times n\)-es külső
  szorzat mátrix rangja \(1\). (b) Ha \(A\) egy olyan \(n \times n\)-es
  mátrix, amelyre \(\operatorname{rank}(A) = 1\), bizonyítsd be, hogy
  léteznek olyan nem nulla \(u\) és \(v\) \(n\)-dimenziós vektorok,
  amelyekre \(A = uv^T\).
\item
  \textbf{2.26.} Egy \(n \times n\)-es \(A\) mátrixot \emph{elemi}nek
  nevezünk, ha az egységmátrixtól egy egyrangú mátrixban tér el, azaz
  \(A = I - uv^T\) valamilyen \(u\) és \(v\) \(n\)-dimenziós
  vektorokkal. (a) Ha \(A\) elemi, milyen feltétel biztosítja \(u\)-ra
  és \(v\)-re, hogy \(A\) nemszinguláris? (b) Ha \(A\) elemi és
  nemszinguláris, bizonyítsd be, hogy \(A^{-1}\) is elemi, megmutatva,
  hogy \(A^{-1} = I - \sigma u v^T\) valamilyen \(\sigma\) skalárral. Mi
  a konkrét \(\sigma\) érték \(u\) és \(v\) függvényében? (c) Elemi-e
  egy a 2.4.3. szakaszban definiált elemi eliminációs mátrix? Ha igen,
  mik \(u\), \(v\) és \(\sigma\) ebben az esetben?
\item
  \textbf{2.27.} Bizonyítsd be, hogy a 2.4.9. szakaszban megadott
  Sherman--Morrison-formula,
\end{itemize}

\[(\boldsymbol{A} - \boldsymbol{u}\boldsymbol{v}^T)^{-1} = A^{-1} + A^{-1}u(1 - v^T A^{-1}u)^{-1}v^T A^{-1},\]

helyes. (\emph{Útmutatás:} Szorozd meg mindkét oldalt \(A - uv^T\)-tal.)

\begin{itemize}
\tightlist
\item
  \textbf{2.28.} Bizonyítsd be, hogy a 2.4.9. szakaszban megadott
  Woodbury-formula,
\end{itemize}

\[(\boldsymbol{A} - \boldsymbol{U}\boldsymbol{V}^T)^{-1} = A^{-1} + A^{-1}U(I - V^T A^{-1}U)^{-1}V^T A^{-1},\]

helyes. (\emph{Útmutatás:} Szorozd meg mindkét oldalt \(A - UV^T\)-tal.)

\begin{itemize}
\tightlist
\item
  \textbf{2.29.} \(p = 1, 2\) és \(\infty\) esetén bizonyítsd be, hogy a
  vektor \(p\)-normák teljesítik a 2.3.1. szakasz végén megadott három
  tulajdonságot.
\item
  \textbf{2.30.} \(p = 1\) és \(\infty\) esetén bizonyítsd be, hogy a
  mátrix \(p\)-normák teljesítik a 2.3.2. szakasz végén megadott öt
  tulajdonságot.
\item
  \textbf{2.31.} Legyen \(A\) egy szimmetrikus pozitív definit mátrix.
  Mutasd meg, hogy az
\end{itemize}

\[\|\boldsymbol{x}\|_A = (\boldsymbol{x}^T \boldsymbol{A} \boldsymbol{x})^{1/2}\]

függvény teljesíti a vektornorma 2.3.1. szakasz végén megadott három
tulajdonságát. Ezt a vektornormát az \(A\) mátrix által
\emph{indukáltnak} mondjuk.

\begin{itemize}
\tightlist
\item
  \textbf{2.32.} Mutasd meg, hogy egy \(m \times n\)-es \(A\) mátrix
  alábbi függvényei teljesítik a mátrixnorma 2.3.2. szakasz végén
  megadott első három tulajdonságát, és ezért mátrixnormák az ott
  említett általánosabb értelemben. (a)
\end{itemize}

\[\|\boldsymbol{A}\|_{\max} = \max_{i,j} |a_{ij}|\]

Figyeljük meg, hogy ez egyszerűen \(A\)-nak az \(\mathbb{R}^{mn}\)-beli
vektorként tekintett \(\infty\)-normája. (b)

\[\|A\|_F = \left(\sum_{i,j} |a_{ij}|^2\right)^{1/2}\]

Figyeljük meg, hogy ez egyszerűen \(A\)-nak az \(\mathbb{R}^{mn}\)-beli
vektorként tekintett \(2\)-normája. Ezt \emph{Frobenius-mátrixnormának}
nevezzük.

\begin{itemize}
\tightlist
\item
  \textbf{2.33.} Bizonyítsd be vagy adj ellenpéldát: Ha \(A\)
  nemszinguláris mátrix, akkor \(\|A^{-1}\| = \|A\|^{-1}\).
\item
  \textbf{2.34.} Tegyük fel, hogy \(\boldsymbol{A}\) pozitív definit
  mátrix. (a) Mutasd meg, hogy \(\boldsymbol{A}\) szükségképpen
  nemszinguláris. (b) Mutasd meg, hogy \(A^{-1}\) szükségképpen pozitív
  definit.
\item
  \textbf{2.35.} Tegyük fel, hogy az \(A\) mátrixnak van \(A = BB^T\)
  alakú felbontása, ahol \(B\) nemszinguláris. Mutasd meg, hogy \(A\)
  szükségképpen szimmetrikus és pozitív definit.
\item
  \textbf{2.36.} Vezess le egy algoritmust, amely egy \(n \times n\)-es
  szimmetrikus pozitív definit \(\boldsymbol{A}\) mátrix
  \(\boldsymbol{L}\boldsymbol{L}^T\) Cholesky-felbontásának
  kiszámításához \(\boldsymbol{A}\) és
  \(\boldsymbol{L}\boldsymbol{L}^T\) megfelelő elemeinek egyenlővé
  tételével kapható.
\item
  \textbf{2.37.} Tegyük fel, hogy az
\end{itemize}

\[\boldsymbol{B} = \begin{bmatrix} \alpha & \boldsymbol{a}^T \\ \boldsymbol{a} & \boldsymbol{A} \end{bmatrix}\]

\((n+1)\)-ed rendű szimmetrikus mátrix pozitív definit. (a) Mutasd meg,
hogy az \(\alpha\) skalár szükségképpen pozitív, és az \(n \times n\)-es
\(\boldsymbol{A}\) mátrix szükségképpen pozitív definit. (b) Mi
\(\boldsymbol{B}\) Cholesky-felbontása az alkotó almátrixok
segítségével?

\begin{itemize}
\tightlist
\item
  \textbf{2.38.} Tegyük fel, hogy az
\end{itemize}

\[\boldsymbol{B} = \begin{bmatrix} \boldsymbol{A} & \boldsymbol{a} \\ \boldsymbol{a}^T & \alpha \end{bmatrix}\]

\((n+1)\)-ed rendű szimmetrikus mátrix pozitív definit. (a) Mutasd meg,
hogy az \(\alpha\) skalár szükségképpen pozitív, és az \(n \times n\)-es
\(\boldsymbol{A}\) mátrix szükségképpen pozitív definit. (b) Mi
\(\boldsymbol{B}\) Cholesky-felbontása az alkotó almátrixok
segítségével?

\begin{itemize}
\tightlist
\item
  \textbf{2.39.} Ellenőrizd, hogy egy \(n\)-ed rendű szimmetrikus
  pozitív definit mátrix Cholesky-felbontásának műveletszámában
  (szorzások vagy összeadások számában) a domináns tag \(n^3/6\).
\item
  \textbf{2.40.} Legyen \(\boldsymbol{A}\) egy \(\beta\) sávszélességű
  sávos mátrix, és tegyük fel, hogy a
  \(\boldsymbol{P}\boldsymbol{A} = \boldsymbol{L}\boldsymbol{U}\)
  LU-felbontást részleges főelemkiválasztással végzett
  Gauss-kiküszöböléssel számítjuk. Mutasd meg, hogy az
  \(\boldsymbol{U}\) felső háromszögű tényező sávszélessége legfeljebb
  \(2\beta\).
\item
  \textbf{2.41.} Legyen \(A\) egy nemszinguláris tridiagonális mátrix.
  (a) Mutasd meg, hogy általában \(A^{-1}\) sűrű. (b) Hasonlítsd össze
  az \(Ax = b\) lineáris rendszer Gauss-kiküszöböléssel és
  visszahelyettesítéssel történő megoldásához, valamint a rendszer
  explicit mátrixinvertálással történő megoldásához szükséges munkát és
  tárhelyet.
\end{itemize}

Ez a példa egy további okot szemléltet arra, hogy az explicit
mátrixinvertálás miért rendszerint rossz ötlet.

\begin{itemize}
\tightlist
\item
  \textbf{2.42.} (a) Készíts algoritmust egy nemszinguláris
  \(n \times n\)-es háromszögmátrix inverzének helyben történő
  kiszámítására, azaz további tömbtárhely nélkül. (b) Lehetséges-e egy
  általános nemszinguláris \(n \times n\)-es mátrix inverzét helyben
  kiszámítani? Ha igen, vázolj egy erre szolgáló algoritmust, ha nem,
  magyarázd meg miért. E feladat céljaira feltételezheted, hogy
  főelemkiválasztásra nincs szükség.
\item
  \textbf{2.43.} Tegyük fel, hogy egy \(Cz = d\) lineáris rendszert kell
  megoldanod, ahol \(C\) egy komplex \(n \times n\)-es mátrix, \(d\) és
  \(z\) pedig komplex \(n\)-dimenziós vektorok, de a lineáris
  egyenletrendszer-megoldód csak valós rendszereket kezel. Legyen
  \(C = A + iB\) és \(d = b + ic\), ahol \(A\), \(B\), \(b\) és \(c\)
  valósak, és \(i = \sqrt{-1}\). Mutasd meg, hogy a \(z = x + iy\)
  megoldást a
\end{itemize}

\[\begin{bmatrix} A & -B \\ B & A \end{bmatrix} \begin{bmatrix} x \\ y \end{bmatrix} = \begin{bmatrix} b \\ c \end{bmatrix}\]

\(2n \times 2n\)-es valós lineáris rendszer adja. Jó módja-e ez a
feladat megoldásának? Miért?

\subsubsection{Számítógépes
feladatok}\label{szuxe1muxedtuxf3guxe9pes-feladatok-1}

\begin{itemize}
\tightlist
\item
  \textbf{2.1.} (a) Mutasd meg, hogy az
\end{itemize}

\[\mathbf{A} = \begin{bmatrix} 0{,}1 & 0{,}2 & 0{,}3 \\ 0{,}4 & 0{,}5 & 0{,}6 \\ 0{,}7 & 0{,}8 & 0{,}9 \end{bmatrix}\]

mátrix szinguláris. Írd le az \(Ax = b\) rendszer megoldásainak
halmazát, ha

\[\boldsymbol{b} = \begin{bmatrix} 0{,}1 \\ 0{,}3 \\ 0{,}5 \end{bmatrix}.\]

\begin{enumerate}
\def\labelenumi{(\alph{enumi})}
\setcounter{enumi}{1}
\tightlist
\item
  Ha részleges főelemkiválasztással végzett Gauss-kiküszöböléssel
  akarnánk megoldani ezt a rendszert egzakt aritmetikával, hol akadna
  meg a folyamat? (c) Mivel \(A\) néhány eleme nem ábrázolható egzaktul
  egy bináris lebegőpontos rendszerben, a mátrix a számítógépbe bevitel
  után már nem egzaktul szinguláris, így a rendszer
  Gauss-kiküszöböléssel történő megoldása nem feltétlenül hiúsul meg.
  Oldd meg ezt a rendszert számítógépen egy Gauss-kiküszöbölést végző
  könyvtári rutinnal. Hasonlítsd össze a számított megoldást az (a)
  részben adott megoldáshalmaz leírásával. Ha a szoftvered tartalmaz
  kondícióbecslőt, mi a \(\operatorname{cond}(A)\) becsült értéke?
  Mennyi értékes jegyet várnál e becslés alapján a megoldásban?
\end{enumerate}

\begin{itemize}
\tightlist
\item
  \textbf{2.2.} (a) Egy Gauss-kiküszöbölést végző könyvtári rutinnal
  oldd meg az \(Ax = b\) rendszert, ahol
\end{itemize}

\[\boldsymbol{A} = \begin{bmatrix} 2 & 4 & -2 \\ 4 & 9 & -3 \\ -2 & -1 & 7 \end{bmatrix}, \quad \boldsymbol{b} = \begin{bmatrix} 2 \\ 8 \\ 10 \end{bmatrix}.\]

\begin{enumerate}
\def\labelenumi{(\alph{enumi})}
\setcounter{enumi}{1}
\tightlist
\item
  Az \(\boldsymbol{A}\) mátrix már kiszámított LU-felbontását
  felhasználva --- a mátrix újrafelbontása nélkül --- oldd meg az
  \(Ay = c\) rendszert, ahol
\end{enumerate}

\[\boldsymbol{c} = \begin{bmatrix} 4 \\ 8 \\ -6 \end{bmatrix}.\]

\begin{enumerate}
\def\labelenumi{(\alph{enumi})}
\setcounter{enumi}{2}
\tightlist
\item
  Ha az \(A\) mátrix úgy változik, hogy \(a_{1,2} = 2\), a
  Sherman--Morrison-féle frissítési technikával --- a mátrix
  újrafelbontása nélkül --- számítsd ki az új \(x\) megoldást az eredeti
  \(b\) jobb oldali vektort használva.
\end{enumerate}

\begin{itemize}
\tightlist
\item
  \textbf{2.3.} Az alábbi ábra egy síkbeli rácsszerkezetet ábrázol,
  amely 13 rúdból (a számozott egyenesek) és 10 csomópontból (a
  számozott körök) áll. A jelzett terhelések, tonnában, a 2., az 5. és a
  6. csomópontra hatnak, és meg szeretnénk határozni a rácsszerkezet
  egyes rúdjaira ható eredő erőt.
\end{itemize}

\pandocbounded{\includegraphics[keepaspectratio]{_page_51_Figure_16.jpeg}}

Ahhoz, hogy a rácsszerkezet statikus egyensúlyban legyen, minden
csomópontban se vízszintes, se függőleges irányú eredő erő ne lépjen
fel. Így a rúderőket úgy határozhatjuk meg, hogy minden csomópontban a
balról és a jobbról ható vízszintes erőket, hasonlóképpen a felfelé és a
lefelé ható függőleges erőket egyenlővé tesszük. A nyolc csomópontra ez
16 egyenletet adna, ami több, mint a 13 meghatározandó ismeretlen erő.
Ahhoz, hogy a rácsszerkezet statikailag határozott legyen, azaz hogy
egyetlen megoldás létezzen, feltételezzük, hogy az 1. csomópont
vízszintesen és függőlegesen is mereven rögzített, a 8. csomópont pedig
függőlegesen rögzített. A rúderőket vízszintes és függőleges
komponensekre felbontva és \(\alpha = \sqrt{2}/2\)-t bevezetve a
rúderőkre, \(f_i\)-re, a következő egyenletrendszert kapjuk:

Egy könyvtári rutinnal oldd meg ezt a lineáris egyenletrendszert az
\(f\) rúderővektorra. Megjegyzendő, hogy e rendszer mátrixa meglehetősen
ritka, így érdemes lehet sávos rendszer-megoldóval vagy általánosabb
ritka megoldóval is kísérletezned, bár ez a konkrét feladatpéldány túl
kicsi ahhoz, hogy ezek egy általános megoldóhoz képest jelentős előnyt
jelentsenek.

\begin{itemize}
\tightlist
\item
  \textbf{2.4.} Írj rutint egy \(\boldsymbol{A}\) mátrix
  kondíciószámának becslésére. Használhatod akár az \(1\)-normát, akár a
  \(\infty\)-normát (vagy próbáld ki mindkettőt, és hasonlítsd össze az
  eredményeket). Ki kell majd számítanod \(\|\boldsymbol{A}\|\)-t, ami
  könnyű, és meg kell becsülnöd \(\|\boldsymbol{A}^{-1}\|\)-t, ami
  nagyobb kihívás. Ahogy a 2.3.3. szakaszban tárgyaltuk,
  \(\|\boldsymbol{A}^{-1}\|\) becslésének egy módja, hogy olyan
  \(\boldsymbol{y}\) vektort választunk, amelyre a
  \(\|\boldsymbol{z}\|/\|\boldsymbol{y}\|\) arány nagy, ahol
  \(\boldsymbol{z}\) az
  \(\boldsymbol{A}\boldsymbol{z} = \boldsymbol{y}\) megoldása. Próbálj
  ki két különböző megközelítést \(\boldsymbol{y}\) választására: (a)
  Válaszd \(\boldsymbol{y}\)-t az \(A^T y = c\) rendszer megoldásaként,
  ahol \(c\) egy olyan vektor, amelynek minden komponense \(\pm 1\), és
  az egyes komponensekhez tartozó előjelet az alábbi heurisztikával
  választjuk. Az \(\mathbf{A} = \mathbf{L}\mathbf{U}\) felbontást
  felhasználva az \(\mathbf{A}^T \mathbf{y} = \mathbf{c}\) rendszert két
  lépésben oldjuk meg: egymás után megoldjuk az
  \(\mathbf{U}^T \mathbf{v} = \mathbf{c}\) és az
  \(\mathbf{L}^T \mathbf{y} = \mathbf{v}\) háromszögű rendszereket. Az
  első háromszögű megoldás minden lépésében \(c\) megfelelő komponensét
  \(1\)-nek vagy \(-1\)-nek választjuk, attól függően, hogy melyik
  választás teszi a kapott \(v\)-komponenst abszolút értékben nagyobbá.
  (Ehhez egy saját háromszögű megoldó rutint kell majd írnod e logika
  megvalósítására.) Ezután oldd meg a második háromszögű rendszert a
  szokott módon \(\boldsymbol{y}\)-ra. Az ötlet az, hogy az \(A\)-ban
  jelen lévő rossz kondicionáltság \(U\)-ban fog tükröződni, viszonylag
  nagy \(v\)-t eredményezve. A viszonylag jól kondicionált
  egységháromszögű \(\boldsymbol{L}\) mátrix aztán megőrzi ezt a
  viszonyt, viszonylag nagy \(\boldsymbol{y}\)-t eredményezve. (b)
  Válassz néhány --- mondjuk öt --- különböző \(\boldsymbol{y}\) vektort
  véletlenszerűen, és azt használd, amelyik a legnagyobb
  \(\|\boldsymbol{z}\|/\|\boldsymbol{y}\|\) arányt adja. (Ehhez
  közönséges háromszögű megoldó rutint használhatsz.)
\end{itemize}

Az \(\boldsymbol{A}\) szükséges LU-felbontásához könyvtári rutint
használhatsz. Teszteld mindkét megközelítést az alábbi mátrixokon:

\[\mathbf{A}_1 = \begin{bmatrix} 10 & -7 & 0 \\ -3 & 2 & 6 \\ 5 & -1 & 5 \end{bmatrix},\]

\[\mathbf{A}_2 = \begin{bmatrix} -73 & 78 & 24 \\ 92 & 66 & 25 \\ -80 & 37 & 10 \end{bmatrix}.\]

Hogyan viszonyulnak egymáshoz a két módszerrel kapott eredmények? A
becsléseid minőségének ellenőrzéséhez számítsd ki \(\mathbf{A}^{-1}\)-et
explicit módon, hogy meghatározd a valódi normáját (ehhez a számításhoz
is felhasználhatod a már kiszámított LU-felbontást). Ha van hozzáférésed
olyan lineáris egyenletrendszer-szoftverhez, amely már tartalmaz
kondícióbecslőt, hogyan viszonyulnak az eredményeid az övéhez?

\begin{itemize}
\tightlist
\item
  \textbf{2.5.} (a) Egy egyszeres pontosságú, Gauss-kiküszöbölést végző
  rutinnal oldd meg az \(Ax = b\) rendszert, ahol
\end{itemize}

\[\boldsymbol{A} = \begin{bmatrix} 21{,}0 & 67{,}0 & 88{,}0 & 73{,}0 \\ 76{,}0 & 63{,}0 & 7{,}0 & 20{,}0 \\ 0{,}0 & 85{,}0 & 56{,}0 & 54{,}0 \\ 19{,}3 & 43{,}0 & 30{,}2 & 29{,}4 \end{bmatrix},\]

\[\boldsymbol{b} = \begin{bmatrix} 141{,}0 \\ 109{,}0 \\ 218{,}0 \\ 93{,}7 \end{bmatrix}.\]

\begin{enumerate}
\def\labelenumi{(\alph{enumi})}
\setcounter{enumi}{1}
\tightlist
\item
  Számítsd ki az \(r = b - Ax\) maradékot dupla pontosságú
  aritmetikával, ha elérhető (de a végső eredményt egyszeres pontosságú
  \(r\) vektorba tárolva). Megjegyzendő, hogy a megoldó rutin
  elpusztíthatja az \(A\)-t tartalmazó tömböt, így külön másolatot kell
  majd megőrizned a maradék számításához. (Ha a használt számítási
  környezetben csak egyetlen precízió érhető el, akkor végezd el a
  teljes feladatot abban a precízióban.) (c) Oldd meg az \(Az = r\)
  lineáris rendszert a „javított" \(x + z\) megoldás megkapásához.
  Megjegyzendő, hogy \(A\)-t nem kell újrafelbontani. (d) Ismételd a (b)
  és (c) lépéseket, amíg további javulást nem tapasztalsz.
\end{enumerate}

\begin{itemize}
\tightlist
\item
  \textbf{2.6.} Egy \(n \times n\)-es \(\boldsymbol{H}\) Hilbert-mátrix
  elemei \(h_{ij} = 1/(i+j-1)\), így az alakja
\end{itemize}

\[\begin{bmatrix} 1 & 1/2 & 1/3 & \cdots \\ 1/2 & 1/3 & 1/4 & \cdots \\ 1/3 & 1/4 & 1/5 & \cdots \\ \vdots & \vdots & \vdots & \ddots \end{bmatrix}.\]

\(n = 2, 3, \ldots\) esetén generáld az \(n\)-ed rendű Hilbert-mátrixot,
valamint a \(\boldsymbol{b} = \boldsymbol{H}\boldsymbol{x}\)
\(n\)-dimenziós vektort, ahol \(\boldsymbol{x}\) az az \(n\)-dimenziós
vektor, amelynek minden komponense \(1\). Egy Gauss-kiküszöbölést (vagy
Cholesky-felbontást, mivel a Hilbert-mátrix szimmetrikus és pozitív
definit) végző könyvtári rutinnal oldd meg a kapott
\(\boldsymbol{H}\boldsymbol{x} = \boldsymbol{b}\) lineáris rendszert,
\(\hat{x}\) közelítő megoldást kapva. Számítsd ki az
\(\mathbf{r} = \mathbf{b} - \mathbf{H}\hat{x}\) maradék és a
\(\Delta \mathbf{x} = \hat{x} - \mathbf{x}\) hiba \(\infty\)-normáját,
ahol \(\mathbf{x}\) a csupa egyes vektor. Milyen nagyra veheted
\(n\)-et, mielőtt a hiba 100 százalékossá válik (azaz a megoldásban
nincs egyetlen értékes jegy sem)? Emellett egy kondícióbecslővel
határozd meg \(\operatorname{cond}(\boldsymbol{H})\)-t minden \(n\)-re.
Próbáld meg jellemezni a kondíciószámot \(n\) függvényében. Ahogy \(n\)
változik, hogyan viszonyul a számított megoldás komponenseiben lévő
helyes számjegyek száma a mátrix kondíciószámához?

\begin{itemize}
\tightlist
\item
  \textbf{2.7.} (a) Mi történik, ha részleges főelemkiválasztással
  végzett Gauss-kiküszöbölést alkalmazunk a következő alakú mátrixra?
\end{itemize}

\[\begin{bmatrix}
1 & 0 & 0 & 0 & 1 \\
-1 & 1 & 0 & 0 & 1 \\
-1 & -1 & 1 & 0 & 1 \\
-1 & -1 & -1 & 1 & 1 \\
-1 & -1 & -1 & -1 & 1
\end{bmatrix}\]

Nőnek-e a transzformált mátrix elemei? Mi történik, ha ehelyett teljes
főelemkiválasztást használunk? (Megjegyzendő, hogy az (a) rész nem
igényel számítógépet.)

\begin{enumerate}
\def\labelenumi{(\alph{enumi})}
\setcounter{enumi}{1}
\tightlist
\item
  Egy részleges főelemkiválasztással végzett Gauss-kiküszöbölést
  alkalmazó könyvtári rutinnal oldj meg ilyen alakú, különböző méretű
  lineáris rendszereket, olyan jobb oldali vektorokat választva, amelyek
  esetén a megoldás ismert. Hogyan viselkedik a hiba, a maradék és a
  kondíciószám, ahogy a rendszerek egyre nagyobbak lesznek? Ez a
  mesterségesen megkonstruált rendszer a 2.4.5. szakaszban idézett,
  legrosszabb esetben fellépő növekedési tényezőt illusztrálja, és nem
  jellemző a részleges főelemkiválasztással végzett Gauss-kiküszöbölés
  szokásos viselkedésére.
\end{enumerate}

\begin{itemize}
\tightlist
\item
  \textbf{2.8.} Egy \(Ax = b\) lineáris rendszer mindkét oldalának egy
  \(D\) nemszinguláris diagonális mátrixszal való szorzása, hogy egy új
  \(DAx = Db\) rendszert kapjunk, egyszerűen átskálázza a rendszer
  sorait, és elméletben nem változtatja meg a megoldást. Az ilyen
  skálázás azonban befolyásolja a mátrix kondíciószámát és a
  Gauss-kiküszöbölésben a főelemek választását, így véges precíziójú
  aritmetikában befolyásolhatja a megoldás pontosságát. Megjegyzendő,
  hogy a skálázás némi kerekítési hibát vihet a mátrixba, hacsak \(D\)
  elemei nem a használt lebegőpontos aritmetikai rendszer alapjának
  hatványai (miért?).
\end{itemize}

Véletlenszerűen választott \(\boldsymbol{A}\) mátrixú lineáris rendszer
és olyan \(\boldsymbol{b}\) jobb oldali vektor használatával, amely
mellett a megoldás ismert, kísérletezz különféle \(D\) skálázó
mátrixokkal, hogy milyen hatással vannak a \(DA\) mátrix kondíciószámára
és a \(DAx = Db\) lineáris rendszer megoldására egy erre szolgáló
könyvtári rutin által adott megoldásra. Mindenképpen próbálj ki elég
torz skálázásokat is, amelyeknél \(D\) diagonális elemeinek
nagyságrendje jelentősen különbözik egymástól (a cél egy rosszul
megválasztott mértékegységekkel adott rendszer szimulálása). Hasonlítsd
össze mind a relatív maradékokat, mind a különféle skálázások által
adott hibát. Találsz-e olyan skálázást, amely nagyon rossz pontosságot
ad? Ebben az esetben a maradék továbbra is kicsi marad-e?

\begin{itemize}
\tightlist
\item
  \textbf{2.9.} (a) Főelemkiválasztás nélküli Gauss-kiküszöböléssel oldd
  meg az
\end{itemize}

\[\begin{bmatrix} \epsilon & 1 \\ 1 & 1 \end{bmatrix} \begin{bmatrix} x_1 \\ x_2 \end{bmatrix} = \begin{bmatrix} 1+\epsilon \\ 2 \end{bmatrix}\]

lineáris rendszert \(\epsilon = 10^{-2k}\) esetén,
\(k = 1, \ldots, 10\). Az egzakt megoldás
\(x = \begin{bmatrix} 1 & 1 \end{bmatrix}^T\), \(\epsilon\) értékétől
függetlenül. Hogyan viselkedik a számított megoldás pontossága, ahogy
\(\epsilon\) értéke csökken?

\begin{enumerate}
\def\labelenumi{(\alph{enumi})}
\setcounter{enumi}{1}
\tightlist
\item
  Ismételd meg az (a) részt, továbbra is főelemkiválasztás nélküli
  Gauss-kiküszöbölést használva, de ezúttal a megoldás javítására egy
  lépés iteratív finomítást is alkalmazz, a maradékot ugyanolyan
  precízióval számítva, mint a számítások többi részét. Most hogyan
  viselkedik a számított megoldás pontossága, ahogy \(\epsilon\) értéke
  csökken?
\end{enumerate}

\begin{itemize}
\tightlist
\item
  \textbf{2.10.} Tekintsük az
\end{itemize}

\[\begin{bmatrix} 1 & 1+\epsilon \\ 1-\epsilon & 1 \end{bmatrix} \begin{bmatrix} x_1 \\ x_2 \end{bmatrix} = \begin{bmatrix} 1+(1+\epsilon)\epsilon \\ 1 \end{bmatrix}\]

lineáris rendszert, ahol \(\epsilon\) egy rögzítendő kis paraméter. Az
egzakt megoldás nyilvánvalóan

\[\boldsymbol{x} = \begin{bmatrix} 1 \\ \epsilon \end{bmatrix}\]

\(\epsilon\) bármely értékére.

Egy Gauss-kiküszöbölésen alapuló könyvtári rutinnal oldd meg ezt a
rendszert. Kísérletezz \(\epsilon\) különféle értékeivel, különösen a
saját számítógéped \(\sqrt{\epsilon_{\text{mach}}}\)-ához közeli
értékekkel. \(\epsilon\) minden kipróbált értékére becsüld meg a mátrix
kondíciószámát és a megoldás minden komponensében lévő relatív hibát.
Milyen pontosan van meghatározva az egyes komponensek? Hogyan viszonyul
az egyes komponensekre kapott pontosság a mátrix kondíciószáma és a
2.3.4. szakaszban megadott hibakorlátok alapján várhatóhoz? Milyen
következtetéseket tudsz levonni ebből a kísérletből?

\begin{itemize}
\tightlist
\item
  \textbf{2.11.} (a) Írj programokat, amelyek megvalósítják a
  főelemkiválasztás nélküli, a részleges főelemkiválasztással és a
  teljes főelemkiválasztással végzett Gauss-kiküszöbölést. (b) Generálj
  több véletlen mátrixú lineáris rendszert (azaz a mátrix elemeinek
  előállítására használj véletlenszám-generátort), valamint olyan jobb
  oldali vektorokat, amelyek mellett a megoldás ismert, és hasonlítsd
  össze a három implementáció pontosságát, maradékait és teljesítményét.
  (c) Tudsz-e alkotni egy (nem véletlen) mátrixot, amelyre a teljes
  főelemkiválasztás lényegesen pontosabb, mint a részleges
  főelemkiválasztás?
\item
  \textbf{2.12.} Írj rutint tridiagonális lineáris egyenletrendszerek
  megoldására a 2.5.3. szakaszban megadott algoritmussal, és teszteld
  néhány mintapéldán. Írd le, hogyan változna a rutinod, ha részleges
  főelemkiválasztást vennél be. Írd le, hogyan változna a rutinod, ha a
  rendszer pozitív definit volna, és LU-felbontás helyett a
  Cholesky-felbontást számítanád ki.
\item
  \textbf{2.13.} Egy háromszögmátrix determinánsa a főátlóbeli elemeinek
  szorzatával egyenlő. Használd ezt a tényt egy olyan rutin
  kidolgozásához, amely egy tetszőleges \(n \times n\)-es
  \(\boldsymbol{A}\) mátrix determinánsát számítja ki az LU-felbontása
  segítségével. Használhatsz egy részleges főelemkiválasztással végzett
  Gauss-kiküszöbölést megvalósító könyvtári rutint az LU-felbontás
  megszerzésére, vagy tervezheted saját rutinodat. Hogyan tudod
  meghatározni a determináns megfelelő előjelét? A túlcsordulás vagy
  alulcsordulás veszélyének elkerülésére érdemes lehet a determináns
  tényleges értéke helyett a logaritmusának kiszámítását fontolóra
  venned.
\item
  \textbf{2.14.} Írj programokat, amelyek megvalósítják a \(C = AB\)
  mátrixszorzást --- ahol \(A\) egy \(m \times n\)-es és \(B\) egy
  \(n \times k\)-s mátrix --- két különböző módon: (a) Számítsd ki az
  \(\boldsymbol{A}\) sorainak és \(\boldsymbol{B}\) oszlopainak \(mk\)
  darab belső szorzatát. (b) \(C\) minden oszlopát \(A\) oszlopainak
  lineáris kombinációjaként állítsd elő.
\end{itemize}

BLAS-terminológiában (lásd a 2.7.2. szakaszt) az első implementáció az
\texttt{sdot}-ot, a második pedig az \texttt{saxpy}-t használja.
Hasonlítsd össze e két implementáció teljesítményét a saját
számítógépeden. Előfordulhat, hogy elég nagy mátrixokat kell
kipróbálnod, mielőtt a teljesítménykülönbségek jelentőssé válnának. Tudj
meg annyit a számítógéprendszeredről, amennyit csak tudsz (pl. a
gyorsítótár mérete és kezelési stratégiája), és használd fel ezt az
információt a megfigyelt eredmények magyarázatára.

\begin{itemize}
\tightlist
\item
  \textbf{2.15.} Valósítsd meg a Gauss-kiküszöbölést a hármas egymásba
  ágyazott ciklus mind a hat különböző sorrendjével, és hasonlítsd össze
  a teljesítményüket a saját számítógépeden. E feladat céljaira a
  numerikus stabilitás érdekében végzett főelemkiválasztást figyelmen
  kívül hagyhatod, de mindenképpen olyan tesztmátrixokat használj,
  amelyek nem igényelnek főelemkiválasztást. Előfordulhat, hogy elég
  nagy rendszert kell kipróbálnod, mielőtt a teljesítménykülönbségek
  jelentőssé válnának. Tudj meg annyit a számítógéprendszeredről,
  amennyit csak tudsz (pl. a gyorsítótár mérete és kezelési
  stratégiája), és használd fel ezt az információt a megfigyelt
  eredmények magyarázatára.
\item
  \textbf{2.16.} Háromszögű lineáris rendszerek megoldásához mind az
  előre-, mind a visszahelyettesítés olyan egymásba ágyazott ciklusokat
  tartalmaz, amelyek két indexe bármelyik sorrendben futhat. Valósítsd
  meg mind az előre-, mind a visszahelyettesítést a két indexsorrend
  mindegyikével (összesen négy algoritmust), és hasonlítsd össze
  teljesítményüket különböző méretű háromszögű tesztmátrixokon.
  Előfordulhat, hogy elég nagy rendszert kell kipróbálnod, mielőtt a
  teljesítménykülönbségek jelentőssé válnának. Az indexsorrendek legjobb
  választása ugyanaz-e mindkét algoritmusnál? Tudj meg annyit a
  számítógéprendszeredről, amennyit csak tudsz (pl. a gyorsítótár mérete
  és kezelési stratégiája), és használd fel ezt az információt a
  megfigyelt eredmények magyarázatára.
\item
  \textbf{2.17.} Tekintsünk egy vízszintes, egyik végén befogott, a
  hossza fennmaradó részén szabad tartót. A tartón ható erők egy
  diszkrét modellje egy \(Ax = b\) lineáris egyenletrendszert ad, ahol
  az \(n \times n\)-es \(A\) mátrix a következő sávos alakú:
\end{itemize}

\[\begin{bmatrix} 9 & -4 & 1 & 0 & \cdots & \cdots & 0 \\ -4 & 6 & -4 & 1 & \ddots & & \vdots \\ 1 & -4 & 6 & -4 & 1 & \ddots & \vdots \\ 0 & \ddots & \ddots & \ddots & \ddots & \ddots & 0 \\ \vdots & \ddots & 1 & -4 & 6 & -4 & 1 \\ \vdots & & \ddots & 1 & -4 & 5 & -2 \\ 0 & \cdots & \cdots & 0 & 1 & -2 & 1 \end{bmatrix},\]

a \(\boldsymbol{b}\) \(n\)-dimenziós vektor a tartóra ható ismert
terhelés (a saját súlyával együtt), a \(\boldsymbol{x}\) \(n\)-dimenziós
vektor pedig a tartó meghatározandó, eredő lehajlását jelöli. A tartót
egyenletesen terheltnek tekintjük, úgy, hogy a terhelésvektor minden
komponensére \(b_i = 1/n^4\).

\begin{enumerate}
\def\labelenumi{(\alph{enumi})}
\item
  \(n = 100\) esetén oldd meg ezt a lineáris rendszert mind egy sűrű
  lineáris rendszerekhez szolgáló standard könyvtári rutinnal, mind egy
  sávos (vagy általánosabb ritka) rendszerekhez tervezett könyvtári
  rutinnal. Hogyan viszonyul egymáshoz a két rutin a megoldás
  kiszámításához szükséges időben? Mennyire egyeznek meg a kapott
  válaszok egymással?
\item
  Ellenőrizd, hogy az \(A\) mátrixnak van \(A = RR^T\) UL-felbontása,
  ahol \(R\) a következő alakú felső háromszögmátrix:
\end{enumerate}

\begin{longtable}[]{@{}
  >{\raggedright\arraybackslash}p{(\linewidth - 10\tabcolsep) * \real{0.2184}}
  >{\raggedright\arraybackslash}p{(\linewidth - 10\tabcolsep) * \real{0.1724}}
  >{\raggedright\arraybackslash}p{(\linewidth - 10\tabcolsep) * \real{0.0460}}
  >{\raggedright\arraybackslash}p{(\linewidth - 10\tabcolsep) * \real{0.0460}}
  >{\raggedright\arraybackslash}p{(\linewidth - 10\tabcolsep) * \real{0.0460}}
  >{\raggedright\arraybackslash}p{(\linewidth - 10\tabcolsep) * \real{0.4713}}@{}}
\toprule\noalign{}
\begin{minipage}[b]{\linewidth}\raggedright
\(\lceil 2 \rceil\)
\end{minipage} & \begin{minipage}[b]{\linewidth}\raggedright
-2
\end{minipage} & \begin{minipage}[b]{\linewidth}\raggedright
1
\end{minipage} & \begin{minipage}[b]{\linewidth}\raggedright
0
\end{minipage} & \begin{minipage}[b]{\linewidth}\raggedright
\end{minipage} & \begin{minipage}[b]{\linewidth}\raggedright
0 {]}
\end{minipage} \\
\midrule\noalign{}
\endhead
\bottomrule\noalign{}
\endlastfoot
0 & -21∴. & -2 & 1 & ٠. & : \\
: & ٠. & ٠ & ٠. & ٠ & 0 \\
: & & ٠. & & -2 & \\
: & & & ٠. & 1 & \(\begin{bmatrix} -2 \\ 1 \end{bmatrix}\) \\
& & & & 0 & 1 \\
\end{longtable}

\(n = 1000\) esetén oldd meg a lineáris rendszert e felbontással (két
háromszögű megoldásra lesz szükség). Emellett oldd meg a rendszert az
eredeti alakjában is egy sávos (vagy általános ritka)
rendszer-megoldóval, ahogy az (a) részben. Mennyire egyeznek meg a
kapott válaszok egymással? Melyik megközelítés tűnik pontosabbnak? Mi
\(\boldsymbol{A}\) kondíciószáma, és milyen pontosságot sugall, hogy
várnunk kellene? Próbáld ki az iteratív finomítást, hogy lássad, a
kevésbé pontos módszernél javul-e a pontosság vagy a maradék.

\end{document}
